\documentclass[11pt,a4paper]{article}
\usepackage[T1]{fontenc}
\usepackage[utf8]{inputenc}
\usepackage{lmodern,amsmath,amssymb,amsthm,mathtools}
\usepackage{geometry,microtype,booktabs,array,longtable,xcolor}
\usepackage{fancyhdr,fancyvrb,fvextra,needspace}
\usepackage{enumitem}
\usepackage{hyperref,bookmark}
\geometry{left=25mm,right=25mm,top=25mm,bottom=25mm,headheight=14pt,headsep=9mm}
\definecolor{Ink}{HTML}{17293C}
\definecolor{Link}{HTML}{235577}
\hypersetup{colorlinks=true,linkcolor=Ink,citecolor=Link,urlcolor=Link,
  pdftitle={Global Proofs, Proof Candidates, and Counterexamples from Twenty-Five Mathematical Research Rounds},
  pdfauthor={Keio and Shaitetsu},bookmarksnumbered=true}
\urlstyle{same}
\setcounter{tocdepth}{1}
\setcounter{secnumdepth}{3}
\setlength{\parindent}{1.2em}
\setlength{\parskip}{3pt}
\setlength{\emergencystretch}{3em}
\allowdisplaybreaks[2]
\setlist{itemsep=2pt,topsep=4pt}
\pagestyle{fancy}
\fancyhf{}
\fancyhead[L]{\small\textsc{Global proofs and counterexamples}}
\fancyhead[R]{\small\textsc{Integrated edition}}
\fancyfoot[C]{\thepage}
\renewcommand{\headrulewidth}{0.3pt}
\newcolumntype{L}[1]{>{\raggedright\arraybackslash}p{#1}}
\newcommand{\SourceID}[1]{\textsf{\footnotesize #1}}
\newtheorem{SourceRbdTtheorem}{Theorem}
\renewcommand{\theSourceRbdTtheorem}{\thesection.\arabic{SourceRbdTtheorem}}
\newtheorem{SourceRbdTlemma}[SourceRbdTtheorem]{Lemma}
\newtheorem{SourceRbdTcorollary}[SourceRbdTtheorem]{Corollary}
\newtheorem{SourceRbgMtheorem}{Theorem}
\renewcommand{\theSourceRbgMtheorem}{\thesection.\arabic{SourceRbgMtheorem}}
\newtheorem{SourceRbgMlemma}[SourceRbgMtheorem]{Lemma}
\newtheorem{SourceRbgMcorollary}[SourceRbgMtheorem]{Corollary}
\newtheorem{SourceRbhMtheorem}{Theorem}
\renewcommand{\theSourceRbhMtheorem}{\thesection.\arabic{SourceRbhMtheorem}}
\newtheorem{SourceRbiKtheorem}{Theorem}
\renewcommand{\theSourceRbiKtheorem}{\thesection.\arabic{SourceRbiKtheorem}}
\newtheorem{SourceRbiKlemma}[SourceRbiKtheorem]{Lemma}
\newtheorem{SourceRbjKtheorem}{Theorem}
\renewcommand{\theSourceRbjKtheorem}{\thesection.\arabic{SourceRbjKtheorem}}
\newtheorem{SourceRbjKcorollary}[SourceRbjKtheorem]{Corollary}
\newtheorem{SourceRcaMtheorem}{Theorem}
\renewcommand{\theSourceRcaMtheorem}{\thesection.\arabic{SourceRcaMtheorem}}
\newtheorem{SourceRcaMlemma}[SourceRcaMtheorem]{Lemma}
\newtheorem{SourceRcbTtheorem}{Theorem}
\renewcommand{\theSourceRcbTtheorem}{\thesection.\arabic{SourceRcbTtheorem}}
\newtheorem{SourceRcbTlemma}[SourceRcbTtheorem]{Lemma}
\newtheorem*{SourceRccStheorem}{Theorem}
\newtheorem*{SourceRccSlemma}{Certificate lemma}
\newtheorem{SourceRcdMtheorem}{Theorem}
\renewcommand{\theSourceRcdMtheorem}{\thesection.\arabic{SourceRcdMtheorem}}
\newtheorem{SourceRcdMlemma}{Lemma}
\renewcommand{\theSourceRcdMlemma}{\thesection.\arabic{SourceRcdMlemma}}
\newtheorem{SourceRceHqtheorem}{Theorem}
\renewcommand{\theSourceRceHqtheorem}{\thesection.\arabic{SourceRceHqtheorem}}
\newtheorem{SourceRceHqlemma}[SourceRceHqtheorem]{Lemma}
\newtheorem{SourceRceHqcorollary}[SourceRceHqtheorem]{Corollary}
\newtheorem{SourceRcfMtheorem}{Theorem}
\renewcommand{\theSourceRcfMtheorem}{\thesection.\arabic{SourceRcfMtheorem}}
\newtheorem{SourceRcfMlemma}[SourceRcfMtheorem]{Lemma}
\newtheorem{SourceRcfMcorollary}[SourceRcfMtheorem]{Corollary}
\newtheorem{SourceRcfKtheorem}{Theorem}
\renewcommand{\theSourceRcfKtheorem}{\thesection.\arabic{SourceRcfKtheorem}}
\newtheorem{SourceRcfKproposition}[SourceRcfKtheorem]{Proposition}
\newtheorem{SourceRcfKlemma}[SourceRcfKtheorem]{Lemma}
\newtheorem{SourceRacKtheorem}{Theorem}
\renewcommand{\theSourceRacKtheorem}{\thesection.\arabic{SourceRacKtheorem}}
\newtheorem{SourceRacKlemma}[SourceRacKtheorem]{Lemma}
\newtheorem{SourceRacKcorollary}[SourceRacKtheorem]{Corollary}
\newtheorem{SourceRbcKtheorem}{Theorem}
\renewcommand{\theSourceRbcKtheorem}{\thesection.\arabic{SourceRbcKtheorem}}
\newtheorem{SourceRbiHtheorem}{Theorem}
\renewcommand{\theSourceRbiHtheorem}{\thesection.\arabic{SourceRbiHtheorem}}
\newtheorem{SourceRbiHlemma}{Lemma}
\renewcommand{\theSourceRbiHlemma}{\thesection.\arabic{SourceRbiHlemma}}
\newtheorem{SourceRbiTtheorem}{Theorem}
\renewcommand{\theSourceRbiTtheorem}{\thesection.\arabic{SourceRbiTtheorem}}
\newtheorem{SourceRbiTlemma}{Lemma}
\renewcommand{\theSourceRbiTlemma}{\thesection.\arabic{SourceRbiTlemma}}
\newtheorem{SourceRbiRtheorem}{Claim}
\renewcommand{\theSourceRbiRtheorem}{\thesection.\arabic{SourceRbiRtheorem}}
\newtheorem{SourceRbjMtheorem}{Theorem}
\renewcommand{\theSourceRbjMtheorem}{\thesection.\arabic{SourceRbjMtheorem}}
\newtheorem{SourceRbjMlemma}{Lemma}
\renewcommand{\theSourceRbjMlemma}{\thesection.\arabic{SourceRbjMlemma}}
\newtheorem{SourceRcaTtheorem}{Theorem}
\renewcommand{\theSourceRcaTtheorem}{\thesection.\arabic{SourceRcaTtheorem}}
\newtheorem{SourceRcaHtheorem}{Theorem}
\renewcommand{\theSourceRcaHtheorem}{\thesection.\arabic{SourceRcaHtheorem}}
\newtheorem{SourceRcaHlemma}[SourceRcaHtheorem]{Lemma}
\newtheorem{SourceRcaRtheorem}{Theorem}
\renewcommand{\theSourceRcaRtheorem}{\thesection.\arabic{SourceRcaRtheorem}}
\newtheorem{SourceRcaStheorem}{Theorem}
\renewcommand{\theSourceRcaStheorem}{\thesection.\arabic{SourceRcaStheorem}}
\newtheorem{SourceRcbRtheorem}{Theorem}
\renewcommand{\theSourceRcbRtheorem}{\thesection.\arabic{SourceRcbRtheorem}}
\newtheorem{SourceRcbRlemma}{Lemma}
\renewcommand{\theSourceRcbRlemma}{\thesection.\arabic{SourceRcbRlemma}}
\newtheorem{SourceRcbStheorem}{Theorem}
\renewcommand{\theSourceRcbStheorem}{\thesection.\arabic{SourceRcbStheorem}}
\newtheorem{SourceRcbSlemma}[SourceRcbStheorem]{Lemma}
\newtheorem*{SourceRccHthm}{Theorem}
\newtheorem{SourceRcdKtheorem}{Theorem}
\renewcommand{\theSourceRcdKtheorem}{\thesection.\arabic{SourceRcdKtheorem}}
\newtheorem{SourceRcdKlemma}[SourceRcdKtheorem]{Lemma}
\newtheorem*{SourceRcdStheorem}{Theorem}
\newtheorem{SourceRceMtheorem}{Theorem}
\renewcommand{\theSourceRceMtheorem}{\thesection.\arabic{SourceRceMtheorem}}
\newtheorem{SourceRceMlemma}{Lemma}
\renewcommand{\theSourceRceMlemma}{\thesection.\arabic{SourceRceMlemma}}
\newtheorem{SourceRceStheorem}{Theorem}
\renewcommand{\theSourceRceStheorem}{\thesection.\arabic{SourceRceStheorem}}
\newtheorem{SourceRceSlemma}{Lemma}
\renewcommand{\theSourceRceSlemma}{\thesection.\arabic{SourceRceSlemma}}
\newtheorem{SourceRceHbtheorem}{Theorem}
\renewcommand{\theSourceRceHbtheorem}{\thesection.\arabic{SourceRceHbtheorem}}
\newtheorem{SourceRceHbcorollary}[SourceRceHbtheorem]{Corollary}
\newtheorem{SourceRceTtheorem}{Theorem}
\renewcommand{\theSourceRceTtheorem}{\thesection.\arabic{SourceRceTtheorem}}
\newtheorem{SourceRceTcorollary}{Corollary}
\renewcommand{\theSourceRceTcorollary}{\thesection.\arabic{SourceRceTcorollary}}
\newtheorem{SourceRbhRtheorem}{Theorem}
\renewcommand{\theSourceRbhRtheorem}{\thesection.\arabic{SourceRbhRtheorem}}
\newtheorem{SourceRbhRlemma}[SourceRbhRtheorem]{Lemma}
\newtheorem{SourceRbbHtheorem}{Theorem}
\renewcommand{\theSourceRbbHtheorem}{\thesection.\arabic{SourceRbbHtheorem}}
\newtheorem{SourceRbcTtheorem}{Theorem}
\renewcommand{\theSourceRbcTtheorem}{\thesection.\arabic{SourceRbcTtheorem}}
\newtheorem{SourceRbcTlemma}[SourceRbcTtheorem]{Lemma}
\newtheorem{SourceRceKtheorem}{Theorem}
\renewcommand{\theSourceRceKtheorem}{\thesection.\arabic{SourceRceKtheorem}}
\newtheorem{SourceRceKlemma}[SourceRceKtheorem]{Lemma}
\begin{document}
\thispagestyle{empty}
\begin{center}
{\small\textsc{Integrated mathematical manuscript \quad | \quad Rounds 1--25}}\\[14pt]
{\LARGE\bfseries Global Proofs, Proof Candidates,\\[4pt]
and Counterexamples}\\[10pt]
{\large From Twenty-Five Mathematical Research Rounds}\\[18pt]
{\large Keio and Shaitetsu}\\[5pt]
{\small Integrated edition, 20 November 2026}
\end{center}
\vspace{9pt}
\begin{abstract}
This paper places the complete mathematical arguments from a selected
twenty-five-round research archive into one continuous manuscript.
It incorporates forty-one original \LaTeX{} manuscripts and one
reconstructed proof checkpoint, grouped into forty-one principal
sections. The subjects include exact graph colorings, prime-power
Diophantine equations, binomial transforms and divisibility, numerical
semigroups, special-function inequalities, total positivity,
set-partition sequences, real-rootedness, modular congruences, and
spectral counterexamples.

Twelve topics have explicit archived internal proof-review records.
Twenty-four additional original manuscripts contain complete-range
arguments or proof candidates, and a code-only checkpoint is reconstructed
with its complete exact certificate. Four further manuscripts supply
negative resolutions of five named assertions. These counts use different
units where stated: the two independently obtained Eulerian-maxima
proofs form one topic, while one Stern-polynomial manuscript and one
gap-semigroup manuscript each address two numbered conjectures.

All theorem statements, proof arguments, substantive auxiliary results,
and source-specific limitations from the selected manuscripts are
retained. Review status is recorded separately from the mathematical
assertions. The integrated edition does not assert external peer review,
formal proof-assistant certification, or universal publication priority.
\end{abstract}
\noindent\textbf{Keywords.} Global inequalities; combinatorial sequences;
Diophantine finiteness; exact certificates; total positivity;
counterexamples; proof provenance.
\clearpage
\tableofcontents
\clearpage
\section*{Introduction and editorial conventions}
\addcontentsline{toc}{section}{Introduction and editorial conventions}
This paper presents the complete selected arguments in one reproducible
document. Each principal section identifies its original target, round,
and source participant, and retains the entire proof exposition. A shared
table of contents, continuous pagination, bibliography, and source
register connect the sections. All text and certificate listings are
embedded in one source file; compilation needs no external
\texttt{\string\input}, image, bibliography, or attached manuscript.

\paragraph{Compilation and certificate replay.}
Compile the supplied single TEX file with pdfLaTeX, repeating the run
until the table of contents and references stabilize. A standard
installation with Latin Modern, AMS, and the listed layout packages is
sufficient. No Python execution is required to typeset the paper.
Optional certificate replay uses the embedded Python listings; two
listings additionally require SymPy. Source-specific filenames in the
preserved prose identify the original archive, while the proof-critical
programs selected for this edition are printed in Appendix~A.

\paragraph{Scope of the selection.}
The four parts correspond exactly to the four requested groups:
internally reviewed positive results; the early Stern and equal-shift
arguments; further full-range proof candidates; and negative resolutions.
The source register at the end records every selected contribution.
The Reiwa and Keio Eulerian-maxima proofs share one principal section,
but neither derivation has been discarded.

\paragraph{Meaning of the status labels.}
An archived internal review is evidence about the review performed in
that research session, not an external publication decision. A candidate
label is retained when no final independent adjudication was recorded.
It does not itself identify a mathematical defect. Conversely, the
presence of a theorem environment or an exact finite test does not
automatically establish a global statement. A finite certificate is
relevant only together with the theorem or analytic argument that
extends it to the full quantified domain.

\paragraph{Notation and numbering.}
Each principal section is mathematically self-contained relative to its
explicitly cited theorem inputs. Symbols are local to the section.
Automatically numbered statements and equations carry the principal
section number. The unqualified statement numbers printed inside a
source quotation or a bibliographic discussion refer to the cited
source, whose target is identified beside the section heading.
All hyperlink labels and bibliographic keys have been made unique.

\paragraph{Proof mechanisms connecting the parts.}
Several mechanisms recur across the otherwise distinct subjects.
Positive finite-path recurrences prove block monotonicity; exact
polynomial coefficient and root-count certificates prove continuum
inequalities; signed Cauchy--Binet expansions preserve determinant
strictness; residue paths and Apery sets control semigroup gaps;
and a ramified lattice converts an infinite matrix recurrence into
uniform bounds on every finite-difference order. These connections
organize the exposition without adding an unproved implication between
the separate targets.

\paragraph{Editorial reconstruction and retained limitations.}
The Brian Yi Sun section was reconstructed from an archived checkpoint
and executable exact verifier, because no original PDF or TEX existed
for that entry. Its derivation, coefficient tables, initial cases, and
entire verifier are included. This reconstruction is distinguished from
an original reviewed submission. Collateral finite computations,
withdrawn duplicate targets, known consequences, and explicitly
unverified clauses remain identified as such wherever they occur in a
selected proof manuscript.

\clearpage

\part{Internally reviewed complete results}

\noindent Twelve topics are grouped here. Eleven were the reviewed round winners; the even Eulerian-maxima result received a separate supplementary review. Both independent source derivations of that last topic are included.


\section{Exact tone coloring of Hamming graphs}
\label{section:R13T}
\noindent{\small\textsf{ROUND 13 \enspace / \enspace Taisho \enspace / \enspace R13T}}\\[3pt]
\noindent\textbf{Target.} Bennett--Nichols Conjecture 2.

\noindent\textbf{Editorial status.} Archived internal proof review passed. The qualifier \emph{for all sufficiently large $n$} belongs to the original conjecture. The proof uses a published decomposition theorem and does not impose a prime-power restriction.

\smallskip\hrule\medskip



\begingroup

\setcounter{equation}{0}
\renewcommand{\theequation}{\thesection.\arabic{equation}}
\renewcommand{\theHequation}{R13T.\arabic{equation}}
\paragraph{Overview.}
We prove Conjecture~2 of Bennett and Nichols: for every pair of positive integers $t,b$, the $t$-tone chromatic number of the Cartesian power $K_n^b$ equals $tn$ for all sufficiently large $n$. Their two-coordinate palette construction from orthogonal Latin squares extends to all dimensions using index-one orthogonal arrays of strength $b$. The needed arrays for all large alphabet sizes follow from Keevash's partite hypergraph decomposition theorem, whose hypotheses hold exactly on complete equal blowups. The palette transfer and the applicability check give an unconditional proof; no new design-existence theorem or effective numerical threshold is claimed.


\subsection{The conjecture and the transfer}
In a $t$-tone coloring of a graph, each vertex $x$ receives a set $C(x)$ of $t$ distinct colors, with $|C(x)\cap C(y)|<d(x,y)$ whenever $x\ne y$. Write $\tau_t$ for the fewest colors in such a coloring. Throughout, $K_n^b$ denotes a \emph{Cartesian} power. Its vertices are $x=(x_1,\ldots,x_b)\in[n]^b$, where $[n]=\{1,\ldots,n\}$, and its distance is Hamming distance: a step changes one coordinate and any one coordinate can be changed in one step.

\begin{SourceRbdTtheorem}\label{R13T:main}
For every pair of positive integers $t,b$ there is an integer $N=N(t,b)$ such that for every integer $n\ge N$,
\[
\tau_t(K_n^b)=tn.
\]
\end{SourceRbdTtheorem}
This is precisely Bennett and Nichols's Conjecture~2 \cite[\S8]{R13T:BN}. They proved an upper bound of $(1+\delta)tn$ for every fixed $\delta>0$ and all sufficiently large $n$ \cite[Theorem~12]{R13T:BN}. Their Theorem~13 and Corollary~2(a) establish the exact large-$n$ result in dimension $b=2$ using mutually orthogonal Latin squares and one separate palette per square. We retain these palettes and replace the squares by an array with stronger coordinate uniqueness.

An index-one orthogonal array of strength $b$, with $b+t$ columns and alphabet $[n]$, is a set $\mathcal A\subseteq[n]^{b+t}$ whose projection onto every set of $b$ columns is a bijection to $[n]^b$. In particular, the first $b$ columns give a unique row $(x,z(x))$ for each $x\in[n]^b$.

\begin{SourceRbdTlemma}\label{R13T:transfer}
If such an array exists, then $\tau_t(K_n^b)=tn$. Conversely, a $t$-tone coloring using a palette partitioned into $t$ disjoint classes of size $n$, with every vertex receiving one color from each class, gives such an array.
\end{SourceRbdTlemma}
\begin{proof}
Label the palette by $[t]\times[n]$, and set
\[
C(x)=\big\{(i,z_i(x)):i\in[t]\big\}.
\]
For vertices $x,y$ at Hamming distance $d>0$, their rows already agree on $b-d$ input positions. If they shared $d$ or more colors, they would agree on at least $d$ output positions as well. This would give two distinct rows agreeing on $b$ columns, contradicting the array property. Thus the coloring uses at most $tn$ colors. A coordinate line is an $n$-clique. Its $n$ labels are pairwise disjoint because adjacent vertices share no color, forcing at least $tn$ colors.

Conversely, number the colors in each class by $[n]$ and append the corresponding $t$ symbols to each input vertex. Two distinct rows associated to vertices at distance $d$ agree on $b-d$ input positions and at most $d-1$ output positions, so they agree on at most $b-1$ positions in total. Every $b$-column projection is injective, and it is bijective because domain and codomain each have cardinality $n^b$.
\end{proof}
The converse concerns only colorings with this fixed palette partition. It is not asserted for arbitrary optimal $t$-tone colorings.

\subsection{All sufficiently large alphabet sizes}
\begin{SourceRbdTlemma}\label{R13T:oa}
For fixed integers $q>b\ge2$, the complete $q$-partite $b$-uniform hypergraph with $n$ vertices in every part, all cross-part $b$-edges present, has a decomposition into transversal copies of the complete $b$-uniform hypergraph on $q$ vertices, for every sufficiently large integer $n$.
\end{SourceRbdTlemma}
\begin{proof}
We apply Keevash \cite[Definition~1.6 and Theorem~1.7]{R13T:Keevash} with the template $H=K_q^{(b)}$ on $[q]$ and with its complete equal $n$-blowup $G$ on the parts $V_1,\ldots,V_q$. A cross-part set contains at most one vertex per part. For any template $b$-edge $f$, the density $d_f(G)$ between the corresponding parts is exactly $1$.

For \emph{balance}, fix a cross-part set $e$ supported on a part set $I\subseteq[q]$. For every template edge $f\supseteq I$, the number of present edges with support $f$ that contain $e$ is $n^{b-|I|}$, independent of $f$. If no such $f$ exists the balance condition is vacuous. Hence $G$ is $H$-balanced in the sense required by the theorem.

For \emph{typicality}, take any list of cross-part $(b-1)$-sets and any target part that extends each of their supports to an edge of $H$. This part is disjoint in support from each set in the list. Every vertex in the target part extends every set to a present edge of $G$. The common extension count is therefore exactly $n$, which equals the target part size times the product of the relevant densities. This works for every permitted list size, hence in particular for the size bound in Theorem~1.7, with arbitrarily small positive tolerance.

All parts have size $n$. Fix a density lower parameter $1/2$, and choose the positive typicality tolerance below the bound required by the theorem for this parameter. For fixed $q$ the prescribed lower bound on the density parameter tends to zero as $n$ grows. Thus the parameter inequalities hold for every sufficiently large integer $n$. Theorem~1.7 gives a partite $H$-decomposition, which means the required copies use exactly one vertex in each $V_i$ and cover each $b$-edge exactly once.
\end{proof}

\begin{proof}[Proof of Theorem~\ref{R13T:main}]
For $b=1$ the claim is immediate for the clique $K_n$. For $b\ge2$, put $q=b+t$ and apply Lemma~\ref{R13T:oa}. Label the vertices in each part by $[n]$ and regard each transversal copy as a row of length $q$. Choosing $b$ columns and their symbols fixes a unique edge of the complete partite $b$-graph. That edge lies in exactly one copy, so each such projection of the rows is bijective. This is the array in Lemma~\ref{R13T:transfer}, which gives equality. The threshold is chosen with $q,b$ fixed and imposes no primality, congruence, or factorization condition on $n$.
\end{proof}

\begin{SourceRbdTcorollary}
For fixed positive $t,b$ and positive integers $n_1,\ldots,n_b$ with $m=\max_i n_i$ sufficiently large in terms of $t,b$,
\[
 \tau_t(K_{n_1}\mathbin{\square}\cdots\mathbin{\square}K_{n_b})=tm.
\]
\end{SourceRbdTcorollary}
\begin{proof}
Restrict the coloring of $[m]^b$ to the subbox $\prod_i[n_i]$. Hamming distances are unchanged. A coordinate line in a largest factor is an $m$-clique, giving the lower bound.
\end{proof}

\subsection{Scope, checks, and provenance}
Over a large finite field one can form the array from evaluations of all polynomials of degree less than $b$ at $b+t$ distinct field elements. More generally this works over $\mathbb Z/n\mathbb Z$ when the evaluation points have pairwise unit differences: each projection matrix is Vandermonde with unit determinant. Such decompositions for orders whose prime factors are sufficiently large are classical; see Fabila-Monroy and Wood \cite{R13T:FMW}. This approach does not establish the required all-integer threshold: large orders can still have small prime factors. For a concrete failure of a naive linear recipe, modulo $6$ the assignments $x+y$ and $x+5y$ are each Latin, but the two inputs $(0,0)$ and $(3,3)$ agree in both outputs at distance two. This is not a general obstruction to unpartitioned colorings. The all-order step in this proof is provided by Keevash, rather than by an unproved extension of the interpolation construction.

The companion Python script checks every $b$-column projection and every pairwise coloring constraint for five examples, including order $35$ at $b=2,t=3$ and order $7$ at $b=4,t=2$. It also checks the modulo-$6$ failure directly. Those finite calculations illustrate the transfer but are not used to infer the large-$n$ claim. The manuscript has only one non-elementary dependency, the cited partite decomposition theorem.

The palette use for $b=2$ and the asymptotic bounds belong to Bennett and Nichols. The partite design existence and its array interpretation are existing theory; the proof contribution is the extension of their fixed-palette construction to every dimension via strength-$b$ projections and the application to their exact conjecture. A bounded source search located no earlier full proof of Conjecture~2 or updated version stating such a proof. This is not a certification of worldwide priority or external review. Conjecture~1 on trees in the original paper is not addressed.
\endgroup


\clearpage

\section{Two-prime products represented by sums of cubes}
\label{section:R16M}
\noindent{\small\textsf{ROUND 16 \enspace / \enspace Meiji \enspace / \enspace R16M}}\\[3pt]
\noindent\textbf{Target.} Tian sub-conjecture (2).

\noindent\textbf{Editorial status.} Archived internal proof review passed. Only the separately stated two-prime sub-conjecture (2) is settled here; the assertion for an arbitrary number of prime factors is outside this section. The reciprocity appendix is retained.

\smallskip\hrule\medskip



\begingroup
\newcommand{\ord}{\operatorname{ord}}
\setcounter{equation}{0}
\renewcommand{\theequation}{\thesection.\arabic{equation}}
\renewcommand{\theHequation}{R16M.\arabic{equation}}
\paragraph{Overview.}
For any two distinct fixed primes, at most two sums of the first
$n$ positive cubes have exactly those prime factors. This proves
Tian's separately stated Sub-conjecture~(2) in
Zeng--Pint\'er--Fu--Tian, \emph{Mathematics} 14 (2026), 127.
A prime-power divisibility argument gives uniqueness within each
sign and orientation; quadratic reciprocity supplies the missing
mutual exclusions. The case involving the prime two is elementary
as well. The arbitrary-number-of-primes Sub-conjecture~(3) is not
addressed.


\subsection{Statement and scope}
Zeng, Pint\'er, Fu and Tian \cite[\S1]{R16M:Tian} record three separate
statements. Their Sub-conjecture~(1), about the fixed pair $3,5$,
is proved in their Theorem~3. Our target is their next statement,
Sub-conjecture~(2), for every pair of distinct fixed primes, with
no upper bound on the primes, the index or the exponents.
Their Theorem~4(1) states the weaker upper bound of four.
Sub-conjecture~(3) concerns arbitrary support size and is separate
from the claim here.

\begin{SourceRbgMtheorem}[Tian's two-prime conjecture]\label{R16M:thm:main}
Let $p,q$ be distinct primes. There are at most two triples of
positive integers $(n,\beta,\gamma)$ satisfying
\[
                     \sum_{i=1}^{n}i^3=p^\beta q^\gamma .
\]
\end{SourceRbgMtheorem}

The standard sum-of-cubes identity gives
$\sum_{i=1}^{n}i^3=T_n^2$, where $T_n=n(n+1)/2$;
the identity also follows by telescoping
$T_i^2-T_{i-1}^2=i^3$. Unique factorization makes both exponents
even. Thus it suffices to count positive triples $(n,a,b)$
with $n(n+1)=2p^a q^b$. The primes are regarded as labeled
throughout; interchanging $p$ and $q$ only interchanges the
corresponding exponents.

\subsection{Uniqueness in a sign and orientation}
For $u$ prime write $v_u(z)$ for its exponent in a nonzero
integer $z$. We use the following standard elementary argument,
included to handle all signs and avoid primitive-divisor
exceptions.

\begin{SourceRbgMlemma}\label{R16M:lem:valuation}
Let $u$ be an odd prime, $Y\equiv1\pmod u$ an integer with
$|Y|>1$, and $k\ge1$. Then
\[
            v_u(1+Y+\cdots+Y^{k-1})=v_u(k).
\]
\end{SourceRbgMlemma}
\begin{proof}
Denote the sum by $S_k(Y)$. For $u\nmid k$ it is $k$ modulo $u$.
For $k=u$, write $Y=1+h$ with $u\mid h$. Then
\[
 S_u(Y)=u+\binom{u}{2}h+\binom{u}{3}h^2+\cdots+h^{u-1}
             \equiv u\pmod{u^2}.
\]
Indeed each intermediate coefficient is divisible by $u$ and
$u-1\ge2$. Use
$S_{kl}(Y)=S_k(Y)S_l(Y^k)$ and $Y^k\equiv1\pmod u$ to factor
off the full power of $u$ in $k$.
\end{proof}

\begin{SourceRbgMlemma}\label{R16M:lem:branches}
For distinct odd primes $P,Q$ and a fixed $\epsilon\in\{1,-1\}$,
\[
                       P^a-\epsilon=2Q^b
\]
has at most one solution $(a,b)$ in positive integers.
\end{SourceRbgMlemma}
\begin{proof}
If there were two solutions, their exponents on $P$ would
differ. Let these be $a<c$, with exponents $b<d$ on $Q$.
Then $P^a-\epsilon$ divides $P^c-\epsilon$. Write $c=ka+r$,
$0\le r<a$, and reduce modulo $P^a-\epsilon$ to see that
$P^a-\epsilon$ divides $\epsilon^k P^r-\epsilon$.
For $r>0$ the latter is nonzero of absolute value at most
$P^{a-1}+1<P^a-1$, a contradiction. Thus $a\mid c$.
When $\epsilon=-1$ the quotient exponent $k$ must be odd:
otherwise the congruence would require $P^a+1$ to divide $2$.
For $\epsilon=1$, this divisibility step is also the usual
gcd identity for $P^a-1$ and $P^c-1$.

Put $X=P^a$, and let $Y=X$ for $\epsilon=1$ and $Y=-X$
for $\epsilon=-1$. In either case $Y\equiv1\pmod Q$ and, using
$k$ odd in the second case,
\[
             Q^{d-b}=\frac{P^{ka}-\epsilon}{P^a-\epsilon}=S_k(Y).
\]
By Lemma~\ref{R16M:lem:valuation}, this $Q$-power is
$Q^{v_Q(k)}\le k$. But $S_k(X)>k$ for $k\ge2$ and $X\ge3$.
If $Y=-X$ and $k\ge3$ is odd, the sum is
$1+\sum_{j=1}^{(k-1)/2}(X^{2j}-X^{2j-1})>k$
as each displayed difference is at least six. This is a
contradiction in both cases.
\end{proof}

\subsection{Odd primes: the missing exclusion}
Assume first that both primes are odd. Since $n,n+1$ are
coprime, each prime power lies in a single factor and the
factor of two lies in the even factor. The two odd prime
powers cannot lie in the same factor. If they both lay in the
odd factor, the even factor would be $2$ and the odd factor
would be $1$ or $3$, which cannot have both odd primes.
If both lay in the even factor, the odd factor would be $1$
and the even factor $2$, also impossible. Thus all solutions
arise from
\[
               P^a-\epsilon=2Q^b,\qquad \epsilon=\pm 1,
\]
in one of the two orientations $(P,Q)=(p,q),(q,p)$, with
$n=\min(P^a,2Q^b)$. Conversely each such solution gives a valid
triple, after matching the exponents to their primes and
doubling both.

We show that any solution having $\epsilon=1$ is the only
solution for its pair. Suppose
\[
                         P^a-1=2Q^b .
\]
Reduction modulo four gives $P^a\equiv3\pmod4$, so
$P\equiv3\pmod4$ and $a$ is odd. Modulo $Q$ one has $P^a=1$;
therefore $P$ is a square modulo $Q$, since
$(P^{(a+1)/2})^2\equiv P\pmod Q$.
Write $(x/u)$ for the Legendre symbol modulo the odd prime $u$.

If $Q\equiv3\pmod4$, quadratic reciprocity gives $(Q/P)=-1$.
The reverse equation $Q^c-1=2P^d$ would require $c$ odd
and hence $(Q/P)=1$, a contradiction. Neither equation with
$\epsilon=-1$ is possible: its left base is $3$ modulo four,
so that $P^e+1=2Q^f$, for example, forces $e$ even by reduction modulo
four. This would make $-1$ a square modulo $Q$, contrary to
$Q\equiv3\pmod4$. The same argument interchanges $P,Q$.

If $Q\equiv1\pmod4$, quadratic reciprocity gives $(Q/P)=1$.
The reverse equation $Q^c-1=2P^d$ is impossible modulo four.
The original equation modulo $P$ gives $2Q^b=-1$, whence
$(2/P)=-1$. The equation $P^c+1=2Q^d$ would give $(2/P)=1$
by the same calculation with right-hand side congruent to $1$.
Finally $Q^c+1=2P^d$ would give $Q^c=-1$ modulo $P$,
impossible because $Q$ is a square and $-1$ is not.

The original sign and orientation have at most one solution
by Lemma~\ref{R16M:lem:branches}. This proves the uniqueness claim
when a solution with $\epsilon=1$ exists. Otherwise only
the two orientations with $\epsilon=-1$ are available.
There is then at most one solution in each, which proves the
bound for odd primes. In particular, if there are two such
solutions for an odd pair, both indices are odd (the odd
factor is smaller in the $\epsilon=-1$ branch).

\subsection{The pair containing two}
Let $p=2$ and let $q$ be odd. The two consecutive factors
are $2^k$ and $q^b$ for $k=a+1\ge2$ and $b\ge1$.
Indeed the even factor carries the full $2$-power, and if it
also carried $q$ then the odd factor would be $1$, leaving
the even factor $2$, which is impossible here.

If $q^b+1=2^k$, then $b$ is odd by reduction modulo four.
For odd $b>1$, the quotient $(q^b+1)/(q+1)$ is an odd
integer greater than one and cannot divide a power of two.
Hence $b=1$.
If $q^b-1=2^k$ and $b$ is odd, the same argument with
$(q^b-1)/(q-1)$ gives $b=1$. If $b$ is even, put
$x=q^{b/2}$. The positive factors $x-1,x+1$ of $2^k$
are both even powers of two and differ by two; they must
be $2,4$. Thus $x=3$ and the only even-$b$ possibility
is $(q,b,k)=(3,2,3)$.

For $b=1$, at most one of $q-1,q+1$ can be a power
of two of size at least four, since their difference is two.
Together with the unique even-$b$ possibility, this proves
the bound and completes Theorem~\ref{R16M:thm:main}. For instance,
the pair $2,3$ has triples $(3,2,2),(8,4,4)$, so the
universal bound is sharp.

\subsection{Exact check for a specified pair}
For odd $P,Q$, let $r=\ord_Q(P)$. Any solution in the
$\epsilon=1$ branch has $r\mid a$, so $P^r-1$ is an even divisor
of $2Q^b$ divisible by $Q$. Consequently $P^r-1=2Q^{b'}$
with $b'>0$, and uniqueness gives $a=r$. In the
$\epsilon=-1$ branch the order is even and $a$ is an odd
multiple of $r/2$. The divisor $P^{r/2}+1$ is again even
and divisible by $Q$, so uniqueness gives $a=r/2$.
Thus there is an unbounded-index decision for every fixed pair
by modular order and exact prime-power tests, with the even-prime
case as in the preceding subsection.

The companion Python verifier checks all 14,028 unordered prime
pairs with both primes at most 997 by this rule. There are 70
pairs with one solution, two pairs with two, and all other pairs
have none. An independent direct factorization sieve of
$T_n$ through $n=10^6$ agrees on their shared range.
The code is not used in the proof.


\subsection*{Appendix: the reciprocity ingredient}
For completeness, let $m=(P-1)/2$ and take the least positive
residues $r_i$ of $Qi$ modulo $P$, $1\le i\le m$.
The numbers $s_i=\min(r_i,P-r_i)$ permute $1,\ldots,m$:
a collision would give $i\equiv\pm j\pmod P$, which requires
$i=j$ in this range. If $N$ of the $r_i$ exceed $P/2$,
multiplication and cancellation of $m!$ show
$Q^m\equiv(-1)^N\pmod P$. The $m$ nonzero squares are exactly
the $m$ roots of $z^m=1$ in the residue field, by Fermat's
little theorem and the degree bound. All other nonzero elements
have $z^m=-1$. Thus the Legendre symbol is multiplicative,
$(-1/P)=(-1)^m$, and $(Q/P)=(-1)^N$.

Since $P,Q$ are odd and $\sum s_i=\sum i$, comparing the parity
of $\sum Qi=P\sum\lfloor Qi/P\rfloor+\sum r_i$ gives
$N\equiv\sum\lfloor Qi/P\rfloor\pmod2$.
Interchange $P,Q$ and add the floor sums. They partition the
lattice points in the rectangle
$1\le i\le(P-1)/2$, $1\le j\le(Q-1)/2$ according as
$Pj<Qi$ or $Pj>Qi$; equality is excluded. Hence their sum
is $(P-1)(Q-1)/4$, yielding the reciprocity law used above.
\endgroup


\clearpage

\section{Block monotonicity of binary binomial transforms}
\label{section:R17M}
\noindent{\small\textsf{ROUND 17 \enspace / \enspace Meiji \enspace / \enspace R17M}}\\[3pt]
\noindent\textbf{Target.} Dilcher--Ulas Conjecture 2 / 5.8.

\noindent\textbf{Editorial status.} Archived internal proof review passed. The target is monotonicity inside every block, including the prescribed final equality. The sign pattern alone is not counted as a new result.

\smallskip\hrule\medskip



\begingroup
\newcommand{\Z}{\mathbb Z}
\setcounter{equation}{0}
\renewcommand{\theequation}{\thesection.\arabic{equation}}
\renewcommand{\theHequation}{R17M.\arabic{equation}}
\paragraph{Overview.}
We prove Conjecture 2 of Dilcher and Ulas (\emph{Integers} 22 (2022), A8), equivalently Conjecture 5.8 of arXiv:2008.13475v2. An odd, antiperiodic extension turns their signed binomial sum into an ordinary walk count on a finite path. Pascal's rule gives a second walk count for each consecutive difference; the endpoint of the path accounts for the unique equality in each block. The argument extends to square-wave binomial transforms of every integer block width at least two.


\subsection{Statement and context}
Let
\[
 f_{m,n}(z)=\sum_{j=0}^n\binom nj z^{\binom jm},\qquad
 \binom jm=0\text{ for }j<m .
\]
Dilcher and Ulas conjectured that for all integers \(k\ge1\) and \(\nu\ge0\), the sequence \((-1)^\nu f_{2^k,n}(-1)\) is positive on
\[
  \nu 2^{k+1}\le n\le (\nu+1)2^{k+1}-2
\]
and strictly increases between consecutive terms there, except that the final two terms are equal~\cite[p.~22, Conjecture~2]{R17M:DU}. They proved the first block in Proposition~8, and their Corollary~6 also gives the terminal equality in all blocks. Independently, the sign and zero pattern already appears in Theorem~1 of Su, Tang and Pott~\cite{R17M:STP}: the value \(f_{2^k,n}(-1)\) equals \(2^n-2\operatorname{wt}(\sigma_{n,2^k})\), since the degree-\(2^k\) elementary symmetric Boolean function takes value \(\binom j{2^k}\bmod 2\) on inputs of weight \(j\). For \(n<2^k\) the sign assertion follows immediately from the definition. The proof below treats the sign and every first difference together; its new step is a nonnegative walk count that certifies the strict inequalities in every block. It does not require a trigonometric estimate.

\subsection{Walk representation and proof}
For an integer \(m\ge2\), set \(L=2m\) and
\[
 B_m(n)=\sum_{j=0}^{n}\binom nj(-1)^{\lfloor j/m\rfloor}.
\]
Let \(A\) be the adjacency matrix of the path with vertices \(1,\dots,L-1\), and let \(v(x)=1\) at odd vertices and \(v(x)=0\) at even vertices. When writing \((A^n v)(L)\), we mean zero by convention.

\begin{SourceRbhMtheorem}\label{R17M:main}
Write \(n=L\nu+r\), with \(0\le r<L\). Then \(B_m(n)=0\) if \(r=L-1\); otherwise
\[
  (-1)^\nu B_m(n)=(A^n v)(r+1)>0.
\]
For \(0\le r\le L-3\) one has
\[
  (-1)^\nu\bigl(B_m(n+1)-B_m(n)\bigr)=(A^n v)(r+3),
\]
which is positive for \(r<L-3\) and zero for \(r=L-3\). In particular this proves Dilcher and Ulas's conjecture when \(m=2^k\).
\end{SourceRbhMtheorem}

\begin{proof}
On the whole integer line define
\[
 g(2j+1)=(-1)^{\lfloor j/m\rfloor}\quad(j\in\Z),\qquad g(2j)=0 .
\]
The definition gives \(g(x+L)=-g(x)\), and \(g(-x)=-g(x)\) follows from
\(\lfloor(-j-1)/m\rfloor=-1-\lfloor j/m\rfloor\). Put
\[
 H_n(x)=\sum_{j=0}^n \binom nj g(x-n+2j) \quad(x\in\Z).
\]
Pascal's identity gives
\begin{equation}\label{R17M:rec}
 H_{n+1}(x)=H_n(x-1)+H_n(x+1).
\end{equation}
With \(H_0=g\), both symmetries propagate under \eqref{R17M:rec}, so \(H_n(0)=H_n(L)=0\) and
\[
 H_n(x+L\nu)=(-1)^\nu H_n(x).
\]
In particular, restricting \eqref{R17M:rec} to \(1\le x\le L-1\) is precisely multiplication by \(A\). As \(g\) restricts to \(v\), it follows that \(H_n(x)=(A^n v)(x)\) on these vertices. This entry counts length-\(n\) walks from \(x\) to any odd vertex. If \(x+n\) is odd, it is strictly positive: choose any edge incident to \(x\) and alternate along it for \(n\) steps. For \(n=0\), use the empty walk at \(x\). Every path vertex has an incident edge because \(m\ge2\).

The original transform is the diagonal sample \(B_m(n)=H_n(n+1)\). After reduction by \(L\nu\), the sample lies at \(r+1\); there \(r+1+n\) is odd, because \(L\) is even. This proves the sign claim, and at \(r=L-1\) the sample lies at the zero boundary.
For successive terms, \eqref{R17M:rec} yields the exact identity
\[
 B_m(n+1)-B_m(n)=H_{n+1}(n+2)-H_n(n+1)=H_n(n+3).
\]
For \(r<L-3\) this gives the positive walk count at \(r+3\), again of the correct parity. At \(r=L-3\) it instead gives the boundary value zero. These are all adjacent pairs in the specified block.

Finally, for \(m=2^k\), the coefficient of \(X^{2^k}\) in \((1+X)^j\) over \(\mathbb F_2\) is the \(k\)-th binary digit of \(j\): if \(j=\sum_i \epsilon_i 2^i\), expand \(\prod_{i:\epsilon_i=1}(1+X^{2^i})\) and use uniqueness of the binary expansion. Therefore \((-1)^{\binom{j}{2^k}}=(-1)^{\lfloor j/2^k\rfloor}\), so \(f_{2^k,n}(-1)=B_{2^k}(n)\). The signs, strict increases, and exact last equality have the original quantifiers.
\end{proof}

\paragraph{A direct cancellation interpretation.}
The identity \(H_n(x)=(A^n v)(x)\) also follows by reflecting walks on \(\Z\). For an interior starting point, count walks with signed endpoint weight \(g\). For walks hitting \(0\) or \(L\), keep the segment up to the first boundary hit and reflect the remaining segment in that boundary. The first hit is unchanged, so this is an involution. The endpoint weight changes sign since \(g(-y)=-g(y)\) and \(g(2L-y)=-g(y)\). Any fixed walk hits the boundary at the last step and has weight zero. The surviving walks stay in the interior and have total weight \((A^n v)(x)\). This is a bijective explanation of why no estimates are needed even when the unfurled walk crosses many sign blocks.

\paragraph{Numerical check.}
An independent exact-integer check compares the original sparse-polynomial sum, the square-wave sum, and the path formula, and tests the first differences and zeros. These checks guard against endpoint and parity errors; no finite check is used in the proof.
\endgroup


\clearpage

\section{Nonmonotone exceptional sets of positive powers}
\label{section:R18K}
\noindent{\small\textsf{ROUND 18 \enspace / \enspace Keio \enspace / \enspace R18K}}\\[3pt]
\noindent\textbf{Target.} Benfield--Lippard Conjecture 10.5.

\noindent\textbf{Editorial status.} Archived internal proof review passed. This concerns the stabilized exceptional-set sequence, not the general Waring problem. The logarithmic limsup and liminf consequences are retained.

\smallskip\hrule\medskip



\begingroup
\newcommand{\N}{\mathbb N}
\newcommand{\B}{\mathbf B}
\setcounter{equation}{0}
\renewcommand{\theequation}{\thesection.\arabic{equation}}
\renewcommand{\theHequation}{R18K.\arabic{equation}}
\paragraph{Overview.}
We prove Benfield and Lippard's Conjecture~10.5 that the number of
stable gaps in the exact-summand Waring problem is not eventually
monotone. Write \(b_k\) for the stable gap count at exponent \(k\).
For every sufficiently large prime \(p\), we prove
\(b_{p-1}>b_p\), with \(b_{p-1}/b_p\longrightarrow\infty\).
The proof compares a Fermat congruence obstruction at exponent
\(p-1\) to a coprimality result for small shifted \(p\)-th powers.
The latter follows by injecting a square of reduced fractions
into the \(p\)-th roots of unity in a finite field.


\subsection{Identification of the stable gap sets}

For \(k\ge2\) let
\[
 S_k=\left\langle x^k-1:x\ge2\right\rangle
 \subseteq\mathbb Z_{\ge0}
\]
be the additive monoid of finite sums, including the empty sum.
For \(j\ge1\), \(j+t\) is a sum of exactly \(j\) positive \(k\)-th
powers if and only if \(t\) is a sum of at most \(j\) of the displayed
generators. For the forward implication, subtract \(1\) from each
power and drop zero terms; for the reverse implication, pad with
copies of \(1^k\). Thus the intersection of the decreasing sets of
positive offsets \(t\) not represented at \(j\) is precisely
\(\N\setminus S_k\). By Zenkin's eventual representability for a
fixed sufficiently large \(j\) \cite{R18K:Zenkin}\cite[Theorem~1.1]{R18K:BL}, one of the
decreasing sets is finite, so the chain stabilizes at its
intersection. Benfield and Lippard define \(\B^k\) to be this
stabilized set \cite[\S10.2]{R18K:BL}. Hence
\[
 b_k=|\B^k|=|\N\setminus S_k|.
\]
No novelty is claimed for stabilization, which also follows directly
from the originating paper's descending-chain argument.

\begin{SourceRbiKlemma}[A residue bound]\label{R18K:lem:residue}
Let \(T\) be generated by positive integers of greatest common divisor
\(1\), including a generator \(m\), with all generators at most \(D\).
Every integer at least \((m-1)D\) belongs to \(T\).
\end{SourceRbiKlemma}

\begin{proof}
In the directed graph of residues modulo \(m\), include all
translations by generators. The residues reachable from \(0\) form a
submonoid of the finite cyclic group, hence a subgroup: the inverse of
any element is a nonnegative multiple of that element.
B\'ezout's identity implies the subgroup is the whole group.
For each residue choose a path from \(0\) with the fewest edges.
Such a path has no repeated vertices and at most \(m-1\) edges.
Summing the corresponding generators gives an element of \(T\) in that
residue of size at most \((m-1)D\). Any larger integer in the same
residue is obtained by adding copies of \(m\).
\end{proof}

\subsection{A short coprimality certificate at prime exponents}

\begin{SourceRbiKlemma}\label{R18K:lem:gcd}
If \(p\ge2500\) is an odd prime and \(M=\lfloor\sqrt{2p}\rfloor\), then
\[
 \gcd\{x^p-1:2\le x\le M\}=1.
\]
\end{SourceRbiKlemma}

\begin{proof}
Suppose a prime \(q\) divides every displayed integer. The
multiplicative order of \(2\bmod q\) divides \(p\), and it is not \(1\),
so it is \(p\). As \(q\) is odd and \(p\mid q-1\), we have
\(q\ge2p+1\). All integers from \(1\) to \(M\) are therefore
nonzero \(p\)-th roots of unity modulo \(q\).
Consider the ordered pairs \(1\le a,b\le M\) with \(\gcd(a,b)=1\).
The map \((a,b)\mapsto a b^{-1}\pmod q\) is injective on them:
an equality of residues implies \(q\mid ad-bc\) for the two pairs
\((a,b),(c,d)\), and \(|ad-bc|<M^2\le2p<q\).
Thus \(ad=bc\), and reduction of both fractions gives equality of
the pairs. Every image is a \(p\)-th root.

On the other hand, the number of pairs that are not coprime is at
most \(M^2\sum_{\ell\,\mathrm{prime}}\ell^{-2}\). Keeping \(2\)
and bounding the other primes by all odd integers at least \(3\)
gives
\[
 \sum_{\ell\,\mathrm{prime}}\ell^{-2}
 \le \frac14+\sum_{i\ge1}\frac1{(2i+1)^2}
 =\frac{\pi^2}{8}-\frac34 .
\]
Consequently the number of coprime pairs is at least
\((7/4-\pi^2/8)M^2>(101/196)M^2>p\).
For the first strict inequality use \(\pi<22/7\);
for the second, \(M^2/p>2-2\sqrt{2/p}>196/101\)
when \(p\ge2500\).
This contradicts the degree bound of \(p\) on the number of roots of
\(X^p-1\) over \(\mathbb F_q\).
\end{proof}

\subsection{Nonmonotonicity and quantitative oscillation}

\begin{SourceRbiKtheorem}\label{R18K:thm:main}
For \(k\ge2\), let \(b_k=|\B^k|\). Then
\[
 \frac{b_{p-1}}{b_p}\longrightarrow\infty
 \quad\text{as }p\text{ tends to infinity through the primes}.
\]
The sequence \(b_k\) has infinitely many strict adjacent rises and
infinitely many strict adjacent falls, proving Conjecture~10.5 of
Benfield and Lippard \cite{R18K:BL}. Moreover,
\[
 \limsup_{k\to\infty}\frac{\log b_k}{k\log k}=1,
 \qquad
 \liminf_{k\to\infty}\frac{\log b_k}{k\log k}\le\frac12.
\]
\end{SourceRbiKtheorem}

\begin{proof}
For prime \(p\ge2500\), apply Lemmas~\ref{R18K:lem:residue}
and~\ref{R18K:lem:gcd} to the generators \(x^p-1\) with \(2\le x\le M\).
This generated monoid is contained in \(S_p\). With
\(m=2^p-1\) and \(D=M^p-1\), it follows that
\[
 b_p\le (m-1)D < (2\sqrt{2p})^p.
\]
For exponent \(p-1\), Fermat's little theorem gives
\(p\mid x^{p-1}-1\) for every \(p\nmid x\).
If \(p\mid x\), then \(x^{p-1}-1\ge p^{p-1}-1\).
Hence every positive nonmultiple of \(p\) strictly below
\(p^{p-1}-1\) is a gap of \(S_{p-1}\).
In particular, the first \(p^{p-2}-1\) complete blocks of \(p\)
give, for \(p\ge5\),
\[
 b_{p-1}\ge(p-1)(p^{p-2}-1)>\tfrac12 p^{p-1}.
\]
This is a conservative version of the obstruction in
\cite[Theorem~10.7]{R18K:BL}. The ratio is bounded below by
\[
 \frac1{2p}\left(\frac{\sqrt p}{2\sqrt2}\right)^p\longrightarrow\infty.
\]

Finally, \(1,\ldots,2^k-2\) are gaps of \(S_k\),
because \(2^k-1\) is its least positive element.
Thus \(b_k\to\infty\); eventual nonincrease is impossible,
so the strict falls already obtained are interspersed with
infinitely many strict rises.

For the upper bound on the limsup, the generators \(x^k-1\)
with \(2\le x\le k+1\) have gcd \(1\).
Indeed, any hypothetical common prime divisor \(q\le k+1\)
is ruled out by \(x=q\); for \(q>k+1\), the \(k+1\) distinct
elements \(1,\ldots,k+1\) would all be roots of \(X^k-1\) modulo \(q\).
Lemma~\ref{R18K:lem:residue} now gives \(b_k<(2(k+1))^k\).
The lower bound at \(p-1\) and upper bound at \(p\) give the
remaining two logarithmic inequalities.
\end{proof}

\subsection{Exact checks and scope}

The proof of Theorem~\ref{R18K:thm:main} does not use computation.
An optional reproducibility script, \texttt{verify.py}, computes
the least element \(d_r\) of \(S_k\) in each residue class \(r\)
modulo \(2^k-1\) by weighted shortest paths. It starts with a
finite generator prefix of gcd \(1\). When the next generator
exceeds the current Frobenius number, all omitted generators
already belong to the current monoid and the computation is final.
For modulus \(m=2^k-1\), the resulting Frobenius number and gap
count are \(\max_r d_r-m\) and
\(\sum_{r=0}^{m-1}(d_r-r)/m\), respectively.
The script also verifies the comparison bounds for
\(13\le p<2500\). The exact comparisons in the two remaining
prime cases are
\[
 b_6=355825>127839=b_7,\qquad
 b_{10}=129689838697>544902765=b_{11}.
\]
Together the analytic theorem and optional finite check prove
\(b_{p-1}>b_p\) for every prime \(p\ge7\).
The corresponding claim is false at \(p=3,5\).
The paper does not assert the parity or symmetry conjectures
adjacent to Conjecture~10.5 in the source.
\endgroup


\clearpage

\section{The diagonal obstruction in universal binomial divisibility}
\label{section:R19K}
\noindent{\small\textsf{ROUND 19 \enspace / \enspace Keio \enspace / \enspace R19K}}\\[3pt]
\noindent\textbf{Target.} Zhi-Wei Sun Conjecture 1.10.

\noindent\textbf{Editorial status.} Archived internal proof review passed. The complete original implication uses the new diagonal argument together with the published He--Yuan cases. Those inputs are explicitly credited.

\smallskip\hrule\medskip



\begingroup
\newcommand{\vp}{v_q}
\setcounter{equation}{0}
\renewcommand{\theequation}{\thesection.\arabic{equation}}
\renewcommand{\theHequation}{R19K.\arabic{equation}}
\paragraph{Overview.}
We prove that for every integer $k\ge3$ there are infinitely many
positive integers $n$ for which $(k-1)n+1$ does not divide
$\binom{k^2n}{kn}$. The witnesses have $(k-1)n+1=q^h$ for a suitable
prime divisor $q$ of $2k-1$. A split of the carry count at digit $h$
gives the required valuation deficit. This supplies the case left open
by He and Yuan and, with their theorems, proves both implications of
Zhi-Wei Sun's Conjecture~1.10.


\subsection{The target}
Sun's Conjecture~1.10 in \cite[p.~193]{R19K:Sun} (Conjecture~1.3 of the
preprint) states the following for integers $k,\ell>1$:
\begin{align}
\left[\binom{kn}{n}\mid\binom{\ell n}{n}\binom{k\ell n}{\ell n-1}
   \text{ for all }n\ge1\right]
   &\Longrightarrow k=\ell\text{ or }\ell=2\text{ or }\{k,\ell\}=\{3,5\},
   \label{R19K:first}\\
\left[\binom{kn}{n}\mid\binom{\ell n}{n-1}\binom{k\ell n}{\ell n}
   \text{ for all }n\ge1\right]
   &\Longrightarrow k=2\text{ and }\ell+1\text{ is a power of two}.
   \label{R19K:second}
\end{align}
He and Yuan \cite[Theorems~1.8--1.9 and the subsequent remark]{R19K:HY}
established \eqref{R19K:first} without restriction and \eqref{R19K:second}
under $k\ne\ell$. On the diagonal, the ratio of the right-hand product
in \eqref{R19K:second} to its proposed divisor is
\begin{equation}
\frac{\binom{kn}{n-1}\binom{k^2n}{kn}}{\binom{kn}{n}}
=\frac{n}{(k-1)n+1}\binom{k^2n}{kn}.
\label{R19K:ratio}
\end{equation}
As $\gcd(n,(k-1)n+1)=1$, the diagonal requires an $n$ for which
$(k-1)n+1\nmid\binom{k^2n}{kn}$. Sun's later sequence entry
\cite{R19K:SunOEIS} records this formulation and asks a stronger question
with a bound on the least coprime witness. We address the 2012
divisibility conjecture, not that stronger question.

\subsection{A prime-power witness}
\begin{SourceRbjKtheorem}\label{R19K:thm}
Let $k\ge3$ be an integer. There is an odd prime $q$ dividing $2k-1$
such that, for every positive integer $h$ with
$q^h\equiv1\pmod{k-1}$ and $q^h>k^2$, the positive integer
$n=(q^h-1)/(k-1)$ satisfies
\[
\vp\!\binom{k^2n}{kn}
\le h-\vp(2k-1)+\vp\bigl(k(k+1)\bigr)<h.
\]
In particular, $(k-1)n+1$ does not divide $\binom{k^2n}{kn}$ for
infinitely many $n$.
\end{SourceRbjKtheorem}

\begin{proof}
Write $q$ for a prime factor of $2k-1$, $s=\vp(2k-1)$, and
$Q=q^h$. In particular, $q$ is odd and $q\nmid k-1$, because
$2k-1=2(k-1)+1$. Multiples of the multiplicative order of $q$ modulo
$k-1$ give arbitrarily large eligible $h$. Fix such an $h$ with
$Q>k^2$; this also gives $h>s$. Set
\[
X=kn,\qquad Y=k(k-1)n,\qquad X+Y=k^2n.
\]
Counting multiples of $q^j$ in the three factorials yields
\[
\vp\!\binom{X+Y}{X}=\sum_{j\ge 1} e_j,
\qquad
e_j=\left\lfloor\frac{X+Y}{q^j}\right\rfloor
   -\left\lfloor\frac{X}{q^j}\right\rfloor
   -\left\lfloor\frac{Y}{q^j}\right\rfloor.
\]
Each $e_j$ is 0 or 1; it is 0 exactly when the least nonnegative
remainders of $X,Y$ modulo $q^j$ sum to less than $q^j$.

For $1\le j\le s$, put $d=q^j$. We have $(k-1)n\equiv-1\pmod d$
and $k\equiv-(k-1)\pmod d$, so the remainder of $X$ is 1.
The remainder of $Y=k(Q-1)$ is the remainder of $-k$, namely
$(d-1)/2$ since $2k\equiv1\pmod d$ and $d$ is odd. The sum
$1+(d-1)/2$ is less than $d$. Thus $e_j=0$ for $1\le j\le s$ and
the contribution up to $h$ is at most $h-s$.

To account for every $j>h$, use the $Q$-blocks
\begin{align*}
X  &=Q+(n-1),\\
Y  &=(k-1)Q+(Q-k),\\
X+Y&=(k+1)Q+(n-k-1).
\end{align*}
All three remainders belong to $[0,Q-1]$: $n>k+1$ follows from
$Q>k^2$, and $n<Q$ follows from $k-1\ge2$.
If a nonnegative integer equals $uQ+r$ with $0\le r<Q$, then for
every $i\ge1$ its quotient by $q^{h+i}$, rounded down, is
$\lfloor u/q^i\rfloor$. Applying this to the three displayed
blocks gives
\begin{align*}
\sum_{j>h}e_j
 &= \sum_{i\ge1}
 \left(\left\lfloor\frac{k+1}{q^i}\right\rfloor
 -\left\lfloor\frac{1}{q^i}\right\rfloor
 -\left\lfloor\frac{k-1}{q^i}\right\rfloor\right)\\
 &= \vp\bigl(k(k+1)\bigr).
\end{align*}

It remains to choose $q$ with $\vp(k(k+1))<s$. If $2k-1$ has a
prime factor $q\ge5$, choose it. Such a $q$ divides neither $k$
nor $k+1$, since from $q\mid 2k-1$ these would imply respectively
$q\mid1$ or $q\mid3$. Otherwise, $2k-1=3^s$ with $s\ge2$; choose
$q=3$. Then $3\nmid k$ and $k+1=3(3^{s-1}+1)/2$ has $3$-adic
valuation 1, again less than $s$. The bound follows, and since
$(k-1)n+1=q^h$, the desired nondivisibility follows as well.
\end{proof}

\begin{SourceRbjKcorollary}
Both implications in Sun's Conjecture~1.10 hold.
\end{SourceRbjKcorollary}

\begin{proof}
Apply He--Yuan \cite[Theorem~1.8]{R19K:HY} to \eqref{R19K:first} and
\cite[Theorem~1.9]{R19K:HY} to \eqref{R19K:second} when $k\ne\ell$.
When $k=\ell\ge3$, Theorem~\ref{R19K:thm} and the ratio~\eqref{R19K:ratio}
disprove the universal divisibility premise of \eqref{R19K:second}.
For $k=\ell=2$, take $n=2$: the proposed divisor is $6$ and the
right-hand product is $4\cdot70=280$, not a multiple of $6$.
This exhausts the cases.
\end{proof}

\subsection{Exact checks}
For illustration, $k=3$, $q=5$, $h=2$ gives $n=12$ and
$v_5\binom{108}{36}=1<2$. The exceptional prime case $k=5$,
$q=3$, $h=4$ gives $n=20$ and $v_3\binom{500}{100}=3<4$.
The accompanying verification script checks the floor sums
against the independent digit-sum form of the factorial valuation,
and checks small instances directly with integer binomial
coefficients. Those checks are not hypotheses or steps in the proof.
\endgroup


\clearpage

\section{A global Turan inequality for the parabolic sine}
\label{section:R20M}
\noindent{\small\textsf{ROUND 20 \enspace / \enspace Meiji \enspace / \enspace R20M}}\\[3pt]
\noindent\textbf{Target.} Jekl Conjecture 6.1.

\noindent\textbf{Editorial status.} Archived internal proof review passed. The periodically extended parabolic sine is used throughout; a positive-half-wave formula is not silently extrapolated to the whole line.

\smallskip\hrule\medskip



\begingroup
\newcommand{\R}{\mathbb R}
\newcommand{\Z}{\mathbb Z}
\setcounter{equation}{0}
\renewcommand{\theequation}{\thesection.\arabic{equation}}
\renewcommand{\theHequation}{R20M.\arabic{equation}}
\paragraph{Overview.}
Jekl conjectured a strict three-point inequality for the periodic parabolic
sine sampled at consecutive integer multiples of a step shorter than its
positive half-wave. We prove the conjecture, allowing any real center.
The positive half-wave is strictly concave. To use this observation without
losing control when the samples cross zeros, we replace each step by the
shorter of two equivalent steps on the period circle.


Write $P=8/3$ for the parabolic analogue of $\pi$ and $S=\sin_{\rm p}$
for Jekl's periodic parabolic sine. We will prove the following exact
criterion, from which the conjecture follows at $t=nx$.

\begin{SourceRcaMtheorem}\label{R20M:thm}
For every $t\in\R$ and $h\in\R$,
\[
 D(t,h):=S(t)^2-S(t-h)S(t+h)\ge0.
\]
Equality holds exactly when $h\in P\Z$. In particular, for every
integer $n\ge1$ and $0<x<P$,
\[
 S(nx)^2>S((n-1)x)S((n+1)x).
\]
The latter assertion is Conjecture~6.1 of Jekl~\cite[p.~736]{R20M:jekl}.
\end{SourceRcaMtheorem}

\paragraph{The functions in the conjecture.}
On $0\le u\le P$, let $C=C(u)$ solve
\[
 C+\tfrac13 C^3=\tfrac43-u,\qquad S(u)=1-C(u)^2.
\]
The cubic is strictly increasing in $C$, so this specifies $C$ uniquely
from $C(0)=1$ to $C(P)=-1$. These are the parabolic cosine and sine
of~\cite[Section~2]{R20M:jekl}: the cubic follows by integrating
$(1+C^2)C'=-1$ with the indicated initial condition. Beyond this
interval take the source's odd, $2P$-periodic extension of $S$; simply
continuing the cubic to all arguments would define a different function.
Since $4/3=P/2$, oddness of the cubic gives $C(P-u)=-C(u)$ and
$S(P-u)=S(u)$. Thus, on $[0,P]$,
\[
 S(u+P)=S(u-P)=-S(P-u)=-S(u).
\]
This relation extends to all real $u$ by periodicity (on the other
half-period shift by $P$ once more). In particular the zeros are $P\Z$
and $|S|$ repeats the first positive half-wave on every nodal interval.


\begin{SourceRcaMlemma}\label{R20M:lemma}
If $0<d\le P/2$ and $t\in\R$, then $D(t,d)>0$.
\end{SourceRcaMlemma}
\begin{proof}
On $0<u<P$, implicit differentiation gives
\[
 C'=-\frac{1}{1+C^2},\qquad S'=\frac{2C}{1+C^2},\qquad
 S''=-\frac{2(1-C^2)}{(1+C^2)^3}<0 .
\]
Hence $|S|$ is strictly concave on each open interval between consecutive
zeros. Put $a=t-d$, $b=t+d$. If $S(a)S(b)>0$, both endpoints
belong to the same nodal interval. Indeed, distinct nonadjacent open
intervals are separated by more than $P$ between interior points,
whereas $b-a\le P$; adjacent nodal intervals have opposite signs.
Strict midpoint concavity, followed by the arithmetic--geometric mean
inequality, now yields
\[
 |S(t)|>\frac{|S(a)|+|S(b)|}{2}\ge\sqrt{S(a)S(b)}.
\]
If $S(a)S(b)<0$, the conclusion is immediate, even if $S(t)=0$.
If $S(a)S(b)=0$, then at least one endpoint lies in $P\Z$ while
its distance $d\in(0,P)$ from $t$ precludes $t\in P\Z$. Thus the
middle square is strictly positive in the last case as well.
\end{proof}

\begin{proof}[Proof of Theorem~\ref{R20M:thm}]
Antiperiodicity of $S$ yields $D(t,h+P)=D(t,h)$, because each
endpoint factor changes sign. Exchanging the endpoints gives
$D(t,-h)=D(t,h)$. Therefore $D(t,h)=D(t,d)$ for
$d=\operatorname{dist}(h,P\Z)\in[0,P/2]$. If $d=0$ the gap is
$D(t,0)=0$ identically. Otherwise Lemma~\ref{R20M:lemma} gives strict
positivity. Taking $h=x$ and $t=nx$ proves the stated conjecture.
\end{proof}

\paragraph{Scope.}
The proof needs neither an addition formula nor a computational check.
In fact it applies verbatim to any real antiperiodic function whose
first half-wave has zero endpoints and is positive and strictly concave
in the interior. The argument concerns the numbered Tur\'an conjecture,
not the separate double-angle approximation in~\cite[Section~6]{R20M:jekl}.
\endgroup


\clearpage

\section{Strict log-concavity for uniform linear hypercombs}
\label{section:R21T}
\noindent{\small\textsf{ROUND 21 \enspace / \enspace Taisho \enspace / \enspace R21T}}\\[3pt]
\noindent\textbf{Target.} Galvin--Sharpe Question 1.10(4).

\noindent\textbf{Editorial status.} Archived internal proof review passed. The quantified family is uniform linear hypercombs. The result does not assert independence-polynomial unimodality for every tree or hypertree.

\smallskip\hrule\medskip



\begingroup
\newcommand{\ii}{\mathrm i}
\newcommand{\Arg}{\operatorname{Arg}}
\setcounter{equation}{0}
\renewcommand{\theequation}{\thesection.\arabic{equation}}
\renewcommand{\theHequation}{R21T.\arabic{equation}}
\paragraph{Overview.}
We prove that the strong independence polynomial of every linear uniform
hypercomb has no zeros in the closed sector $|\Arg x|\le 2\pi/3$ about
the positive real axis. Its coefficients are therefore strictly log-concave
at every interior index. This eliminates the entire possible gap in the
unimodality result of Galvin and Sharpe, with no lower bound on the length
and no upper bound on the uniformity. The argument represents choices of
overlap vertices as dimers of a path and uses two complementary nonsingularity
criteria for a single tridiagonal matrix.


\subsection{The question and the result}
For $n\ge0$ the hypercomb $C_{n,\ell}$ has hubs $v_0,\ldots,v_n$.
Each of its $n$ spine edges contains two consecutive hubs and $s=\ell-2$
private vertices. Each of its $n+1$ tooth edges contains one hub and
$t=\ell-1$ private vertices. All private vertices are distinct. A
\emph{strong independent set} meets every edge at most once. Write
$C_{n,\ell}(x)=\sum_k c_k x^k$ for its size enumerator.
Galvin and Sharpe identified the hypercomb question in 2024~\cite{R21T:vone}.
Their final paper~\cite{R21T:gs} proves unimodality outside an interval of
length $O_\ell(\sqrt n)$ and asks whether that interval can be
eliminated (Question 1.10(4)). Note 1.9 of that paper explains why a
stronger assertion in an intermediate draft was withdrawn. A subsequent
preprint~\cite{R21T:eventual} establishes strict log-concavity for sufficiently
large lengths at each fixed pair of positive weights. Here the conclusion
holds at \emph{every} length for the full uniform family.
\begin{SourceRcbTtheorem}\label{R21T:main}
For $n\ge0$ and $\ell\ge3$, all zeros of $C_{n,\ell}(x)$ lie in the
open sector $|\Arg(-x)|<\pi/3$. All coefficients from $0$ to $2n+1$
are positive, and
\[
 c_k^2>c_{k-1}c_{k+1}\qquad(1\le k\le 2n).
\]
In particular the sequence is unimodal.
\end{SourceRcbTtheorem}
This theorem concerns the uniform hypercomb in Question 1.10(4); it does
not assert unimodality of all linear hypertrees.

\subsection{The path determinant}
Put $A=1+tx$ and $B=1+sx$. With no hubs selected, the private vertices
have enumerator $A^{n+1}B^n$. The selected hubs must be nonconsecutive.
Relative to the no-hub weight, each selection contributes $x/A$, as its
tooth is lost; it also removes a factor $B$ for each incident spine
edge. No spine edge is removed twice, since adjacent hubs cannot both
be selected.

Index the edges of an ordinary path on $n+2$ nodes by $0,\ldots,n$.
A hub selection corresponds precisely to a matching of these edges.
Each interior node represents one spine edge, and endpoint nodes represent
no spine edges. In particular this convention works for $n=0$, when
there are just two endpoint nodes. For $x\ne0$ with $|\Arg x|\le
2\pi/3$, define
\[
 d=\sqrt{A/x}=\sqrt{t+1/x},\qquad h=Bd,
\]
using the square root with positive real part (the radicand avoids the
nonpositive real axis). Assign $d$ to each endpoint node and $h$ to every
interior node; denote the resulting diagonal matrix by $D$. If an edge
$uv$ is selected, deleting its two monomers gives relative weight
$(D_{uu}D_{vv})^{-1}$. This supplies exactly $x/A$ and one inverse
factor of $B$ per covered interior node. It follows that
\begin{equation}\label{R21T:det}
 C_{n,\ell}(x)=\frac{A^{n+1}B^n}{d^2h^n}
 \det(D+\ii H),\qquad H=\text{adjacency matrix of the path}.
\end{equation}
Indeed, in the determinant a nonzero permutation term consists only of
fixed points and disjoint transpositions of neighboring indices. Each
such transposition contributes $(-1)\ii^2=1$, and a fixed point contributes
its monomer. Dividing by $d^2h^n$ normalizes the empty matching to one.
The prefactor is finite and nonzero throughout the indicated punctured
sector.

\subsection{Two complementary matrix bounds}
By conjugation it suffices to take $0\le\theta=\Arg x\le2\pi/3$.
The arguments of $1+tx$ and $1+sx$ both belong to $[0,\theta]$.
Consequently
\[
 -\theta/2\le\Arg d=\tfrac12(\Arg(1+tx)-\theta)\le0.
\]
In particular $\Re d>0$. If $\Re h>0$, the Hermitian part of
$D+\ii H$ is the positive definite diagonal matrix $\Re D$. No nonzero
vector can lie in its kernel, by taking the real part of its quadratic
form.

Suppose instead $\Re h\le0$. The unwrapped argument $\Arg B+\Arg d$
belongs to $[-\theta/2,\theta]$, so it can have nonpositive cosine
only if it is at least $\pi/2$. Hence $\Arg B\ge\pi/2$ and
\[
 \theta>\pi/2,\qquad s|x|\ge-1/\cos\theta,\qquad
 |B|\ge \Im B=s|x|\sin\theta\ge-\tan\theta\ge\sqrt3 .
\]
Completing the square with $r=1/|x|$ gives the independent bound
\[
 |d|^4=|t+r e^{-\ii\theta}|^2
       =(r+t\cos\theta)^2+t^2\sin^2\theta
       \ge t^2\sin^2\theta.
\]
As $t\ge2$ and $\sin\theta\ge\sqrt3/2$, we have
$|d|\ge3^{1/4}>1$ and $|h|\ge3^{3/4}>2$.
The path has one off-diagonal entry of modulus one in each endpoint
row, and two in each interior row. Thus the matrix in~\eqref{R21T:det} is
strictly diagonally dominant. For completeness, if it annihilated a
nonzero vector, a coordinate of largest modulus would violate the
strict row bound. This proves nonsingularity also in the second case.
Since $C_{n,\ell}(0)=1$, the claimed zero exclusion follows.

\subsection{From the sector to strict log-concavity}
A real root must be negative. For a nonreal conjugate pair write the two
roots as $-\rho e^{\pm\ii\phi}$ with $0<\phi<\pi/3$. Its real monic
factor has coefficients $(\rho^2,2\rho\cos\phi,1)$, all positive,
and $(2\rho\cos\phi)^2>\rho^2$. Real linear factors likewise have
positive coefficients and no interior inequality to check.

Here is the strict convolution argument. Extend each finite positive
coefficient sequence $a=(a_0,\ldots,a_m)$ by zero outside its support.
If it is log-concave, its Toeplitz matrix $T(a)_{ij}=a_{j-i}$ has all
order-two minors nonnegative: for nonzero entries this follows by
comparing products of the successive nonincreasing ratios
$a_k/a_{k-1}$; if the product on the negative diagonal is positive,
contiguity of support forces the other product to be positive as well.
For two such sequences, $T(a*b)=T(a)T(b)$. Cauchy--Binet gives each
order-two minor as a sum of nonnegative products (only finitely many
are nonzero). If both sequences are strictly log-concave, the adjacent
minor $a_i^2-a_{i-1}a_{i+1}$ is positive at \emph{every} support index,
including endpoints. Given a support index $k$ of $a*b$, choose $i,j$ in
the respective supports with $i+j=k$. In the minor with rows $(0,1)$ and
columns $(k,k+1)$, intermediate columns $(i,i+1)$ then yield a strictly
positive summand. Inducting over the factors gives strict log-concavity
of their product. Finally, choosing one private vertex from each of the
$2n+1$ hyperedges realizes the maximum possible size; all smaller
sizes occur by taking subsets. This proves Theorem~\ref{R21T:main}.

\paragraph{Numerical independence of the proof.}
The argument above uses no finite search. The companion integer check
compares hub enumeration with the published recurrence on 130 parameter
pairs, includes seven direct enumerations over actual vertex subsets,
and compares the scaled continuant against 3,672 recurrence evaluations.
\endgroup


\clearpage

\section{An inverse Gudermannian upper bound on its full domain}
\label{section:R22S}
\noindent{\small\textsf{ROUND 22 \enspace / \enspace Showa \enspace / \enspace R22S}}\\[3pt]
\noindent\textbf{Target.} Sun--Chen Conjecture 7.

\noindent\textbf{Editorial status.} Archived internal proof review passed. The exact polynomial certificate establishes positivity on the entire real interval. Its finite coefficient checks are not a sampling argument.

\smallskip\hrule\medskip



\begingroup
\newcommand{\gd}{\operatorname{gd}}
\setcounter{equation}{0}
\renewcommand{\theequation}{\thesection.\arabic{equation}}
\renewcommand{\theHequation}{R22S.\arabic{equation}}
\paragraph{Overview.}
A fractional-power approximation proposed by Sun and Chen is an upper
bound on its entire stated interval. A sixth-degree Taylor minorant for
cosine reduces the derivative comparison to the positivity of a sparse
polynomial. An explicit integer coefficient certificate proves the sign
on the continuum rather than at a mesh of arguments.


\subsection*{The conjecture and the derivative comparison}
For $0\le x<\pi/2$ let $G(x)=\gd^{-1}(x)=\int_0^x\sec s\,ds$.
A local expansion led Sun and Chen~\cite[p.\,140]{R22S:SC} to propose the following
inequality as their Conjecture~7.
\begin{SourceRccStheorem}
For every real $x$ with $0<x<\sqrt{105/47}$,
\[
 G(x)<B(x):=\frac{59x}{24+35(1-47x^2/105)^{59/94}}.
\]
\end{SourceRccStheorem}
\begin{proof}
On this interval $x<3/2<\pi/2$. Set $t=(1-47x^2/105)^{1/94}$
and $w=(105/47)(1-t^{94})=x^2$. Thus $0<t<1$. Introduce
\[
 A=24+35t^{59},\qquad
 M=24t^{35}+35t^{94}+70\frac{59}{94}(1-t^{94})>0 .
\]
Differentiation gives $B'(x)=59 M/(t^{35}A^2)$. Since $G'(x)=\sec x$,
it suffices to show $59M\cos x-t^{35}A^2>0$. By the alternating
cosine expansion,
\[
 \cos x > C(w):=1-\frac w2+\frac{w^2}{24}-\frac{w^3}{720};
\]
the omitted tail starts positive and its terms decrease in absolute
value since $x^2<9/4$. The certificate on the next page shows
$59M C(w)-t^{35}A^2>0$. Therefore $B'(x)>G'(x)$ throughout the
interval; integrating from $0$, where both functions vanish, proves
the strict inequality.
\end{proof}


\subsection*{An exact integer certificate}
\begin{SourceRccSlemma}
For $0<t<1$ the polynomial
$H(t)=\tfrac{78074896}{35}\big(59M C(w)-t^{35}A^2\big)$ is positive.
\end{SourceRccSlemma}
\begin{proof}
The nonzero coefficients $h_i=[t^i]H(t)$, obtained by multiplication,
are as follows. Columns are paired solely to save space.
\begin{center}
\begin{tabular}{rr@{\hspace{16mm}}rr}
\toprule
$i$ & $h_i$ & $i$ & $h_i$ \\
\midrule
0&436284173 &188&54646095\\
35&$-1046570808$ &223&510121080\\
94&486496863 &282&$-100389975$\\
129&2361331512 &317&48915720\\
153&$-2732621360$ &376&$-18213300$\\
\bottomrule
\end{tabular}
\end{center}
Put $u=t^{-1}-1>0$. By the binomial theorem the coefficient of $u^j$ in
$(1+u)^{376}H(1/(1+u))$ is
\[
 b_j=\sum_{i\le 376-j}h_i\binom{376-i}{j}\qquad(0\le j\le376).
\]
Direct integer evaluation gives
\[
 b_0=b_1=b_2=b_3=0,\quad
 b_4=738388790816670,\quad
 \min_{4\le j\le376}b_j=b_{376}=436284173.
\]
For complete reproducibility the following exact loop verifies all
these finite arithmetic statements from the table (Python 3).
\begin{Verbatim}[fontsize=\footnotesize,breaklines,breakanywhere]
from math import comb
I = [0,35,94,129,153,188,223,282,317,376]
h = [436284173,-1046570808,486496863,2361331512,-2732621360,
     54646095,510121080,-100389975,48915720,-18213300]
b = [sum(v*comb(376-i,j) for i,v in zip(I,h) if j<=376-i)
     for j in range(377)]
assert b[:4]==[0]*4 and b[4]==738388790816670
assert min(b[4:])==b[376]==436284173
\end{Verbatim}
Thus the transformed polynomial is strictly positive for every $u>0$,
and division by $(1+u)^{376}$ gives $H(t)>0$.
\end{proof}
\endgroup


\clearpage

\section{A sharp fractional-power lower bound for the inverse sine}
\label{section:R23M}
\noindent{\small\textsf{ROUND 23 \enspace / \enspace Meiji \enspace / \enspace R23M}}\\[3pt]
\noindent\textbf{Target.} Xu--Sun--Chen Conjecture 1.

\noindent\textbf{Editorial status.} Archived internal proof review passed. The root-count certificate excludes roots on the complete comparison domain. This is the inverse-sine conjecture, distinct from the inverse-tangent assertion elsewhere.

\smallskip\hrule\medskip



\begingroup

\setcounter{equation}{0}
\renewcommand{\theequation}{\thesection.\arabic{equation}}
\renewcommand{\theHequation}{R23M.\arabic{equation}}
\paragraph{Overview.}
Xu, Sun and Chen proposed a fractional-power lower bound for the inverse sine
that matches its Taylor expansion through degree seven. We prove the
inequality on the full interval and show more: the derivative of the proposed
bound, divided by the derivative of the inverse sine, is strictly decreasing.
The certificate is a polynomial of degree 8423 with only seven terms.
Three positive roots, counted with multiplicity, are forced at one point,
and another root is forced outside the comparison region. Descartes' rule
then excludes any root in that region.


\subsection*{The comparison}
Write
\[
a=\frac{1459}{595},\quad b=\frac{864}{595},\quad
c=\frac{1741}{1785},\quad r=\frac{1459}{3482},\qquad
L(x)=\frac{a x}{b+(1-cx^2)^r}.
\]
The conjecture is not merely an asymptotic statement: an expansion at zero
alone gives no control close to the other endpoint.
\begin{SourceRcdMtheorem}\label{R23M:thm}
For $0<x<1$, the function $K(x)=\sqrt{1-x^2}\,L'(x)$ is strictly
decreasing. In particular $L(x)<\arcsin x$ throughout this interval, as
asserted by Xu--Sun--Chen in Conjecture~1, Eq.~(23) of \cite{R23M:XSC}.
\end{SourceRcdMtheorem}
\begin{SourceRcdMlemma}\label{R23M:lem}
Suppose a real polynomial $P$ has $P(0)>0$ and its nonzero coefficients
have at most $m+1$ sign changes, with $m$ odd. If $P(1+s)=A s^m+O(s^{m+1})$
with $A>0$, then $P(q)>0$ for all $q>1$.
\end{SourceRcdMlemma}
\begin{proof}
Immediately to the left of 1 the polynomial is negative, so it has a root
in $(0,1)$. This root and the root at 1 already account for $m+1$ positive
roots counting multiplicity, the maximum allowed by Descartes' rule of signs.
Its positive sign immediately to the right of 1 therefore persists.
\end{proof}

\subsection*{Proof of the comparison}
Set $t=1-cx^2$, $d=1-c$, $Y=t^r$ and $D=b+Y$. On the comparison interval
$d<t<1$; in particular all powers below are real, and all factors we clear
are strictly positive. By the quotient rule,
\[
L'(x)=a\frac{N}{D^2},\qquad
N=b+(1-2r)Y+2rY/t .
\]
Consequently $K(x)=a\sqrt{(t-d)/c}\,N/D^2$. Using primes for total
derivatives in $t$, define
\[
H=t^2\bigl[DN+2(t-d)(DN'-2ND')\bigr].
\]
The derivative of $K$ with respect to $t$ equals
$aH/(2\sqrt{c}\sqrt{t-d}\,t^2D^3)$, with $Y'=rY/t$.
The key step is to retain only the sparse terms of $H$, not to expand around
the boundary to full degree. Set $q=t^{-1/3482}$ and form
\[
P(q)=\frac{1073078451025}{36}\,q^{9882}H .
\]
Substitution of $t=q^{-3482}$ and $Y=q^{-1459}$ gives the following complete
coefficient table for $P(q)=\sum_j B_j q^j$:
\[
\begin{array}{r|r@{\qquad}r|r}
j&B_j&j&B_j\\ \hline
0&782041225&4941&-4231917040\\
1459&-16375689960&6964&1747495365\\
2918&62852495616&8423&1038948064\\
3482&-45813373270&&
\end{array}
\]
In particular the signs in increasing order of degree are $+,-,+,-,-,+,+$.
For $k=0,1,2,3$ respectively, one obtains
\[
\sum_j B_j\binom{j}{k}
\;=\;0,\quad 0,\quad 0,\quad 46206946595031268620 .
\]
These are the first four Taylor coefficients of $P(1+s)$. There are exactly
four sign changes, so Lemma~\ref{R23M:lem} applies with $m=3$. Since $q>1$
throughout $d<t<1$, we have $P(q)>0$, hence $H>0$. Thus $K$ is strictly
increasing in $t$ and strictly decreasing in $x$. As $a=b+1$, its value
at zero is $K(0)=1$. This gives $L'(x)<1/\sqrt{1-x^2}=(\arcsin x)'$.
Integration from zero proves Theorem~\ref{R23M:thm}. \hfill$\square$
\endgroup


\clearpage

\section{Strict sign regularity of the reciprocal Jackson gamma kernel}
\label{section:R24Hq}
\noindent{\small\textsf{ROUND 24 \enspace / \enspace Heisei \enspace / \enspace R24Hq}}\\[3pt]
\noindent\textbf{Target.} Karp--Vishnyakova--Zhang, Section 3, Example 8.

\noindent\textbf{Editorial status.} Archived internal proof review passed. The target is the reciprocal Jackson gamma kernel at all orders; total positivity of the positive gamma kernel alone would not imply this conclusion.

\smallskip\hrule\medskip



\begingroup
\newcommand{\epsr}{\varepsilon_r}
\setcounter{equation}{0}
\renewcommand{\theequation}{\thesection.\arabic{equation}}
\renewcommand{\theHequation}{R24Hq.\arabic{equation}}
\paragraph{Overview.}
For each $0<q<1$, every order-$r$ collocation determinant of
$1/\Gamma_q(x+y)$ on the positive half-line has strict sign
$(-1)^{r(r-1)/2}$. This proves the reciprocal-kernel conjecture in
Example 8 of Section 3 of Karp, Vishnyakova and Zhang. The proof is a
$q$-binomial factorization: a totally positive shifted-factorial matrix
is composed with a reversed monomial matrix. A fixed positive summand
in each signed Cauchy--Binet expansion certifies strictness even in
the infinite limit.


\subsection{The determinant and the scope}
Fix $0<q<1$. Write $(u;q)_n=\prod_{j=0}^{n-1}(1-uq^j)$, with empty product
one, and $(u;q)_\infty=\prod_{j=0}^\infty(1-uq^j)$. For $s>0$ set
\[
 \Gamma_q(s)=(1-q)^{1-s}\frac{(q;q)_\infty}{(q^s;q)_\infty},
 \qquad \epsr=(-1)^{\binom r2}.
\]
A kernel is strictly sign regular of all orders if, at every order $r$,
all its determinants formed from strictly increasing row and column nodes
are nonzero and have the same sign.
\begin{SourceRceHqtheorem}\label{R24Hq:main}
For $r\ge1$, $0<x_1<\cdots<x_r$ and $0<y_1<\cdots<y_r$,
\[
 \epsr\det\left[\frac{1}{\Gamma_q(x_i+y_j)}\right]_{i,j=1}^r>0 .
\]
In fact $x_i>0$ and $y_j>0$ may be replaced by $x_1+y_1>0$.
\end{SourceRceHqtheorem}
The conjecture here is the unnumbered one about the \emph{reciprocal}
kernel in \cite[Section 3]{R24Hq:KVZ}; it is not the separately numbered
Bessel Conjecture 1. The positive kernel $\Gamma_q(x+y)$ was treated
by Salazar \cite{R24Hq:S}, who distinguishes the reciprocal question. Strict
total positivity of a kernel does not generally establish sign
regularity of its entrywise reciprocal. Subsection~\ref{R24Hq:bridge} gives the
additional factorization that is needed here.

\subsection{A positive auxiliary kernel}
We recall two facts and give the short arguments that fix all signs.
If $0<t_1<\cdots<t_r$ and $0\le k_1<\cdots<k_r$ are integers, then
$\det[t_i^{k_j}]>0$. Indeed, by Descartes' rule of signs a nonzero linear
combination of $r$ such monomials has at most $r-1$ positive roots, so
the determinant never vanishes. Its sign is constant on the connected
set of increasing nodes. At $t_i=t^{i-1}$, $t>1$, the determinant is the
ordinary Vandermonde on the increasing numbers $t^{k_j}$, hence positive.
For decreasing positive nodes, reversing the rows multiplies the sign by
$\epsr$.

For $0\le A<1$ and $0\le z<1$ the $q$-binomial identity reads
\begin{equation}\label{R24Hq:qb}
 \frac{(Az;q)_\infty}{(z;q)_\infty}
    =\sum_{n\ge0}\frac{(A;q)_n}{(q;q)_n}z^n .
\end{equation}
For completeness, the product ratio $F(z)$ is analytic on $|z|<1$,
$F(0)=1$, and $(1-z)F(z)=(1-Az)F(qz)$. Equating Taylor coefficients
recursively proves \eqref{R24Hq:qb}.

\begin{SourceRceHqlemma}\label{R24Hq:pos}
Choose strictly increasing lists $x_i>0$ and $n_j\in\mathbb Z_{\ge0}$. Then
$\det[(q^{x_i};q)_{n_j}]>0$.
\end{SourceRceHqlemma}
\begin{proof}
Using \eqref{R24Hq:qb} with $A=0$, write
\[
 H(x,n)=\frac{1}{(q^{x+n};q)_\infty}
 =\sum_{k\ge0}\frac{(q^x)^k(q^n)^k}{(q;q)_k}.
\]
In the Cauchy--Binet expansion of its order-$r$ determinant truncated at
$k=M\ge r-1$, every term has positive sign: the determinants on
$q^{x_i}$ and $q^{n_j}$ each have sign $\epsr$ and all weights are positive.
In particular the term on $k=0,\ldots,r-1$ is fixed and positive.
The entries converge and hence so do the determinants, which retain
this positive lower bound. Finally
$(q^x;q)_n=(q^x;q)_\infty H(x,n)$, a positive row scaling.
\end{proof}

\subsection{The bridge to the reciprocal}\label{R24Hq:bridge}
\begin{proof}[Proof of Theorem~\ref{R24Hq:main}]
Set $z_j=q^{y_j}$; these nodes are strictly decreasing in $(0,1)$.
Using \eqref{R24Hq:qb} with $A=q^{x_i}$, put
\[
 T_{ij}=\frac{(q^{x_i}z_j;q)_\infty}{(z_j;q)_\infty}
       =\sum_{n\ge0}(q^{x_i};q)_n\frac{z_j^n}{(q;q)_n}.
\]
Let $T^{(M)}$ use only $0\le n\le M$. For $M\ge r-1$,
\[
 \epsr\det T^{(M)}
 =\sum_{0\le n_1<\cdots<n_r\le M}
  \det[(q^{x_i};q)_{n_h}]_{i,h}
  \frac{\epsr\det[z_j^{n_h}]_{h,j}}{\prod_{h=1}^r(q;q)_{n_h}} .
\]
Lemma~\ref{R24Hq:pos} and the decreasing-node Vandermonde sign show that
\emph{every} summand is strictly positive. The contribution
$(n_1,\ldots,n_r)=(0,1,\ldots,r-1)$ is independent of $M$. Passing to
the limit therefore gives $\epsr\det T>0$. The target matrix is obtained by
scaling row $i$ by $(1-q)^{x_i}$ and column $j$ by
$(1-q)^{y_j-1}(z_j;q)_\infty/(q;q)_\infty$, all strictly positive.
If merely $x_1+y_1>0$, replace each $x_i$ by $x_i+c$ and each $y_j$
by $y_j-c$, where $-x_1<c<y_1$. All sums are unchanged.
\end{proof}
\begin{SourceRceHqcorollary}
Choose strictly increasing lists $x_i>0$ and $n_j\in\mathbb Z_{\ge0}$. Then
$\epsr\det[1/(q^{x_i};q)_{n_j}]>0$. Thus both finite
$q$-factorial conclusions of \cite[Appendix A]{R24Hq:KVZ} hold at every order,
with positive sign in Lemma~\ref{R24Hq:pos} and reversed-column sign for the reciprocal.
\end{SourceRceHqcorollary}
\begin{proof}
Use $1/(q^x;q)_n=(1-q)^{-n}\Gamma_q(x)/\Gamma_q(x+n)$
and the extended domain in Theorem~\ref{R24Hq:main}.
\end{proof}
\paragraph{Verification and status.}
The proof is independent of numerical testing. The accompanying exact-arithmetic script checks the final Cauchy--Binet expansion
and a grid of finite shifted-factorial minors. The all-order conclusion is a manuscript result, not
a claim of peer review or established priority.
\endgroup


\clearpage

\section{A global Gudermannian bound by signed quadrature}
\label{section:R25M}
\noindent{\small\textsf{ROUND 25 \enspace / \enspace Meiji \enspace / \enspace R25M}}\\[3pt]
\noindent\textbf{Target.} Sun--Chen Conjecture 10.

\noindent\textbf{Editorial status.} Archived internal proof review passed. The proof covers every positive real argument. Both the explicit error bounds and the companion quadrature lower bound are retained.

\smallskip\hrule\medskip



\begingroup
\newcommand{\sech}{\operatorname{sech}}
\newcommand{\gd}{\operatorname{gd}}



\newcommand{\I}{\mathcal I}
\setcounter{equation}{0}
\renewcommand{\theequation}{\thesection.\arabic{equation}}
\renewcommand{\theHequation}{R25M.\arabic{equation}}
\paragraph{Overview.}
Sun and Chen conjectured an upper bound for the mean value of the
hyperbolic secant, and proved it for $0<x<6/5$. We prove the bound
on the entire positive half-line. The two terms in the conjectured
bound are the endpoint and interior terms of a Gauss--Radau rule
after changing to the squared integration variable. The product
for the hyperbolic cosine fixes the sign of the quadrature error.
We also give an explicit, asymptotically sharp upper bound for that error.


\subsection{The inequality and its constants}
Write $\gd(x)=\int_0^x\sech s\,ds$.
Conjecture~10 of Sun and Chen~\cite[p.\,142, (28)]{R22S:SC} is the following.
\begin{SourceRcfMtheorem}\label{R25M:thm}
For every $x>0$,
\[
  \frac{\gd(x)}{x}<\frac{4+5\sech(\sqrt{3/5}\,x)}{9}.
\]
\end{SourceRcfMtheorem}
Their Theorem~12 proves the inequality locally, for $0<x<6/5$.
We establish the sign globally, without a cutoff or numerical
evaluation. Put $f_x(u)=\sech(x\sqrt{u})$ and define
\[
 \I(h)=\int_0^1 h(t^2)\,dt,\qquad
 Q(h)=\tfrac49 h(0)+\tfrac59 h(3/5).
\]
Then $\I(f_x)=\gd(x)/x$. The moments $\I(u^k)=1/(2k+1)$
explain all the constants: a two-node rule with nodes $0,a$ and
interior weight $w$ integrates $1,u,u^2$ exactly precisely when
$wa=1/3$ and $wa^2=1/5$, hence $a=3/5$, $w=5/9$. This
is the standard left Gauss--Radau rule for $du/(2\sqrt{u})$ on
$[0,1]$. The point is to apply it to $f_x$, rather than to
differentiate $\sech(xt)$ three times in $t$.

\begin{SourceRcfMlemma}\label{R25M:sign}
For each $x>0$, $f_x'''(u)<0$ on $[0,1]$, and $-f_x'''(u)$
is strictly decreasing there.
\end{SourceRcfMlemma}
\begin{proof}
The classical product for $\cosh$~\cite[(4.36.2)]{R25M:dlmf} gives
\[
 f_x(u)=\prod_{n\ge1}(1+c_nu)^{-1},\qquad
 c_n=\frac{4x^2}{(2n-1)^2\pi^2}>0.
\]
Set $S_j(u)=\sum_{n\ge1}\bigl(c_n/(1+c_nu)\bigr)^j$.
Since $\sum c_n<\infty$, both the series of logarithms and
its first three derivatives converge uniformly on $[0,1]$
(the terms are bounded by $c_n$ and $(j-1)!c_n^j$ respectively).
The logarithmic derivatives of $f_x$ are therefore $-S_1,S_2,-2S_3$.
It follows that
\begin{equation}
 -f_x'''(u)=f_x(u)\bigl(S_1(u)^3+3S_1(u)S_2(u)+2S_3(u)\bigr).
 \label{R25M:deriv}
\end{equation}
Every factor and every $S_j$ on the right is positive and strictly
decreasing with $u$. This proves both assertions.
\end{proof}

\subsection{Interpolation gives the global sign}
\begin{proof}[Proof of Theorem~\ref{R25M:thm}]
Let $a=3/5$. For $h\in C^3[0,1]$ let $H$ be the quadratic
(degree at most two) matching $h$ at $0,a$ and $h'$ at $a$.
For $u\notin\{0,a\}$ the Hermite remainder reads
\[
 h(u)-H(u)=\frac{h'''(\xi_u)}{6}u(u-a)^2,\qquad 0<\xi_u<1.
\]
Indeed one subtracts $C v(v-a)^2$ from $h(v)-H(v)$ to make
the result vanish at $v=u$; the four zeros counted with multiplicity
give a zero of its third derivative by Rolle's theorem.
For $h=f_x$ the right-hand side is strictly negative, apart
from the interpolation nodes where the difference is zero.
Now $\I(H)=Q(H)=Q(f_x)$ by moment exactness. Integration gives
$\I(f_x)<Q(f_x)$, as required.
\end{proof}

\begin{SourceRcfMcorollary}\label{R25M:error}
The gap $D(x)=Q(f_x)-\gd(x)/x$ satisfies, for all $x>0$,
\[
 \frac{x^3\tanh^3(x)\sech(x)}{2100}
 <D(x)<\frac{61}{31500}x^6.
\]
The coefficient of $x^6$ on the right cannot be decreased.
\end{SourceRcfMcorollary}
\begin{proof}
The integral of the remainder factor is $\I(u(u-a)^2)=4/175$.
By Lemma~\ref{R25M:sign}, we have
\[
 (2/525)(-f_x'''(1))<D(x)<(2/525)(-f_x'''(0)).
\]
At $u=1$ use
$S_1(1)=x\tanh(x)/2$ in \eqref{R25M:deriv} and discard the other
positive terms. At zero the term of degree six in $\sech z$ is
$-61z^6/720$, whence $-f_x'''(0)=61x^6/120$.
Since $Q$ matches the moments up to degree two,
the same Taylor expansion yields
$D(x)=(61/31500)x^6+O(x^8)$ as $x\to0$.
\end{proof}

\noindent A companion lower bound comes from fixing the right
endpoint instead of the left. The rule
$Q_R(h)=(5h(1/5)+h(1))/6$ again matches the first three moments.
Interpolating at $1/5$ twice and at $1$ once replaces the remainder
factor by $(u-1/5)^2(u-1)$, so its sign reverses:
\[
 \frac{5\sech(x/\sqrt{5})+\sech x}{6}<\frac{\gd(x)}{x}
 \qquad(x>0).
\]
\endgroup


\clearpage

\section{Strict log-convexity of the even Eulerian maxima}
\label{section:R25K}
\noindent{\small\textsf{ROUND 25 \enspace / \enspace Keio \enspace / \enspace R25K}}\\[3pt]
\noindent\textbf{Target.} Ambrus--Gargyan log-convexity conjecture.

\noindent\textbf{Editorial status.} The ratio-bound proof passed the archived supplementary review. Both the ratio proof and the separate signed-integral proof, the finite certificate, and the higher fixed-order Hankel consequence are reproduced. The archived supplementary audit specifically certifies the ratio route; it does not certify all of the integral extensions.

\smallskip\hrule\medskip



\begingroup
\newcommand{\sinc}{\operatorname{sinc}}
\setcounter{equation}{0}
\renewcommand{\theequation}{\thesection.\arabic{equation}}
\renewcommand{\theHequation}{R25K.\arabic{equation}}
\paragraph{Overview.}
For the largest Eulerian number $M_m$ of order $m$, let
$a_p=M_{2p}/(2p)!$. We prove that $a_{p-1}a_{p+1}>a_p^2$
exactly for the centers $p\ge3$ among the indices $p\ge2$.
This proves the conjecture of Ambrus and G\'argy\'an stated after their
Proposition~4.4. The first proof compares a lower bound and an upper
bound at different indices. The second separates the positive central
lobe of a signed sinc integral from its tail. The latter argument also
proves eventual strict positivity of every fixed-order consecutive
Hankel determinant of the same sequence.


\subsection{Statement and the first proof}
Write $A(m,k)$ for the number of permutations of $m$ letters with
$k-1$ ascents, $1\le k\le m$, and $f(m)=\max_k A(m,k)/m!$.
The exact published conjecture concerns the sequence
$(f(2p))_{p\ge2}$ \cite[end of Section~4]{R25K:AG}. Its first three
terms come from rows $4,6,8$. The conjecture therefore begins with
a three-term inequality centered at row 6, not row 4.
\begin{SourceRcfKtheorem}\label{R25K:t:main}
Set $a_p=f(2p)$ for $p\ge1$. For every integer $p\ge3$,
$D_p:=a_{p-1}a_{p+1}-a_p^2>0$, whereas $D_2=-1/2880$.
\end{SourceRcfKtheorem}
\begin{proof}
Let $R_p=a_p/a_{p-1}$. The classical ratio estimates of Lesieur and
Nicolas \cite[Theorem~2, p.~395]{R25K:LN}, as transferred to even rows
in the unnumbered display immediately after (4.11) of \cite{R25K:AG}, give,
for the indices $p\ge 3$ used below,
\[
 L(p):=\frac{2p^2}{(p+1)(2p-1)}
   <R_p<
 U(p):=\frac{p(2p+1)^2}{2(p+1)^2(2p-1)}.
\]
Specifically, the bounds in (4.8) of \cite{R25K:AG} are on
$f(2p+1)/f(2p-1)$; multiply both by
$p(2p+1)/((p+1)(2p-1))$, using their (4.10).
The authors' weaker upper bound in (4.11) does not suffice here.
Comparing at successive indices yields
\[
 \frac{D_p}{a_p^2}=\frac{R_{p+1}}{R_p}-1
 >\frac{L(p+1)}{U(p)}-1
 =\frac{2p^3-5p^2-10p-4}{p(p+2)(2p+1)^3}.
\]
With $p=h+4$, the numerator is $2h^3+19h^2+46h+4>0$ for
$h\ge0$. The first four values $a_1,\ldots,a_4$ are
$1/2$, $11/24$, $151/360$, and $15619/40320$; substitution gives
$D_3=23423/14515200>0$ and the asserted value of $D_2$.
\end{proof}
Thus Theorem~\ref{R25K:t:main} is a short consequence of a classical
ratio theorem, not a new estimate on those ratios. The next proof
does not use that theorem.
\subsection{An independent estimate for the signed moments}
Let $s(x)=\sin(x)/x$, with $s(0)=1$, and put $z=s^2$,
$w=s\cos$. The integral representation \cite[(1.1), (4.7)]{R25K:AG} gives
\[
 a_k=\frac2\pi\int_0^\infty w(x)z(x)^k\,dx
 =\frac2\pi(B_k+E_k),\qquad
 B_k=\int_0^{\pi/2} w(x)z(x)^k\,dx .
\]
The factors in $B_k$ are nonnegative, and $B_k\le\pi/2<2$.
Since $|s(x)|\le1/x$ for $x\ge\pi/2$, for $k\ge1$ we have
\[
 |E_k|\le \int_{\pi/2}^\infty x^{-2k-1}\,dx
 =\frac{(2/\pi)^{2k}}{2k}<(2/3)^{2k}.
\]
Set $T_p=B_{p-1}B_{p+1}-B_p^2$. Expanding the two products, the
four mixed linear terms have total absolute value at most $8q$
and the two quadratic terms at most $2q^2$, where
$q=(2/3)^{2p-2}$. Consequently
\begin{equation}\label{R25K:eq:tail}
 \big|(\pi/2)^2D_p-T_p\big|<12(2/3)^{2p-2}\qquad (p\ge2).
\end{equation}
On the square $[0,\pi/2]^2$, the integrand in the identity
\[
 T_p=\tfrac12\iint w(x)w(y)[z(x)z(y)]^{p-1}
                  [z(x)-z(y)]^2\,dx\,dy
\]
has a definite sign. To obtain a lower bound, take $p\ge100$, and
restrict to the rectangle $0\le x\le 1/(2\sqrt p)$,
$1/\sqrt p\le y\le2/\sqrt p$ and its transpose.
For $0\le t\le1$ the elementary bounds
\[
 1-t^2/3\le z(t)\le 1-t^2/4,\qquad w(t)>1/3
\]
hold. Indeed $1-t^2/6\le s(t)\le1-t^2/6+t^4/120
\le1-t^2/7$ by alternating Taylor bounds, and
$(1-t^2/7)^2\le1-t^2/4$; also $s(t)\ge5/6$ and
$\cos t\ge 1-t^2/2\ge1/2$. Hence on the chosen rectangle
$z(x)-z(y)\ge 1/(6p)$. Moreover $z(t)^{p-1}\ge
(1-2/p)^p\ge 1/16$ there: $u\log(1-2/u)$ is increasing for
$u>2$ since $\log(1+v)\le v$ for $v=2/(u-2)$, and the
value of $(1-2/u)^u$ at $u=4$ is $1/16$.
The area of one rectangle is $1/(2p)$, so
\[
 T_p\ge \frac{1}{2p}\left(\frac{1}{48}\right)^2
                   \left(\frac{1}{6p}\right)^2
     = \frac1{165888\,p^3}.
\]
This exceeds the error in \eqref{R25K:eq:tail} for $p\ge100$.
At $100$ the comparison reduces to
$12\cdot165888\cdot100^3\cdot 2^{198}<3^{198}$.
The ratio of successive values of $p^3(2/3)^{2p-2}$ is at most
$(101/100)^3(4/9)<1$ for $p\ge100$.
For completeness, the following entire integer certificate verifies the
remaining finite range. It uses the usual Eulerian recurrence with
zero-based column indices; the last inequality clears all factorials.
\begin{Verbatim}[fontsize=\footnotesize,breaklines,breakanywhere]
from math import factorial as fact
r = [1]; M = {}
for n in range(1,201):
    r = [(j+1)*(r[j] if j<len(r) else 0)
         + (n-j)*(r[j-1] if j else 0) for j in range(n)]
    if n%2 == 0: M[n//2] = max(r)
for p in range(3,100):
    assert (M[p-1]*M[p+1]*fact(2*p)**2 >
            M[p]**2*fact(2*p-2)*fact(2*p+2))
\end{Verbatim}
Together these estimates give a second proof of Theorem~\ref{R25K:t:main}.
\subsection{Higher consecutive Hankel determinants}
\begin{SourceRcfKproposition}
For each fixed $d\ge1$, there exists $N_d$ such that, for every
integer $n\ge N_d$, one has
$\det(a_{n+i+j})_{0\le i,j<d}>0$.
\end{SourceRcfKproposition}
\begin{proof}
Expanding the determinant under the integrals and symmetrizing gives
\[
 \det(B_{n+i+j})_{0\le i,j<d}
 =\frac1{d!}\int_{[0,\pi/2]^d}\prod_{j=1}^d w(x_j) z(x_j)^n
     \prod_{i<j}(z(x_i)-z(x_j))^2\,d\mathbf{x}.
\]
This follows also by expanding the square of the Vandermonde determinant.
Restrict to the box $\sqrt n x_j\in [j,j+1/4]$, $1\le j\le d$,
for $n$ sufficiently large. Uniformly for $u$ in these fixed intervals,
$z(u/\sqrt n)=1-u^2/(3n)+O_d(n^{-2})$; therefore
$z(u/\sqrt n)^n\to e^{-u^2/3}$ and $w(u/\sqrt n)\to1$ uniformly.
The squared coordinates $u^2$ on distinct intervals are bounded apart.
Every factor $|z(x_i)-z(x_j)|$ is thus at least $c_d/n$,
for some $c_d>0$ and all large $n$. The box has volume
$4^{-d}n^{-d/2}$, and the full integrand is nonnegative. The determinant
on the left is accordingly at least
$C_d n^{-d/2-d(d-1)}$ with $C_d>0$.
Replacing each entry $B_{n+i+j}$ by $B_{n+i+j}+E_{n+i+j}$
alters the determinant by $O_d((2/3)^{2n})$ by the finite determinant
expansion and the uniform bound on its entries. Multiplication by
$(2/\pi)^d>0$ gives the claimed determinant.
\end{proof}
The threshold here may depend on $d$. No uniform total positivity
or positive moment representation for the full sequence is asserted.
\paragraph{Attribution.}
A separate Reiwa analysis in the same research session independently
identified the central-lobe versus signed-tail approach.
The ratio theorem and the integral representation are from the
literature; no priority claim is made for the general determinant
identity or for the analytic method.
\endgroup


\subsection{Independent derivation of the ratio proof}

\label{section:R25R}

\noindent\textbf{Source.} Reiwa, Round 25. This is an independent derivation of the same main Eulerian result, not a second solved topic. The first three-term center is $p=3$; extending it to $p=2$ is false.



\begingroup

\setcounter{equation}{0}
\renewcommand{\theequation}{\thesection.\arabic{equation}}
\renewcommand{\theHequation}{R25R.\arabic{equation}}
Let $M_m$ be the largest coefficient of the Eulerian polynomial of order $m$, and
put $f(m)=M_m/m!$. At the end of Section~4 of their paper \cite{R25K:AG}, Ambrus and
G\'argy\'an conjectured that $(f(2p))_{p\ge2}$ is log-convex. They proved
ordinary convexity and noted that their estimate (4.11) falls just short of
log-convexity. The sharper estimate printed directly below (4.11), however,
suffices after a shift of index.

\textbf{Theorem.} $f(2p-2)f(2p+2)>f(2p)^2$ for every integer $p\ge3$.
Thus the conjectured log-convexity holds strictly, starting with the triple
$f(4),f(6),f(8)$.

\textit{Proof.} Set $r_p=f(2p)/f(2p-2)$ for $p\ge2$. The strict Lesieur--Nicolas
estimate, as stated in \cite[Theorem~4.2]{R25K:AG}, and the even-to-odd conversion
\cite[equation~(4.10)]{R25K:AG} read
\[
 \frac{2p}{2p+1}<\frac{f(2p+1)}{f(2p-1)}<\frac{2p+1}{2p+2},
 \qquad f(2p)=\frac{2p+1}{2p+2}f(2p+1).
\]
Applying the second identity twice and the first inequality once gives
\[
 L_p:=\frac{2p^2}{(p+1)(2p-1)}<r_p<
 \frac{p(2p+1)^2}{2(p+1)^2(2p-1)}=:U_p.
\]
It suffices to put the lower bound at the \emph{next} index above the upper
bound at the current one. In fact
\[
 L_{p+1}-U_p=
 \frac{2p^3-5p^2-10p-4}{2(p+1)^2(p+2)(2p-1)(2p+1)}.
\]
For $u=p-4\ge0$, the numerator is $2u^3+19u^2+46u+4>0$.
Thus $r_{p+1}>r_p$ for $p\ge4$.
The remaining case follows from the exact central Eulerian numbers
$M_4=11$, $M_6=302$, $M_8=15619$:
\[
 f(4)f(8)-f(6)^2=\frac{23423}{14515200}>0.\qquad\square
\]

\textbf{Endpoint and credit.} The sequence in \cite{R25K:AG} starts at $p=2$, so its
first interior center is $p=3$. In fact $f(2)f(6)-f(4)^2=-1/2880$.
The estimates are classical; the contribution here is the shifted
comparison. No new bound on the Laplace--P\'olya integral is needed.
The Keio seat independently considered the same original estimates and a
signed-integral approach. No external verification or priority is asserted.
\endgroup


\clearpage

\part{Additional complete all-parameter arguments}

\noindent These two manuscripts contain complete-range arguments that were not counted as tournament victories. The first addresses two Ulas conjectures; the second treats a separately posed equal-shift extension. The scope of a mathematical argument and the outcome of a contest are different records.


\section{Uniform real-root bounds and positivity of Stern polynomials}
\label{section:R02K}
\noindent{\small\textsf{ROUND 2 \enspace / \enspace Keio \enspace / \enspace R02K}}\\[3pt]
\noindent\textbf{Target.} Ulas Conjectures 4.2(2) and 4.4.

\noindent\textbf{Editorial status.} Complete-range argument preserved; no final tournament adjudication or priority certification. Both numbered Ulas conjectures are retained. The original contest rejected the new-method qualification, not the universal scope of the displayed arguments; priority and then-open status were not certified.

\smallskip\hrule\medskip



\begingroup
\newcommand{\R}{\mathbb R}
\newcommand{\N}{\mathbb N}
\setcounter{equation}{0}
\renewcommand{\theequation}{\thesection.\arabic{equation}}
\renewcommand{\theHequation}{R02K.\arabic{equation}}
\paragraph{Overview.}
We give complete elementary arguments for two statements posed by Ulas:
Conjecture 4.2(2), concerning a uniform bound on real roots in a
parameterized Stern family, and Conjecture 4.4, concerning the absence of
real roots in another family. A sparse numerator gives the explicit bound
$17$, simultaneously for every $k\ge2$ and $n\ge1$. A corrected quadratic
identity gives strict positivity on the whole real line for every member
of the second family. Both arguments cover the original quantifiers and
use no computational inference or unproved hypothesis. These consequences are not submitted as historically novel breakthroughs.
The requested corpus already explicitly screens out square-completion and
fewnomial consequences of Ulas's formulas as new methods. Thus the
all-parameter mathematical arguments pass, but the tournament's new-method
requirement does not. Priority and current open status are not certified.


\subsection{Precise targets and attribution}
The Stern polynomials are
\begin{equation}\label{R02K:def:B}
 B_0(t)=0,\quad B_1(t)=1,\qquad
 B_{2m}(t)=tB_m(t),\quad B_{2m+1}(t)=B_m(t)+B_{m+1}(t).
\end{equation}
We use the definition in the body of Ulas~\cite{R02K:Ulas}; its abstract has a
subscript typo. Our exact target text is the accessible version
\texttt{arXiv:1909.10844v1}, page 16. The publisher identifies the later
2020 article~\cite{R02K:Publisher}, but its full text was subscription-blocked
in this session; no assertion about silent changes to that version is made.

For integers $k\ge2$ and $n\ge1$, put
\[
p_{k,n}=2^{2n+k}-3\cdot2^{n+k-1}+2^k-3,\qquad
N_k(n)=\#\{t\in\R:B_{p_{k,n}}(t)=0\}.
\]
\textbf{Target A, Ulas Conjecture 4.2(2).} For every odd $k\ge3$, find a
constant $C_k$ that bounds $N_k(n)$ for all $n\ge1$.
The relevant clause in the source reads:
``there is a number $C_k$ such that for each $n\in\N_+$ we have
$N_k(n)\le C_k$''.

\textbf{Target B, Ulas Conjecture 4.4.} For every integer $n\ge0$, define
\[
h_n=\frac23(2^{2n}-1)(2^{2n+1}+1)+1.
\]
The assertion is that $B_{h_n}$ has no real zero, stated in the source as
``the equation $B_{h_n}(t)=0$ has no real roots''.
The case $n=0$, where $B_{h_0}=B_1=1$, is included.

These are two distinct numbered assertions. In particular, Target A is
not Conjecture 4.2(1), whose monotonicity clause is discussed in corpus
manuscript [071]~\cite{R02K:Corpus}. We make no claim about the irreducibility
clauses or the remaining assertions of Conjecture 4.3.

\begin{SourceRacKtheorem}[Uniform real-root bound]\label{R02K:thm:bound}
For every $k\ge2$ and every $n\ge1$, $N_k(n)\le17$.
The same bound holds even when roots are counted with multiplicity.
Thus Target A holds with $C_k=17$, independent of $k$.
\end{SourceRacKtheorem}
\begin{SourceRacKtheorem}[Strict positivity]\label{R02K:thm:positive}
For every $n\ge0$ and every real $t$, $B_{h_n}(t)>0$.
In particular, Target B holds.
\end{SourceRacKtheorem}

\subsection{Elementary Stern identities}
The following standard addition identity is also recorded in
\cite[Lemma 2.1]{R02K:Ulas}; we include its proof.
\begin{SourceRacKlemma}\label{R02K:lem:add}
For integers $r,m\ge0$ and $0\le b\le2^r$,
\begin{align}
 B_{m2^r+b}&=B_{2^r-b}B_m+B_bB_{m+1},\label{R02K:eq:plus}\\
 B_{m2^r-b}&=B_{2^r-b}B_m+B_bB_{m-1}\quad(m\ge1).\label{R02K:eq:minus}
\end{align}
All polynomials are evaluated at the same variable $t$.
\end{SourceRacKlemma}
\begin{proof}
For $r=0$, \eqref{R02K:eq:plus} follows at $b=0,1$ from $B_0=0,B_1=1$.
Induct on $r$. At level $r+1$, if $b=2c$, use $B_{2j}=tB_j$ on both
sides and the level-$r$ identity. If $b=2c+1$, add the level-$r$
identities for $c$ and $c+1$. The coefficients combine by
$B_c+B_{c+1}=B_{2c+1}$ and
$B_{2^r-c-1}+B_{2^r-c}=B_{2^{r+1}-2c-1}$. This proves the plus identity
including the endpoints. Substitute $m-1$ and $2^r-b$ into it to obtain
the minus identity.
\end{proof}
The recursion also gives, for $r\ge0$ and $s\ge2$,
\begin{equation}\label{R02K:eq:geometric}
 B_{2^r}=t^r,\quad B_{2^r-1}=\frac{t^r-1}{t-1},\quad
 B_{2^s-3}=\frac{2t^{s-1}-t-1}{t-1}.
\end{equation}
The middle identity includes $r=0$. These are polynomial identities:
the quotients are geometric sums. They follow by induction from
\eqref{R02K:def:B}; in the last identity, use
$B_{2^s-3}=B_{2^{s-1}-2}+B_{2^{s-1}-1}$ for $s\ge3$, starting at $s=2$.

Every coefficient of every $B_j$ is nonnegative. Moreover, $B_j(0)=1$
for odd $j$: in the odd-index recursion exactly one of the two smaller
indices is odd. These facts follow simultaneously by induction, with
$B_0(0)=0$ and $B_1(0)=1$. Consequently an odd-index Stern polynomial
is strictly positive for $t\ge0$.

\subsection{Proof of the uniform bound}
Write $V_{k,n}=B_{p_{k,n}}$, and put $u=t^{n-1}$. Applying
\eqref{R02K:eq:plus} to
\[
p_{k,n}=2^{n+k-1}(2^{n+1}-3)+(2^k-3)
\]
and then \eqref{R02K:eq:minus} to the complementary index gives
\begin{align*}
V_{k,n}
&=\left((t+1)t^{n-1}+B_{2^k-3}B_{2^{n-1}-1}\right)B_{2^{n+1}-3}
   +B_{2^k-3}B_{2^{n+1}-2}.
\end{align*}
This calculation is valid also when $n=1$, because $B_{2^{n-1}-1}=B_0=0$.
Using \eqref{R02K:eq:geometric} and
$B_{2^{n+1}-2}=t(t^n-1)/(t-1)$ yields the identity already appearing in
\cite[proof of Theorem 2.3]{R02K:Ulas}:
\begin{align*}
V_{k,n}(t)
&=\left((t+1)t^{n-1}+
 \frac{(2t^{k-1}-t-1)(t^{n-1}-1)}{(t-1)^2}\right)
 \frac{2t^n-t-1}{t-1}\\
&\qquad+\frac{t(2t^{k-1}-t-1)(t^n-1)}{(t-1)^2}.
\end{align*}
Clear the denominator and collect powers of $u$:
\begin{equation}\label{R02K:eq:sparse}
F_{k,n}(t):=(t-1)^3V_{k,n}(t)=C_0(t)+C_1(t)t^{n-1}+C_2(t)t^{2n-2},
\end{equation}
where
\begin{align*}
C_0(t)&=-2t^{k+1}+4t^k+2t^{k-1}+t^3-t^2-3t-1,\\
C_1(t)&=-2t^4+6t^2+4t+2t^{k+2}-2t^{k+1}-6t^k-2t^{k-1},\\
C_2(t)&=2t^4-2t^3-4t^2+4t^k.
\end{align*}
There are at most $7+7+4=18$ nonzero monomials. All exponents are
nonnegative for $k\ge2,n\ge1$. Coincident exponents must be combined;
combination or cancellation can only decrease the support size.

\begin{SourceRacKlemma}[Classical sparse-polynomial bound]\label{R02K:lem:sparse}
A nonzero real polynomial with $s$ nonzero monomials has at most $s-1$
positive roots, counted with multiplicity.
\end{SourceRacKlemma}
\begin{proof}
Write it as $\sum_{j=1}^s a_jx^{e_j}$ with
$0\le e_1<\cdots<e_s$ and all $a_j\ne0$. Divide by $x^{e_1}$ on
$x>0$. For $s=1$ there are no roots. For $s>1$, the derivative of the
quotient has $s-1$ nonzero monomials and, by induction, at most $s-2$
positive roots counted with multiplicity. If the quotient has positive roots
of multiplicities $m_1,\ldots,m_q$, its derivative has at least
$\sum_i(m_i-1)+(q-1)=\sum_i m_i-1$ positive roots counted with
multiplicity, by the local multiplicity rule and Rolle's theorem. Hence
$\sum_i m_i\le s-1$.
\end{proof}

\begin{proof}[Proof of Theorem~\ref{R02K:thm:bound}]
The index $p_{k,n}$ is positive and odd: its minimum at $n=1$ is
$2^{k+1}-3\ge5$. Hence $V_{k,n}(0)=1$ and $V_{k,n}(t)>0$ for $t\ge0$.
All its real zeros are strictly negative. Under $t=-x$, $x>0$,
\[
F_{k,n}(-x)=(-x-1)^3V_{k,n}(-x).
\]
The multiplier is nonzero, so the negative zeros to be counted correspond
exactly to the positive zeros of $F_{k,n}(-x)$. This nonzero polynomial
has at most $18$ monomials and therefore at most $17$ positive roots by
Lemma~\ref{R02K:lem:sparse}. Notice that the extraneous triple zero $t=1$
introduced in \eqref{R02K:eq:sparse} is on the opposite half-line and is never
counted. No assumption about the relative size of $k,n$ was used.
\end{proof}

\subsection{Proof of strict positivity and a source correction}
Fix $n\ge0$, and set
\[
a=\frac{4^n-1}{3},\quad X=B_a(t),\quad Z=B_{a+1}(t),\quad
Y=B_{4a+1}(t)=(t+1)X+Z.
\]
For $n=0$, $X=0,Z=Y=1$ and $B_{h_0}=1$. Suppose $n\ge1$.
Since $4^n=3a+1$ and
$h_n=2^{2n+2}a+2a+1$, Lemma~\ref{R02K:lem:add} gives
\[
B_{h_n}=B_{2^{2n+2}-2a-1}X+B_{2a+1}Z.
\]
The complementary index is $10a+3$. Splitting its odd-index recursion
and applying the minus identity at scale $2^{2n}$ yields
\begin{align*}
 B_{5a+1}&=tB_{2a}+B_{a+1}=t^2X+Z,\\
 B_{5a+2}&=tB_{2a+1}+B_a=(t+1)X+tZ,\\
 B_{10a+3}&=(t^2+t+1)X+(t+1)Z.
\end{align*}
As $B_{2a+1}=X+Z$, we obtain the exact identity
\begin{align}
B_{h_n}&=(t^2+t+1)X^2+(t+2)XZ+Z^2\label{R02K:eq:Q}\\
       &=t^2X^2-tXY+Y^2\label{R02K:eq:F}\\
       &=\left(Y-\frac t2X\right)^2+\frac{3t^2}{4}X^2.\label{R02K:eq:sos}
\end{align}
The first line agrees with an identity in the proof of
\cite[Theorem 3.1, p.~13]{R02K:Ulas}. For the same stated $X,Y$, the next
quadratic form printed in that preprint is
$(t^2+t+1)X^2-(t+1)XY+Y^2$; it is not the substitution of
$Z=Y-(t+1)X$ into \eqref{R02K:eq:Q}. Formula \eqref{R02K:eq:F} is the corrected
substitution, derived above from the defining recursion. At $n=1$ the
recursion gives $B_{19}=1+3t+3t^2$; the printed $W_1=1+t+3t^2$ on
preprint page 14 also disagrees. These observations concern that accessible
version, not an unchecked published correction.

\begin{SourceRacKlemma}[No simultaneous adjacent zeros]\label{R02K:lem:pair}
For any $m\ge0$ and any real $t\ne0$, the pair
$(B_m(t),B_{m+1}(t))$ is not $(0,0)$.
\end{SourceRacKlemma}
\begin{proof}
For the column vector $P_m=(B_m,B_{m+1})^{\mathsf T}$, the recursion gives
\[
P_{2m}=\begin{pmatrix}t&0\\1&1\end{pmatrix}P_m,\qquad
P_{2m+1}=\begin{pmatrix}1&1\\0&t\end{pmatrix}P_m.
\]
Both matrices are invertible when $t\ne0$. Reading the binary digits of
$m$ expresses $P_m$ as their product applied to $P_0=(0,1)^{\mathsf T}$,
so it cannot vanish.
\end{proof}
\begin{proof}[Proof of Theorem~\ref{R02K:thm:positive}]
For $t\ne0$, Lemma~\ref{R02K:lem:pair} shows that $X,Z$ are not both zero.
The transformation $Y=(t+1)X+Z$ is invertible as a transformation of pairs,
so $X,Y$ are not both zero. The positive-definite form in \eqref{R02K:eq:sos}
is therefore strictly positive. At $t=0$, the oddness of $h_n$ gives
$B_{h_n}(0)=1$. This proves the statement for every real $t$, including
$n=0$.
\end{proof}

\subsection{Verification, prior work, and exact scope}
The proofs above are all-parameter arguments. The accompanying
standard-library script is a reproducibility check, not a substitute for
those proofs. It computes Stern pairs directly from their binary indices,
independently of the asserted closed forms.
\begin{center}
\begin{tabular}{@{}p{.68\linewidth}r@{}}
\toprule
Check & Passed\\\midrule
Sparse identity, $2\le k\le60$, $1\le n\le60$ &3540\\
Exact multiplicity three of $t=1$ in the cleared numerator &3540\\
Quadratic and sum-of-squares identities, $0\le n\le100$ &101\\
Formal expansion in independent variables $t,u,v$ &18 terms\\\bottomrule
\end{tabular}
\end{center}
The formal check treats $u=t^{n-1}$ and $v=t^{k-1}$ as independent
indeterminates before substitution; it confirms the universal expansion,
not just sampled parameter values. No numerical root finding is used.
Run \texttt{python3 -I reproduce.py}; the exact output is
\texttt{results.json} and \texttt{results.txt}.

\paragraph{What is old and what is asserted here.}
The Stern recursion, addition formulas, geometric sums, the unsimplified
formula for $V_{k,n}$, and \eqref{R02K:eq:Q} are prior. The sparse-root lemma is
classical. We do not rename any of them as discoveries. The claims of this
note are their explicit all-parameter consequences for Targets A and B,
with the source correction in \eqref{R02K:eq:F}. The bound $17$ is not asserted
to be sharp.

\paragraph{Historical audit.}
The search covered the original preprint, its publisher metadata, targeted
queries by conjecture numbers and root terminology, the later Stern-zero
paper of Altizio~\cite{R02K:Altizio}, and the corpus snapshot
\texttt{dd9c67cb72c5991cce5816de540db8c2132ca343}.
No separately titled earlier resolution was located in the targeted web
search. That absence is not evidence of novelty. More importantly, corpus
[071] already contains the same three-block sparse numerator and explicitly
states: ``Direct square-completion and fewnomial consequences of Ulas's
formulas are likewise not promoted as new methods.'' Corpus [035] also
records the sparse positive-root bound used here. Thus Target A is an
immediate combination of displayed prior ingredients. Target B follows
by square completion of a quadratic form already printed in Ulas.
The corpus warning strongly overlaps precisely these two routes, even
though those consequences are not separately promoted there as theorems.

\paragraph{Tournament qualification.}
\textbf{New-method criterion G5: not satisfied.} The proofs are
self-contained, unconditional, and global, but they introduce no new
methodological principle. We do not certify these targets as genuinely
still open or claim historical originality. This packet is an explicit
mathematical-consequence and overlap audit, not a victory submission.
The independent mathematical review requested by the coordinator remains
separate from both the exact computational checks and the novelty issue.
\endgroup


\clearpage

\section{Finiteness of equal-shift repeated binomial coefficients}
\label{section:R12K}
\noindent{\small\textsf{ROUND 12 \enspace / \enspace Keio \enspace / \enspace R12K}}\\[3pt]
\noindent\textbf{Target.} Jenkins equal-shift extension.

\noindent\textbf{Editorial status.} Complete-range argument preserved; no final tournament adjudication or priority certification. The all-$d\ge2$ equal-shift extension is a separate target. It does not prove Singmaster's uniform multiplicity conjecture.

\smallskip\hrule\medskip



\begingroup
\newcommand{\Z}{\mathbb Z}
\setcounter{equation}{0}
\renewcommand{\theequation}{\thesection.\arabic{equation}}
\renewcommand{\theHequation}{R12K.\arabic{equation}}
\paragraph{Overview.}
Singmaster's uniform multiplicity conjecture is not proved here. We prove that the equal-shift equation
\(\binom nk=\binom{n-d}{k+d}\) has finitely many valid integer solutions for every fixed \(d\ge2\), in fact establishing finiteness for the entire associated algebraic curve over \(\Z\). This addresses the equal-shift extension proposed by Jenkins. The proof uses centered finite products: a quadratic expression on a discrete lattice approaches a rational limit with a nonzero first-order term. An exact algebraic argument also rules out all solutions with \(k\ge1\) for \(d=2,3\).


\subsection{Target and scope}
For an integer \(t>1\), let \(N(t)\) count positions \((n,k)\) with \(1\le k<n\) and \(\binom nk=t\). The round target is Singmaster's 1971 conjecture \(N(t)=O(1)\) \cite{R12K:singmaster1971}. It is not the sharper assertion \(N(t)\le8\). The interior theorem of Matom\"aki, Radziwi\l\l, Shao, Tao and Ter\"av\"ainen bounds the number of representations only in a specified interior region \cite[Theorem 1.3]{R12K:mrstt}. Our argument does not control the complementary near-edge region.

Jenkins considers \(\binom nk=\binom{n-a}{k+b}\) with \(a,b\) fixed. His Theorem 1 treats \(a\ne b\); Section~3 already obtains finiteness for \(a=b=2\) via Siegel's theorem on a smooth quartic. In Section~4, p.~9, he proposes \textquotedblleft to extend the theorem to the case when \(a=b\ne1\)\textquotedblright{} \cite{R12K:jenkins}. Theorem~\ref{R12K:thm:finite} gives a proof of that proposed extension, independently of the algebraic-geometry argument for \(d=2\). The case \(d=1\) is the known infinite Lind--Singmaster family \cite{R12K:lind,R12K:singmaster1975}.

\subsection{The centered equation}
\begin{SourceRbcKtheorem}\label{R12K:thm:finite}
For every fixed integer \(d\ge2\), the equation
\[
\prod_{i=0}^{2d-1}(n-k-i) =
 \prod_{i=0}^{d-1}(n-i)\prod_{i=1}^d(k+i)
\]
has finitely many solutions \((n,k)\in\Z^2\). In particular, this holds for \(\binom nk=\binom{n-d}{k+d}\) in the valid domain \(k\ge0\), \(n\ge k+2d\).
\end{SourceRbcKtheorem}

\begin{proof}
The polynomial equation is obtained by factorial cancellation in the stated binomial domain. For real \(s\), put
\[
P_m(s)=\prod_{j=0}^{m-1}\left(s+\frac{m-1}{2}-j\right).
\]
Set \(x=n-k-(2d-1)/2\), \(y=k+(d+1)/2\), \(z=n-(d-1)/2\). Then the equation and the relation between the centered variables are
\begin{equation}\label{R12K:eq:cent}
 P_{2d}(x)=P_d(y)P_d(z), \qquad z-y=x-\tfrac12 .
\end{equation}
Suppose there are infinitely many integral solutions, and choose a sequence with \(M=\max(|x|,|y|,|z|)\to\infty\). Pass to a subsequence on which \((x,y,z)/M\to(a,b,c)\). The monic leading terms and the linear relation give \(a^{2d}=(bc)^d\) and \(c-b=a\). If any of \(a,b,c\) vanishes, the equations force all three to vanish, contrary to \(\max(|a|,|b|,|c|)=1\). If \(bc<0\), then \(d\) is even and \(bc=-a^2\), and \((b+c)^2=(c-b)^2+4bc\) would be \(-3a^2<0\). Thus \(bc=a^2>0\). On this subsequence the absolute values of \(x,y,z\) are comparable, tend to infinity, and \(t_0:=zy>0\) eventually.

The centered factors occur in opposite pairs (with a central factor \(s\) when \(m\) is odd). Thus
\[
P_m(s)=s^m\left(1-\frac{m(m^2-1)}{24s^2}+O_m(|s|^{-4})\right),
\]
since the coefficient of \(s^{m-2}\) is minus half the sum of the squares of the offsets. Write \(A=(4d^2-1)/12\) and \(B=(d^2-1)/24\). Take the positive \(d\)-th root on each side of \eqref{R12K:eq:cent} and use the first-order expansion of \((1+u)^{1/d}\). Along the subsequence,
\[
 x^2-A+O_d(x^{-2})=
 t_0-B\!\left(\frac{z}{y}+\frac{y}{z}\right)+O(x^{-2}).
\]
On the right, \(t_0\) multiplies \(1-B(y^{-2}+z^{-2})+O(M^{-4})\). Comparability justifies the displayed remainders, with constants allowed to depend on the chosen limiting direction.

Write \(w=x^2-t_0\). The last display first gives \(w=O(1)\), as the ratios are bounded. From \eqref{R12K:eq:cent},
\[
 \frac{z}{y}+\frac{y}{z}
   =2+\frac{(x-1/2)^2}{t_0}
   =3+\frac{w-x+1/4}{t_0}
   =3-\frac1x+O(x^{-2}),
\]
where \(t_0=x^2-w\). It follows that
\begin{equation}\label{R12K:eq:drift}
 w=\frac{5d^2+1}{24}+\frac{d^2-1}{24x}+O(x^{-2}).
\end{equation}
Because all three variables are half-integers, \(w\in\frac14\Z\). A converging sequence in this lattice must eventually equal its limit. But \(d^2-1\ne0\), and multiplying the deviation in \eqref{R12K:eq:drift} by \(x\) gives a nonzero limit, contradicting eventual equality. This also works for \(x<0\). Hence an unbounded sequence of integral solutions is impossible.
\end{proof}

The special case \(d=1\) gives \(P_2(x)=x^2-1/4\) and \(P_1(y)P_1(z)=zy\). There \(w=1/4\) holds exactly: the drift coefficient vanishes. This is consistent with the classical infinite family, not an exception imposed by the argument.

\subsection{Exact exclusions at offsets two and three}
\begin{SourceRbcKtheorem}\label{R12K:thm:exact}
For \(d=2\) or \(3\), the only solution of \(\binom nk=\binom{n-d}{k+d}\) in \(k\ge0,\ n\ge k+2d\) is \((n,k)=(2d,0)\), where both coefficients equal \(1\).
\end{SourceRbcKtheorem}
\begin{proof}
For fixed \(k\ge1\) and \(u=n-k>2d-1\), the original product ratio, left over right, equals
\[
  \frac{1}{\prod_{i=1}^d(k+i)}
  \prod_{i=0}^{d-1}\frac{u-i}{u+k-i}
  \prod_{i=d}^{2d-1}(u-i).
\]
It is positive and strictly increasing in \(u\); its limiting values at the two ends of the interval are \(0\) and \(+\infty\). Thus there is precisely one real root in that interval. In the centered variables, fix \(y=k+(d+1)/2\), set \(J=z-y=x-1/2\), and use \(z=x+y-1/2\). At any trial value \(w_0>0\), there is exactly one real \(x>y\) with \(w=w_0\): the quadratic \(x^2-y(x+y-1/2)\) is strictly increasing on \(x\ge y\), is negative at \(x=y\), and tends to \(+\infty\). Here \(x>y\ge d-1/2\) for \(d=2,3\), hence the trials lie in the domain \(u>2d-1\). They have \(J>0\) and
\[
 t_0=y(y+J),\qquad t_0-J=y^2+(y-1)J\ge y^2.
\]

Let \(F_d=P_{2d}(x)-P_d(y)P_d(z)\). We have \(x^2=t_0+w\), \(z^2+y^2=3t_0+w-J-1/4\). For \(d=2,3\), expansion and these substitutions give the following trial evaluations:
\[
\begin{array}{ccl}
\toprule
d&w_0&F_d\text{ at }w=w_0\\
\midrule
2&1/2&-J/4-3t_0/4-7/16\\
2&1&(t_0-J)/4-13/16\\
3&5/4&-Jt_0-2t_0^2-t_0+5\\
3&9/4&t_0(t_0-J-7)\\
\bottomrule
\end{array}
\]
For example, \(P_4(x)=x^4-5x^2/2+9/16\) and \(P_2(s)=s^2-1/4\). In the sextic case \(P_6(x)=x^6-35x^4/4+259x^2/16-225/64\) and \(P_3(s)=s^3-s\). These yield the table without approximations.

The trial signs in each case are negative and positive, respectively, since \(y\ge5/2\) for \(d=2\), and \(t_0\ge y^2\ge9\) for \(d=3\). Strict monotonicity of the product ratio puts the only real root between the trial values of \(x\), hence between the two listed values of \(w\). For integral input with \(d=2\), \(x,y,z\) are half-integers of the form integer plus \(1/2\), so \(x^2-zy\in\frac12\Z\). For \(d=3\), \(y,z\) are integers and \(x\) is an integer plus \(1/2\), so \(x^2-zy\in\Z+1/4\). The respective open intervals contain no points of those lattices. The case \(k=0\) reduces to \(1=\binom{n-d}{d}\), which holds exactly at \(n=2d\) in the valid domain.
\end{proof}

\subsection{What this does not establish}
Theorem~\ref{R12K:thm:finite} has quantifiers \emph{for each fixed \(d\), finitely many points}. It neither bounds the union over \(d\), nor controls how many distinct offsets can occur from one position. Thus even combining it with Jenkins's unequal-shift theorem leaves Singmaster's conjecture unresolved. That theorem fixes the two shifts, rather than fixing the two lower indices; known finiteness for the latter kind of equation does not cover the equal-shift family. Existing large finite searches \cite[Theorem 1]{R12K:blokhuis} do not replace infinite arguments. Our smaller census through \(10^8\), kept with the replication files, was a control rather than a records claim; it reproduces \(N(3003)=8\) as its unique maximum and provides no uniform bound. Fixed-precision congruences and log-concavity on individual antidiagonals were also insufficient: the former cannot even remove residue classes containing the boundary point \((2d,0)\), and the latter does not compare distant antidiagonals.
\endgroup


\clearpage

\part{Further proof candidates and a reconstructed certificate}

\noindent This part preserves twenty-two additional source manuscripts and one editorial reconstruction. The heading of every section states the target and the archive's review status. Inclusion preserves a full submitted argument; it does not silently upgrade an unadjudicated claim to an independently certified theorem.


\section{An exact Hilbert series for reflected harmonic coefficient spaces}
\label{section:R18H}
\noindent{\small\textsf{ROUND 18 \enspace / \enspace Heisei \enspace / \enspace R18H}}\\[3pt]
\noindent\textbf{Target.} Au Conjecture 4.3.

\noindent\textbf{Editorial status.} Full-range proof candidate; not independently adjudicated at the tournament freeze. 

\smallskip\hrule\medskip



\begingroup
\newcommand{\Q}{\mathbb Q}
\newcommand{\orddeg}{\operatorname{deg}}
\setcounter{equation}{0}
\renewcommand{\theequation}{\thesection.\arabic{equation}}
\renewcommand{\theHequation}{R18H.\arabic{equation}}
\paragraph{Overview.}
Kam Cheong Au's Conjecture 4.3 predicts the dimension in every degree of
a space of coefficient sequences associated with a Wilf--Zeilberger
seed. We prove the exact Hilbert series
$1/((1-T)^2(1-T^2))$. The duplication identity for rising factorials reveals
a reflection symmetry giving the upper bound. For the lower bound, a
filtration by the number of harmonic-sum factors yields independent
leading forms, while discrete residues rule out any algebraic relation
among the shifted harmonic sequences that could invalidate the
filtration after evaluation.


\subsection{The conjecture}
For $n\geq 0$, write $(v)_n=\prod_{j=0}^{n-1}(v+j)$, with empty product $1$.
Au's Example VIII \cite[Conjecture 4.3]{R18H:au} specifies the summand
\begin{equation}
S(a,b,c;n)=
\frac{3^{-2n}(2a-6b-6c-8n-1)
 (\frac12-a-3b)_n(\frac12-a+3b)_n
 (\frac12+a+3b+3c)_{2n}}
 {(1-2a)_{2n}(1+c)_n(1+2b+c)_n}.
\label{R18H:source}
\end{equation}
Let $V_N(S)$ be the $\Q$-span of the sequences
$n\mapsto[a^i b^j c^k]S(a,b,c;n)$ with $i+j+k=N$. This is the
coefficient-function space of \cite[\S2.2]{R18H:au}, not the span of the
corresponding summed numerical constants.

\begin{SourceRbiHtheorem}
For every integer $N\geq 0$,
\[
\dim_{\Q}V_N(S)=\sum_{r=0}^{\lfloor N/2\rfloor}(N-2r+1),
\qquad
\sum_{N\geq0}\dim_{\Q}V_N(S)T^N=\frac{1}{(1-T)^2(1-T^2)}.
\]
\end{SourceRbiHtheorem}

The proof does not require a Wilf--Zeilberger boundary limit or an evaluation
of any of the harmonic series displayed after the conjecture.

\subsection{An injective harmonic-sequence model}
For $\gamma\in\Q\cap(0,1]$ and $r\geq1$, define
\[
h_{\gamma,r}(n)=\sum_{j=0}^{n-1}\frac1{(j+\gamma)^r}.
\]
To evaluate rational functions of $n$ at integers, we use sequences modulo
equality for all sufficiently large $n$. Thus finitely many poles in a
rational coefficient are harmless.

\begin{SourceRbiHlemma}\label{R18H:ind}
For any finite set $I$ of distinct pairs $(\gamma,r)$ as above, no nonzero
polynomial in $\Q(n)[X_{\gamma,r}:(\gamma,r)\in I]$ becomes the zero sequence
under $X_{\gamma,r}\mapsto h_{\gamma,r}(n)$.
\end{SourceRbiHlemma}

\begin{proof}
We first record the rational summability obstruction. Suppose
$R(n)\in\Q(n)$ and $c_{\gamma,r}\in\Q$ satisfy
\[
R(n+1)-R(n)=\sum_{(\gamma,r)\in I}c_{\gamma,r}(n+\gamma)^{-r}.
\]
For each orbit $-\gamma+\mathbb Z$, sum all coefficients of
$(n-\alpha)^{-r}$ in the partial fractions of the left-hand side over
$\alpha$ in that orbit, with missing coefficients set to zero. All but
finitely many terms vanish, and the difference has sum zero since
$n\mapsto n+1$ shifts poles by $-1$ without changing their coefficients.
The orbit and order sum on the right is exactly $c_{\gamma,r}$, since
the $\gamma$'s lie in $(0,1]$. Hence all coefficients are zero.

If a nonzero polynomial relation existed, choose one, $P$, whose leading
monomial in the $X$'s is minimal in graded lexicographic order, and make
its leading coefficient $1$. Put $f_{\gamma,r}(n)=(n+\gamma)^{-r}$
and define $\sigma P=P(n+1,X+f(n))$. The relations form a set stable
under this shift by the recurrence for $h_{\gamma,r}$. The leading
monomial and its coefficient are unchanged by the shift. Therefore
$\sigma P-P$, being a relation with smaller leading monomial if nonzero,
is identically zero as a polynomial.

Let $D$ be the degree in $X$ of $P$; $D=0$ is impossible since a nonzero
rational function cannot vanish at all large integers. In
$\sigma P=P$ compare the homogeneous part $P_D$ of degree $D$. Its
coefficients are rational functions fixed by $n\mapsto n+1$, hence
constants (a periodic rational function has no finite poles and a
periodic polynomial is constant). Comparing degree $D-1$ gives
\[
P_{D-1}(n+1,X)-P_{D-1}(n,X)
  =-\sum_{(\gamma,r)\in I}f_{\gamma,r}(n)\,\partial_{X_{\gamma,r}}P_D(X).
\]
Since $P_D$ has positive degree in characteristic zero, some monomial on
the right has a nontrivial constant coefficient vector in the $f$'s.
The corresponding coefficient on the left violates the rational
summability obstruction.
\end{proof}

This elementary proof is a special case of Schneider's parameterized
telescoping criterion for algebraic independence of sums \cite{R18H:schneider}.

\subsection{Reflection symmetry and leading forms}
Change variables by
\[
x=-a+3b,\quad y=-a-3b,\quad z=\frac{a+3b+3c}{2},
\qquad s=\frac{x+y}{2},\quad d=\frac{x-y}{2}.
\]
This is invertible over $\Q$ and preserves coefficient spans in each
total degree, because on degree-$N$ forms the substitution is an invertible
linear map over $\Q$. Divide the whole summand by $S(0,0,0;n)$: this
sequence never vanishes and multiplication by its reciprocal preserves
ranks of sequence families.

Write $A_\gamma(v)=(\gamma+v)_n/(\gamma)_n$. By
$(v)_{2n}=2^{2n}(v/2)_n((v+1)/2)_n$, the normalized summand is
\begin{align}
F(s,z,d;n)={}&\left(1+\frac{4(s+z)}{1+8n}\right)
 \frac{A_{1/2}(s+d) A_{1/2}(s-d) A_{1/4}(z) A_{3/4}(z)}
 {A_{1/2}(s) A_1(s)
  A_1((2z+s+d)/3) A_1((2z+s-d)/3)}.
\label{R18H:factor}
\end{align}
The rational factor and every paired product are unchanged by $d\mapsto-d$.
Put $u=d^2$; then $s,z,u$ have weights $1,1,2$ in the original degree.
Consequently the only possibly nonzero coefficients in degree $N$
are indexed by $i+j+2r=N$, an upper bound for the dimension.

For the lower bound, use the formal identity
\[
\log A_\gamma(v)=\sum_{m\geq1}
 \frac{(-1)^{m-1}}m h_{\gamma,m}(n) v^m.
\]
In this calculation temporarily replace $h_{\gamma,m}$ by independent
indeterminates $X_{\gamma,m}$. Only $m\leq N$ matters in weight $N$.
Let $E$ be the formal logarithm of the Pochhammer quotient in
\eqref{R18H:factor}. All its coefficients are linear in the $X$'s, and it has
zero constant term. Its ordinary degree-one part in $s,z,u$ (giving $u$
ordinary degree one for this purpose) is $L_s s+L_z z+M u$ where
\begin{align*}
L_s &= X_{1/2,1}-\tfrac53 X_{1,1},\
L_z &= X_{1/4,1}+X_{3/4,1}-\tfrac43 X_{1,1},\
M   &= -X_{1/2,2}+\tfrac19 X_{1,2}.
\end{align*}
For example, the $d^2$ coefficient only comes from the quadratic
log terms of the two half-integer factors in the numerator and the two
shifted integer factors in the denominator.

If $i+j+2r=N$, consider the polynomial
\[
C_{i,j,r}(n,X)=[s^i z^j u^r]\left(1+\frac{4(s+z)}{1+8n}\right)e^E.
\]
It has degree in $X$ at most $i+j+r$. Indeed, each factor of $E$ uses
at least one unit of ordinary degree in the parameters, and supplies
one $X$ factor. For equality one must use the constant part of the
prefactor, and only ordinary degree-one terms of $E$. Thus the top part is
exactly
\begin{equation}
\frac{L_s^i L_z^j M^r}{i!j!r!}.
\label{R18H:leading}
\end{equation}
The displayed forms are independent linear forms, so the polynomial
monomials in them are linearly independent over $\Q(n)$.

Suppose a nontrivial $\Q$-linear combination of the degree-$N$ coefficient
sequences of $F$ vanishes. Among the terms with nonzero coefficient take
the minimum $r$. Those terms have the strictly highest $X$-degree,
namely $N-r$. Their top parts \eqref{R18H:leading} cannot cancel, so the
corresponding combination is a nonzero polynomial in the
harmonic indeterminates over $\Q(n)$. Lemma~\ref{R18H:ind} prevents this
polynomial from vanishing on the harmonic sequences, a contradiction.
This argument includes $N=0$, for which $F$ has constant term $1$.
It even gives independence of the selected coefficients over $\Q(n)$
in the eventual-sequence model. Counting the indices $i+j+2r=N$
now proves the theorem. \hfill$\square$

\subsection{Reproducibility and scope}
The companion \texttt{check.py} constructs the truncated power series
of the normalized version of \eqref{R18H:source} from the original variables,
by successive rational factor multiplication over the finite field
with $1{,}000{,}000{,}007$ elements. It separately computes the normalized
series in \eqref{R18H:factor}, checks the reflection, and compares exact
modular ranks in the two coordinates. The output \texttt{results.txt}
is only a finite cross-check, not a premise of the all-degree proof.

Conjecture 4.3 asserts a Hilbert series for coefficient \emph{sequences}.
It does not assert independence of the numerical values of their
associated infinite sums; the proof makes no claim about such
independence. It establishes the full conjecture as stated in version 4
of Au's preprint. A targeted search on 19 November 2026 found no prior
proof of this specific conjecture. This is not a priority claim beyond
the sources checked.
\endgroup


\clearpage

\section{Cyclotomic supercongruences and prime-power specializations}
\label{section:R18T}
\noindent{\small\textsf{ROUND 18 \enspace / \enspace Taisho \enspace / \enspace R18T}}\\[3pt]
\noindent\textbf{Target.} Guo--Schlosser Conjectures 6.3 and 6.1.

\noindent\textbf{Editorial status.} Full-range proof candidate; not independently adjudicated at the tournament freeze. Both the cyclotomic and prime-power assertions are reproduced as candidates. Arithmetic regression tests do not certify the all-level Taylor-jet interpolation argument.

\smallskip\hrule\medskip



\begingroup
\newcommand{\Zp}{\mathbb Z_p}
\newcommand{\ord}{\operatorname{ord}}
\newcommand{\cut}{\delta}
\setcounter{equation}{0}
\renewcommand{\theequation}{\thesection.\arabic{equation}}
\renewcommand{\theHequation}{R18T.\arabic{equation}}
\paragraph{Overview.}
We record a proof candidate for Guo and Schlosser's Conjectures 6.3
and 6.1. The argument interpolates each of three local Taylor
coefficients independently at terminating very-well-poised
specializations. The degree bound comes from cancellation of a
quadratic parameter moment under a residue reflection. This is a
frozen, not independently reviewed candidate, not an adjudicated
tournament result.


\subsection{Exact target and notation}
Let \((a;q)_k=\prod_{j=0}^{k-1}(1-aq^j)\) and \([v]_q=(1-q^v)/(1-q)\).
For \(A\equiv1\pmod3\) and \(\cut\in\{1,3\}\) put
\[
 F_{A,\cut}(q,a)=\sum_{k=0}^{(A-1)/\cut}
 [6k+1]_q\frac{(aq;q^3)_k(q/a;q^3)_k(q;q^3)_k}
 {(aq^3;q^3)_k(q^3/a;q^3)_k(q^3;q^3)_k},
 \quad R_A(q)=\frac{(q^3;q^3)_{(A-1)/3}}{(q^2;q^3)_{(A-1)/3}}.
\]
For \(n>1\) and \(L\equiv1\pmod3\) with \(n\equiv1\pmod3\), define
\[
 K_{nL}(q)=q^{2(1-n)/3}[n]_q\frac{R_{nL}(q)}{R_L(q^n)}.
\]
\begin{SourceRbiTtheorem}[Candidate: Guo--Schlosser, Conjecture 6.3]
For every \(n>1\) with \(n\equiv1\pmod3\), every \(r\ge1\), and both
\(\cut=1,3\), one has
\[
 F_{n^r,\cut}(q,1)-K_{n^r}(q)F_{n^{r-1},\cut}(q^n,1)
 \equiv0\pmod{\prod_{i=1}^r\Phi_{n^i}(q)^3}.
\]
Congruences of rational functions are understood in the localization
at the indicated cyclotomic polynomials.
\end{SourceRbiTtheorem}
This is the exact modulus and normalization in \cite[Conjecture 6.3]{R18T:GS}.
The cited paper proves the one-level cyclotomic statement (Theorem 1.2),
and separately gives an all-level lower-order divisibility (equation (6.1)).
Its signed companion (6.3) does not provide the unsigned congruence here.

\subsection{Local normalization to first order}
Fix \(s=n^i\), \(1\le i\le r\), and let \(\zeta\) be a primitive \(s\)-th root.
Use \(q=\zeta\exp(t)\). For any positive \(N\equiv1\pmod3\) divisible by \(s\),
\[
 J_N(q):=\frac{R_N(q)}{R_{N/n}(q^n)}
   =\frac{\prod_{k\le (N-1)/3,\ n\nmid 3k}(1-q^{3k})}
          {\prod_{k\le (N-1)/3,\ n\nmid (3k-1)}(1-q^{3k-1})}.
\]
Indeed the omitted numerator exponents are \(3nh\) and the omitted
denominator exponents are \(n(3h-1)\), both for
\(1\le h\le (N/n-1)/3\). Thus \(J_N\) is a unit near \(\zeta\).
We claim that \(J_N/J_{N'}=1+O(t^2)\) for any two such \(N,N'\).

It suffices to compare \(N+3s\) to \(N\).
The added numerator exponents are \(N+2+3h\) and the added denominator
exponents \(N+1+3h\), \(0\le h<s\), in each case omitting multiples of \(n\).
Match numerator index \(h\) to denominator index
\(h+(2s+1)/3\pmod s\). The denominator exponent minus the numerator
exponent is \(2s\) if \(h<b_*=(s-1)/3\), and \(-s\) otherwise.
The value of the quotient at \(t=0\) is 1. Write
\(\varphi(w)=-w/(1-w)\), so that
\(\varphi(w)+\varphi(w^{-1})=1\) for \(w\ne1\).
The non-omitted residues in the full range are stable under negation
and number \(s-s/n\). The corresponding residues for \(h<b_*\)
are \(2,5,\dots,s-2\), excluding multiples of \(n\); they are also
stable under negation and number \((s-s/n)/3\).
The logarithmic derivative of the quotient is therefore
\[
 s\left(\frac{s-s/n}{2}-3\frac{s-s/n}{6}\right)=0.
\]
This proves the claim. In particular
\begin{equation}\label{R18T:eq:K}
 K_N(q)-K_{n^r}(q)=
 \begin{cases}O(t^3)&i=1,\\ O(t^2)&i\ge2,\end{cases}
\end{equation}
since \([n]_q\) has a simple zero for \(i=1\), and is a unit otherwise.

\subsection{A degree bound for the three jets}
Now put \(a=\exp(sty)\), treating \(y\) as an indeterminate.
Let
\[
 U=\frac{n^{r-i}-1}{\cut},\qquad
 D_U(y)=\prod_{h=1}^{U}(9h^2-y^2).
\]
Consider the Taylor coefficients through degree two in \(t\)
of \(F_{n^r,\cut}(q,a)\) and \(F_{n^{r-1},\cut}(q^n,a)\).
We show that multiplying each by \(D_U\) gives a polynomial in
\(y\) of degree at most \(2U\). We also show that the constant
coefficient vanishes on any side whose base has order greater than 1.

Here are the details for a base \(Q=\omega\exp(\lambda t)\) of order
\(m\) at \(t=0\), with \(\lambda m=s\). Use
\((Q,m,\lambda)=(q,s,1)\) or \((q^n,s/n,n)\).
Write \(k=mu+b\), \(0\le b<m\), and let \(c=(m-1)/3\).
Both cutoffs have complete groups with \(0\le u<U\), followed by
\(0\le b\le c\) for \(\cut=3\) or all \(0\le b<m\) for \(\cut=1\).
The factors with vanishing root-of-unity constant term from the complete
groups contribute the same ratio
\begin{equation}\label{R18T:eq:B}
 B_u(t,y)=\prod_{h=0}^{u-1}
 \frac{(1-e^{s(3h+1)t})(1-e^{s(3h+1+y)t})(1-e^{s(3h+1-y)t})}
 {(1-e^{s(3h+3)t})(1-e^{s(3h+3+y)t})(1-e^{s(3h+3-y)t})}.
\end{equation}
For \(b>c\) there are three additional zeros in the numerator and
none in the denominator, so these terms are \(O(t^3)\) over
\(\mathbb C(y)\). For \(m=1\) only \(b=0\) occurs, and the
\(Q\)-integer is taken by its removable value.

The constant factor \(B_u(0,y)\) given by the ratio of the linear
coefficients in \eqref{R18T:eq:B} is a rational function with numerator
degree \(2u\) and denominator, up to a scalar, \(D_u(y)\).
Moreover,
\[
 B_u(t,y)/B_u(0,y)=1+\beta_{u,1}t+\beta_{u,2}t^2+O(t^3)
\]
with \(\beta_{u,1},\beta_{u,2}\) independent of \(y\).
To see this, use \(\log((e^z-1)/z)=z/2+z^2/24+O(z^3)\);
the opposite shifts cancel linearly, and the quadratic \(y^2\)
contributions cancel between numerator and denominator since each
contains \(u\) opposite-shift pairs.

Now take \(m>1\) and \(b\le c\). After factoring out \(B_u\),
the nonvanishing factors' constant value, with the weight, is
\[
 W_b=[6b+1]_\omega \rho_b^3,\qquad
 \rho_b=\frac{(\omega;\omega^3)_b}{(\omega^3;\omega^3)_b}.
\]
Every completed group has identical numerator and denominator
nonzero residue multisets, so only the remainder contributes here.
The remaining factors are independent of \(y\) through degree one
in \(t\). Their only possible additional term in degree two is
\(s^2 y^2 W_b H_b\), where
\[
 H_b=\sum_{h=0}^{b-1}\psi(\omega^{1+3h})
       -\sum_{h=1}^b\psi(\omega^{3h}),\qquad
 \psi(w)=-\frac{w}{(1-w)^2}.
\]
Indeed the coefficient follows by twice differentiating
\(\log(1-we^z)\), and the completed residue groups again cancel.
This description of the \(y^2\) coefficient remains valid when
\(W_b=0\), because the weight itself is independent of \(y\).
Since \(m-(1+3h)=3(c-h)\) and \(\psi(w)=\psi(w^{-1})\),
we have \(H_c=0\) and \(H_{c-b}=H_b\).

The same reflection gives
\[
 \rho_{c-b}=\rho_c\omega^{2b}\rho_b,\qquad
 \rho_c=(-1)^c\omega^{-3c(c+1)/2},\qquad
 \rho_c^3=\omega .
\]
For the last equality note that
\(-9c(c+1)/2-1=-m(m+1)/2\); the resulting phase is \(1\) for
odd \(m\) and \(-1\) for even \(m\), exactly \((-1)^c\).
The weights satisfy
\([6(c-b)+1]_\omega=-\omega^{-(6b+1)}[6b+1]_\omega\).
Hence \(W_{c-b}=-W_b\), and both \(\sum_{b=0}^c W_b\) and
\(\sum_{b=0}^c W_bH_b\) vanish. For each \(u\), the coefficient
through degree two is therefore \(B_u(0,y)\) times a scalar
independent of \(y\), and the constant term is zero.
For \(m=1\), that same degree conclusion holds by \eqref{R18T:eq:B}
and the \(y\)-independent weight, with no assertion of constant
vanishing. Multiplication by \(D_U\) proves the claimed polynomial bound.

\subsection{Interpolation and descent}
Let \(\mathcal E(t,y)\) denote
\[
 F_{n^r,\cut}(q,a)-K_{n^r}(q)F_{n^{r-1},\cut}(q^n,a),
 \qquad q=\zeta e^t,\quad a=e^{sty}.
\]
By the degree bound, \([t^h]D_U(y)\mathcal E(t,y)\) is a polynomial
of degree at most \(2U\) for \(h=0,1,2\).
For \(0\le j\le U\) set \(y=\pm(3j+1)\) and \(N=s(3j+1)\).
The parameter is then \(a=q^{\pm N}\). Both sums terminate:
the first at \((N-1)/3=sj+(s-1)/3\), the second at
\((N/n-1)/3\), within their stated cutoffs for either \(\cut\).
All denominators at this substitution are nonzero rational functions
because \(N,N/n\equiv1\pmod3\).
The terminating very-well-poised \({}_6\phi_5\) identity,
in the substitution already used by \cite[proof of Theorem 1.2]{R18T:GS},
evaluates the first sum as
\(q^{2(1-N)/3}[N]_qR_N(q)\), and the second by replacing \(q,N\)
with \(q^n,N/n\). Their ratio is \(K_N(q)\).
Consequently \(\mathcal E(t,y)=O(t^3)\) at each such node:
for \(i=1\) use \eqref{R18T:eq:K} and regularity of the second sum;
for \(i\ge2\) the second sum also has zero constant coefficient.
No node is a pole of any jet since \(D_U(\pm(3j+1))\ne0\).

There are \(2U+2\) distinct nodes, so each of the three polynomials
vanishes identically. In particular \(\mathcal E(t,0)=O(t^3)\).
Since \(q-\zeta\) has order one in \(t\) and this holds for every
primitive \(n^i\)-th root, the desired congruence follows.

\begin{SourceRbiTtheorem}[Candidate: Guo--Schlosser, Conjecture 6.1]
For every prime \(p\equiv1\pmod3\), every \(r\ge1\), and both
\(\cut=1,3\),
\[
 \sum_{k=0}^{(p^r-1)/\cut}(6k+1)\frac{(1/3)_k^3}{k!^3}
 \equiv p\Gamma_p(2/3)^3
 \sum_{k=0}^{(p^{r-1}-1)/\cut}(6k+1)\frac{(1/3)_k^3}{k!^3}
 \pmod{p^{3r}}.
\]
\end{SourceRbiTtheorem}
\begin{proof}
Take \(n=p\) and then \(q\to1\). For each summand, replacing
factors \(1-q^e\) by \([e]_q\) removes all factors at \(q=1\).
The denominator after cyclotomic cancellation evaluates to a
\(p\)-adic unit: factors \(\Phi_{p^j}\) for \(1\le j\le r\)
do not occur in the reduced denominator, by the same residue
count leading to \eqref{R18T:eq:B}; factors with \(j>r\) cannot occur
since all denominator exponents are below \(3p^r<p^{r+1}\).
The smaller sum with base \(q^p\) uses the same count, with the
case \(j=1\) removable by equal numbers of factors \(1-q^p\).
The product defining \(J_{p^r}\) has denominator exponents prime
to \(p\). Monic cyclotomic factors can be canceled over
\(\mathbb Z[q]\); the only monomial factor in a denominator is a
power of \(q\). Thus specializing the proven divisibility gives
modulus \(\prod_{j=1}^r\Phi_{p^j}(1)^3=p^{3r}\).

Put \(X=p^r/3\). Removing the multiples of \(p\) from the two
rising products and using the Morita functional equation gives
\[
 J_{p^r}(1)=
 \frac{\Gamma_p(2/3+X)\Gamma_p(2/3)}
 {\Gamma_p(1)\Gamma_p(1/3+X)}
 =\Gamma_p(2/3+X)\Gamma_p(2/3-X)\Gamma_p(2/3).
\]
In the first equality, the powers of 3 in the numerator and
denominator cancel because each has the same number of unit
factors. In the second equality use reflection, with
\(a_0(1/3+X)=(2p+1)/3\) odd, and \(\Gamma_p(1)=-1\).
Taylor expansion with integral first-order remainder
\cite[Theorem 14]{R18T:LR} gives
\(J_{p^r}(1)\equiv\Gamma_p(2/3)^3\pmod{p^{2r}}\).
The smaller undeformed sum is divisible by \(p^{r-1}\) for both
cutoffs by \cite[(6.2)]{R18T:GS}, with the convention that at level zero
it is 1. Since \(K_{p^r}(1)=pJ_{p^r}(1)\), substituting the claimed
Gamma factor does not change the congruence modulo \(p^{3r}\).
\end{proof}

\subsection*{Verification status}
Candidate derived at 04:53:58 UTC on 19 November 2026, before notice
that another candidate had passed independent review. This checkpoint
contains the entire proposed argument; no specific unproved sublemma
is presently identified by the proposer. It has not received
independent verification. In particular, a reviewer should verify
the common-denominator degree bound, the normalization comparison
at all terminating nodes, and the arithmetic specialization
rather than relying on numerical checks.
\endgroup


\clearpage

\section{Log-concavity of restricted Stirling rows and restricted Bell roots}
\label{section:R18R}
\noindent{\small\textsf{ROUND 18 \enspace / \enspace Reiwa \enspace / \enspace R18R}}\\[3pt]
\noindent\textbf{Target.} Moll--Ramirez--Villamizar 6.3; Arnoczki--Nyul 3.4.

\noindent\textbf{Editorial status.} Full-range candidate; the separate restricted Bell-root assertion is identified as a classical consequence. Both original claims are retained. The restricted Stirling row claim is distinct from the restricted $r$-Bell root claim, whose classical provenance is stated explicitly.

\smallskip\hrule\medskip



\begingroup
\newcommand{\lc}{\text{log-concave}}
\setcounter{equation}{0}
\renewcommand{\theequation}{\thesection.\arabic{equation}}
\renewcommand{\theHequation}{R18R.\arabic{equation}}
\paragraph{Overview.}
We give two candidate proofs: one for the row log-concavity of restricted
Stirling numbers (Moll--Ram\'{i}rez--Villamizar, Conjecture 6.3), and
one for the simplicity of the negative zeros of restricted
$r$-Bell polynomials (Arn\'{o}czki--Nyul, Conjecture 3.4).
These are frozen research claims without independent review or a priority verdict.
The second assertion also follows from a classical matching-polynomial theorem.


\subsection{Restricted Stirling rows}
Let $S_m(n,k)$ be the number of partitions of an $n$-element labelled
set into $k$ nonempty blocks, each of size at most $m$.
Let $S_m(0,0)=1$ and use zero for impossible choices. The exact target is
Conjecture~6.3 of \cite{R18R:MRV}, which asserts log-concavity in $k$ for every
$n\ge 0$ and $m\ge 1$. The argument below gives a quantitative strengthening.

\begin{SourceRbiRtheorem}
For every $n\ge 0$, $m\ge 1$ and $k\ge 1$,
\[
 k\,S_m(n,k)^2\ge (k+1)S_m(n,k-1)S_m(n,k+1).
\]
Moreover, the sequence $k!S_m(n,k)$ is log-concave in $k$.
\end{SourceRbiRtheorem}

\begin{proof}
For any nonnegative log-concave sequence $b$ with interval support,
\begin{equation}\label{R18R:near}
 b(u-s)b(u-t)\ge b(u)b(u-s-t)\qquad(s,t\ge0).
\end{equation}
If the right product vanishes this is immediate. Otherwise all indices
lie between two positive endpoints. Products of the decreasing successive
ratios $b(v)/b(v-1)$ show that
$b(u-t)/b(u-s-t)\ge b(u)/b(u-s)$, proving \eqref{R18R:near}.

Write $A(z)=\sum_{j=1}^{m} z^j/j!$, $a_j=[z^j]A(z)$ and
$F_k(u)=[z^u]A(z)^k$. The sequence $a$ is log-concave with
interval support, so each coefficient sequence $F_j$ is too, by ordinary
convolution closure \cite{R18R:Menon}. For completeness, this closure follows
because the Toeplitz matrix $T_a=(a(i-j))_{i,j\in\mathbb Z}$ has
nonnegative $2\times 2$ minors by the same ratio balancing as
\eqref{R18R:near}; convolution is multiplication of these matrices and
Cauchy--Binet gives nonnegative $2\times2$ minors for the product.
Supports add as intervals. All sums here are finite; the convolution
identity starts at $F_0=\delta_0$.

For $b(u)=F_{k-1}(u)$ one now has the exact expansion
\begin{align*}
F_k(u)^2 - F_{k-1}(u)F_{k+1}(u)
 &=\sum_{s,t\ge0}a_s a_t
    \bigl(b(u-s)b(u-t)-b(u)b(u-s-t)\bigr)\ge0.
\end{align*}
The expansion uses one convolution in each factor on the positive side
and two convolutions in the second factor on the negative side. Thus
this proves log-concavity in the \emph{number} of convolutions, not just
in the coefficient index.

Finally, $F_k(n)=k!S_m(n,k)/n!$, because labelling the $k$ blocks
in order gives the multinomial coefficients $n!/(j_1!\cdots j_k!)$.
Substitution and cancellation yield the stated factor $(k+1)/k$.
In the nonzero support the requested ordinary log-concavity follows,
and the support is precisely $\lceil n/m\rceil\le k\le n$ for $n>0$.
\end{proof}

The same argument works with any nonnegative gapless log-concave block
coefficient sequence on nonnegative indices, and with any additional
log-concave factor $G(z)$ of finite support and nonnegative indices:
replace $F_k(u)$ by $[z^u]G(z)A(z)^k$. In particular, a version with $r$
distinguished labels constrained to different blocks of size at most $m$
uses $G(z)=(\sum_{j=0}^{m-1}z^j/j!)^r$. This counts the ordinary labels
placed in distinguished blocks while retaining the same normalization
for the $k$ nondistinguished blocks.

\subsection{Simple roots of restricted $r$-Bell polynomials}
Let $B_{n,r}$ be the restricted $r$-Bell polynomial of \cite{R18R:AN}:
the coefficient of $x^k$ counts partitions of $n+r$ labels into $k+r$
blocks of size at most $2$, with $r$ distinguished labels in different blocks.
Theorem~3.1 of \cite{R18R:AN} gives
\[
B_{0,r}=1,\qquad B_{n+1,r}=(x+r-n)B_{n,r}+2xB_{n,r}'.
\]
Conjecture~3.4 asks for all roots other than zero to be simple and negative,
with zero multiplicity $e_{n,r}=\max(0,\lceil(n-r)/2\rceil)$.
Theorem~3.5 of \cite{R18R:AN} advances the conjecture only at one parity.

\begin{SourceRbiRtheorem}
For all $n,r\ge0$, $B_{n,r}=x^{e_{n,r}}H_{n,r}(x)$, where $H_{n,r}$
is monic and all its roots are simple and negative.
\end{SourceRbiRtheorem}
\begin{proof}
Fix $r$ and induct on $n$, starting with $H_{0,r}=1$.
Put $e=e_{n,r}$ and $c=r-n+2e$. Then $c\ge0$, and direct examination
of the ceiling shows that $e_{n+1,r}=e+1$ if $c=0$ and $e_{n+1,r}=e$
if $c>0$. Substitute $B_{n,r}=x^e H$ into the displayed recurrence.
After removing $x^{e_{n+1,r}}$, the next polynomial is
\[
   H_{n+1,r}=
   \begin{cases}
     (x+c)H+2xH',&c>0,\\
     H+2H',&c=0.
   \end{cases}
\]
When $H=1$, the outputs are $x+c$ and $1$, respectively. Otherwise write
$\rho_1<\cdots<\rho_d<0$ for the roots of monic $H$;
$\operatorname{sgn}H'(\rho_i)=(-1)^{d-i}$.

For $c>0$ the output is monic of degree $d+1$ and its sign at $\rho_i$
is $(-1)^{d-i+1}$. There is a sign change on each successive-root interval,
on $(-\infty,\rho_1)$ by the leading term, and on $(\rho_d,0)$ because
the value at zero is $cH(0)>0$. These $d+1$ distinct roots exhaust the degree.
For $c=0$ the output is monic of degree $d$ and the sign at $\rho_i$ is
$(-1)^{d-i}$. Each successive-root interval and the leftmost interval
$(-\infty,\rho_1)$ contains one root. These $d$ roots exhaust the degree.
In both cases no root can be zero or multiple, completing the induction.
\end{proof}

\paragraph{Priority caveat.}
There is also a direct classical implication. The graph in \cite[Thm.~3.2]{R18R:AN}
is the join of an $n$-vertex clique with an $r$-vertex independent set.
Deleting a path traversing the clique leaves only the independent set,
whose matching polynomial has no nonzero root. Godsil's path deletion
inequality \cite{R18R:Godsil} therefore forces all nonzero matching roots of the
original graph to have multiplicity at most one. Under the relation
$\mu(G,y)=y^{r-n}(-1)^n B_{n,r}(-y^2)$ for $y\ne0$, this immediately implies
the claimed simplicity. Consequently only the alternate deflation/interlacing
argument is offered here, not a claim that simplicity was absent as a consequence
of earlier theorems.
\endgroup


\clearpage

\section{Coefficientwise Turan inequalities for cubic Franel polynomials}
\label{section:R19M}
\noindent{\small\textsf{ROUND 19 \enspace / \enspace Meiji \enspace / \enspace R19M}}\\[3pt]
\noindent\textbf{Target.} Sun cubic Franel assertion in 1.1(i).

\noindent\textbf{Editorial status.} Full-range proof candidate; not independently adjudicated at the tournament freeze. Only the cubic Franel assertion is targeted. The sequence starts at index $8$, so its first three-term center is $9$; the preceding center fails. The polynomial certificate recipe is retained.

\smallskip\hrule\medskip



\begingroup
\newcommand{\supp}{\operatorname{supp}}
\setcounter{equation}{0}
\renewcommand{\theequation}{\thesection.\arabic{equation}}
\renewcommand{\theHequation}{R19M.\arabic{equation}}
\paragraph{Overview.}
We give a proof candidate for the cubic Franel assertion in Z.-W. Sun's
Conjecture 1.1(i): the sequence of polynomials
$F_n(q)=\sum_{j=0}^n\binom nj^3q^j$, starting at $n=8$, is
coefficientwise log-convex. In fact, all interior coefficients of the
Tur\'an difference are strictly positive and the boundary index is sharp.
The argument cancels a constant term using exact one-step balance,
dominates two size-biased moments by lattice Gaussian moments,
and finishes with a finite polynomial coefficient certificate
valid throughout the infinite parameter domain.


\begin{SourceRbjMtheorem}\label{R19M:thm:main}
For all $n\ge9$ and $1\le k\le 2n-1$,
\[
  [q^k](F_{n+1}F_{n-1}-F_n^2)>0.
\]
The coefficients at $k=0,2n$ are zero. The $q^2$ coefficient at
$n=8$ is $-84$.
\end{SourceRbjMtheorem}
Sun independently posed exactly the cubic Franel assertion in 2024
\cite{R19M:MO} and included it as the first assertion of Conjecture 1.1(i)
in his 2026 paper \cite{R19M:Sun}. Our result concerns only this assertion,
not the separately stated rising-binomial or double-sum predictions.

\subsection{One-step balance}
Write $D_{n,k}=[q^k](F_{n+1}F_{n-1}-F_n^2)$.
Reciprocity of each $F_n$ gives $D_{n,k}=D_{n,2n-k}$ and
$D_{n,0}=0$. Direct expansion gives
\[
 D_{n,1}=6n,\quad D_{n,2}=\frac34 n(n-1)(n^2-9n+6),\quad
 D_{n,3}=\frac{n(n-1)^3}{6}(2n^3+n^2+20n-24).
\]
These are positive for $n\ge9$ and give the stated failure at $n=8$.
Hence assume $4\le k\le n$. Put $N=n+1$, $h=2N-k$ and, for
$0\le j\le k$, set
\[
 \ell=k-j,\quad d=j-\ell,\quad
 A=N-j,\quad B=N-\ell,\qquad
 w(d)=\left(\binom Nj\binom N\ell\right)^3.
\]
Extend $w$ by zero to the remaining integers of the same parity as $k$.
It is an even weight. Coefficient extraction using
$\binom{N-1}{j}=(A/N)\binom Nj$ and
$\binom{N-2}{j}=[A(A-1)/(N(N-1))]\binom Nj$ yields
\begin{align}
 N^6(N-1)^3 D_{n,k}&=\sum_d w(d)P(d),\
 P(d)&=\frac{N^3}{2}\bigl((A(A-1))^3+(B(B-1))^3\bigr)
                   -(N-1)^3(AB)^3.\label{R19M:eq:P}
\end{align}
Next, the adjacent ratios give
$w(d)(A\ell)^3=w(d+2)((j+1)(B+1))^3$ for $j<k$.
Both boundary terms in the following summation vanish
($\ell=0$ at the right, $j=0$ at the left), so
\[
 \sum_d w(d)\bigl((A\ell)^3(d+1)-(jB)^3(d-1)\bigr)=0.
\]
For $c=(kh+d^2)/4$ and $t=Nd/2$, we have $A\ell=c-t$ and $jB=c+t$.
Half the bracket is therefore
\[
 H(d)=c^3+3ct^2-d(3c^2t+t^3).
\]
As $H(0)=(kh/4)^3>0$, set
\begin{equation}\label{R19M:eq:C}
 C(d)=P(d)-\frac{P(0)}{H(0)}H(d)=a d^2+b d^4+e d^6.
\end{equation}
The $w$-sum of $C$ equals that of $P$. This replacement removes
the constant exactly, so only two ratios of even moments are needed.

\subsection{Two size-biased moment bounds}
Set
\[
 M=\frac{(k+1)(h+1)}{2N+2}.
\]
It is at least $85/22$: for fixed $N$ its numerator increases over
$4\le k\le N-1$, and at $k=4$ the quotient increases with $N\ge10$.
For $0\le d<k$ on the parity lattice,
\[
 \frac{w(d+2)}{w(d)}
 =\left(\frac{k-d}{k+d+2}\frac{h-d}{h+d+2}\right)^3
 \le \exp[-6(d+1)/M].
\]
Here $\log((1-v)/(1+v))\le-2v$ for $0\le v<1$,
applied with $v=(d+1)/(k+1)$ and $(d+1)/(h+1)$.
The right side is the ratio of the weights
$g(d)=\exp(-3d^2/(2M))$. Thus $w(d)/g(d)$ is
nonincreasing on the positive half-lattice (including the zero tail).
The ratio is unchanged by multiplying both weights by $d^2$.
Under the probability with weights $d^2g(d)$, a nonincreasing function
has nonpositive covariance with an increasing one: this follows by
averaging $(f(x)-f(y))(u(x)-u(y))\le0$ over independent $x,y$.
Apply this to $u=w/g$ and $f(d)=d^2,d^4$.
Evenness of the weights and vanishing of the $d=0$ contribution give
\begin{equation}\label{R19M:eq:dom}
 \frac{\sum w(d)d^{2r}}{\sum w(d)d^2}
 \le \frac{S_r}{S_1},\qquad r=2,3,\qquad
 S_r=\sum_{d\equiv k\, (2)} d^{2r}g(d).
\end{equation}

\begin{SourceRbjMlemma}\label{R19M:lem:gauss}
For $M\ge85/22$, on either parity lattice,
$S_2/S_1<M+3/4$ and $S_3/S_1<2M^2$.
\end{SourceRbjMlemma}
\begin{proof}
Put $v=M/3$, $z=\pi^2 v$, and let $p\in\{0,1\}$ be the parity.
Poisson summation on $p+2\mathbb Z$ gives
\[
 S_r=\tfrac12\sqrt{2\pi v}\,v^r(2r-1)!!\,(1+E_r),\qquad
 E_r=2\sum_{j\ge1}(-1)^{pj}R_r(zj^2)e^{-zj^2/2},
\]
where
\[
 R_1(u)=1-u,\qquad R_2(u)=1-2u+u^2/3,\qquad
 R_3(u)=1-3u+u^2-u^3/15.
\]
The polynomials follow by differentiating the Fourier transform
$\sqrt{2\pi v}\exp(-v\xi^2/2)$ respectively two, four and six
times with respect to $\xi$, multiplying by $(-1)^r$,
and evaluating at frequencies $\xi=\pi j$.

Here are explicit bounds, valid without numerical approximations. Let
$z_0=1269/100$. Since $\pi>157/50$ and $M\ge85/22$, we have
$z>z_0$. The elementary bound on $\pi$ follows, for example, from
Machin's formula and the alternating series for the arctangent.
We have $e^3>20$ by its Taylor sum through degree $8$;
hence $e^{z_0/2}>400(1+0.345+0.345^2/2)>550$.
For $1\le r\le3$ write
\[
 B_r=\frac{z_0^r}{550\cdot999999}.
\]
Since $j^2-1\ge3(j-1)$, $j^{2r}\le64^{j-1}$ for $j\ge2$,
and $64e^{-18}<10^{-6}$, we obtain
\[
 \sum_{j\ge2}(zj^2)^r e^{-zj^2/2}
 \le z^r e^{-z/2}\sum_{j\ge2}10^{-6(j-1)}<B_r.
\]
The last inequality uses the decrease of $z^r e^{-z/2}$ on
$[z_0,\infty)$. Consequently
\[
 |E_1|\le 2((z_0-1)/550+B_1)<1/20
\]
using also that $(z-1)e^{-z/2}$ decreases there.
Since $R_2(u)-R_1(u)=u^2/3-u$ and $z\le zj^2$, we have
\[
 M|E_2-E_1|
 <\frac{6}{9}\left(\frac{z_0^3/3-z_0^2}{550}+\frac{B_3}{3}\right).
\]
This uses $M=3z/\pi^2$, $\pi^2>9$, and decrease of
$(u^3/3-u^2)e^{-u/2}$ for $u\ge z_0$ (its derivative has
sign opposite to $u^2-9u+12$). Multiplying this bound by $20/19$
gives a rational number $<3/4$. Since
$S_2/S_1=M(1+E_2)/(1+E_1)$, the first bound follows.

For the other difference, let $\alpha=(15+\sqrt{105})/2$
and $\beta=(15-\sqrt{105})/2$. For $u\ge z_0>\alpha>25/2$,
\[
 |R_3(u)-R_1(u)|=u(u-\alpha)(u-\beta)/15
 \le u^2(u-25/2)/15.
\]
For the leading term, bound
$u^2e^{-u/4}\le (25/2)^2e^{-(25/2)/4}$ by monotonicity and
$(u-25/2)e^{-u/4}\le (4/e)e^{-(25/2)/4}$ by maximization.
Together, after division by $15$, the bound is $<1/24$:
use $e>2$ and $e^{25/4}>400(1+1/4)=500$.
For the tail use $|R_3(u)-R_1(u)|\le u^3/15$.
It follows that $|E_3-E_1|<2(1/24+B_3/15)<19/100$.
Dividing by $1+E_1>19/20$ gives $(1+E_3)/(1+E_1)<6/5$.
Now $S_3/S_1=(5/3)M^2(1+E_3)/(1+E_1)$ proves the claim.
\end{proof}

\subsection{Polynomial certificate and conclusion}
Put $L=N+1$, $T=(k+1)(2N-k+1)$ and define
\[
 U=\frac{32k^2a}{N},\quad V=-\frac{32k^2b}{N},\quad
 W=-\frac{32k^2e}{N},\qquad
 J=4L^2U-(2LT+3L^2)V-2T^2W.
\]
All are integer polynomials in $N,k$, by direct expansion from
\eqref{R19M:eq:P} and \eqref{R19M:eq:C}. Table~\ref{R19M:tab:cert} gives
the total degree, number of nonzero terms, and smallest nonzero
coefficient after the indicated substitutions. Every constant
coefficient is positive. Thus $V,W,J>0$ on the required domain.
For reproducibility, the coefficient extraction can be carried out
using just these formulas; a short exact check is included below.

\begin{table}[ht]
\centering
\begin{tabular}{lrrr}
\toprule
Substitution $(k,N)$ & $V$ & $W$ & $J$\\
\midrule
$(4,10+y)$ & $(6,7,48)$ & $(4,5,36)$ & $(8,9,1200)$\\
$(5,10+y)$ & $(6,7,48)$ & $(4,5,36)$ & $(8,9,1584)$\\
$(6,10+y)$ & $(6,7,48)$ & $(4,5,36)$ & $(8,9,1968)$\\
$(7,10+y)$ & $(6,7,48)$ & $(4,5,36)$ & $(8,9,2352)$\\
$(8,10+y)$ & $(6,7,48)$ & $(4,5,36)$ & $(8,9,2736)$\\
$(9+x,10+x+y)$ & $(6,28,24)$ & $(4,15,8)$ & $(9,54,12)$\\
\bottomrule
\end{tabular}
\caption{Exact coefficient certificate on the six orthants $x,y\ge0$.}
\label{R19M:tab:cert}
\end{table}

In particular $b,e<0$, and because $M=T/(2L)$,
\[
 a+b(M+3/4)+2eM^2=\frac{N}{128k^2L^2}J>0.
\]
By \eqref{R19M:eq:dom} and Lemma~\ref{R19M:lem:gauss}, the sum of $C$ against $w$,
divided by $\sum_d w(d)d^2>0$, is greater than the left side of this
identity. This proves $D_{n,k}>0$ for the remaining $k$ and completes
the argument.

\subsubsection*{Exact recipe for Table~\ref{R19M:tab:cert}}
The following code expands only the polynomials defining the six
orthants. It is not a search over the conjecture's indices.
\begin{Verbatim}[fontsize=\footnotesize,breaklines,breakanywhere]
import sympy as s
N,k,d,x,y=s.symbols("N k d x y")
A,B=N-(k+d)/2,N-(k-d)/2
c,t=(k*(2*N-k)+d*d)/4,N*d/2
P=N**3*((A*(A-1))**3+(B*(B-1))**3)/2-(N-1)**3*(A*B)**3
H=c**3+3*c*t*t-d*(3*c*c*t+t**3)
C=s.Poly(P-P.subs(d,0)*H/H.subs(d,0),d)
U,V,W=[s.cancel(32*k*k*C.nth(j)*sgn/N)
       for j,sgn in [(2,1),(4,-1),(6,-1)]]
L,T=N+1,(k+1)*(2*N-k+1)
J=4*L**2*U-(2*L*T+3*L**2)*V-2*T**2*W
for K in range(4,10):
    sub={k:K,N:10+y} if K<9 else {k:9+x,N:10+x+y}
    for f in (V,W,J):
        p=s.Poly(f.subs(sub,simultaneous=True),x,y)
        assert p.nth(0,0)>0 and min(p.coeffs())>0
        print(K,p.total_degree(),len(p.terms()),min(p.coeffs()))
\end{Verbatim}
\endgroup


\clearpage

\section{Log-concavity for two consecutive permutation patterns}
\label{section:R20T}
\noindent{\small\textsf{ROUND 20 \enspace / \enspace Taisho \enspace / \enspace R20T}}\\[3pt]
\noindent\textbf{Target.} Bach--Remmel {1324,1423} conjecture.

\noindent\textbf{Editorial status.} Full-range proof candidate; not independently adjudicated at the tournament freeze. 

\smallskip\hrule\medskip



\begingroup
\newcommand{\GG}{\Gamma}
\setcounter{equation}{0}
\renewcommand{\theequation}{\thesection.\arabic{equation}}
\renewcommand{\theHequation}{R20T.\arabic{equation}}
\paragraph{Overview.}
For the consecutive patterns $1324$ and $1423$, Bach and Remmel conjectured strict log-concavity of the coefficients of every sign-substituted reciprocity polynomial. We prove the stated all-orders result. Reversing the polynomials exposes a conserved linear combination, which reduces the three-step recurrence to a two-step recurrence with a correction in its two lowest degrees. Two families of positive adjacent minors propagate jointly, with explicit controls at the exceptional degrees.


\subsection{The problem and a conserved difference}
Set $\GG=\{1324,1423\}$, with patterns occurring in consecutive positions. Write $P_n(y)=U_{\GG,n}(-y)$ in the notation of Bach and Remmel~\cite{R20T:BR}. Their $U_{\GG,n}(y)$ are the exponential-series coefficients of the multiplicative inverse of
\[
1+\sum_{n\ge1}\frac{t^n}{n!}\sum_{\substack{\sigma\in S_n\\\text{$\sigma$ avoids $\GG$ consecutively}}} y^{1+\operatorname{des}(\sigma)} .
\]
The conjecture immediately below Table~1 on p.~208 of~\cite{R20T:BR} concerns this specific pair of patterns for all $n$, not the larger set with an additional increasing pattern.
\begin{SourceRcaTtheorem}
For every $n\ge1$, $P_n(y)$ has positive coefficients in degrees $1,\ldots,n$, strictly log-concave on the interior of this support. In particular its coefficient sequence is unimodal.
\end{SourceRcaTtheorem}
Theorem~3 of~\cite{R20T:BR} specialized at $k_1=k_2=2$ and $y\mapsto-y$ gives
\[
 P_m=(1+y)P_{m-1}+(m-2)yP_{m-2}-(m-2)y^2P_{m-3}\quad(m\ge4).
\]
The last sign follows from the specialization on p.~207; the plus in (9) on p.~208 is inconsistent with that formula and with Table~1 already at $m=4$.
The initial rows are $P_1=y$, $P_2=y+y^2$, $P_3=y+2y^2+y^3$.
For $R_n(x)=x^nP_n(1/x)$ the recurrence reads
\[
 R_m=(1+x)R_{m-1}+(m-2)x(R_{m-2}-R_{m-3}).
\]
Consequently $R_{n+1}-xR_n-nxR_{n-1}$ is independent of $n\ge2$ by subtraction of consecutive expressions. Evaluation at $n=2$ gives
\begin{equation}
 R_{n+1}=1-x+x(R_n+nR_{n-1}) .\label{R20T:first}
\end{equation}
Together with the initial rows, this shows inductively that $R_n$ is positive on degrees $0,\ldots,n-1$ with constant coefficient 1. In particular the linear coefficient of $R_n$ is $n-1$; its quadratic coefficient is $(n-2)+(n-1)(n-3)=n^2-3n+1$ for $n\ge3$.

\subsection{Two strictly positive minors}
Write $s_i=[x^i]S$ for any polynomial $S$, using zero outside its support, and set
\[
 M_i(S,T)=s_i t_{i+1}-s_{i+1}t_i.
\]
In particular $M_i(S,xS)=s_i^2-s_{i-1}s_{i+1}$. For every $n\ge3$, set $A=R_{n-1}$ and $B=R_n$. We establish the simultaneous assertions
\begin{align}
 M_i(A,B)>0,\qquad M_i(B,xA)>0\qquad (0\le i\le n-2),\label{R20T:ord}\\
 M_i(B,xB)>0\qquad(0\le i\le n-1).\label{R20T:conc}
\end{align}
At $n=3$, $A=1+x$ and $B=1+2x+x^2$, and all three assertions follow directly. Now put $C=R_{n+1}$. For $2\le i\le n-1$, the low-degree correction $1-x$ in~\eqref{R20T:first} does not affect $M_i$. Thus
\begin{align*}
 M_i(B,C)&=M_i(B,xB)+n M_i(B,xA)>0,\\
 M_i(C,xB)&=n M_{i-1}(A,B)>0 .
\end{align*}
Here $M_{n-1}(B,xA)=0$ since both operands have degree $n-1$. For $i=0,1$, the explicit first entries are
\[
 (b_0,b_1,b_2)=(1,n-1,n^2-3n+1),\qquad
 (c_0,c_1,c_2)=(1,n,n^2-n-1).
\]
They yield
\[
\begin{array}{c|cc}
 i & M_i(B,C)&M_i(C,xB)\\ \hline
 0 &1&1\\
 1 &n^2-n+1&1
\end{array}
\]
and hence prove~\eqref{R20T:ord} for the pair $B,C$.
For every $1\le i\le n-1$, $M_{i-1}(B,C)>0$ and $M_i(C,xB)>0$ now give
\[
 b_{i-1}c_i>b_ic_{i-1},\qquad c_ib_i>c_{i+1}b_{i-1}.
\]
Multiplying and cancelling $b_{i-1}b_i>0$ proves $c_i^2>c_{i-1}c_{i+1}$. The two endpoint minors are squares of positive coefficients. This proves~\eqref{R20T:conc} for $C$, completing induction. Reversing the row and restoring the factor $y$ preserves the inequalities, proving the theorem.\hfill$\square$
\endgroup


\clearpage

\section{Gaussian limits for four q-Stirling coefficient families}
\label{section:R20H}
\noindent{\small\textsf{ROUND 20 \enspace / \enspace Heisei \enspace / \enspace R20H}}\\[3pt]
\noindent\textbf{Target.} Sagan--Swanson Conjecture 7.8.

\noindent\textbf{Editorial status.} Full-range proof candidate; not independently adjudicated at the tournament freeze. The proportional regime $k/n\to\alpha\in(0,1)$ is part of the target. No stronger local limit theorem for the exponent distribution is asserted.

\smallskip\hrule\medskip



\begingroup
\newcommand{\E}{\mathbb{E}}\newcommand{\PP}{\mathbb{P}}
\newcommand{\Var}{\operatorname{Var}}\newcommand{\Cov}{\operatorname{Cov}}
\newcommand{\ee}{\mathrm e}
\setcounter{equation}{0}
\renewcommand{\theequation}{\thesection.\arabic{equation}}
\renewcommand{\theHequation}{R20H.\arabic{equation}}
\paragraph{Overview.}
Sagan and Swanson ask for coefficient normality in both types and both kinds when the column index is proportional to the row index. We give a single proof by conditioning independent marked occupancy variables. Passing from elementary to complete homogeneous functions replaces Bernoulli occupations by geometric ones. An exponential tilt of the marks preserves those occupation laws, permitting the conditional moment generating function to be evaluated by a one-dimensional lattice limit theorem. The variance of the exponent is asymptotic to an explicit positive constant times $n^3$.


\subsection{Statement and exact encoding}
For a nonzero polynomial $P(q)$ with nonnegative coefficients, let $X(P)$ have law $\PP(X=j)=[q^j]P(q)/P(1)$. Take the unshifted unsigned $q$-Stirling polynomials in \cite{R20H:ss}, not the distribution of the column variable $k$ along a row.
\begin{SourceRcaHtheorem}\label{R20H:main}
Let $k=k(n)$ with $k/n\to\alpha\in(0,1)$. Choose $P$ from $S[n,k]$, $S_B[n,k]$, $c[n,k]$, or $c_B[n,k]$. Then
\[
 \frac{X(P)-\E X(P)}{\sqrt{\Var X(P)}}\ \Longrightarrow\ N(0,1).
\]
This is Conjecture~7.8 of \cite{R20H:ss}. In each family $\Var X(P)\sim V n^3$, with $V$ as in \eqref{R20H:V} below.
\end{SourceRcaHtheorem}
Write $m=n-k$, $[a]=1+q+\cdots+q^{a-1}$ for $a\ge1$. With $a_1,\ldots,a_J$ the listed lengths, all four polynomials are $f_m([a_1],\ldots,[a_J])$:
\[
\begin{array}{c|c|c|c}
 P & f & (a_1,\ldots,a_J)& \beta=\lim J/n\\ \hline
 c[n,k]   &e&1,2,\ldots,n-1 &1\\
 c_B[n,k] &e&1,3,\ldots,2n-1 &1\\
 S[n,k]   &h&1,2,\ldots,k &\alpha\\
 S_B[n,k] &h&1,3,\ldots,2k+1 &\alpha
\end{array}
\]
Here $e_m$ and $h_m$ are elementary and complete homogeneous symmetric polynomials. The table follows from the symmetric-function formulas in \cite{R20H:ss}; equivalently, take the coefficient of $z^m$ in $\prod(1+z[a_j])$ or $\prod(1-z[a_j])^{-1}$, respectively. The omitted zero length in the first row changes no positive-degree factor.

Let $M=a_J$, $x_j=a_j/M$, $d=1$ in type A and $d=2$ in type B. We have $M/n\to d\beta$, and $\max_j|x_j-j/J|=O(1/J)$. For $r>0$ in the elementary case or $0<r<1$ in the homogeneous case write $s_j=r x_j$. Define independent $K_j$ by
\[
\begin{array}{c|c|c|c}
 &\PP(K_j=l)&H(s_j)=\E K_j&D(s_j)=\Var K_j \\ \hline
 e&s_j^l/(1+s_j),\ l=0,1 &s_j/(1+s_j)&s_j/(1+s_j)^2\\
 h&(1-s_j)s_j^l,\ l\ge0&s_j/(1-s_j)&s_j/(1-s_j)^2 .
\end{array}
\]
Conditional on the occupancies, take $K_j$ independent uniform marks on $\{0,1,\ldots,a_j-1\}$ at each $j$, all marks independent. Let $Y_j$ be the sum of the marks at $j$ divided by $n$, and let $K=\sum K_j$, $T=\sum Y_j$. For a particular occupancy vector with $\sum K_j=m$ and a specified sequence of marks at each site, the probability equals a common normalizing factor times
$\prod_j(s_j/a_j)^{K_j}=(r/M)^m$. Summing the marks' generating monomials therefore proves the exact identity
\begin{equation}\label{R20H:encode}
 \mathcal L(X(P))=\mathcal L(nT\mid K=m).
\end{equation}
The unbounded occupancy in the $h$ case allows ordered lists of marks at each site, as required by the ordinary power $[a_j]^{K_j}$.

\subsection{The exact saddle and its covariances}
Choose $r=r_n$ so that $\sum_j H(r x_j)=m$. For all sufficiently large $n$, $0<m<J$ in the elementary case; in the homogeneous case the term with $x_J=1$ diverges at the right endpoint. Hence existence and uniqueness follow by strict monotonicity. At each fixed interior $r$ the left side divided by $n$ tends to
\[
 F(r)=\beta\int_0^1 H(rx)\,dx.
\]
This strictly increasing function starts at zero. At the right endpoint it tends to $1$ in the elementary case ($\beta=1$) and to infinity in the homogeneous case (logarithmic divergence). Thus a unique interior $r_*$ satisfies $F(r_*)=1-\alpha$. Bracket $r_*$ by any two nearby interior points, and use pointwise convergence and monotonicity to obtain $r_n\to r_*$. In particular, in the $h$ case the parameters stay uniformly away from 1. The $K_j$ have a uniformly bounded exponential moment at some positive argument, and $0\le Y_j\le 2K_j$.

For a single scaled mark write
\[
 b_j=(a_j-1)/(2n),\qquad w_j=(a_j^2-1)/(12n^2).
\]
Put $A_n=\Var K/n$, $B_n=\Cov(K,T)/n$, $C_n=\Var T/n$. By conditioning on $K_j$ these are exactly the sums divided by $n$ of
\[
 D(s_j),\qquad D(s_j)b_j,\qquad D(s_j)b_j^2+H(s_j)w_j,
\]
respectively. The Riemann sums converge to $A,B,C$, obtained by integrating over $[0,1]$ against $\beta\,dx$ with substitutions
\[
 s_j\rightsquigarrow r_*x,\quad b_j\rightsquigarrow b(x)=d\beta x/2,\quad
 w_j\rightsquigarrow w(x)=(d\beta x)^2/12.
\]
In particular $A>0$, and the variance after conditioning will be
\begin{equation}\label{R20H:V}
 V=C-B^2/A\ \ge\ \beta\int_0^1 H(r_*x)w(x)\,dx>0.
\end{equation}
The inequality is weighted Cauchy--Schwarz applied to the $D$ terms. This ensures that conditioning on occupancy cannot extinguish the mark fluctuations.

\subsection{A local estimate stable under a real tilt}
We give the short lattice estimate used below rather than presupposing a conditional limit theorem. Consider an array of $O(n)$ independent integer-valued variables $L_j$ with uniformly bounded third absolute moments, with sum $L$ of variance $v_n\asymp n$. Suppose that for at least $c n$ indices, both $\PP(L_j=0)$ and $\PP(L_j=1)$ are at least $c>0$. Uniformly in integers $b$ one has
\begin{equation}\label{R20H:local}
 \sqrt{v_n}\,\PP(L=b)=
 (2\pi)^{-1/2}\exp\!\left(-\frac{(b-\E L)^2}{2v_n}\right)+o(1).
\end{equation}
Indeed, with $L_j'$ an independent copy, the squared modulus of the characteristic function is $\E\cos(u(L_j-L_j'))\le 1-2\PP(L_j=0)\PP(L_j=1)(1-\cos u)$. Multiplying over the indicated indices bounds the centered characteristic function of $L$ on $[-\pi,\pi]$ by $\exp(-c_1 n u^2)$. At $u=z/\sqrt{v_n}$ the log characteristic function of the centered sum equals $-z^2/2+O(n|u|^3)$ for fixed $z$, by Taylor expansion and the uniform third moment bound. After this rescaling, the bound $\exp(-c_2 z^2)$ holds on the integration interval. Dominated convergence gives $L^1$ convergence of the rescaled integrand (without its unit-modulus Fourier phase and extended by zero) to the Gaussian characteristic function. Fourier inversion gives \eqref{R20H:local}, uniformly because the omitted phase has modulus one.

For fixed real $t$, let $h=t/\sqrt n$, and tilt the joint distribution by $\exp(h T)/\E\exp(h T)$. Denote this law by $\PP_t$. Independence between sites is preserved. If $U_j$ is a single integer mark at $j$, the marginal distribution of $K_j$ under this law is of the original kind with $s_j$ replaced by
\[
 s'_j=s_j\E\exp(h U_j/n)=s_j(1+O_t(n^{-1/2})).
\]
The latter formula holds uniformly since the scaled marks are bounded. For $J/3\le j\le 2J/3$, the parameters stay in a compact subinterval of the positive parameter domain under either law. Hence the probabilities of 0 and 1 are bounded below on linearly many sites. Also each tilted $K_j$ still has uniformly bounded exponential moments for large $n$.

We record the small-tilt expansions, where all unadorned expectations are under the unconditioned original law:
\begin{align}
 \log\E\exp\big(h(T-\E T)\big)&=\tfrac12 t^2 C_n+O_t(n^{-1/2}),\label{R20H:t1}\\
 \E_t K&=m+t\sqrt n\,B_n+O_t(1),\label{R20H:t2}\\
 \Var_t K&=n A_n+O_t(\sqrt n).\label{R20H:t3}
\end{align}
For details, put $g_j(u,v)=\log\E\exp(uK_j+vY_j)$. On a fixed real neighborhood of zero its derivatives through order three are bounded uniformly: derivatives before taking the logarithm are bounded by the common exponential moment and $Y_j\le2K_j$; the generating functions are bounded away from zero by Jensen's inequality. Sum the Taylor expansions of $g_j(0,h)$, $\partial_u g_j(0,h)$, and $\partial_u^2 g_j(0,h)$, using $\sum\partial_{uv}g_j(0,0)=n B_n$. The respective errors are $O(n|h|^3)$, $O(nh^2)$, and $O(n|h|)$, proving \eqref{R20H:t1}--\eqref{R20H:t3}.

Apply \eqref{R20H:local} at $b=m$ under both laws. Equations \eqref{R20H:t2}--\eqref{R20H:t3} put the tilted evaluation point a bounded number of standard deviations from its mean, so the absolute error in the scaled estimate is $o(1)$ relative to a nonzero Gaussian limit. Taking the ratio gives
\[
 \frac{\PP_t(K=m)}{\PP(K=m)}
 \longrightarrow \exp\big(-t^2 B^2/(2A)\big).
\]
By the definition of the tilt and \eqref{R20H:t1}, this yields, for every fixed real $t$,
\[
 \E\left[\left.\exp\!\left(t\frac{T-\E T}{\sqrt n}\right)\right|K=m\right]
 =\E\exp\big(h(T-\E T)\big)\,
 \frac{\PP_t(K=m)}{\PP(K=m)}
 \longrightarrow\exp(t^2 V/2).
\]
The moment generating function convergence theorem identifies the limit as $N(0,V)$. In addition the conditional laws of $Z=(T-\E T)/\sqrt n$ have uniformly bounded $\E(\exp(|Z|)\mid K=m)$ by the two tilts $t=\pm1$. Thus every fixed power is uniformly integrable; in particular the conditional mean of $Z$ tends to 0 and the conditional variance tends to $V$. Equation \eqref{R20H:encode} now gives $\Var X(P)\sim n^3V$ and replacement of the centering and scale by the true mean and standard deviation proves Theorem~\ref{R20H:main}.

\subsection*{Scope and prior work}
The theorem concerns convergence in distribution within each fixed row and column pair as the pair varies. It makes no assertion of a uniform local limit in the exponent, coefficient log-concavity, or the stronger limit shape mentioned after Conjecture~7.8. The fixed-column second-kind limit in \cite[Thm.~7.7]{R20H:ss} does not cover the proportional regime. Kasraoui's type-A moment formulas \cite{R20H:k} and the general conditional-limit method exemplified by Janson \cite{R20H:j} are related; the pairs here are independent but not identically distributed. The nearby Conjecture~7.6 concerns bottom heaviness, not normality. A search through 19 November 2026 located no prior resolution of the four-family proportional-regime assertion. This is a research manuscript, not an externally refereed result; priority is subject to further literature review.
\endgroup


\clearpage

\section{Real-rootedness of fixed-peak Fuss--Dyck polynomials}
\label{section:R20R}
\noindent{\small\textsf{ROUND 20 \enspace / \enspace Reiwa \enspace / \enspace R20R}}\\[3pt]
\noindent\textbf{Target.} Alexandersson fixed-peak conjecture.

\noindent\textbf{Editorial status.} Full-range proof candidate; not independently adjudicated at the tournament freeze. 

\smallskip\hrule\medskip



\begingroup

\setcounter{equation}{0}
\renewcommand{\theequation}{\thesection.\arabic{equation}}
\renewcommand{\theHequation}{R20R.\arabic{equation}}
\paragraph{Overview.}
Fix the number of north runs in a path with north and east steps in a
Fuss rectangle. The distribution of the number of runs of length at least
two does not depend on the integer slope parameter. A cyclic enumeration
followed by the Gauss--Lucas theorem proves Alexandersson's
real-rootedness conjecture. The enumeration belongs to the
classical Fuss--Kreweras framework; our proof does not require the
interlacing theorem for ordinary Dyck paths.


\subsection{The target and the formula}
For integers $m,n\ge1$, let $\mathcal D_n^{(m)}$ consist of paths from
$(0,0)$ to $(mn,n)$, with steps $N=(0,1)$ and $E=(1,0)$, contained
in $x\le my$. Count consecutive factors $NE$ and $NNE$ by $K$ and
$J$, respectively. Set
\[
 W^{(m)}_{n,k}(t)=\sum_{P\in\mathcal D_n^{(m)}:\,K(P)=k}t^{J(P)} .
\]
Alexandersson's conjecture \cite{R20R:A} asserts real-rootedness for every
nonempty slice. An ordinary Dyck path is the case $m=1$.

\begin{SourceRcaRtheorem}
Let $L=mn+1$ and $1\le k\le n$. When $k=n$ the slice is the constant
$\binom{L}{n}/L$. When $k<n$ it is
\begin{equation}\label{R20R:eq:kernel}
 W^{(m)}_{n,k}(t)=\frac{1}{L}\binom{L}{k}
 \sum_{i=1}^{\min(k,n-k)}
 \binom{k}{i}\binom{n-k-1}{i-1}t^i.
\end{equation}
Every zero of every slice is real and nonpositive. For fixed $n,k$,
the polynomials at different $m$ are positive scalar multiples.
\end{SourceRcaRtheorem}

\begin{proof}[Enumeration]
The height $my-x$ changes by $m$ on $N$ and by $-1$ on $E$.
Append one extra $E$. The resulting word is written uniquely as
$N^{a_1}E\cdots N^{a_L}E$, where $a_j\ge0$ and $\sum a_j=n$.
The slot increments $d_j=ma_j-1$ have sum $-1$. Admissibility
means precisely that the first $L-1$ partial sums are nonnegative:
within each slot height increases until the final east step.
Conversely the final increment of such a word must be $-1$,
since it is at least $-1$ and the total is $-1$. Thus the last
slot is empty and deleting its $E$ recovers a path. Every original
path ends in $E$, so appending $E$ has not created a marked factor.

Exactly one cyclic rotation of any slot tuple is admissible.
Indeed, if $S_j=d_1+\cdots+d_j$ with $S_0=0$, start just after
the first minimum of $S_0,\ldots,S_{L-1}$. Before wrapping,
all relevant later partial sums are at least this minimum. After
wrapping, each earlier sum is strictly greater, hence greater by at
least one, which compensates for $S_L=-1$. Conversely, a good
cut must satisfy these same weak and strict inequalities and hence
is this first minimum. The tuples have no proper cyclic period:
otherwise their sum in increments would be an integer multiple
of the number of repeated pieces, contrary to $-1$.

The number $K$ is the number of positive slots, and $J$ counts
slots of size at least two. Both are invariant under slot rotation.
Choose the $k$ positive slots in $\binom{L}{k}$ ways. Remove one
unit from each. If $n=k$, there is no remainder. Otherwise choose
the $i$ slots receiving the $n-k$ remaining units in
$\binom{k}{i}$ ways and give each chosen slot a positive share in
$\binom{n-k-1}{i-1}$ ways. Division by $L$ proves the formulas.
\end{proof}

\begin{proof}[Location of the zeros]
Suppose $k<n$ and put $b=n-k-1\ge0$. The sum in
\eqref{R20R:eq:kernel} is
\[
 [x^{k-1}](x+1)^b(x+t)^k
 =\frac{1}{(k-1)!}
 \left.\frac{d^{k-1}}{dx^{k-1}}\big((x+1)^b(x+t)^k\big)\right|_{x=0}.
\]
Indeed, the coefficient of $t^i$ in the expression on the left is
$\binom{k}{i}\binom{b}{i-1}$. For nonreal $t$ the zeros in $x$ of
the undifferentiated polynomial lie on the segment joining $-1$ to
$-t$. This segment does not contain $0$: its only real point is
$-1$. By iterated Gauss--Lucas the same segment contains the zeros of
the displayed derivative, which is not the zero polynomial.
Its value at zero is therefore nonzero. Thus $t$ cannot be a
nonreal zero of the slice. The coefficients are nonnegative and not all
zero, ruling out positive zeros. The constant boundary case is immediate.
\end{proof}

\subsection{Scope and prior work}
The fixed-peak law also identifies the zeros, including multiplicity,
with those at $m=1$. Hence interlacing statements from
Wang--Zhang \cite{R20R:WZ} transfer wherever the ordinary fixed-size or
fixed-peak slices interlace: multiplying individual polynomials by
positive constants does not move their zeros. For example, for
$m=2,n=5$ the five slices are $t,10t+10t^2,45t+45t^2,120t,42$.

The ordinary slice formula and its real-rootedness were already obtained
by B\'ona et al.\ \cite{R20R:B}. Enumeration of Fuss paths by the
whole multiset of run lengths is also classical \cite{R20R:Arm}.
In fact, summing the known type formula over multiplicities yields the
enumeration above. These facts are not claimed as discoveries. The
point is that the stated Fuss real-rootedness problem is settled
directly, with the dependence on $m$ separated. The target source still
lists this statement as a conjecture; we make no claim of exhaustive
bibliographic priority.
\endgroup


\clearpage

\section{Infinitely many strict binomial-difference contractions}
\label{section:R20S}
\noindent{\small\textsf{ROUND 20 \enspace / \enspace Showa \enspace / \enspace R20S}}\\[3pt]
\noindent\textbf{Target.} Bobby Ocean, MathOverflow 171394.

\noindent\textbf{Editorial status.} Full-range proof candidate; not independently adjudicated at the tournament freeze. 

\smallskip\hrule\medskip



\begingroup
\newcommand{\norm}[1]{\left\|#1\right\|_2}
\newcommand{\dd}{\mathrm d}
\setcounter{equation}{0}
\renewcommand{\theequation}{\thesection.\arabic{equation}}
\renewcommand{\theHequation}{R20S.\arabic{equation}}
\paragraph{Overview.}
A 2014 question asks whether two binomial differences must have opposite signs for some order whenever a positive strictly log-concave sequence first rises and then falls. We give an affirmative argument requiring only geometric decay. The proof uses two expansions of one Cauchy transform. Their common singularity region confines the angular mass of an entire function, while its coefficients confine its high-degree mass on a subsequence of circles. This is a candidate proof, separate from the tournament's winning submission; no external refereeing or global priority is claimed.


\subsection{Problem and reduction}
For a real sequence $a=(a_j)_{j\ge0}$ write
\[
 b_n=\sum_{j=0}^n(-1)^{n-j}\binom nj a_j,\qquad
 c_n=\sum_{j=0}^n(-1)^{n-j}\binom nj a_{j+1}=b_{n+1}+b_n .
\]
The last equality is Pascal's identity. The product in Ocean's question~\cite{R20S:ocean} is $(c_n-b_n)(c_n+b_n)$, so strict negativity is equivalent to $|c_n|<|b_n|$.
\begin{SourceRcaStheorem}\label{R20S:main}
Let $a$ be a nonzero real sequence and suppose $|a_j|\le C q^j$ for some $C<\infty$ and $0<q<1$. Then $|c_n|<|b_n|$ for infinitely many $n$.
\end{SourceRcaStheorem}
Indeed, in Ocean's setting the positive successive ratios $a_{j+1}/a_j$ strictly decrease; at the mode $m$ the ratio $q=a_{m+1}/a_m$ is less than one. From that index onward $a_{j+1}\le q a_j$. Increasing $C$ absorbs the finite initial segment, so the theorem answers the original existential question for every allowed sequence.

Put $\rho=\limsup_{n\to\infty}|b_n|^{1/n}$. The binomial formula gives $\rho\le 1+q$. Write $F(z)=\sum_{n\ge0}b_nz^n/n!$. On $z=re^{i\theta}$, $\norm{\cdot}$ denotes the norm with respect to $\dd\theta/(2\pi)$. The function $F+F'$ has coefficients $c_n/n!$. We shall find radii $r\to\infty$ for which
\[
 \limsup\frac{\norm{F+F'}}{\norm{F}}\le q.
\]

\subsection{The two disks and angular localization}
Consider the exterior Cauchy transform
\[
 B(t)=\sum_{n\ge0}\frac{b_n}{t^{n+1}}
      =\sum_{j\ge0}\frac{a_j}{(1+t)^{j+1}}.
\]
For $|t|>1+q$ one can sum the binomial formula absolutely (using the bound $Cq^j$), and the inner sum is $(1+t)^{-j-1}$. The two series are analytic on $|t|>\rho$ and $|1+t|>q$, respectively. Their overlap is connected: the exterior of the union of two closed disks is path-connected (from any point outside both, move vertically away from the real axis, then join around a large circle). Thus they glue to an analytic function on the complement of the compact lens
\[
 K=\{t:|t|\le\rho,\ |1+t|\le q\}.
\]
The lens cannot be empty: otherwise $B$ would be entire, vanish at infinity, and hence be zero. Its Laurent coefficients $b_n$ would all vanish, contradicting invertibility of the triangular binomial transform. In particular $\rho\ge 1-q>0$.

For $\epsilon>0$ put $K_\epsilon=\{|t|\le\rho+\epsilon,\ |1+t|\le q+\epsilon\}$. This inflated lens has interior, contains $K$ in its interior, and has a rectifiable, piecewise circular boundary. Deforming a large-circle contour, then expanding $e^{zt}$ there, gives
\[
 F(z)=\frac{1}{2\pi i}\int_{\partial K_\epsilon}B(t)e^{zt}\,\dd t.
\]
Here and below contours have positive orientation. Fix $\eta>0$. For the closed set of angles
\[
 E_\eta=\{\theta:|1+\rho e^{-i\theta}|\ge q+\eta\},
\]
one has $\max_{t\in K}\operatorname{Re}(e^{i\theta}t)<\rho$ uniformly if $E_\eta$ is nonempty. Indeed equality on $|t|\le\rho$ forces $t=\rho e^{-i\theta}$, which violates the other disk constraint. By compactness and $K_\epsilon\downarrow K$, the uniform bound remains at most $\rho-d$ for some $d>0$ on a sufficiently small fixed inflated lens. The contour formula now gives
\begin{equation}\label{R20S:angular}
 \sup_{\theta\in E_\eta}|F(re^{i\theta})|\le C_\eta e^{(\rho-d)r}.
\end{equation}
If $E_\eta$ is empty, its contribution throughout is zero.

\subsection{Choice of radii}
Select $N\to\infty$ such that $|b_N|^{1/N}\to\rho$, and set $r=N/\rho$. Stirling's formula and Parseval yield
\begin{equation}\label{R20S:peak}
 \norm{F}\ge |b_N|\frac{r^N}{N!}=\exp\{\rho r+o(r)\}.
\end{equation}
Thus (\ref{R20S:angular}) bounds an exponentially negligible fraction of the norm on these circles. The remaining issue is the angular phase in the derivative.

\subsection{Localization of the derivative}
For $v_n=b_n r^n/n!$, Parseval gives the exact identity
\[
 \norm{F'-\rho e^{-i\theta}F}^{2}
       =\sum_{n\ge0}(n/r-\rho)^2|v_n|^2.
\]
For $0<\delta<\rho$ the contribution with $|n/r-\rho|\le\delta$ is at most $\delta^2\norm{F}^2$. We justify exponential negligibility of the complementary terms. Given $\alpha>0$, $|b_n|\le C_\alpha(\rho+\alpha)^n$. Set $\lambda=\rho+\alpha$. For the upper tail $n/r\ge\rho+\delta$, multiplying the majorant $\lambda^n r^n/n!$ by $\exp(h(n-r(\rho+\delta)))\ge1$ gives
\[
 \sum_{n/r\ge\rho+\delta}(1+n/r)|v_n|
 \le C_\alpha(1+\lambda e^h)
 \exp\{r(\lambda e^h-h(\rho+\delta))\}.
\]
At $\alpha=0$, the exponent per unit $r$ equals $\rho$ at $h=0$ and has derivative $-\delta$ there. Choose small $h>0$, then small $\alpha>0$, to make it strictly less than $\rho$. On the lower tail $n/r\le\rho-\delta$, take $h<0$ and replace $\rho+\delta$ by $\rho-\delta$ in the same bound. The derivative at zero is now $+\delta$ and one can again achieve a strict gap. These $\ell^1$ estimates bound the corresponding $\ell^2$ norms as well. Using (\ref{R20S:peak}) and then $\delta\downarrow0$, we conclude that along our chosen radii
\[
 \norm{F'-\rho e^{-i\theta}F}=o(\norm F).
\]
Off $E_\eta$ the multiplier $1+\rho e^{-i\theta}$ has modulus at most $q+\eta$; its total contribution on $E_\eta$ is $o(\norm F)$ by (\ref{R20S:angular}) and (\ref{R20S:peak}). The triangle inequality proves
\[
 \limsup\frac{\norm{F+F'}}{\norm F}\le q+\eta .
\]
Let $\eta\downarrow0$. If $|c_n|\ge|b_n|$ for all but finitely many $n$, Parseval instead gives $\norm{F+F'}^2\ge\norm F^2-P(r)$ for a fixed polynomial $P$ collecting the exceptional degrees. By (\ref{R20S:peak}) the resulting lower limit of the norm ratio is at least one. This contradiction proves Theorem~\ref{R20S:main}.\hfill$\square$
\endgroup


\clearpage

\section{A global fractional-power bound for the arctangent}
\label{section:R21M}
\noindent{\small\textsf{ROUND 21 \enspace / \enspace Meiji \enspace / \enspace R21M}}\\[3pt]
\noindent\textbf{Target.} Xu--Sun--Chen Conjecture 2.

\noindent\textbf{Editorial status.} Full-range proof candidate; not independently adjudicated at the tournament freeze. 

\smallskip\hrule\medskip



\begingroup

\setcounter{equation}{0}
\renewcommand{\theequation}{\thesection.\arabic{equation}}
\renewcommand{\theHequation}{R21M.\arabic{equation}}
\paragraph{Overview.}
An algebraic upper bound for the arctangent was proposed by Xu, Sun and
Chen after matching its Maclaurin expansion through degree seven.
We derive a six-term polynomial for the derivative of the error.
Its fourth-order zero at one exhausts its possible positive roots by
Descartes' rule of signs. This makes the approximation error strictly
increasing on the entire positive half-line.


\subsection*{The inequality}
\textbf{Theorem.} For every real $x>0$,
\begin{equation}
  \arctan x< U(x):=
  \frac{53x}{18+35\bigl(1+\frac{109}{105}x^2\bigr)^{53/109}}.
  \label{R21M:eq:th}
\end{equation}
In fact $E(x):=U(x)-\arctan x$ has positive derivative for all $x>0$.

The inequality is exactly Conjecture 2, equation (27), of Xu, Sun and
Chen~\cite{R23M:XSC}, with the numerator and denominator multiplied by 35.
Their Theorem 2 concerns a different approximant with exponent $7/15$.
Agreement of the local power series does not by itself determine the sign
of the global error in~\eqref{R21M:eq:th}.

\textbf{Proof.} Put $s=(1+109x^2/105)^{1/109}$ and $D=18+35s^{53}$.
For $x>0$ we have $s>1$ and $D>0$. Using $x^2=105(s^{109}-1)/109$
and differentiating $s^{109}=1+109x^2/105$ gives
\[
  1+x^2=\frac{105s^{109}+4}{109},\qquad
  xD'=\frac{3710}{109}(s^{53}-s^{-56}).
\]
It follows that
\[
  D-xD'=\frac{1962s^{56}+105s^{109}+3710}{109s^{56}}.
\]
Subtracting $(\arctan x)'=1/(1+x^2)$ from
$U'(x)=53(D-xD')/D^2$ now yields
\begin{equation}
  E'(x)=\frac{P(s)}
  {109s^{56}(105s^{109}+4)D^2},\label{R21M:eq:der}
\end{equation}
where the polynomial in the numerator is
\begin{align*}
 P(s)&=53(1962s^{56}+105s^{109}+3710)(105s^{109}+4)
              -109^2s^{56}(18+35s^{53})^2.
\end{align*}

\subsection*{The sign certificate}
Writing $P(s)=\sum_d c_ds^d$, its complete list of nonzero coefficients is
\[
\begin{array}{c|rrrrrr}
d &0&56&109&162&165&218\\ \hline
c_d&786520&-3433500&5698350&-14554225&10918530&584325
\end{array}
\]
The coefficients of $(s-1)^k$ in the Taylor expansion at one are
$C_k=\sum_d c_d\binom{d}{k}$. For $k=0,1,2,3,4$ this gives
\[
 (C_0,C_1,C_2,C_3,C_4)=(0,0,0,0,6623016520050).
\]
Consequently $s=1$ is a root of multiplicity four. The coefficient signs
in increasing degree are $+,-,+,-,+,+$, hence there are four sign changes.
Descartes' rule bounds the number of positive real roots,
\emph{with multiplicity}, by four. Thus $P$ has no root in $(1,\infty)$.
Its leading coefficient is positive, so $P(s)>0$ throughout that interval.

The denominator in~\eqref{R21M:eq:der} is positive there. Hence $E'(x)>0$
for $x>0$. Since $E(0)=0$, this proves~\eqref{R21M:eq:th}.\hfill$\square$

\textbf{Verification and scope.}
All coefficients in the argument are integers. The accompanying
standard-library Python script reconstructs the numerator from its two
products and checks both the Taylor coefficients and the sign variations.
As an alternative positivity check independent of root counting, it
also verifies every coefficient of $P(1+t)/t^4$ is positive.
No numerical sampling is needed. Only Conjecture 2 is asserted here.
This is a proof submission, not an externally peer-reviewed result;
this tournament round stopped before its independent adjudication.
\endgroup


\clearpage

\section{Overcolored partition congruences from a finite modular certificate}
\label{section:R21H}
\noindent{\small\textsf{ROUND 21 \enspace / \enspace Heisei \enspace / \enspace R21H}}\\[3pt]
\noindent\textbf{Target.} Thejitha--Fathima Conjecture 7.1.

\noindent\textbf{Editorial status.} Full-range proof candidate; not independently adjudicated at the tournament freeze. The global step uses a finite modular-form criterion. The complete supplied certificate reaches coefficient $1795$; the original separate check through coefficient $500$ was only corroborative.

\smallskip\hrule\medskip



\begingroup
\newcommand{\Fp}{\mathbb F_5}
\newcommand{\bb}{\bar b}
\setcounter{equation}{0}
\renewcommand{\theequation}{\thesection.\arabic{equation}}
\renewcommand{\theHequation}{R21H.\arabic{equation}}
\paragraph{Overview.}
We prove all nine progression congruences of Thejitha and Fathima's Conjecture 7.1. Reduction of three Euler exponents modulo 25 makes every cusp inequality automatic in the finite verification theorem of Radu and Sellers. The resulting tests range only up to dissection index 71. For each listed triple the argument gives two progressions, and its conclusion persists whenever the two effective color counts are unchanged modulo 25.


\subsection{Statement}
Write $f_d(q)=\prod_{j\ge1}(1-q^{dj})$. Let $\bb^k_{r,s}(n)$ count $k$-tuples of overpartitions with $r$ colors for the even parts and $s$ for the odd parts; each color has its own overlining choice. With positive integer parameters, their generating function is
\[
 F_{r,s,k}(q)=\sum_{n\ge0}\bb^k_{r,s}(n)q^n
 =f_1^{-2sk}f_2^{(3s-2r)k}f_4^{(r-s)k}.
\]
Indeed, each positive part $j$ contributes $((1+q^j)/(1-q^j))^{rk}$ for even $j$ and the same factor with exponent $sk$ for odd $j$.
\medskip

\noindent\textbf{Theorem.} Take a triple from the indicated row. For every $n\ge0$ and every indicated $j$, one has $\bb^k_{r,s}(25n+j)\equiv0\pmod5$.
\[
\begin{array}{cl}
j & (r,s,k)\\ \hline
10,15 &(5,2,2),\ (4,4,5),\ (1,2,2),\ (3,1,6),\ (4,1,5)\\
5,20 &(5,3,2),\ (4,5,5),\ (1,3,2),\ (4,1,9).
\end{array}
\]
The same congruences hold when $(rk,sk)$ is replaced by any effective color pair congruent to it modulo $25$ componentwise.

The $j=10$ and $j=20$ statements, respectively, are exactly Conjecture 7.1 of the maintained November 15 version of \cite{R21H:TF}. The September version of a separate, since withdrawn preprint \cite{R21H:TFold} listed two extra triples: $(2,2,1)$ for $j=10$, and $(4,4,4)$ for $j=20$. We certify these separately in the last two rows of the verification table; no result of the withdrawn paper is used. We do not claim novelty separately for each specialization or for the modular criterion.

\subsection{Reduction to a positive Euler product}
Let $e=(-2sk,(3s-2r)k,(r-s)k)$ and let $E_d\in\{0,\ldots,24\}$ be the least residue of $e_d$ modulo $25$, for $d\in\{1,2,4\}$. Form $H=\prod_{d\mid4}f_d^{E_d}=\sum h(n)q^n$. In $\Fp[[q]]$, Frobenius gives $f_d(q)^{25}=f_d(q^{25})$, hence
\[
 F_{r,s,k}(q)=H(q)D(q^{25}),\qquad
 D(q)=\prod_{d\mid4} f_d(q)^{(e_d-E_d)/25}.
\]
Every Euler factor has constant term one. Thus $D$ is a unit, even for negative exponents. For any $0\le j<25$ the $j$-dissection of $F$ equals $D(q)\sum_{n\ge0}h(25n+j)q^n$ and vanishes if and only if the latter sum vanishes. This also proves persistence under changing effective color counts by multiples of 25.

\subsection{Finite criterion and its hypotheses}
Put $W=\sum E_d$, $S=\sum dE_d$. Since $\sum d e_d=0$, we have $25\mid S$. For $t=10,20$ define $P=\{t,25-t\}$, $t_0=\min P$ and
\[
 B=\left\lfloor\frac{36W}{24}-\frac{S}{600}-\frac{t_0}{25}\right\rfloor .
\]
We claim that $5\mid h(25n+j)$ for all $n\ge0$, $j\in P$, follows from the same assertion for $0\le n\le B$. Apply Radu--Sellers \cite[Lemma 2.4]{R21H:RS} with $(m,M,N,t)=(25,4,20,t)$, Euler vector $E$, and identically zero auxiliary vector.

Here are the admissibility checks in detail. The only prime dividing $m$ divides $N$, and every $d\mid4$ divides $mN$. With $\kappa=\gcd(m^2-1,24)=24$, the two required divisibilities are
\[
24\mid\kappa N\sum_{d\mid4}E_d\,mN/d,
\qquad 8\mid\kappa N W .
\]
Both hold. The remaining required divisor is
$600/\gcd(24(-24t-S),600)=5$, since $t$ is a multiple of $5$ but not of $25$ and $25\mid S$. It divides $N$. The parity proviso for even $m$ does not apply. For any cusp representative with first column $(a,c)$, the lower bound is the minimum over $0\le\lambda<25$ of
\[
 \frac1{600}\sum_{d\mid4}\frac{E_d}{d}\gcd\big(d(a+24\lambda c),25c\big)^2 .
\]
Every summand is nonnegative and the auxiliary contribution is zero. This argument covers every cusp, without a squarefree-level assumption.

The residue orbit consists of $st+(s-1)S/24$ modulo $25$ for square units $s$ modulo $600$. Such $s$ are $1$ modulo $24$, so the second term vanishes modulo $25$. The possible values of $s$ modulo 5 are $1,4$, both attained, for example by the squares of $1,7$. This gives exactly $P$. The index $[\mathrm{SL}_2(\mathbb Z):\Gamma_0(20)]$ equals $20(1+1/2)(1+1/5)=36$. Substitution into the lemma now yields precisely the displayed $B$.

\subsection{Exact finite test}
The table gives, in order, the nonnegative exponents, the weighted sum, the bound, the largest coefficient index $L=25B+\max P$, and the number $z$ of nonzero residues among the $2(B+1)$ tested coefficients.
\begin{center}
\begin{tabular}{ccrrrr}
\toprule
$(r,s,k);t$ & $(E_1,E_2,E_4)$ & $S$& $B$& $L$&$z$\\ \midrule
$(5,2,2);10$ & $(17,17,6)$ &75&59&1490&0\\
$(4,4,5);10$ & $(10,20,0)$ &50&44&1115&0\\
$(1,2,2);10$ & $(17,8,23)$ &125&71&1790&0\\
$(3,1,6);10$ & $(13,7,12)$ &75&47&1190&0\\
$(4,1,5);10$ & $(15,0,15)$ &75&44&1115&0\\
$(5,3,2);20$ & $(13,23,4)$ &75&59&1495&0\\
$(4,5,5);20$ & $(0,10,20)$ &100&44&1120&0\\
$(1,3,2);20$ & $(13,14,21)$ &125&71&1795&0\\
$(4,1,9);20$ & $(7,5,2)$ &25&20&520&0\\ \midrule
$(2,2,1);10$ & $(21,2,0)$ &25&34&865&0\\
$(4,4,4);20$ & $(18,16,0)$ &50&50&1270&0\\ \bottomrule
\end{tabular}
\end{center}
For reproducibility here is the exact procedure to obtain $z$ using Euler's pentagonal theorem. Array slices in the pseudocode are replaced by explicit coefficient updates; no floating-point operations or unproved pattern detection are involved.
\begin{Verbatim}[fontsize=\footnotesize,breaklines,breakanywhere]
a = [1, 0, ..., 0]             # positions 0 through L
for d in (1, 2, 4):
  repeat E[d] times:
    b = copy(a)
    for integers h >= 1 with d*h*(3*h-1)/2 <= L:
      for v in (d*h*(3*h-1)/2, d*h*(3*h+1)/2):
        for i = v,...,L:
          b[i] = b[i] + (-1)^h * a[i-v]
    a = [x mod 5 for x in b]
z = count(a[25*n+j] != 0 for 0 <= n <= B, j in P)
\end{Verbatim}
Since one update multiplies $a$ by $f_d$ through degree $L$, the counts certify exactly the finite hypotheses, not a sample beyond which a pattern is presumed. The finite modular criterion yields the infinite assertions for $H$, and the Frobenius reduction transfers them to $F$. \hfill$\square$
\endgroup


\clearpage

\section{Divisibility and square structure of Jacobi-symbol determinants}
\label{section:R21R}
\noindent{\small\textsf{ROUND 21 \enspace / \enspace Reiwa \enspace / \enspace R21R}}\\[3pt]
\noindent\textbf{Target.} Sun 2024 Conjecture 5.1(ii),(iii).

\noindent\textbf{Editorial status.} Full-range proof candidate; not independently adjudicated at the tournament freeze. The remaining parts (ii) and (iii) are the targets; part (i) has prior credit and is rederived. Composite-modulus rank defects and prime-modulus block arguments are both retained.

\smallskip\hrule\medskip



\begingroup
\newcommand{\Z}{\mathbb Z}
\newcommand{\F}{\mathbb F}
\newcommand{\adj}{\operatorname{adj}}
\newcommand{\cof}{\operatorname{cof}}
\setcounter{equation}{0}
\renewcommand{\theequation}{\thesection.\arabic{equation}}
\renewcommand{\theHequation}{R21R.\arabic{equation}}
\paragraph{Overview.}
For odd $n\equiv3\pmod4$ the determinant obtained by removing $0,1,-1$ from the Jacobi-symbol matrix of $(j+k)(j+2k)$ is divisible by $\varphi(n)/2$. When $n\equiv3\pmod8$, the quotient is an integer square. This proves the last two assertions of Sun's Conjecture 5.1; we also recover its previously proved vanishing assertion. At prime moduli the choice of deleted indices makes two circulant cofactors appear. The odd part becomes a symmetric zero-sum circulant; in the square case its companion becomes skew-symmetric. We include the integral argument that supplies the full totient factor, even when it is not square-free.


\subsection{The determinant and the claim}
For odd $n>1$ let $\chi_n(a)=(\frac an)$ denote the Jacobi symbol and set
\[
 D_n=\det[\chi_n((j+k)(j+2k))]_{2\le j,k\le n-2}.
\]
The determinant of the empty matrix is one. The deleted pair $1,-1$ is essential to the proof: it deletes exactly one coordinate from each parity block.
\begin{SourceRcbRtheorem}[Sun's Conjecture 5.1 \cite{R21R:Sun}]
\leavevmode
\begin{enumerate}
\item If $n\equiv1\pmod4$ is not a sum of two squares, then $D_n=0$.
\item If $n\equiv3\pmod4$, then $\varphi(n)/2$ divides $D_n$.
\item If $n\equiv3\pmod8$, then $D_n=(\varphi(n)/2)a^2$ for some $a\in\Z$.
\end{enumerate}
\end{SourceRcbRtheorem}
Li and Yan \cite{R21R:LY} already proved the first assertion. The vanishing of the \emph{full} determinant for $n\equiv3\pmod4$ is older \cite{R21R:K}; it does not itself give the divisibility of a minor. The contributions claimed here are the second and third assertions, with the first included for completeness. This is an unrefereed proof submission.

\subsection{Composite moduli}
Let $M_n$ denote the matrix of the same form with all residues modulo $n$ as indices. If $n$ is not square-free, its entries depend on the row index only modulo $\operatorname{rad}(n)$, since the Jacobi symbol is a product of powers of Legendre symbols. Thus $\operatorname{rank} M_n\le \operatorname{rad}(n)\le n/3<n-3$ for $n\ge9$.

For a prime $q\equiv3\pmod4$ write $\chi=\chi_q$, and $f(t)=\chi((t+1)(t+2))$. The sum
\begin{equation}
 \sum_{t\in\F_q}\chi(t)f(t)=0 \label{R21R:cubic}
\end{equation}
vanishes because $t\mapsto-2-t$ negates $t(t+1)(t+2)$ and $\chi(-1)=-1$. The vector $(\chi(k))_{k\in\F_q}$ is in the right kernel of $M_q$. For row $j\ne0$, the substitution $t=j/k$ over nonzero $k$ gives $\chi(j)$ times \eqref{R21R:cubic}; for row zero the product is $\chi(2)\sum_{k\ne0}\chi(k)=0$.

For square-free composite $n$ divisible by such a $q$, the Chinese remainder theorem identifies $M_n$ with $M_q\otimes M_{n/q}$ up to simultaneous reordering. Hence its corank is at least $n/q$. One can choose $q$ so that $n/q\ge5$: if the initial complementary factor is $3$, instead take the prime $3$, whose complementary factor is then at least $7$. It follows that $D_n=0$ for every odd composite $n$ having a prime factor $3\bmod4$ (and also for every odd nonsquare-free composite).

Every $n\equiv3\pmod4$ has such a prime factor. A positive integer not representable as a sum of two squares has, by the classical two-square theorem, a prime $3\bmod4$ to an odd exponent; it cannot itself be a prime $1\bmod4$. This proves assertion 1 and disposes of the composite cases of assertions 2 and 3.

\subsection{Two cofactor facts}
For a matrix indexed by a group with identity, let $\cof C$ be its principal cofactor at the identity. A circulant on a cyclic group $H$ has entries $h(x/y)$.
\begin{SourceRcbRlemma}[Recentering]
Suppose $|H|$ is odd and $\sum_{x\in H}h(x)=0$. Replacing the symbol $h(t)$ by $h(bt)$ for $b\in H$ does not change $\cof C$.
\end{SourceRcbRlemma}
\begin{proof}
This replacement multiplies $C$ on the left by a translation permutation $P$. Its cycles have odd length, so $\det P=1$. Row and column sums of $C$ vanish; hence $\adj C$ is a constant matrix. Indeed if its rank is less than $|H|-1$ the adjugate vanishes, and otherwise the left and right kernels are both the constant line. Consequently $\adj(PC)=(\adj C)P^{-1}=\adj C$. This also covers the zero matrix of size one, whose cofactor is one.
\end{proof}
\begin{SourceRcbRlemma}[Symmetric cofactor]
If $C$ is an integral symmetric circulant of odd size $m$ with zero row sum, then $\cof C=m u^2$ for some $u\in\Z$.
\end{SourceRcbRlemma}
\begin{proof}
The case $m=1$ is immediate. In cyclic coordinates put $s=(m-1)/2$. The Laurent symbol has a symmetric representative
\[
 L(X)=a_0+\sum_{r=1}^{s}a_r(X^r+X^{-r})=P(X+X^{-1}).
\]
Here $P\in\Z[T]$ by the two-term recurrence for $X^r+X^{-r}$ (multiply by $X+X^{-1}$ and subtract the preceding term). Since $P(2)=0$, write $P(T)=(2-T)Q(T)$ with $Q\in\Z[T]$. Write $\zeta=e^{2\pi i/m}$, $t_j=\zeta^j+\zeta^{-j}$. The eigenvalues of $C$ are $L(\zeta^j)$, by the Fourier basis. Comparing the coefficients of $z$ in $\det(zI+C)$ in the Fourier basis and by cofactor expansion gives, even for additional zero eigenvalues,
\[
 m\cof C=\prod_{j=1}^{m-1}L(\zeta^j)
 =\left(\prod_{j=1}^{s}(2-t_j)Q(t_j)\right)^2 .
\]
The first half-product $\prod_{j=1}^{s}(2-t_j)$ equals $\prod_{j=1}^{m-1}(1-\zeta^j)=m$. The other half-product $u=\prod_{j=1}^s Q(t_j)$ is an algebraic integer fixed by every cyclotomic automorphism, because multiplication by a unit modulo $m$ permutes the nonzero indices up to sign. Therefore $u\in\Z$, as needed.
\end{proof}
The unsigned graph-Laplacian analogue is familiar from the enumeration of circulant spanning trees \cite{R21R:Trees}. The proof above permits arbitrary signed integral entries.

\subsection{A split on quadratic residues}
It remains to handle a prime $p\equiv3\pmod4$. For $p=3$ the assertion is $1=1^2$. For $p>3$, let $H=(\F_p^\times)^2$ and $m=|H|=(p-1)/2$. Then $m$ is odd and $\F_p^\times=H\sqcup(-H)$. Put $f(t)=\chi_p((t+1)(t+2))$ as before and
\[
 h_+(t)=f(t)+f(-t),\qquad h_-(t)=f(t)-f(-t)\quad(t\in H).
\]
Write $C_+,C_-$ for the circulants on $H$ with these symbols. For nonzero $j,k$, the matrix entry equals $f(j/k)$. Order the surviving indices as $H\setminus\{1\}$ followed by their negatives, in matching order. Simultaneous permutation does not affect the determinant. The resulting matrix has blocks $\left(\begin{smallmatrix}A&B\\ B&A\end{smallmatrix}\right)$ with entries of $A,B$ given by $f(j/k), f(-j/k)$ respectively. The subspaces of pairs $(v,v)$ and $(v,-v)$ decompose its action over $\mathbb Q$; thus, with no scalar or sign factor,
\begin{equation}
 D_p=\cof C_+\, \cof C_- . \label{R21R:split}
\end{equation}
Since $\chi_p(t)=1$ on $H$ and $-1$ on $-H$, \eqref{R21R:cubic} says $\sum_{t\in H} h_-(t)=0$.

Set $\varepsilon=\chi_p(2)$ and $\delta=2\varepsilon$. We have $\delta\in H$ because $\chi_p(-1)=-1$. Directly, $f(2/t)=\varepsilon f(t)$, since $(2/t+1)(2/t+2)=2(t+1)(t+2)/t^2$. Applying this identity at $t$ and $-t$ gives
\[
 h_-(\delta/t)=h_-(t), \qquad h_+(\delta/t)=\varepsilon h_+(t).
\]
Choose $b\in H$ such that $b^2=\delta$, possible since $H$ has odd order. Hence $t\mapsto h_-(bt)$ is invariant under $t\mapsto t^{-1}$, and its circulant is symmetric. By both lemmas, $\cof C_-=m u^2$. This proves $m\mid D_p$ regardless of $C_+$.

If $p\equiv3\pmod8$, the supplementary law for $2$ gives $\varepsilon=-1$. Now $h_+(bt^{-1})=-h_+(bt)$, so the centered plus circulant is skew-symmetric. In particular the symbol sums to zero, and the recentering lemma applies. Its principal deletion is an even-dimensional integral skew-symmetric matrix, whose determinant is the square $v^2$ of its integral Pfaffian. Thus \eqref{R21R:split} gives $D_p=m(uv)^2$, proving assertion 3. \hfill $\square$

\paragraph{Exact checks.}
The included script uses fraction-free integer elimination on the original matrix for every applicable odd $3\le n\le127$ (39 determinants). Independently it verifies both block cofactors and each required centered symmetry in all 16 prime cases. For example $D_{11}=25\cdot5=5\cdot5^2$ and $D_{19}=361\cdot25281=9\cdot1007^2$. In contrast $D_{23}=-187590821$ is divisible by $11$, as required, but is negative. These checks are not used in the proof.
\endgroup


\clearpage

\section{Consecutive interlacing for type D Eulerian polynomials}
\label{section:R21S}
\noindent{\small\textsf{ROUND 21 \enspace / \enspace Showa \enspace / \enspace R21S}}\\[3pt]
\noindent\textbf{Target.} Chow consecutive-rank question.

\noindent\textbf{Editorial status.} Full-range proof candidate; not independently adjudicated at the tournament freeze. The cited fixed-rank oriented refinement is an external theorem input. Strict interlacing starts at rank three, while rank two has a shared root.

\smallskip\hrule\medskip



\begingroup
\newcommand{\DD}{\mathcal D}
\setcounter{equation}{0}
\renewcommand{\theequation}{\thesection.\arabic{equation}}
\renewcommand{\theHequation}{R21S.\arabic{equation}}
\paragraph{Overview.}
Chow asked whether the Eulerian polynomials of consecutive type $D$ Coxeter groups interlace. We give a proof from Savage and Visontai's fixed-rank refinement. Its first entry is the preceding rank's total, and differences of neighboring cutoff sums recover every entry of the preceding row. The latter observation rules out cancellation at a shared zero and gives strict interlacing from rank three. The rank-two exception is explicit.


Let $\DD_n$ be the signed permutations of $[n]$ with evenly many minus signs. For $w\in\DD_n$ write
\[
 d(w)=\mathbf1_{w_1+w_2<0}+\#\{1\le i<n:w_i>w_{i+1}\},
 \qquad D_n(x)=\sum_{w\in\DD_n}x^{d(w)} .
\]
Chow \cite[Sec.~5]{R21S:chow} singled out interlacing with the next rank, beyond the known real-rootedness of each $D_n$.
\begin{SourceRcbStheorem}\label{R21S:thm}
For every $n\ge3$, $D_n$ has $n$ distinct negative zeros. Between each two successive zeros of $D_{n+1}$ lies exactly one zero of $D_n$. The pair $D_2,D_3$ interlaces weakly, sharing the zero $-1$.
\end{SourceRcbStheorem}

\paragraph{Known fixed-rank input.}
We use the refined Eulerian polynomials $T_{n,i}$ of Savage and Visontai
\cite[Secs.~4--5 and proof of Thm.~7.1]{R21S:sv}. Their results give
\[
 \sum_{i=0}^{2n-1}T_{n,i}=2D_n,\qquad
 (T_{2,0},T_{2,1},T_{2,2},T_{2,3})=(2,2x,2x,2x^2).
\]
For $0\le i\le2n+1$, put $c_i=\lceil ni/(n+1)\rceil$. With
\begin{equation}\label{R21S:cut}
 H_c(x)=x\sum_{j<c}T_{n,j}(x)+\sum_{c\le j<2n}T_{n,j}(x),
 \qquad 0\le c\le2n,
\end{equation}
where all indices $j$ are nonnegative, their recurrence is $T_{n+1,i}=H_{c_i}$. At each rank $m\ge4$, the ordered row $T_{m,i}$ is pairwise weakly interlacing in index order. In particular $F=T_{m,0}$ interleaves each $G=T_{m,i}$ with the rightmost zero belonging to $G$ (ties allowed). The degree of $G$ is $\deg F$ or $\deg F+1$. These statements about the refinement are the inputs, not new claims.


\paragraph{The first entry and a nonvanishing test.}
Since $c_0=0$, $T_{n+1,0}=H_0=2D_n$. Also $c_{2n+1}=2n$, and consecutive $c_i$ differ by at most one. Thus every cutoff in \eqref{R21S:cut} occurs in the new row.
\begin{SourceRcbSlemma}\label{R21S:noc}
No complex number is a common root of an entire row $(T_{m,i})_{i=0}^{2m-1}$ for $m\ge2$.
\end{SourceRcbSlemma}
\begin{proof}
At $x=1$ every entry is positive by induction using \eqref{R21S:cut} and the initial row. If $z\ne1$ is a common zero at rank $n+1$, then $H_c(z)=0$ for every $0\le c\le2n$. But
\[
 H_{c+1}(z)-H_c(z)=(z-1)T_{n,c}(z)\qquad(0\le c<2n).
\]
So $z$ is a common zero of the preceding row. Descent reaches the nonzero constant $T_{2,0}=2$, a contradiction.
\end{proof}

\begin{proof}[Proof of Theorem~\ref{R21S:thm}]
The initial polynomials are $(1+x)^2$ and $(1+x)(x^2+10x+1)$. The quadratic has one zero on either side of $-1$, both negative, giving the exceptional weak interlacing and the simple rank-three base.

Suppose $n\ge3$ and $F=T_{n+1,0}=2D_n$ has simple zeros $r_1<\cdots<r_n<0$. For $G=T_{n+1,i}$, oriented weak interlacing and the positive leading coefficient give
\[
 (-1)^{n-j+1}G(r_j)\ge0\qquad(1\le j\le n).
\]
Indeed, unless $G(r_j)=0$, there are exactly $n-j+1$ zeros of $G$ to the right of $r_j$, counted with multiplicity; this holds for both allowed degrees. Lemma~\ref{R21S:noc} makes at least one inequality strict at each $r_j$. Hence $2D_{n+1}=\sum_i T_{n+1,i}$ has strictly alternating signs there, ending with a negative value at $r_n$. It has a positive constant term. Its degree is $n+1$ with positive leading coefficient: $c_{2n+1}=2n$ gives the last entry $xF$, and no entry has larger degree or negative coefficients. Thus its sign at $-\infty$ is opposite its sign at $r_1$. There is a zero in each gap $(r_j,r_{j+1})$, as well as in $(-\infty,r_1)$ and $(r_n,0)$. These $n+1$ gaps account for the entire degree, proving simplicity and strict interlacing.
\end{proof}

\small
\endgroup


\clearpage

\section{The optimal lower coefficient in a Gudermannian Huygens inequality}
\label{section:R22K}
\noindent{\small\textsf{ROUND 22 \enspace / \enspace Keio \enspace / \enspace R22K}}\\[3pt]
\noindent\textbf{Target.} Xu--Chen Conjecture 4.1, remaining lower bound.

\noindent\textbf{Editorial status.} Full-range proof candidate; not independently adjudicated at the tournament freeze. The remaining sharp lower inequality is proved in the submitted argument; the upper inequality is credited to the original article. This section also retains the exact enclosure identifying the minimizer.

\smallskip\hrule\medskip



\begingroup
\newcommand{\gd}{\operatorname{gd}}
\newcommand{\sech}{\operatorname{sech}}
\setcounter{equation}{0}
\renewcommand{\theequation}{\thesection.\arabic{equation}}
\renewcommand{\theHequation}{R22K.\arabic{equation}}
\paragraph{Overview.}
Xu and Chen conjectured the sharp coefficient of the inverse Gudermannian in a two-sided Huygens inequality. Their upper bound was already proved. We derive the lower bound by showing that the numerator minus a suitable multiple of the denominator is strictly convex in the square of the argument. The proof gives uniqueness of the equality point. This is a nonwinning, independently unreviewed tournament submission.


\subsection*{The sharp constant}
For $0<x<\pi/2$, let $h(x)=\gd(x)=\int_0^x\sech t\,dt$ and $a(x)=\gd^{-1}(x)=\int_0^x\sec t\,dt$. Put
\[
 G(x)=\frac{a(x)+h(x)-2x}{x^4 a(x)} .
\]
There is a unique point $x_*$ minimizing $G$ on this interval. It satisfies
\[
 1.12840812792<x_*<1.12840812794,\qquad
 0.0745497347<G(x_*)<0.0745497348.
\]
In particular, with $\lambda=G(x_*)$ and $\mu=(2/\pi)^4$ the optimal constants,
\[
2+\lambda x^3a(x)\ \leq\ \frac{h(x)+a(x)}x\ <\ 2+\mu x^3a(x).
\]
The right side is Theorem 4.2 of Xu and Chen~\cite{R22K:xc}. What remains of their Conjecture 4.1 is the asserted sharp left side, including the identification of the proposed point of equality. We do not assume that $G$ is unimodal.

\subsection*{A convex comparison}
Write $y=x^2$, $R=\pi^2/4$ and $u_j=|E_{2j}|/(2j+1)!$, where $E_{2j}$ are the signed Euler numbers from $\sech t$. The convergent secant and hyperbolic secant series imply
\[
 G(x)=N(y)/D(y),\qquad
 D(y)=\sum_{k\ge0}u_k y^k,\qquad
 N(y)=2\sum_{j\ge0}u_{2j+2}y^{2j}.
\]
Indeed, $a(x)/x=D(y)$, and the odd powers of $y$ cancel from $(a+h)/x$. Put $c=3/40$, $F_s=N-sD$. We prove $F_s''(y)>0$ for every $s\le c$ and $0<y<R$. Since $D''>0$, it suffices to use $s=c$.


Write $b_k=[y^k]F_c=2\mathbf1_{2\mid k}u_{k+2}-cu_k$ and $P(y)=\sum_{k=2}^5k(k-1)b_k y^{k-2}$. The coefficients of $(1+t)^3 P(5t/(2+2t))$, in increasing order, are
\[
 \frac{409}{45360},\qquad \frac{325}{24192},\qquad
 \frac{1387667}{31933440},\qquad \frac{1197337}{127733760}.
\]
All are positive, so $P(y)>0$ on $[0,5/2)$, which contains $[0,R]$.

For even $k\ge6$, the terms of degrees $k-2,k-1$ in $F_c''$ have combined sign given by
\[
 2u_{k+2}-c\left(u_k+\frac{k+1}{k-1}yu_{k+1}\right).
\]
Using the classical Euler-number formula (equivalently the secant partial fractions),
\[
 u_j=\frac4\pi R^{-j}\frac{\beta(2j+1)}{2j+1},\quad
 \beta(s)=\sum_{m\ge0}\frac{(-1)^m}{(2m+1)^s},
\]
we have $1-3^{-s}<\beta(s)<1$ for $s>0$. Since $R<5/2$ and $y<R$,
\[
 \frac{u_{k+2}}{u_k}>\frac4{25}\frac{13}{17}(1-3^{-17})>\frac3{25},\qquad
 y\frac{u_{k+1}}{u_k}<\frac1{1-3^{-13}}<\frac{101}{100}.
\]
As $(k+1)/(k-1)\le7/5$ and $6/25>c(1+\tfrac75\tfrac{101}{100})$, each pair is positive. Absolute convergence on compact subintervals justifies pairing. This proves the convexity assertion.

\subsection*{Existence, uniqueness and sharpness}
At $z=5/4$ exact substitution gives
\[
 \sum_{k=0}^8 b_kz^k<-59/100000 .
\]
For a short reproducible check, generate $E_0=1$ and
$E_{2j}=-\sum_{i<j}\binom{2j}{2i}E_{2i}$, which follows from $\sech t\cosh t=1$. Discarding the negative part of the omitted terms, using $u_k<(4/3)(4/9)^k/(2k+1)$, bounds the tail above by
\[
 \frac83\left(\frac49\right)^2\frac{(5/9)^9}{23(1-5/9)}<26/100000 .
\]
Thus $G(\sqrt z)<c$. The endpoint limits are $G(0+)=2u_2=1/12$ and $G(\pi/2-)=1/R^2$ (since $a(x)\to\infty$ and $h$ is bounded). Both exceed $c$. Continuity now yields an interior minimum $s<c$. The strictly convex function $F_s$ is nonnegative and vanishes at each minimizing point in the $y$ variable. Two zeros would force negative values between them. Hence the minimizer is unique. Equality there proves sharpness of the lower coefficient; the right endpoint limit proves sharpness of the known upper coefficient.


\subsection*{Identification with the proposed equality point}
The decimals in the conjecture identify a particular point rather than define a real constant exactly. Here are finite rational bounds that distinguish the minimizer. Take
\[
 L=1.12840812792,\quad U=1.12840812794,\quad M=120,\quad r=26/45 .
\]
Throughout $1\le y\le13/10$, both series $H=N,D$ and their first derivatives have nonnegative coefficients. Truncate at degree $M$ before differentiating. For each of the four resulting expressions, the error is at most
\[
 T=3r^{M+1}\left(\frac{M+1}{1-r}+\frac r{(1-r)^2}\right).
\]
Indeed the coefficient bounds $[y^k]H\le3(4/9)^k$, together with $y\ge1$ for the differentiated series, give an omitted sum bounded by $3\sum_{k>M}kr^k$. Thus on $[v,w]$, lower and upper bounds are the corresponding truncated expressions at $v,w$, with $T$ added only to the upper bound.

Write $W=N'D-ND'$. Substituting $y=L^2,U^2$ separately and multiplying positive bounds as above gives
\[
 W(L^2)<0<W(U^2).
\]
Dividing bounds for $N,D$ on $[L^2,U^2]$ gives
\[
 0.0745497347<N(y)/D(y)<0.0745497348<c.
\]
All these comparisons are integer tests after clearing denominators. The supplied check evaluates the finite sums by Horner's rule with exact fractions. For clarity, if subscripts $-,+$ denote the two bounds, the interval for $W$ is $[N'_-D_- -N_+D'_+,\ N'_+D_+ -N_-D'_-]$. There exists a stationary point in the indicated interval. At that point, with $s=N/D<c$, one has $F_s=F_s'=0$. Strict convexity forces $F_s$ to have its unique global minimum there, so this point must be the minimizer already found. This finishes the identification required by the conjecture.
\endgroup


\clearpage

\section{The sharp integral-order cutoff in a modified Bessel inequality}
\label{section:R22H}
\noindent{\small\textsf{ROUND 22 \enspace / \enspace Heisei \enspace / \enspace R22H}}\\[3pt]
\noindent\textbf{Target.} Qi Remark 5.9.

\noindent\textbf{Editorial status.} Full-range proof candidate; not independently adjudicated at the tournament freeze. The statement asks for a violating positive argument at every integral order $k\ge6$. It is distinct from Qi's neighboring shape conjectures.

\smallskip\hrule\medskip



\begingroup
\newcommand{\ii}{\mathrm{i}}
\newcommand{\R}{\operatorname{Re}}
\setcounter{equation}{0}
\renewcommand{\theequation}{\thesection.\arabic{equation}}
\renewcommand{\theHequation}{R22H.\arabic{equation}}
\paragraph{Overview.}
Feng Qi proved a lower bound for a normalized modified Bessel function in
orders one through five, and asked whether it fails in each higher order.
We give a common construction of violations: the argument is chosen to
align the first pair of conjugate poles of the differentiated Bernoulli
kernel. The remaining poles admit a uniform absolute bound.


\subsection{Statement and choice of phase}
Extend $f(u)=u/(1-e^{-u})$ by continuity at zero, and write
\[
 B_k(u)=\frac{I_k(2\sqrt u)}{u^{k/2}}
       =\sum_{j=0}^\infty\frac{u^j}{j!(k+j)!},\qquad u\ge0,
\]
with the series defining $B_k(0)$. Qi's Theorem 1.2 establishes
$B_k(u)\ge f^{(k-1)}(u)$ for all $u>0$ when $1\le k\le5$.
His Remark 5.9 conjectures that no other positive integral order
has this property.
\begin{SourceRccHthm}
For each integer $k\ge6$, some $u>0$ satisfies
$f^{(k-1)}(u)>B_k(u)$.
\end{SourceRccHthm}
Put $a=2\pi$. Let $r_k$ be the least nonnegative residue of $3k-1$
modulo 4, and put $\theta_k=r_k\pi/(2k)$ and
$u_k=a\tan\theta_k$. The strict inequality holds at $u_k$ itself,
even if $u_k=0$. In that last case continuity gives the assertion on
a nonempty positive interval.

The partial-fraction expansion of $\coth$ gives
\[
 f(u)=1+\tfrac12u+2\sum_{j\ge1}\frac{u^2}{u^2+a^2j^2}.
\]
For completeness, the expansion follows by taking the logarithmic
derivative of Euler's product
$\sinh(u/2)/(u/2)=\prod_{j\ge1}(1+u^2/(a^2j^2))$
and using $f(u)=u/2+(u/2)\coth(u/2)$.
For each summand, $2u^2/(u^2+b^2)=2+\ii b/(u-\ii b)-\ii b/(u+\ii b)$.
The rational-function series is locally normally convergent away from
its poles; differentiation on such compact subsets gives
\[
 f^{(k-1)}(u)=2(-1)^{k-1}(k-1)!
    \sum_{j\ge1}\R\frac{\ii a j}{(u-\ii a j)^k}.
\]

\subsection{Separating the first pole pair}
Since $u_k-\ii a=a\sec\theta_k\exp(\ii(\theta_k-\pi/2))$,
the signed real part of the first term is
\[
 a^{1-k}\cos^k\theta_k\,
 \cos\big((3k-1)\pi/2-k\theta_k\big)
 =a^{1-k}\cos^k\theta_k .
\]
For $j\ge2$ the absolute value of the corresponding real part is at most
$(aj)^{1-k}$. Therefore
\[
 f^{(k-1)}(u_k)\ge\frac{2(k-1)!}{a^{k-1}}\left(
     \cos^k\theta_k-\sum_{j=2}^\infty j^{1-k}\right).
\]
The series in parentheses that is subtracted is bounded above by
$2^{1-k}+2^{2-k}/(k-2)\le 3/64$ for $k\ge6$.

For $k\ge8$, put $c=3\pi/2$. Concavity of $\log\cos x$ on
$[0,\pi/2)$ and its value zero at the origin give
$k\log\cos(c/k)\ge8\log\cos(c/8)$. Thus
\[
 \cos^k\theta_k\ge\cos^8(3\pi/16)>
 (4/5)^8>1/8.
\]
Here $\cos(3\pi/16)\ge1-(3\pi/16)^2/2>4/5$, using $\pi<22/7$.
For $k=7$ the cosine is 1. For $k=6$ it is
$\cos(\pi/12)>\sqrt3/2$, giving $\cos^6\theta_6>27/64$.
Consequently the parenthesized difference in the derivative estimate
is at least $5/64$ for $k\ge7$, and greater than $3/8$ for $k=6$.

In all cases $0\le\theta_k<\pi/4$, hence $u_k<a<7$.
Directly from its positive series and
$(k+j)!\ge k!(k+1)^j$, we have
\[
 k! B_k(u_k)\le e^{u_k/(k+1)}<3 .
\]
In the opposite direction $k! f^{(k-1)}(u_k)>3$ follows for $k=6$
from $(3/4)5!6!/7^5>3$. For $k\ge7$ use
\[
 \frac{5}{32}\frac{(k-1)!k!}{7^{k-1}}>3 .
\]
This holds at $k=7$ and the quotient of successive left-hand sides
is $k(k+1)/7>1$. This proves the theorem.\hfill$\square$

\paragraph{Source and scope.} See \cite{R22H:qi}. This result concerns all
integral orders of that inequality, not the separate parameterized
shape conjectures in Remark 5.8. Tournament status: nonwinning submission,
not independently refereed.
\endgroup


\clearpage

\section{Dyadic valuations of singleton-free Bell numbers}
\label{section:R23K}
\noindent{\small\textsf{ROUND 23 \enspace / \enspace Keio \enspace / \enspace R23K}}\\[3pt]
\noindent\textbf{Target.} Moll--Ramirez--Villamizar Conjecture 7.4.

\noindent\textbf{Editorial status.} Full-range proof candidate; not independently adjudicated at the tournament freeze. The dyadic-lattice bound applies to all forward-difference orders. The implicit $2$-adic root supplements the displayed arithmetic progressions; no closed digit formula for that root is claimed.

\smallskip\hrule\medskip



\begingroup
\newcommand{\Z}{\mathbb Z}
\newcommand{\val}{v_2}
\setcounter{equation}{0}
\renewcommand{\theequation}{\thesection.\arabic{equation}}
\renewcommand{\theHequation}{R23K.\arabic{equation}}
\paragraph{Overview.}
The centered Poisson moment operator falls into independent pairs of states
after a change of lattice and reduction modulo $\sqrt 2$. Its 24th power is
therefore the identity modulo 8. We use the resulting bound on finite
differences to give a proof candidate for all the formulas in
Moll--Ram\'irez--Villamizar's Conjecture~7.4. In the progression that
their formulas do not exhaust, valuations are distances to a single
2-adic integer.


Write $D_n$ for the number of partitions of an $n$-element set with no
singleton blocks, including $D_0=1$. Write $\val$ for the 2-adic valuation,
with $\val(0)=+\infty$.
\begin{SourceRcdKtheorem}\label{R23K:thm}
The following formulas hold, with $m\ge0$ except where indicated:
\begin{align*}
 \val(D_{12m+4})&=2, &
 \val(D_{12m+7})=\val(D_{12m+10})&=1,\\
 \val(D_{24m+1})&=5+\val(m) \quad(m\ge1),\\
 \val(D_{48m+37})&=5, &
 \val(D_{96m+61})&=6.
\end{align*}
More precisely, there is a unique $\alpha\in\Z_2$ with
$\val(\alpha)=2$ such that
$\val(D_{24m+13})=5+\val(m-\alpha)$ for every $m\ge0$.
\end{SourceRcdKtheorem}
The three displayed lines are exactly equations (26)--(28) of the
published conjecture~\cite{R23K:MRV}. The last assertion specifies the rest of
this progression by an implicit 2-adic parameter, not by a formula for
its digits.

\subsection*{A tridiagonal realization}
On the basis $e_j$ indexed by $j\ge0$, define
\[
 J e_j=e_{j+1}+j e_j+j e_{j-1},\qquad e_{-1}=0.
\]
Then $D_n=[e_0]J^n e_0$. To see this, set
$F_j(t)=(e^t-1)^j\exp(e^t-1-t)/j!$.
One has $F'_j=F_{j-1}+jF_j+(j+1)F_{j+1}$, where $F_{-1}=0$,
and $F_j(0)=\mathbf1_{j=0}$. Induction gives
$[e_j]J^n e_0=F_j^{(n)}(0)$.
The exponential formula for partitions into blocks of size at least two
gives $F_0(t)=\sum_{n\ge0}D_n t^n/n!$.


\subsection*{The lattice and its difference bound}
Put $\pi=\sqrt2$, $R=\Z[\pi]$, and take the free $R$-module
$L$ spanned by $f_j=\pi^{-\lfloor j/2\rfloor} e_j$. These are
finite linear combinations: no convergence or bounded-operator
assertion for an infinite matrix is needed. The operator $J$
preserves $L$. In fact, for an even $j$, the upward coefficient in
$J f_j$ is $1$ and the downward coefficient is $j/\pi$
(absent at $j=0$). For odd $j$ the coefficients are $\pi$ and $j$,
respectively. The middle coefficient is always $j$.

Over $R/(\pi)=\mathbb F_2$, all transitions between pairs
$(f_{2i},f_{2i+1})$ vanish. Each pair supports the matrix
$\left(\begin{smallmatrix}0&1\\1&1\end{smallmatrix}\right)$, whose cube
is the identity. Thus $(J^3-I)L\subseteq\pi L$. If
$(A-I)L\subseteq \pi^b L$, then
$A^2-I=2(A-I)+(A-I)^2$ gives exponent at least
$\min(b+2,2b)$. Starting at $b=1$ and squaring three times gives
exponents $2,4,6$:
\[
 (J^{12}-I)L\subseteq4L,\qquad (J^{24}-I)L\subseteq8L.
\]
Let $c_k(a,h)=\sum_{i=0}^k(-1)^{k-i}\binom ki D_{a+ih}$; thus
$c_0(a,h)=D_a$. These are the $e_0$ coefficients of
$J^a(J^h-I)^k e_0$. Since $f_0=e_0$ and $2^rR\cap\Z=2^r\Z$,
\begin{equation}\label{R23K:bound}
 \val(c_k(a,12))\ge 2k,\qquad \val(c_k(a,24))\ge 3k
 \qquad(a,k\ge0).
\end{equation}

\subsection*{The finite arithmetic inputs}
The distinguished-element recurrence
\[
 D_0=1,\qquad D_{n+1}=\sum_{s=1}^n\binom nsD_{n-s}
\]
chooses the $s$ other members of the block containing that element.
Using it through index 61 gives the following entries modulo 256:
\[
\begin{array}{rrrrr}\toprule
 a&h&c_0(a,h)&c_1(a,h)&c_2(a,h)\\ \midrule
 1&24&0&160&128\\
 13&24&128&96&128\\
 4&12&4&232&96\\
 7&12&162&24&0\\
 10&12&58&240&0\\ \bottomrule
\end{array}
\]
The table is reproducible by initializing the list with $[1]$,
appending the recurrence reduced modulo 256 at each step, and taking
successive differences at the indicated strides. A second recurrence
using $J$ is included in the accompanying verification script.


\subsection*{From finite differences to exact valuations}
For integers $m\ge0$ the binomial theorem applied to $J^h=I+(J^h-I)$ gives
\[
 D_{a+hm}=\sum_{k=0}^m\binom mk c_k(a,h).
\]
With $h=12$, all terms of order $k\ge2$ vanish modulo 16 by
\eqref{R23K:bound}. The three bottom rows of the table now give the
first line of Theorem~\ref{R23K:thm}.

For $a\in\{1,13\}$ set $b_k=c_k(a,24)/32$. These are 2-adic integers;
$b_1$ is odd and $\val(b_2)=2$. For $k\ge3$, by~\eqref{R23K:bound},
\[
 \val(b_k)-\val(k!)\ge 3k-5-\val(k!)\ge 2k-4\ge2.
\]
Here $\val(k!)\le k-1$. The polynomial
$P_k(X)=X(X-1)\cdots(X-k+1)$ has integer coefficients, so
$P_k(x)-P_k(y)$ is divisible by $x-y$ as a polynomial identity.
Consequently each term
$b_k\bigl(\binom{x}{k}-\binom{y}{k}\bigr)$ for $k\ge2$
has valuation at least $1+\val(x-y)$ at distinct nonnegative
integers $x,y$. In particular, the normalized function
$H_a(x)=D_{24x+a}/32$ satisfies
\[
 \val\bigl(H_a(x)-H_a(y)\bigr)=\val(x-y).
\]
Indeed the $k=1$ term has precisely this valuation and the sum is finite
after $k=\max(x,y)$.

This isometry extends uniquely from nonnegative integers to
$\Z_2$ by completeness. It permutes every quotient
$\Z/2^s\Z$: isometry makes its induced self-map injective and the
quotient is finite. The compatible preimages of zero give a unique
root $\alpha_a\in\Z_2$. By isometry,
$\val(D_{24m+a})=5+\val(m-\alpha_a)$.
For $a=1$, $D_1=0$ hence $\alpha_1=0$.
For $a=13$, $D_{13}\equiv128\pmod{256}$ gives
$\val(\alpha_{13})=\val(H_{13}(0))=2$. Now
$m=2u+1$ and $m=4u+2$ have distances to $\alpha_{13}$ of valuations
0 and 1 respectively. Substitution gives the final line of the theorem.
Uniqueness of the parameter in the integer valuation formula follows
also by approximating a putative parameter modulo arbitrarily high
powers of two.\hfill$\square$
\endgroup


\clearpage

\section{A real-parameter cosine quotient and its strict monotonicity}
\label{section:R23T}
\noindent{\small\textsf{ROUND 23 \enspace / \enspace Taisho \enspace / \enspace R23T}}\\[3pt]
\noindent\textbf{Target.} Sandor Conjecture 1.

\noindent\textbf{Editorial status.} Full-range proof candidate; not independently adjudicated at the tournament freeze. 

\smallskip\hrule\medskip



\begingroup
\newcommand{\sinc}{\operatorname{sinc}}
\setcounter{equation}{0}
\renewcommand{\theequation}{\thesection.\arabic{equation}}
\renewcommand{\theHequation}{R23T.\arabic{equation}}
\noindent
S\'andor \cite{R23T:plain1} proved a trigonometric inequality for integral parameters. At
the end of his paper he asks whether the underlying monotonicity holds
for every real parameter $n\ge3$. The following argument proves it on
an interval longer than the one in the question. This note records an
unreviewed tournament submission, not a verified claim of priority.

\paragraph{Theorem.} For every real $n\ge3$, the function
\[
   g_n(t)=\frac{\cos t-\cos(nt)}{t^2\cos t}
\]
has negative derivative on $0<t<2\pi/(n+1)$. In particular this holds
throughout the interval $0<t<\pi/(2n)$ in \cite{R23T:plain1}.

\paragraph{Proof.} Put $a=(n+1)/2$, $b=(n-1)/2$, and
$S(z)=\sin(z)/z$. Then $a\ge2$, $b\ge1$, and $0<t<\pi/a$.
All four arguments $t,bt,2t,at$ belong to $(0,\pi)$. The cosine
difference identity and $S(2t)=S(t)\cos t$ give the exact factorization
\[
          g_n(t)=2ab\,S(t)S(bt)\,\frac{S(at)}{S(2t)}.
\]
Write $H(z)=z\cot z$. Its derivative has numerator
$\sin z\cos z-z=-2\int_0^z\sin^2 u\,du<0$ and denominator $\sin^2 z>0$.
Thus $H$ decreases from $H(0+)=1$ on $(0,\pi)$.
Since $zS'(z)/S(z)=H(z)-1$, logarithmic differentiation gives
\[
 t\frac{g_n'(t)}{g_n(t)}=
 (H(t)-1)+(H(bt)-1)+(H(at)-H(2t)).
\]
The first two terms are strictly negative and the third is nonpositive.
The quotient is positive by the factorization, so $g_n'(t)<0$. Finally
$\pi/(2n)<2\pi/(n+1)$ for $n\ge3$. \hfill$\square$

\paragraph{Source and scope.} See \cite{R23T:plain1}. The auxiliary checks are not needed for the proof.
\endgroup


\clearpage

\section{Sharp weighted power bounds for the hyperbolic sine mean}
\label{section:R23S}
\noindent{\small\textsf{ROUND 23 \enspace / \enspace Showa \enspace / \enspace R23S}}\\[3pt]
\noindent\textbf{Target.} Hui-Zuo Xu Conjecture 1.

\noindent\textbf{Editorial status.} Full-range proof candidate; not independently adjudicated at the tournament freeze. 

\smallskip\hrule\medskip



\begingroup
\newcommand{\sh}{M_{\sinh}}
\newcommand{\csch}{\operatorname{csch}}
\setcounter{equation}{0}
\renewcommand{\theequation}{\thesection.\arabic{equation}}
\renewcommand{\theHequation}{R23S.\arabic{equation}}
\paragraph{Overview.}
A product expansion reduces a sharp mean comparison proposed by Hui-Zuo Xu to a threshold function that decreases from $16/5$ to $2$. The two endpoint orders give the upper bounds; an intermediate order at which the endpoint weights agree gives the common lower bound. The argument covers both conjectured inequalities and the optimality of their exponents.


\subsection*{The problem}
For distinct $a,b>0$, write $A=(a+b)/2$, $G=\sqrt{ab}$ and
\[
 \sh=\frac{a-b}{2\sinh((a-b)/(a+b))},\qquad
 W_p(w)=\big[w A^p+(1-w)G^p\big]^{1/p}.
\]
Use the continuous geometric-mean interpretation when $p=0$. Put $c=1/\sinh 1$, $p_*=\log(2/3)/\log c$ and $\beta_p=c^p$ for $p>0$, with $\beta_p=1$ otherwise.
\begin{SourceRcdStheorem}[Xu's Conjecture 1, Eqs.\ (4.1)--(4.2) in \cite{R23S:xu}]
For all such $a,b$,
\begin{align*}
 W_{p_*}(2/3)&<\sh<W_{16/5}(2/3),\\
 W_{p_*}(\beta_{p_*})&<\sh<W_2(\beta_2).
\end{align*}
All exponents in the two comparisons are best possible in their respective families.
\end{SourceRcdStheorem}
The two lower bounds are identical since $c^{p_*}=2/3$. The exponent two in the second upper bound also occurs in the fixed-order work of Nowicka and Witkowski \cite{R23S:nw}. The issue is how to optimize the orders, in particular how to prove the lower bound at the nonrationally specified endpoint order.
\subsection*{Product and threshold}
Set $x=((a-b)/(a+b))^2$. Then $x$ ranges over $(0,1)$,
$G/A=\sqrt{1-x}$ and $\sh/A=m(x)=\sqrt x/\sinh\sqrt x$, with $m(0)=1$ and $m(1)=c$.
For $j=1,2,3$ set
\[
 S_j(x)=\sum_{n=1}^\infty(\pi^2 n^2+x)^{-j}.
\]
Euler's product for $\sinh$ \cite{R23S:nist} gives $-(\log m)'=S_1$ and $S_j'=-jS_{j+1}$ for $j=1,2$. Differentiation is justified by uniform convergence on $[0,1]$, by comparison with the corresponding sums at zero.

For $p>0$, let $f=m^p$, $g=(1-x)^{p/2}$, and define
\[
 R=\frac{f'}{g'}=2m^pS_1(1-x)^{1-p/2},\qquad H=\frac{1-f}{1-g}.
\]
Thus $H(x)$ is the average of $R$ on $[0,x]$ with positive weight $-g'$.
At the endpoints $H(0)=R(0)=1/3$ (by continuity) and $H(1)=1-c^p$.
The key identity is
\[
 (1-x)\frac{R'}R=
 \big[\tfrac12-(1-x)S_1\big]\big[p-B(x)\big],\qquad
 B(x)=\frac{1+(1-x)S_2/S_1}{1/2-(1-x)S_1}.
\]
The expansion $\log(\sinh\sqrt x/\sqrt x)=x/6-x^2/180+O(x^3)$ gives $S_1(0)=1/6$ and $S_2(0)=1/90$. The denominator in $B$ is therefore positive and strictly increasing. The ratio $S_2/S_1$ strictly decreases: its derivative has numerator $S_2^2-2S_1S_3<0$ by Cauchy--Schwarz. Consequently the numerator of $B$ is positive and strictly decreasing. In particular
\[
 B(0)=\frac{1+(1/90)/(1/6)}{1/2-1/6}=\frac{16}{5},\qquad B(1)=2.
\]
\subsection*{The three critical orders}
For $p=2$, $R$ strictly decreases throughout $(0,1)$, so its running average $H$ strictly decreases as well. Thus $H(x)>H(1)=1-c^2$, or
$f(x)<c^2+(1-c^2)g(x)$. This is the second upper bound.
For $p=16/5$, $R$ strictly increases, giving $H(x)>H(0)=1/3$ and $f(x)<2/3+g(x)/3$, the first upper bound. The strict endpoint comparisons in these two cases give
\[
 1-c^2<1/3<1-c^{16/5},\quad\text{so}\quad 2<p_*<16/5.
\]
For $p=p_*$, $H(0)=H(1)=1/3$. Since $B$ strictly decreases across $p_*$, $R$ first strictly decreases, then strictly increases. Consider
\[
 D(x)=\int_0^x(R(s)-1/3)(-g'(s))\,ds=(1-g(x))(H(x)-1/3).
\]
The integrand is negative near zero and can switch sign at most once, to positive. It must switch, as $D(1)=D(0)=0$. Hence $D(x)<0$ for all $0<x<1$: before the switch this follows by integrating from zero and afterwards by integrating backwards from one. Thus $f(x)>2/3+g(x)/3$, as required.

\subsection*{Why none of the exponents can be improved}
\textit{Fixed arithmetic weight $2/3$.}
For $p>p_*$ the limit at $x=1$ of the normalized proposed lower bound is $(2/3)^{1/p}>c$, so the lower bound fails nearby. Conversely, for any $0<p<16/5$, $R$ initially decreases, and $H(x)<1/3$ near zero. This reverses the desired upper bound there. For nonpositive orders the same failure follows from the monotonicity of the power mean in its order (comparison with $p=1$).

\medskip\noindent
\textit{Endpoint-calibrated arithmetic weight $\beta_p$.}
At $0<p<p_*$ one has $H(0)=1/3>1-c^p$; a lower bound in this family would instead require $H(x)<1-c^p$ for every $x\in(0,1)$. For nonpositive $p$ the proposed lower bound is $A>\sh$. For $p>2$, expansion about $x=1$ gives
\[
 f(x)=c^p+p c^p S_1(1)(1-x)+O((1-x)^2),\qquad
 (1-c^p)g(x)=o(1-x).
\]
The linear coefficient is positive. Hence $f(x)>c^p+(1-c^p)g(x)$ close to one, ruling out any larger upper-bound order. These obstructions go in the improving directions: ordinary fixed-weight power means increase in their order, whereas $W_p(c^p)$ decreases in $p>0$ by Xu's Lemma~2 \cite{R23S:xu}. This proves the claimed sharpness.\hfill$\square$
\endgroup


\clearpage

\section{Nonreal zeros in two normalized Stirling polynomial families}
\label{section:R23R}
\noindent{\small\textsf{ROUND 23 \enspace / \enspace Reiwa \enspace / \enspace R23R}}\\[3pt]
\noindent\textbf{Target.} Wang--Zhang Remark 3.5.

\noindent\textbf{Editorial status.} Full-range proof candidate; not independently adjudicated at the tournament freeze. 

\smallskip\hrule\medskip



\begingroup
\newcommand{\St}[2]{\left\{\!\begin{smallmatrix}#1\\#2\end{smallmatrix}\!\right\}}
\setcounter{equation}{0}
\renewcommand{\theequation}{\thesection.\arabic{equation}}
\renewcommand{\theHequation}{R23R.\arabic{equation}}
\noindent Let $\St{N}{K}$ count partitions of a labeled $N$-element set into $K$ nonempty
blocks. Wang and Zhang \cite[Remark~3.5]{R23R:plain1} asked whether both of the following
polynomials have nonreal zeros for every $n\ge3$:
\[
  a_k=\frac{\St{n+1}{k+1}}{\binom{n+1}{k+1}},\qquad
  F_n(x)=\sum_{k=0}^n a_kx^k,\qquad H_n(x)=\sum_{k=0}^n\frac{a_k}{k!}x^k.
\]
\textbf{Theorem.} $F_n$ has a nonreal conjugate pair for every $n\ge2$;
$H_n$ has one for every $n\ge3$. Both lower cutoffs are sharp.
\par\medskip\noindent
\textit{Proof.} Write $P(x)=\sum_{k=0}^n\binom nk b_kx^k$ with $b_n\ne0$. By Rolle's
theorem, if all zeros of $P$ are real, then so are all zeros of each derivative
(repeated zeros retain their appropriate multiplicity). In particular,
\[
 \frac{2}{n!}P^{(n-2)}(x)= b_nx^2+2b_{n-1}x+b_{n-2}
 \quad\Longrightarrow\quad b_{n-1}^2-b_nb_{n-2}\ge0
\]
under this hypothesis. Put $N=n+1$. A partition into $N-1$ blocks has one
pair. A partition into $N-2$ blocks has either one triple or two disjoint pairs;
there are three ways to pair four prescribed elements. Consequently
\[
 a_n=1,\qquad a_{n-1}=\frac n2,\qquad
 a_{n-2}=\frac{\binom N3+3\binom N4}{\binom N2}
        =\frac{(n-1)(3n-2)}{12}.
\]
For $P=F_n$ the three normalized coefficients, in descending degree, are
$1,1/2,(3n-2)/(6n)$, so the required difference is $(4-3n)/(12n)<0$ for
$n\ge2$. For $P=n!H_n$ those coefficients are
$1,n/2,(n-1)(3n-2)/6$. Now the difference is
\[
 -\frac{3n^2-10n+4}{12}
 =-\frac{3(n-3)^2+8(n-3)+1}{12}<0 \qquad(n\ge3).
\]
Neither polynomial can have all roots real in the stated range. Real coefficients
supply the conjugate root. Degrees zero and one are immediate, while
$H_2(x)=1/3+x+x^2/2$ has positive discriminant.\hfill$\square$
\par\medskip\noindent\small
This argument addresses neither the neighboring Boros--Moll conjecture nor the number of nonreal pairs.
\par\medskip\noindent\textit{Tournament status: complete proof candidate; not independently adjudicated.}
\endgroup


\clearpage

\section{Shape sums of recursively filtered matchings}
\label{section:R24M}
\noindent{\small\textsf{ROUND 24 \enspace / \enspace Meiji \enspace / \enspace R24M}}\\[3pt]
\noindent\textbf{Target.} Coulter--Do Conjecture 5.7(a), recursive interpretation.

\noindent\textbf{Editorial status.} Full-range proof candidate; not independently adjudicated at the tournament freeze. The claim is the shape sum for the explicitly recursive tableau vectors. It does not establish a simultaneous eigenbasis, commutativity, part (b), or the entire conjectural package.

\smallskip\hrule\medskip



\begingroup
\newcommand{\ch}{\operatorname{ch}}
\newcommand{\J}{\mathcal J}
\newcommand{\cA}{\mathcal A}
\newcommand{\cR}{\mathcal R}
\setcounter{equation}{0}
\renewcommand{\theequation}{\thesection.\arabic{equation}}
\renewcommand{\theHequation}{R24M.\arabic{equation}}
\paragraph{Overview.}
Coulter and Do conjectured that summing their tableau vectors over a fixed shape gives an explicit vector of Jack coefficients. We prove this for precisely the vectors generated by their stated last-box recursion, without presuming that individual vectors form a common eigenbasis. Keeping the component containing the last identity edge distinguished identifies the last matching operator with the Nazarov--Sklyanin Lax operator. The shift theorem then groups the spectral vectors from all predecessor diagrams into the required class vector. This argument does not assert commutativity of the matching operators.


\subsection{The precise assertion}
Work over $K=\mathbb C(b)$, with $\alpha=1+b$. On perfect matchings of $[2k]$ use the identity matching $e_k=(12|34|\cdots|2k-1,2k)$. The superposition of a matching with $e_k$ is a union of alternating cycles, counting doubled edges. Their half-lengths give its coset-type. For each output matching, charge the largest vertex of each cycle positively and alternate charges along its edges. If $T_a=(a,2k-1)$, the last Coulter--Do operator on coefficient functions is
\[
  (\J_k v)(m)=\sum_{a=1}^{2k-2}
    \omega(m,T_am) v(T_am),\qquad
  \omega(m,T_am)=\begin{cases}1& q_m(a)=q_m(2k-1),\\
                                   b&\text{otherwise}.\end{cases}
\]
This is their Definition 5.1 written as an incoming row, not the charge rule at the input matching.

If a vector has coefficient $v_\nu$ on each matching of type $\nu$, its characteristic is
\[
 \ch(v)=\sum_\nu v_\nu\,p_\nu / a_\nu,
 \qquad a_\nu=\alpha^{\ell(\nu)}z_\nu,\qquad
 z_\nu=\prod_{h\ge1}h^{m_h(\nu)}m_h(\nu)! .
\]
Let $J_\lambda$ be the integral Jack symmetric function at parameter $\alpha$, and $N_\lambda=\langle J_\lambda,J_\lambda\rangle_\alpha$. Here the power sums are orthogonal with squared norms $a_\nu$. The Jack hook norm formula and Coulter--Do (18)--(19) identify their vector $\mathbf w_\lambda$ with the unique class vector for which
\begin{equation}
  \ch(\mathbf w_\lambda)=J_\lambda/N_\lambda .
  \label{R24M:jack}
\end{equation}

Write $\cA(\mu)$ for the addable boxes. A box in row $r$ and column $c$ (one-based) has content $h(s)=\alpha(c-1)-(r-1)$. For $s\in\cA(\mu)$ put
\[
 E_{\mu,s}(X)=
  \prod_{t\in\cA(\mu)\setminus\{s\}}
   \frac{X-h(t)}{h(s)-h(t)} .
\]
Embed matching vectors by appending the last identity edge, extending coefficients by zero elsewhere; denote this by $\iota$. Define $u_\varnothing=1$ and, for a standard tableau ending in box $s$ of shape $\mu+s$, recursively set
\begin{equation}
  u_T= E_{\mu,s}(\J_k)\iota u_{\bar T}.
  \label{R24M:native}
\end{equation}
This is exactly the polynomial expression of Coulter--Do Conjecture 5.4(d) with the root normalization of 5.4(c). The factors here all involve just \emph{one} operator, so they have an unambiguous order.

\begin{SourceRceMtheorem}\label{R24M:main}
For every $k\ge0$ and $\lambda\vdash k$, the sum $\sum_{T:\,\operatorname{shape}(T)=\lambda}u_T$ equals $\mathbf w_\lambda$.
\end{SourceRceMtheorem}
This proves the shape identity of Conjecture 5.7(a) for the recursive meaning of its tableau vectors. It does not repair the simultaneous-eigenbasis assertion of 5.4(b), and it asserts neither 5.7(b) nor all of 5.7.

\subsection{The pointed operator}
Distinguish the cycle of the last identity edge. Denote its half-length by $j$ and the ordinary partition of the remaining half-lengths by $\nu$. All pairs $(j,\nu)$, $j+|\nu|=k$, $j\ge1$, occur. On vectors constant on each such pointed type define the isomorphism
\[
 \Phi_k(f)=\sum_{j,\nu} f(j,\nu)\, w^{j-1}p_\nu/a_\nu
 \quad\text{into }K[p_1,p_2,\ldots,w]_{k-1}.
\]
Give $p_i$ degree $i$ and $w$ degree one. A lifted ordinary class vector satisfies $\Phi_k(\iota v)=\ch(v)$.

\begin{SourceRceMlemma}\label{R24M:iso}
The space of pointed-type functions is invariant under $\J_k$, and
\[
 \Phi_k\J_k=L\Phi_k,\qquad
 L=\pi_{\ge0}\sum_{h\ge1}\left(w^{-h}p_h+w^h\alpha h\partial_{p_h}\right)
               +b w\partial_w,
\]
where $\pi_{\ge0}$ discards negative powers of $w$.
If $v$ is an ordinary class vector at rank $k$ then
\[
    \Phi_k(v)= w^{-1}L\,\ch(v).
\]
\end{SourceRceMlemma}
\begin{proof}
The cut, flip and join enumeration in the proof of Coulter--Do Proposition 4.7 gives the following incoming row:
\begin{align*}
 (\J_k f)(j,\nu)={}&
  \sum_{h=1}^{j-1} f(j-h,\nu\cup\{h\})
  + b(j-1)f(j,\nu)\\
 &+\alpha\sum_{h\ge1} h m_h(\nu) f(j+h,\nu\setminus\{h\}) .
\end{align*}
In the last sum use only parts actually present. Indeed, on the distinguished $2j$-cycle the allowed equal-charge vertices other than the pivot give one cut of each ordered positive pair $(j-h,h)$, with the first component distinguished. This can also be read by walking in steps of two starting from the pivot; the surviving segment containing its identity edge has each half-length $1,\ldots,j-1$ once. The opposite-charge choices in that cycle, excluding the fixed partner $2k$, give $j-1$ flips. Each other $2h$-cycle has $h$ choices of each charge, all joining it to the distinguished component. Inserting the row in $\Phi_k$ and using $a_{\nu\cup h}=\alpha h(m_h(\nu)+1)a_\nu$ gives the stated operator. For an ordinary class vector $v$ its value is $v_{\nu\cup h}$ when $j=h$. Thus the coefficient of $w^{h-1}$ in its pointed image is $\alpha h\partial_{p_h}\ch(v)$, proving the last assertion.
\end{proof}

\subsection{Identifying the two shape sums}
We recall the precise Jack--Lax inputs. Use Mickler's conventions $\varepsilon_1=-1$, $\varepsilon_2=\alpha$, so $\hbar=\alpha$, $\bar\varepsilon=b$, $V_h=p_h$ and the homogeneous Jack polynomial is $J_\mu$. The contents in the horizontal and vertical directions are $\alpha$ and $-1$, respectively. Generic rational identities in these parameters give the displayed one-parameter normalization; the denominators between addable contents are nonzero in $K$.

For each $\mu$ there are Lax eigenfunctions $\psi_\mu^s$ for $s\in\cA(\mu)$ with eigenvalue $h(s)$ and $w^0$ coefficient $J_\mu$. From the spectral theorem of Mickler--Moll, and the one-cell Pieri relation as recorded in Mickler (44), (49), (52)--(53), there are scalars $\tau_\mu^s$ with
\begin{equation}
 J_\mu=\sum_s\tau_\mu^s\psi_\mu^s,
 \qquad p_1J_\mu=\sum_s\tau_\mu^s J_{\mu+s}.
 \label{R24M:exp}
\end{equation}
In particular $E_{\mu,s}(L)J_\mu=\tau_\mu^s\psi_\mu^s$. The additional, essential input is that the shift theorem of Mickler (Theorem 3.4, (70), together with (54) or (74)) gives
\begin{equation}
  w^{-1}LJ_\lambda
  \ \in\ \operatorname{span}_K\{\psi_{\lambda-s}^s:s\text{ removable from }\lambda\}.
  \label{R24M:shift}
\end{equation}
Indeed, $LJ_\lambda$ is a linear combination of compressed Lax eigenfunctions indexed by removable corners; the shift theorem identifies each with $w\psi_{\lambda-s}^s$. No spectral separation between unrelated parent diagrams is being assumed.

\begin{proof}[Proof of Theorem \ref{R24M:main}]
Suppose the result holds at rank $k-1$; it holds at rank zero by normalization. Grouping the terms in \eqref{R24M:native} by their last box, the pointed image of the proposed sum at rank $k$ is
\begin{equation}
  \sum_{s\text{ removable from }\lambda}
  \frac{\tau_{\lambda-s}^s}{N_{\lambda-s}}\,\psi_{\lambda-s}^s .
  \label{R24M:goal}
\end{equation}
Here Lemma \ref{R24M:iso} applies \emph{after} the grouping, because each parent shape sum is the ordinary class vector in \eqref{R24M:jack}. Individual tableaux need not themselves be pointed-type functions.

On the other hand the pointed image of the right side of the theorem equals $w^{-1}LJ_\lambda/N_\lambda$. By \eqref{R24M:shift}, it is in the same span as \eqref{R24M:goal}. The constant-in-$w$ coefficient is injective on this particular span: the vectors $\psi_{\lambda-s}^s$ have distinct Jack polynomials $J_{\lambda-s}$ as constant terms. For the target this coefficient is $\alpha\partial_{p_1}J_\lambda/N_\lambda$. Multiplication by $p_1$ is adjoint to $\alpha\partial_{p_1}$ under the specified inner product, so the coefficient of $J_\mu$ therein is
\[
  \frac{\langle J_\lambda,p_1J_\mu\rangle_\alpha}{N_\lambda N_\mu}
  =\frac{\tau_\mu^s}{N_\mu}\qquad (\lambda=\mu+s),
\]
by \eqref{R24M:exp}. These are exactly the constant coefficients of \eqref{R24M:goal}. Equality follows by injectivity, first of the constant term on this span and then of the pointed transform. This proves the induction step.
\end{proof}

\paragraph{Status and checks.}
The finite computation supplied with this note confirms native class constancy of the recursively defined shape sums at $b=2$ modulo $10\,000\,019$ through $k=6$; this is a diagnostic, not the argument for the theorem. The prior public certificate at \cite{R24M:prior} refutes commutativity on the full cyclic space and explicitly distinguishes 5.7(a) from that refutation. The old genus-asymptotic route from this round duplicates earlier work and is not part of the theorem. No priority is inferred from the absence of another proof in a literature search, and no external peer review is asserted.
\endgroup


\clearpage

\section{All-order absolute monotonicity of a logarithmic Cusa--Huygens quotient}
\label{section:R24S}
\noindent{\small\textsf{ROUND 24 \enspace / \enspace Showa \enspace / \enspace R24S}}\\[3pt]
\noindent\textbf{Target.} Ling Zhu Conjecture 4.1.

\noindent\textbf{Editorial status.} Full-range proof candidate; not independently adjudicated at the tournament freeze. The finite negative-coefficient calculation is coupled to a proved bound for the entire infinite tail. The conclusion concerns every derivative order on $(0,\pi)$.

\smallskip\hrule\medskip



\begingroup
\newcommand{\sinc}{\operatorname{sinc}}
\setcounter{equation}{0}
\renewcommand{\theequation}{\thesection.\arabic{equation}}
\renewcommand{\theHequation}{R24S.\arabic{equation}}
\paragraph{Overview.}
Ling Zhu conjectured that the logarithm of the sinc function divided by the logarithm of the Cusa--Huygens upper bound is absolutely monotonic on $(0,\pi/2)$. We give a computer-assisted proof on $(0,\pi)$. After normalization at the first sinc zero, the problem reduces to bounding one signed Taylor series in a weighted coefficient norm. An exact finite rational calculation and a Cauchy estimate bound this norm for all indices. The verification code is included. This is a nonwinning research submission; it has not undergone independent peer review.


\subsection{The target and the coefficient comparison}
For $0<x<\pi$ let
\[
 G(x)=\frac{\log(\sin(x)/x)}{\log((2+\cos x)/3)}.
\]
Both logarithms are real on this interval, and the denominator is nonzero. Zhu's Conjecture~4.1 \cite{R24S:zhu} asks whether every derivative of $G$ is nonnegative on $(0,\pi/2)$. Mere monotonicity of $G$ is not the conjecture.
\begin{SourceRceStheorem}
For every integer $r\ge0$ and every $0<x<\pi$ one has $G^{(r)}(x)>0$.
More precisely, $G(\pi\sqrt{s})=\sum_{n\ge1}c_{n-1}s^{n-1}$ on $|s|<1$, and $n c_{n-1}>3/20$ for every $n\ge1$.
\end{SourceRceStheorem}

Write $z=\pi\sqrt{s}$, $B(s)=-\log((2+\cos z)/3)$, and
$D(s)=s/B(s)=\sum_{j\ge0}d_j s^j$. The even power series make these definitions independent of the local square root, and $d_0=6/\pi^2>3/5$. Euler's sine product gives
\[
 A(s)=-\log(\sinc z)=\sum_{k=1}^\infty \frac{\zeta(2k)}{k}s^k
 \qquad(|s|<1),
\]
by expanding the individual logarithms and summing absolutely. Hence $G(\pi\sqrt{s})=A(s)D(s)/s$ and
\[
 n c_{n-1}=\sum_{j=0}^{n-1}d_j\frac{n}{n-j}\zeta(2(n-j)).
\]
The term $j=0$ is at least $d_0$. As $n/(n-j)\le j+1$ and $1<\zeta(2k)\le\zeta(2)<2$, all possible losses together are bounded by $2S$, where
\[
 S=\sum_{j\ge1}(j+1)\max\{-d_j,0\}.
\]
The bound $S<1/5+1/40$ proved below gives $n c_{n-1}>3/20$. Termwise differentiation in $x$ inside $|x|<\pi$ then proves the assertion for every derivative: for any order there is an even monomial of sufficiently large degree and positive coefficient.

\subsection{Controlling the infinite tail}
\begin{SourceRceSlemma}
$D$ is analytic on a neighborhood of $|s|\le 11/10$, and $|d_j|\le20(10/11)^j$.
\end{SourceRceSlemma}
\begin{proof}
Set $u=(1-\cos z)/3$. On this disk, $|z|^2\le (11/10)\pi^2<11$. The value $1-u=(2+\cos z)/3$ avoids $(-\infty,0]$. Indeed if $z=a+ib$ and $\cos z\le-2$ is real, $b$ cannot be zero; hence $\sin a=0$, $|a|\ge\pi$ and $\cosh b\ge2$, whence $|b|>1$. This implies $|z|^2>\pi^2+1>(11/10)\pi^2$. The principal logarithm therefore exists on a neighborhood of the disk. Its only zero there on the disk is where $\cos z=1$, namely at $z=0$. The corresponding singularity of $D$ is removable.

Using the principal power, write
\[
 D(s)=\frac{s}{u}\int_0^1(1-u)^t\,dt,\qquad
 \frac{s}{u}=\frac{6}{\pi^2\sinc^2(z/2)} .
\]
The first identity follows by integrating the exponential; both sides have the same limiting value at zero. Euler's product and $|1-w|\ge1-|w|$ give
$|\sinc(z/2)|\ge\sinc(|z|/2)\ge 1-|z|^2/24>13/24$; all real product factors are positive. Moreover $|1-u|\le(2+\cosh\sqrt{11})/3<6$. For the last inequality the hyperbolic cosine series is bounded by $1+11/2+(121/24)/(1-11/30)<16$. Thus
\[
 |D(s)|\le\frac{6}{9}(24/13)^2\cdot6<20,
\]
since $|(1-u)^t|=|1-u|^t\le\max(1,|1-u|)$ on $0\le t\le1$. Cauchy's coefficient estimate gives the claim.
\end{proof}
For $N=150$ and $q=10/11$, the tail in $S$ is at most
\[
 V=20\sum_{j>N}(j+1)q^j
   =20q^{N+1}\left(\frac{N+2}{1-q}+\frac{q}{(1-q)^2}\right)<\frac1{40}.
\]
It remains to bound the finite negative contribution. In the unscaled variable $y=x^2$ let $f(y)=(2+\cos\sqrt y)/3$, $b(y)=-\log f(y)$, and $h(y)=y/b(y)$. Then $d_j=\pi^{2j-2}[y^j]h(y)$, so $\pi^2<10$ gives
\[
 \sum_{j=1}^{150}(j+1)\max\{-d_j,0\}\le
 U:=\sum_{j=1}^{150}(j+1)10^{j-1}\max\{-[y^j]h(y),0\}<\frac15.
\]
The final strict bound is a finite rational assertion. The following certificate subsection derives and verifies it.

\subsection{The finite certificate}
Writing subscripts for coefficients, $fb'=-f'$ and $(b/y)h=1$ give
\[
 b_n=-f_n-\frac1n\sum_{k=1}^{n-1} k b_k f_{n-k},
 \qquad h_n=-\frac1{b_1}\sum_{i=1}^n b_{i+1}h_{n-i}\quad(n\ge1).
\]
Here $f_0=1$, $b_0=0$, $h_0=1/b_1=6$, and $f_n=(-1)^n/(3(2n)!)$ for $n\ge1$. The following standalone Python program uses exact rational arithmetic. It verifies the aggregate bounds, not just a prefix of the desired positivity assertion.
\begin{Verbatim}[fontsize=\footnotesize,breaklines,breakanywhere]
from fractions import Fraction as Q
from math import factorial
N = 150
f = [Q(1)] + [Q((-1)**k,3*factorial(2*k)) for k in range(1,N+2)]
b = [Q(0)]
for n in range(1,N+2):
    b.append(-f[n]-sum((k*b[k]*f[n-k] for k in range(1,n)),Q(0))/n)
h = [1/b[1]]
for n in range(1,N+1):
    h.append(-sum(b[i+1]*h[n-i] for i in range(1,n+1))/b[1])
U = sum((j+1)*(-h[j])*10**(j-1) for j in range(1,N+1) if h[j]<0)
q = Q(10,11)
V = 20*q**(N+1)*((N+2)/(1-q)+q/(1-q)**2)
assert U < Q(1,5) and V < Q(1,40)
def ceil6(r): return -((-10**6*r.numerator)//r.denominator)
assert (ceil6(U), ceil6(V)) == (182200, 20029)
\end{Verbatim}

\paragraph{Scope and priority.}
The germ of $G$ at zero is even. Bounds on $G$ or its first derivative, or absolute monotonicity of a different Cusa--Huygens quotient, do not prove the stated conjecture. A bounded search found the author's 2020 conjecture but no prior proof of this exact derivative assertion. This is not a worldwide priority claim. The sine product and Cauchy estimate are classical; no new theorem is claimed for either tool. The proof-specific ingredient is the negative-mass estimate after rescaling by the first sinc zero.
\endgroup


\clearpage

\section{A full shape classification for scaled modified Bessel ratios}
\label{section:R24Hb}
\noindent{\small\textsf{ROUND 24 \enspace / \enspace Heisei \enspace / \enspace R24Hb}}\\[3pt]
\noindent\textbf{Target.} Karp--Vishnyakova--Zhang Conjecture 1, unimodality clause.

\noindent\textbf{Editorial status.} Full-range proof candidate; not independently adjudicated at the tournament freeze. The unimodality clause is affirmatively addressed. The additional log-concavity clause is false and is kept as a separately credited negative statement. The Nuttall-ratio corollary is retained.

\smallskip\hrule\medskip



\begingroup

\setcounter{equation}{0}
\renewcommand{\theequation}{\thesection.\arabic{equation}}
\renewcommand{\theHequation}{R24Hb.\arabic{equation}}
\paragraph{Overview.}
A weighted Wronskian classifies the monotonicity of
$I_{\nu_1}(a_1x)/I_{\nu_2}(a_2x)$ for all $\nu_1\ge\nu_2>-1$ and
$0<a_1\le a_2$. It proves the unimodality assertion in Conjecture 1 of
Karp, Vishnyakova and Zhang and removes their even-order-gap
restriction for ratios of Nuttall functions. The conjecture's further
log-concavity assertion is false; a two-term expansion shows why an
entire interval of equal positive orders violates it at large scales.


\subsection{Classification}
Fix $\nu_1\ge\nu_2>-1$ and $0<a_1\le a_2$. Set
$y_i(x)=I_{\nu_i}(a_ix)$ and $R(x)=y_1(x)/y_2(x)$ for $x>0$.
The defining series has positive terms:
\[
 I_\nu(z)=\sum_{k\ge0}\frac{(z/2)^{2k+\nu}}{k!\Gamma(k+\nu+1)}\qquad(z>0,\ \nu>-1).
\]
Here a peak means strictly increasing up to a unique interior maximum
and strictly decreasing thereafter. This convention keeps monotone
and constant cases separate.
\begin{SourceRceHbtheorem}\label{R24Hb:main}
The possible shapes of $R$ are as follows.
\begin{enumerate}
\item If both parameters agree, $R\equiv1$. If $\nu_1=\nu_2$ and $a_1<a_2$, then $R$ strictly decreases.
\item If $\nu_1>\nu_2$ and $a_1<a_2$, then $R$ has a peak.
\item If $\nu_1>\nu_2$ and $a_1=a_2$, then $R$ has a peak if
$\nu_1+\nu_2<0$ and strictly increases otherwise.
\end{enumerate}
In particular $R$ is unimodal in the sense of \cite[Conjecture 1]{R24Hq:KVZ}.
\end{SourceRceHbtheorem}
\begin{proof}
Write $W=x(y_1'y_2-y_1y_2')$, so $R'=W/(xy_2^2)$.
The scaled Bessel equations $(xy_i')'=(a_i^2x+\nu_i^2/x)y_i$ yield
\[
 W'=\frac{y_1y_2}{x}(A-Bx^2),\qquad
 A=\nu_1^2-\nu_2^2,\quad B=a_2^2-a_1^2\ge0 .
\]
Except when constant, $W$ is monotone or first increasing then
decreasing. When $\nu_1>\nu_2$ the series gives
$xy_i'/y_i=\nu_i+O(x^2)$, so $W>0$ near zero.
Consequently, $W$ either stays positive or crosses zero exactly once,
from positive to negative. In the degenerate case $A=B=0$, $W$ is a
positive constant when $\nu_1>\nu_2$. If the orders agree and $a_1<a_2$,
$W'<0$ throughout, and the coefficient $a_i^2/[4(\nu_i+1)]$ of
$x^2$ in the normalized series gives
\[
  xR'/R=-\frac{B}{2(\nu_1+1)}x^2+O(x^4)<0
\]
near zero. This proves the first case.

For fixed real order the classical expansion \cite[10.40.1]{R24Hb:DLMF} gives
\[
 I_\nu(z)=\frac{e^z}{\sqrt{2\pi z}}\left(1-\frac{4\nu^2-1}{8z}+O(z^{-2})\right).
\]
If $a_1<a_2$, then $R(x)\to0$, so a positive initially increasing
ratio must eventually have a point with negative derivative. If
$a_1=a_2=a$ and $A\ge0$, then $W'\ge0$ and $R$ increases. If $A<0$,
then $R=1-A/(2ax)+O(x^{-2})$ approaches one from above. It cannot
remain increasing, and the single downward crossing gives a peak.
Since $\nu_1>\nu_2$ here, $A$ has the sign of $\nu_1+\nu_2$.
\end{proof}

\subsection{Nuttall ratios and the distinct log-concavity clause}
For $b\ge0$, $a>0$, $\mu>0$ and $\nu>-1$ define
\[
 Q_{\mu,\nu}(a,b)=\int_b^\infty t^\mu e^{-(t^2+a^2)/2}I_\nu(at)\,dt .
\]
\begin{SourceRceHbcorollary}
For every fixed $b\ge0$ the function
$\mu\mapsto Q_{\mu,\nu_1}(a_1,b)/Q_{\mu,\nu_2}(a_2,b)$ is unimodal
on $(0,\infty)$. If not monotone, it first increases and then decreases.
\end{SourceRceHbcorollary}
\begin{proof}
Apply the ratio-of-integrals theorem \cite[Theorem 2.6]{R24Hq:KVZ} as in that
paper's Theorem 4.3: the kernel $t^\mu=e^{\mu\log t}$ is strictly
totally positive, the Gaussian weight is positive, and the ratio of
the remaining two integrands is a positive constant times $R(t)$.
Restriction to $(b,\infty)$ preserves its shape. On a compact
$\mu$-interval $[m,M]\subset(0,\infty)$, integrability at zero follows
from $m+\nu_i>-1$, and at infinity from a Gaussian times
$t^M e^{a_i t}$ for $t\ge1$. This verifies the uniform convergence
hypothesis. The positive kernel preserves the peak orientation.
\end{proof}
The source proved the above corollary for positive even integer order
gaps; Theorem~\ref{R24Hb:main} supplies its unrestricted Bessel input.

For $0<\nu_1=\nu_2=\nu<1/2$ and $a_1<a_2$, let $L=\log R$.
The same expansion, without differentiation of any remainder, gives
\[
 L(x)=(a_1-a_2)x+\tfrac12\log(a_2/a_1)+c/x+O(x^{-2}),\qquad
 c=\tfrac{1-4\nu^2}{8}(a_1^{-1}-a_2^{-1})>0 .
\]
Then $L(x)+L(3x)-2L(2x)=c/(3x)+O(x^{-2})>0$ eventually. Thus the
log-concavity assertion is false, despite full unimodality. A previous
explicit midpoint counterexample at $\nu=1/10$ is recorded in \cite{R24Hb:TT};
no priority for refuting this clause is claimed.
\endgroup


\clearpage

\section{Rectangular endpoint fibers in linear hyperpaths}
\label{section:R24T}
\noindent{\small\textsf{ROUND 24 \enspace / \enspace Taisho \enspace / \enspace R24T}}\\[3pt]
\noindent\textbf{Target.} Galvin--Sharpe combinatorial explanation request.

\noindent\textbf{Editorial status.} A complete combinatorial explanation of a previously known divisibility; no claim that the numerical divisibility is new. The numerical divisibility was already known; the completed request is a direct combinatorial explanation. The final restricted-Bell paragraph remains explicitly finite evidence, not an additional theorem.

\smallskip\hrule\medskip



\begingroup

\setcounter{equation}{0}
\renewcommand{\theequation}{\thesection.\arabic{equation}}
\renewcommand{\theHequation}{R24T.\arabic{equation}}
\paragraph{Overview.}
Galvin and Sharpe asked for a direct combinatorial explanation of the
factor $(\ell-1)^2$ in the number of strong independent $k$-sets of a
linear $\ell$-uniform hyperpath, for $n,k,\ell\ge2$. We partition the
independent sets into rectangles indexed by two edge positions and
the interior vertices. This proves the requested divisibility without
a formula and gives a variable-edge-size version. The numerical
divisibility itself was already known. Neither this result nor our
separate Bell-polynomial checks are a new solution of an unproved
numerical conjecture.


A linear hyperpath consists of successive edges $e_1,\ldots,e_n$,
each consecutive pair meeting in one vertex and nonconsecutive edges
being disjoint. The vertices are their union; let $s_i=|e_i|\ge2$.
A strong independent set meets each edge in at most one vertex.
For a vertex $v$, write $J(v)=\{i:v\in e_i\}$, a singleton or a
pair of consecutive positions. Strong independence is exactly
pairwise disjointness of these supports.

\begin{SourceRceTtheorem}
For $n,k\ge2$, the strong independent $k$-sets of a linear hyperpath
are partitioned into rectangular fibers. For each $1\le L<R\le n$,
one such fiber occurs for each strong independent $(k-2)$-set $M$
whose vertices have support strictly between $L$ and $R$.
The corresponding fiber has size $(s_L-1)(s_R-1)$.
\end{SourceRceTtheorem}
\begin{proof}
For an independent $k$-set, order its vertices by their supports
from left to right; disjoint intervals give a unique order.
Let $u$ and $v$ be its first and last vertices, let $M$ be the remaining
set, and let $L=\max J(u)$ and $R=\min J(v)$. Then $L<R$ and all the
supports in $M$ lie strictly between $L$ and $R$.

Conversely, given such $(L,R,M)$, choose independently
\[
   u\in e_L\setminus e_{L+1},\qquad v\in e_R\setminus e_{R-1}.
\]
The first condition is equivalent to $\max J(u)=L$, the second to
$\min J(v)=R$. Hence $M\cup\{u,v\}$ is independent, including in
the adjacent case $R=L+1$, in which $M$ must be empty. These two
constructions are inverse. The consecutive intersections each remove
one vertex, giving the stated size.
\end{proof}
\begin{SourceRceTcorollary}
For $s_i=\ell$ this gives fibers of size $(\ell-1)^2$. For general
edge sizes the count is a multiple of
$\gcd_{L<R}((s_L-1)(s_R-1))$.
\end{SourceRceTcorollary}
This answers the question following Question 1.10 in \cite[p.~101]{R24T:gs}.
It concerns hyperpaths, not the hypercomb unimodality problem in their
Question 1.10(4). The divisibility in the uniform case was known from
their formula; it is the natural factorization of its combinatorial
fibers that is at issue.

\paragraph{Separate checkpoint on restricted Bell polynomials.}
Let $B_{n,\le m}(x)$ count set partitions by block number with block size
at most $m$. Conjecture 6.2 of \cite{R24T:mrv} asks for real nonpositive
zeros when $m\ne3,4$. By separating the block containing element $n$ one gets
\[
 B_{n,\le m}(x)=x\sum_{j=1}^{\min(n,m)}
       {n-1\choose j-1} B_{n-j,\le m}(x),\qquad B_{0,\le m}=1 .
\]
For $5\le m\le 9$, our recorded exact root-isolation tests gave the
expected number of real roots at $15\le n\le60$ and at the multiples
of five from $65$ to $150$. This is finite evidence only, not a
proof. A second lead on restricted $r$-Bell polynomials
(Arn\'oczki--Nyul, Conjecture 3.4) was discarded because the same
target and proof route had already been recorded in round 18.
No external priority claim or competitive closure is made.
\endgroup


\clearpage

\section{Ratio log-concavity from exact quadratic-field envelopes}
\label{section:R17R}
\noindent{\small\textsf{ROUND 17 \enspace / \enspace Reiwa \enspace / \enspace R17R}}\\[3pt]
\noindent\textbf{Target.} Brian Yi Sun Conjecture 4.7.

\noindent\textbf{Editorial status.} Editorial reconstruction of an archived proof outline and exact verifier; not an original adjudicated submission. The original archive contained no PDF or TEX for this entry. The present reconstruction supplies the derivation, all coefficient tables, and the exact verifier. The first admissible three-term center is five, and center four is explicitly false.

\smallskip\hrule\medskip



\begingroup
\newcommand{\Z}{\mathbb Z}
\newcommand{\Q}{\mathbb Q}
\setcounter{equation}{0}
\renewcommand{\theequation}{\thesection.\arabic{equation}}
\renewcommand{\theHequation}{R17R.\arabic{equation}}
\paragraph{Overview.}
The Round 17 archive preserves an exact algebraic certificate and a
proof outline for the strict log-concavity of
$\bigl(R_{n+1}/R_n\bigr)_{n\ge4}$, where
$R_n=\sum_{k=0}^n\binom nk\binom{n+k}k/(2k-1)$.
We reconstruct the argument as a self-contained proof. A four-term
recurrence propagates two rational bounds over $\Q(\sqrt2)$.
Four fixed polynomials have strictly positive coefficients after the
shift $n=15+t$, certifying the induction and all subsequent Tur\'an
inequalities. Exact initial values cover the remaining centers.
Every certificate is specified in the text and reproduced by an
embedded standard-library Python program.


\paragraph{Editorial and archival status.}
This is a reconstruction of the checkpoint reported on
19 November 2026 at 04:42:02 UTC. The original entry contained
\texttt{CHECKPOINT.md}, \texttt{verify.py}, and \texttt{results.txt},
but no mathematical manuscript. It was not independently adjudicated
before the tournament stopped. The present proof and exact replay do
not alter that historical status, and no publication-priority claim
is made. The reconstruction fills in the generating-function derivation,
the recurrence monotonicity, the denominator signs, and the induction
indices. No unresolved mathematical step is left in the certificate
argument below. No assertion about Sun's distinct Conjecture 4.8 is made.

\subsection{Statement and the indexing convention}
For integers $n\ge0$ put
\[
 R_n=\sum_{k=0}^n\binom nk\binom{n+k}k\frac1{2k-1}.
\]
The term $k=0$ equals $-1$. For $n\ge1$, the term $k=1$
is $n(n+1)$ and all later terms are nonnegative, so
$R_n\ge n(n+1)-1>0$. Define $r_n=R_{n+1}/R_n$ for $n\ge1$.

\begin{SourceRbhRtheorem}\label{R17R:bys:main}
For every integer $n\ge5$,
\[
 r_n^2>r_{n-1}r_{n+1}.
\]
Consequently the indexed sequence $(r_n)_{n\ge4}$ is strictly
log-concave, which implies the log-concavity assertion in
Conjecture 4.7 of Brian Yi Sun~\cite{R17R:bys:sun}.
The inequality at center $n=4$ is false.
\end{SourceRbhRtheorem}

A sequence beginning at index $4$ has its first three-term inequality
at center $5$: its first triple is $(r_4,r_5,r_6)$.
This convention is essential here. Set
\begin{equation}\label{R17R:bys:T}
 T_n=R_{n+1}^3R_{n-1}-R_n^3R_{n+2}.
\end{equation}
For $n\ge2$ all denominators are positive and
\[
 r_n^2-r_{n-1}r_{n+1}
 =\frac{T_n}{R_n^2R_{n-1}R_{n+1}}.
\]
In particular, the exact value $T_4=-4623352$ rules out extending
the conclusion to center $4$.

\subsection{An independently derived recurrence}
We use formal power series, so no convergence interchange is needed.
The identity
\[
 \binom nk\binom{n+k}k
 =\binom{2k}k\binom{n+k}{2k}
\]
and the binomial-series formula give
\[
 \sum_{n\ge k}\binom nk\binom{n+k}k t^n
 =\binom{2k}k\frac{t^k}{(1-t)^{2k+1}}.
\]
Also the binomial theorem yields
\[
 \sum_{k\ge0}\frac{\binom{2k}k}{2k-1}z^k
 =-\sqrt{1-4z},
\]
including the constant term $-1$; for $k\ge1$ the coefficient identity
follows directly by expanding the square root. It follows that
\begin{equation}\label{R17R:bys:gf}
 G(t):=\sum_{n\ge0}R_nt^n
 =-\frac{\sqrt{1-6t+t^2}}{(1-t)^2}.
\end{equation}
The square root is the unique formal series with constant term $1$.
Logarithmic differentiation, or direct differentiation of
\eqref{R17R:bys:gf}, gives
\[
 (1-t)(1-6t+t^2)G'(t)=(t^2-8t-1)G(t).
\]
Taking the coefficient of $t^{n+2}$ proves, for every $n\ge0$,
\begin{equation}\label{R17R:bys:rec}
 (n+3)R_{n+3}-(7n+13)R_{n+2}
 +(7n+15)R_{n+1}-(n+1)R_n=0.
\end{equation}
The initial values are $R_0=-1$, $R_1=1$, and $R_2=7$.
This reproduces the recurrence stated in~\cite[Eq.~(1.4)]{R17R:bys:sun}
without assuming that recurrence as an external input.

For $n\ge1$, division by $R_{n+2}>0$ gives
\begin{equation}\label{R17R:bys:update}
 r_{n+2}=F_n(r_n,r_{n+1}),\qquad
 F_n(x,y)=\frac{7n+13-(7n+15)/y+(n+1)/(xy)}{n+3}.
\end{equation}
For $n\ge0$ and $x,y>1$,
\begin{align*}
 \frac{\partial F_n}{\partial x}
   &=-\frac{n+1}{(n+3)x^2y}<0,\\
 \frac{\partial F_n}{\partial y}
   &=\frac{7n+15-(n+1)/x}{(n+3)y^2}>0.
\end{align*}
The second numerator is positive since it is greater than
$7n+15-(n+1)=6n+14$.

\subsection{The enclosure and its four certificates}
For $n>0$ set $B(n)=64n^4$ and
\begin{align*}
 A_\pm(n)={}&(192+128\sqrt2)n^4-(288+192\sqrt2)n^3\\
 &+(192+102\sqrt2)n^2-(519+255\sqrt2)n\ \pm4096,\\
 L_n={}&A_-(n)/B(n),\qquad U_n=A_+(n)/B(n).
\end{align*}
Thus $U_n-L_n=128/n^4>0$.
For $e,f,g\in\{-1,+1\}$, where $A_{-1}=A_-$ and
$A_{+1}=A_+$, define the polynomial
\begin{align}
 \Delta_{efg}(n)={}&(7n+13)A_e(n)A_f(n+1)B(n+2)\notag\\
 &-(7n+15)A_e(n)B(n+1)B(n+2)\notag\\
 &+(n+1)B(n)B(n+1)B(n+2)\notag\\
 &-(n+3)A_e(n)A_f(n+1)A_g(n+2).
 \label{R17R:bys:Delta}
\end{align}
The four certificates are
\begin{align}
 Q_0(n)&=A_-(n)-B(n),\notag\\
 Q_L(n)&=\Delta_{+,-,-}(n),\notag\\
 Q_U(n)&=-\Delta_{-,+,+}(n),\notag\\
 Q_T(n)&=A_-(n+1)^2B(n)B(n+2)
       -A_+(n)A_+(n+2)B(n+1)^2.\label{R17R:bys:Q}
\end{align}

\begin{SourceRbhRlemma}\label{R17R:bys:cert}
For every real $n\ge15$, all four polynomials in
\eqref{R17R:bys:Q} are strictly positive. Consequently
\begin{align*}
 1<L_n<U_n,\qquad
 F_n(U_n,L_{n+1})>L_{n+2},\qquad
 F_n(L_n,U_{n+1})<U_{n+2},\\
 L_{n+1}^2>U_nU_{n+2}.
\end{align*}
\end{SourceRbhRlemma}

\begin{proof}
Put $n=15+t$. Each $Q(15+t)$ has the form
$\sum_{j=0}^d(a_j+b_j\sqrt2)t^j$, with every integer $a_j,b_j$
strictly positive. The full coefficient lists are in
Subsection~\ref{R17R:bys:data}; the following summary records the degree
and the minimum in each of the two integer components.
\begin{center}
\begin{tabular}{lrrr}
\toprule
Polynomial & Degree & $\min a_j$ & $\min b_j$\\\midrule
$Q_0$ & 4 & 128 & 128\\
$Q_L$ & 9 & 125762560 & 163347456\\
$Q_U$ & 9 & 209781760 & 239305728\\
$Q_T$ & 13 & 855638016 & 603979776\\\bottomrule
\end{tabular}
\end{center}
All identities are finite expansions in the ring
$\Z[\sqrt2][t]$ and can be checked by the complete program in
Subsection~\ref{R17R:bys:code}. Thus each polynomial is positive for $t\ge0$.
Since $B(n)>0$, positivity of $Q_0$ gives $L_n>1$, and
$A_+(n)>A_-(n)>0$ follows.

For any signs $e,f,g$, subtraction in \eqref{R17R:bys:update} gives
\[
 F_n\!\left(\frac{A_e(n)}{B(n)},
              \frac{A_f(n+1)}{B(n+1)}\right)
 -\frac{A_g(n+2)}{B(n+2)}
 =\frac{\Delta_{efg}(n)}
 {(n+3)A_e(n)A_f(n+1)B(n+2)}.
\]
The denominator is positive for $n\ge15$. This proves both propagation
inequalities. Finally $Q_T(n)>0$, divided by
$B(n+1)^2B(n)B(n+2)>0$, gives the last displayed inequality.
\end{proof}

\subsection{Base cases, induction, and all centers}
For $m=15,16$, exact calculation gives
\[
 B(m)R_{m+1}-A_-(m)R_m>0,\qquad
 A_+(m)R_m-B(m)R_{m+1}>0.
\]
For clarity, these four quantities, written as $a+b\sqrt2$, are
\begin{center}
\begin{tabular}{llrr}
\toprule
$m$ & Endpoint & $a$ & $b$\\\midrule
15 & lower & 18862838954118761 & -13333752423991125\\
15 & upper & -18844170732220009 & 13333752423991125\\
16 & lower & 129384610949751184 & -91465704019757040\\
16 & upper & -129286333588046224 & 91465704019757040\\\bottomrule
\end{tabular}
\end{center}
The lower rows have $a>0>b$ and $a^2-2b^2>0$; the upper rows have
$b>0>a$ and $2b^2-a^2>0$. These are exact integer comparisons, not
floating-point approximations. Their positive square differences, in
row order, are
\begin{align*}
 &228786000441750009097934643871,\\
 &475136823521213199746577571169,\\
 &8427526958875688778470250278656,\\
 &16993971020029074293775698465024.
\end{align*}
Dividing by $B(m)R_m>0$ establishes $L_m<r_m<U_m$ at both base indices.

Now suppose the enclosure is known at consecutive indices $n,n+1$,
where $n\ge15$. The monotonicity of $F_n$ on $x,y>1$, together with
Lemma~\ref{R17R:bys:cert}, gives
\[
 L_{n+2}<F_n(U_n,L_{n+1})
 \le F_n(r_n,r_{n+1})
 \le F_n(L_n,U_{n+1})<U_{n+2}.
\]
Equation~\eqref{R17R:bys:update} therefore propagates the enclosure.
Starting from $15,16$ proves $L_m<r_m<U_m$ for every integer $m\ge15$.
For $n\ge15$, all involved quantities are positive, so
\[
 r_{n+1}^2>L_{n+1}^2>U_nU_{n+2}>r_nr_{n+2}.
\]
These are exactly all required inequalities with center at least $16$.

For the remaining centers, the exact values of $T_n$ from
\eqref{R17R:bys:T} are as follows. They can be obtained directly from the
binomial definition; the program also checks agreement with
\eqref{R17R:bys:rec} for every initial term used.
\begin{center}
\begin{tabular}{rr}
\toprule
$n$ & $T_n$\\\midrule
5 & 4658935984\\
6 & 1426910728568\\
7 & 393897783470272\\
8 & 119802351175061384\\
9 & 42058621023318837712\\
10 & 16992443478932799051512\\
11 & 7732754708894786339153952\\
12 & 3872046511651699311391784584\\
13 & 2093303776078940228058135063536\\
14 & 1204760462362780543122917625300216\\
15 & 730563686670982239259274106754908544\\
\bottomrule
\end{tabular}
\end{center}
Every displayed integer is positive. These centers $5$ through $15$
join the infinite range proved above, completing
Theorem~\ref{R17R:bys:main}.\hfill$\square$

\paragraph{Why the finite certificates cover an infinite domain.}
The initial checks are only the two enclosure bases and eleven
remaining Tur\'an centers. The four polynomial identities establish
inequalities for every real $n\ge15$, so the induction does not stop at
a computational cutoff. The coefficient lists are part of the proof;
finite sampling at large $n$ would not replace them.


\subsection{Full algebraic certificate data}\label{R17R:bys:data}
For each certificate $Q$, the rows below give the complete expansion
$Q(15+t)=\sum_j(a_j+b_j\sqrt2)t^j$. All components are positive
integers. Terms not listed are zero. The definitions
\eqref{R17R:bys:Delta}--\eqref{R17R:bys:Q} uniquely determine every entry.

\subsubsection*{$Q_0$}
{\small
\begin{longtable}{rrr}
\toprule
$j$ & $a_j$ & $b_j$\\\midrule
\endhead
0 & 5539319 & 5851125\\
1 & 1538841 & 1601205\\
2 & 160032 & 164262\\
3 & 7392 & 7488\\
4 & 128 & 128\\
\bottomrule
\end{longtable}
}

\subsubsection*{$Q_L$}
{\small
\begin{longtable}{rrr}
\toprule
$j$ & $a_j$ & $b_j$\\\midrule
\endhead
0 & 1651060960399963616 & 8016289336174125728\\
1 & 1328908915455442382 & 4667872835307781050\\
2 & 434419882217826919 & 1209326587928491801\\
3 & 78582022181739466 & 182923165973801116\\
4 & 8826335639173425 & 17800104822217471\\
5 & 644869171091976 & 1155413171724986\\
6 & 30839436440520 & 50021568072264\\
7 & 934712026752 & 1392637661904\\
8 & 16339090176 & 22622299136\\
9 & 125762560 & 163347456\\
\bottomrule
\end{longtable}
}

\subsubsection*{$Q_U$}
{\small
\begin{longtable}{rrr}
\toprule
$j$ & $a_j$ & $b_j$\\\midrule
\endhead
0 & 7170637559529179680 & 13224814771590028640\\
1 & 4466948423504945714 & 7613431526408488006\\
2 & 1226426152202737049 & 1948681949445081319\\
3 & 195035522999616566 & 291042244307310436\\
4 & 19818359697704655 & 27951583756608769\\
5 & 1335517300549112 & 1790060272065862\\
6 & 59723232087096 & 76439830672824\\
7 & 1709964417408 & 2098686307632\\
8 & 28456076544 & 33614928896\\
9 & 209781760 & 239305728\\
\bottomrule
\end{longtable}
}

\subsubsection*{$Q_T$}
{\small
\begin{longtable}{rrr}
\toprule
$j$ & $a_j$ & $b_j$\\\midrule
\endhead
0 & 44384753889664723582976 & 163829972812558332067840\\
1 & 248871109381232741711872 & 275651291623062572892160\\
2 & 174594778146145309503488 & 157824377880393453608960\\
3 & 58977918797802694082560 & 48881821816470068183040\\
4 & 12233470351217752522752 & 9661655319217785241600\\
5 & 1722080270122376626176 & 1318904414037146927104\\
6 & 172857575570866774016 & 129603023686161580032\\
7 & 12671998275168043008 & 9354233590479093760\\
8 & 683204653015953408 & 498367983599132672\\
9 & 26877236667973632 & 19422953845538816\\
10 & 752266603069440 & 539530281091072\\
11 & 14217523494912 & 10133484601344\\
12 & 162898378752 & 115499597824\\
13 & 855638016 & 603979776\\
\bottomrule
\end{longtable}
}

\subsection{Complete exact certificate program}\label{R17R:bys:code}
The following program is included in full. It uses only Python's
standard library and exact integer, rational, and quadratic-field
arithmetic. It reconstructs all four coefficient lists from their
mathematical definitions, verifies every component sign, obtains the
initial values independently from the binomial sum, verifies the
base enclosures, and checks all remaining centers. Assertions must be
enabled (run ordinary \texttt{python3}, without the \texttt{-O} option).
No external data file is required. Its output agrees with the archived
\texttt{results.txt}; the recurrence derivative and coefficient lists
were additionally checked by an independent symbolic expansion for
this reconstruction.

{\footnotesize
\begin{Verbatim}[fontsize=\footnotesize,breaklines,breakanywhere]
#!/usr/bin/env python3
"""Exact certificate for Brian Yi Sun's 2018 Conjecture 4.7."""
from fractions import Fraction
from math import comb

# Pairs represent a+b sqrt(2), with a,b integers.
def qm(x, y):
    a, b = x
    c, d = y
    return a*c + 2*b*d, a*d + b*c

def qpos(x):
    a, b = x
    if a >= 0 and b >= 0:
        return bool(a or b)
    if a <= 0 and b <= 0:
        return False
    return (a*a > 2*b*b) if a > 0 else (2*b*b > a*a)

# Low-to-high coefficient lists. Keep two integer components throughout.
def c(a, b=0):
    return [(a, b)]

def add(*args):
    return [tuple(sum(p[i][j] for p in args if len(p) > i) for j in (0, 1))
            for i in range(max(map(len, args)))]

def mul(*args):
    out = c(1)
    for p in args:
        z = [[0, 0] for _ in range(len(out) + len(p) - 1)]
        for i, u in enumerate(out):
            for j, v in enumerate(p):
                a, b = qm(u, v)
                z[i+j][0] += a
                z[i+j][1] += b
        out = list(map(tuple, z))
    return out

def sub(p, q):
    return add(p, mul(c(-1), q))

def A(x, e):
    return add(mul(c(192,128), x,x,x,x),
               mul(c(-288,-192), x,x,x),
               mul(c(192,102), x,x),
               mul(c(-519,-255),x), c(e*4096))

def B(x):
    return mul(c(64),x,x,x,x)

def evalp(p, t):
    return tuple(sum(co[j]*t**i for i, co in enumerate(p)) for j in (0, 1))

n = [(15,0), (1,0)]
x = [add(n, c(j)) for j in range(3)]
N = {e: [A(v,e) for v in x] for e in (-1,+1)}
D = [B(v) for v in x]

def delta(e, f, g):
    p, q, r = N[e][0], N[f][1], N[g][2]
    return add(mul(add(mul(c(7),n), c(13)),p,q,D[2]),
               mul(add(mul(c(-7),n),c(-15)),p,D[1],D[2]),
               mul(add(n,c(1)),*D),
               mul(sub(c(-3),n),p,q,r))

checks = {
    'A_- - B': sub(N[-1][0], D[0]),
    'lower propagation': delta(+1,-1,-1),
    'upper propagation': mul(c(-1),delta(-1,+1,+1)),
    'Turan separation': sub(mul(N[-1][1],N[-1][1],D[0],D[2]),
                            mul(N[+1][0],N[+1][2],D[1],D[1])),
}
print('Certificate polynomials in t = n - 15:')
for label, p in checks.items():
    while len(p) > 1 and p[-1] == (0,0):
        p.pop()
    assert all(a > 0 and b > 0 for a,b in p), label
    print(label, 'degree', len(p)-1, 'component minima',
          tuple(min(co[j] for co in p) for j in (0,1)))
    print('  coefficients [t^0, t^1, ...] =', p)

R = [-1,1,7]
for m in range(15):
    q,r = divmod((7*m+13)*R[m+2]-(7*m+15)*R[m+1]+(m+1)*R[m],m+3)
    assert r==0
    R.append(q)
assert R == [sum((Fraction(comb(m,k)*comb(m+k,k),2*k-1)
                 for k in range(m+1)), Fraction()) for m in range(len(R))]
print('Exact initial terms R_0 through R_17:', R)
for m in (15,16):
    t = m-15
    d = B(x[0])
    bn = evalp(d,t)[0]
    al, bl = evalp(N[-1][0],t)
    au, bu = evalp(N[+1][0],t)
    low = (R[m+1]*bn-R[m]*al, -R[m]*bl)
    high = (R[m]*au-R[m+1]*bn, R[m]*bu)
    assert qpos(low) and qpos(high)
    # Optional stronger check: lower bound each by replacing sqrt(2)
    # by its endpoint giving the least value on (7/5, 3/2).
    print(f'Base at {m}: lower gap pair = {low}, upper gap pair = {high}')
T = lambda m: R[m+1]**3*R[m-1] - R[m]**3*R[m+2]
for m in range(5,16):
    assert T(m)>0
print('Centers 5 through 15: Turan integer gaps =',[T(m) for m in range(5,16)])
assert T(4) == -4623352
print('Excluded center 4, outside original indexed tail:', T(4))
print('All exact checks passed.')
\end{Verbatim}
}
\endgroup


\clearpage

\part{Counterexamples and negative resolutions}

\noindent Four manuscripts refute five named assertions. Their constructive and asymptotic replacement theorems are retained along with the counterexamples. A negative resolution is not a positive proof of the original claim.


\section{A spectral counterfamily to Cazenave Conjecture 2}
\label{section:R11H}
\noindent{\small\textsf{ROUND 11 \enspace / \enspace Heisei \enspace / \enspace R11H}}\\[3pt]
\noindent\textbf{Target.} Cazenave Conjecture 2, identifier 20829.

\noindent\textbf{Editorial status.} Negative resolution: a counterexample, with the additional full arguments preserved. This refutes the specific adjacency--Laplacian inequality, not the different Cazenave problem considered in another round. The submitted smallest example is not claimed globally minimal.

\smallskip\hrule\medskip



\begingroup

\setcounter{equation}{0}
\renewcommand{\theequation}{\thesection.\arabic{equation}}
\renewcommand{\theHequation}{R11H.\arabic{equation}}
\paragraph{Overview.}
Cazenave's Conjecture 2 (corpus identifier \#20829) asserts
\(\lambda_2(A)+2\leq\mu_1(D-A)\) for every simple connected graph
on at least two vertices. For each \(k\geq6\), take two copies
of \(K_k\) minus one edge, and connect one endpoint of the
missing edge in each copy to a common new vertex. These connected
graphs violate the conjecture. We give exact rational bounds, an
independent matrix certificate for a 13-vertex example, and show
that the spectral gap on this family tends to the golden ratio.
No priority claim is made.


\subsection{Statement and counterexample}
Let \(A\) be the ordinary adjacency matrix, \(D\) the diagonal matrix
of vertex degrees, and \(L=D-A\). Let \(\lambda_2\) denote the second
largest eigenvalue of \(A\) and \(\mu_1\) the largest eigenvalue
of \(L\), each ordered over the reals with multiplicity. Cazenave's
Conjecture 2 (\#20829) in Section 7 of \cite{R11H:cazenave} states that
\[
  \lambda_2 \leq \mu_1-2
\]
for all finite simple undirected connected graphs of order at least
two. We use exactly those hypotheses and conventions.

For \(k\geq 3\), take two disjoint \(k\)-vertex cliques. In each clique
delete the edge between two distinguished vertices \(a_i,b_i\), where
\(i=1,2\). Finally add a single vertex \(c\) and edges \(ca_1,ca_2\).
Denote the resulting graph by \(G_k\). Its order is \(2k+1\);
it is simple and connected. The \(a_i\) and the \(k-2\) untouched
vertices in each clique have degree \(k-1\), each \(b_i\) has degree
\(k-2\), and \(c\) has degree two.

\begin{SourceRbbHtheorem}\label{R11H:thm:main}
For every \(k\geq6\),
\[
   \lambda_2(G_k)>k-\frac{13}{10},\qquad
   \mu_1(G_k)<k+\frac7{10}.
\]
In particular \(\mu_1(G_k)-\lambda_2(G_k)<2\), refuting the conjecture.
Moreover,
\[
  \lim_{k\to\infty}\bigl(\mu_1(G_k)-\lambda_2(G_k)\bigr)
  =\frac{1+\sqrt{5}}{2}.
\]
\end{SourceRbbHtheorem}

\subsection{Proof of the strict inequalities}
\paragraph{Adjacency.}
The principal submatrix after deleting \(c\) is the adjacency matrix
of two identical components \(K_k-a_ib_i\). Its top two eigenvalues
coincide, both equaling the top eigenvalue \(\rho\) of one component.
Interlacing above and below forces \(\lambda_2(G_k)=\rho\).
In a component, a positive eigenvector can be chosen constant with
value \(u\) on \(a_i,b_i\) and value \(v\) on the other \(k-2\)
vertices. This follows by averaging a positive Perron eigenvector
over the automorphisms preserving those classes. The equations are
\[
   \rho u=(k-2)v,\qquad
   \rho v=2u+(k-3)v.
\]
Thus \(\rho\) is the positive root of
\(p_k(t)=t^2-(k-3)t-2(k-2)\).
For \(t=k-13/10>0\),
\[
  p_k(t)=\frac{179-30k}{100}<0\qquad(k\geq6).
\]
Since the other root is negative, this proves \(\lambda_2>k-13/10\).

\paragraph{Laplacian.}
Within either clique block, order \(a_i,b_i\) first. Write \(J\)
for the all-ones matrix, \(e_a\) for the first coordinate vector,
and let \(P\) have leading block
\(\left(\begin{smallmatrix}0&1\\1&-1\end{smallmatrix}\right)\)
and zeros elsewhere. The diagonal block of \(L\) is \(kI-J+P\).
The \(c\)-diagonal entry is 2 and its column contains \(-e_a\)
in each clique block. We prove that \((k+s)I-L\) is positive
definite for \(s=7/10\). Set \(h=k-2\geq4\); both clique blocks
of the shifted matrix are \(C=sI+J-P\). The largest eigenvalue of
\(P\) is \(\gamma=(\sqrt5-1)/2<s\), hence \(C\) is positive
definite. The central Schur complement is
\[
    h+s-2(C^{-1})_{aa}.
\]
The untouched \(h\) vertices of \(C\) have block \(sI_h+J_h\).
Taking the Schur complement of that block changes the all-ones
term on the leading pair to \(\beta J_2\), where
\(\beta=s/(s+h)\); indeed the eliminated term is
\(hJ_2/(s+h)\). The resulting block is
\[
   \begin{pmatrix}s+\beta&\beta-1\\ \beta-1&s+1+\beta\end{pmatrix}.
\]
Set \(\delta=s^2+s-1+(2s+3)\beta\), its positive determinant.
The inverse formula gives \((C^{-1})_{aa}=(s+1+\beta)/\delta\).
Substituting \(s=7/10\) and multiplying the numerator
\((h+s)\delta-2(s+1+\beta)\) by \(1000(10h+7)\) yields
\[
  1900h^2-540h-15309
   =1900(h-4)^2+14660(h-4)+12931>0.
\]
The central Schur complement is positive, proving
\(\mu_1<k+7/10\). This completes the strict comparison without
approximating either eigenvalue.

\subsection{Limiting gap}
For every \(k\), the equal-coordinate Rayleigh test on one almost-clique
gives \(k-1-2/k\leq\rho\leq k-1\), with the upper bound following
from the maximum degree. Consequently \(\lambda_2=k-1+o(1)\).
The two clique blocks of \(L\) are bounded above by
\((k+\gamma)I\) in quadratic form. Combining their coordinates
into one vector of norm \(u\), and writing \(v\) for the absolute
value of the central coordinate, gives the quadratic-form bound
\[
   (k+\gamma)u^2+2\sqrt2\,uv+2v^2.
\]
The largest eigenvalue of the corresponding \(2\times2\) matrix
is at most \(k+\gamma+2/(k+\gamma-2)\): its increment over
\(k+\gamma\) multiplied by its increment over 2 is exactly 2.
For a lower bound, use a vector supported on the first clique,
with entries \(1,\gamma\) at the distinguished vertices and
\(-(1+\gamma)/h\) at each other vertex. Its sum vanishes, so
\(J\) contributes zero. Its distinguished coordinate pair is a
\(\gamma\)-eigenvector of the leading block of \(P\). The resulting
Rayleigh quotient is
\[
   k+\frac{\gamma(1+\gamma^2)}
     {1+\gamma^2+(1+\gamma)^2/h},
\]
which is \(k+\gamma+o(1)\). Thus
\(\mu_1-\lambda_2\to 1+\gamma\), as asserted. This does not
claim an optimal universal constant replacing 2.

\subsection{Exact replay and scope}
The smallest member of the proved range is \(G_6\), with thirteen
vertices and thirty edges. Number the vertices \(0,\ldots,12\).
It has all within-block edges in \(\{0,\ldots,5\}\) and
\(\{6,\ldots,11\}\) except \(01\) and \(67\), plus \(0\,12\)
and \(6\,12\). The following is an exact check of that particular
case using the unreduced Laplacian. Ordinary symmetric elimination
on \((669/100)I-L\), in label order, has the six successive
pivots below, repeated once, followed by \(59547021/380120900\):
\begin{center}
\renewcommand{\arraystretch}{1.3}
\begin{tabular}{c@{\quad}c@{\quad}c@{\quad}c@{\quad}c@{\quad}c}
\(\frac{169}{100}\)&\(\frac{269}{100}\)&\(\frac{3302909}{4546100}\)&
\(\frac{239361621}{330290900}\)&\(\frac{250822521}{346900900}\)&
\(\frac{87427807}{121170300}\)
\end{tabular}
\end{center}
All are positive, certifying \(\mu_1<669/100\); the exact
adjacency comparison \(p_6(47/10)=-1/100\) certifies
\(\lambda_2>47/10\), so the gap is strictly below \(199/100\).
We do not assert minimum counterexample order. The attached
standard-library verifier constructs actual vertex-edge matrices,
checks positivity using rational elimination and checks the
adjacency polynomial at its rational threshold for seven values
of \(k\). No finite run is an assumption of Theorem~\ref{R11H:thm:main}.

The cited version lists \#20829 among its unrefuted conjectures,
not only as a restricted-family problem. The bounded literature
search found settlements of other numbered conjectures in that
paper, including an exact prior settlement of Conjecture 5
(\#239534) \cite{R11H:prior}; an independent proof of that other
inequality is not submitted as a new resolution here. The linked
proof/refutation ZIP was not readable in this environment, so its
internal row for \#20829 could not be checked. This paper supplies
a full mathematical counterexample to the exact written assertion.
Independence of the construction is reported; worldwide priority
is not established.
\endgroup


\clearpage

\section{A power-64 counterexample to two gap-symmetry conjectures}
\label{section:R12T}
\noindent{\small\textsf{ROUND 12 \enspace / \enspace Taisho \enspace / \enspace R12T}}\\[3pt]
\noindent\textbf{Target.} Benfield--Lippard Conjectures 10.1 and 10.2.

\noindent\textbf{Editorial status.} Negative resolution: a counterexample, with the additional full arguments preserved. One exponent-$64$ construction refutes two numbered conjectures. The complete elementary Apery certificate and all bounds on omitted generators are preserved.

\smallskip\hrule\medskip



\begingroup
\newcommand{\N}{\mathbb N_0}
\newcommand{\B}{\mathbf B}
\setcounter{equation}{0}
\renewcommand{\theequation}{\thesection.\arabic{equation}}
\renewcommand{\theHequation}{R12T.\arabic{equation}}
\paragraph{Overview.}
Benfield and Lippard conjectured that when $k$ is a power of two, the stabilized exceptional offsets for sums of positive $k$th powers are symmetric about half their maximum. We disprove this at $k=64$. The shifted powers generate a numerical semigroup with four explicitly given generators. A compact Ap\'ery certificate and two residue decompositions produce a pair of gaps summing to the largest gap; they also show that twice the number of gaps minus one exceeds the largest gap by exactly $118657108319850903491736959827380951772312320$. The proof uses only elementary integer arithmetic after identifying the stabilized offsets. Neither a search over all residues modulo $2^{64}-1$ nor a theorem classifying symmetric semigroups is needed.


\subsection{Original statements and their semigroup form}
In \cite[\S10.1]{R12T:BL2}, Benfield and Lippard let $A_k=\max\B^k$ and $G_k=|\B^k|$, where $\B^k$ is their stabilized set of positive offsets $s$ for which $s+j$ cannot be expressed as a sum of exactly $j$ positive $k$th powers. Their Conjecture~10.1 asserts $A_k=2G_k-1$ for every power-of-two exponent $k\ge2$. Conjecture~10.2 asserts the corresponding gap symmetry: exactly one of $s$ and $A_k-s$ lies in $\B^k$, for $0\le s\le A_k$. We use different letters for $A_k,G_k$ to distinguish the two quantities from the generators below.

Put $S_k=\langle n^k-1:n\ge2\rangle$, where brackets mean all finite nonnegative integer sums, including zero. For a fixed $s>0$, a representation of $s+j$ by $j$ positive powers gives $s=\sum_{i=1}^j(n_i^k-1)\in S_k$. Conversely, if $s\in S_k$ with a representation using $r$ shifted terms, then $s+j$ is representable for every $j\ge r$ by padding with $j-r$ ones. Thus the permanent offset exceptions are precisely $\B^k=\N\setminus S_k$. To connect this pointwise argument to the stabilized set in \cite[\S10.2]{R12T:BL2}, its Theorem~1.1, credited to Zenkin, makes the offset exception sets finite from some $j$ on; padding makes them decreasing, so they eventually equal their intersection. For a cofinite additive submonoid $S$ of $\N$, write $F(S)$ for its largest gap and $g(S)$ for its gap count. Then $A_k=F(S_k)$ and $G_k=g(S_k)$.

\begin{SourceRbcTtheorem}\label{R12T:main}
At $k=64=2^6$, one has
\[
S_{64}=\langle 2^{64}-1,3^{64}-1,5^{64}-1,17^{64}-1\rangle.
\]
There are complementary gaps $P,Q\in\B^{64}$ with $P+Q=A_{64}$, and
\[
2G_{64}-1-A_{64}=118657108319850903491736959827380951772312320>0.
\]
In particular, Conjectures~10.1 and~10.2 of \cite{R12T:BL2} are both false.
\end{SourceRbcTtheorem}

\subsection{A three-generator certificate}
Let
\[
a=(2^{64}-1)/85,\qquad b=(3^{64}-1)/85,\qquad
c=(5^{64}-1)/17,\qquad T=\langle a,b,c\rangle.
\]
These quotients are integers and $\gcd(a,b)=1$. Use the parameters
\[
(\ell,h,w,\eta)=(18766804522483320,\ 12,\ 743739614142711,\ 11).
\]
Direct substitution gives $0<w<\ell$, $0<\eta<h$ and the three exact identities
\begin{align*}
\ell h-w\eta&=a=217020518514230019,\\
\ell b-\eta c&=ua,\qquad u=1876955259190588010127036672,\\
hc-wb&=va,\qquad v=1624801557131671352222750464.
\end{align*}
Let $L$ consist of integer points in $[0,\ell)\times[0,h)$ except those in $[\ell-w,\ell)\times[h-\eta,h)$, and let the weight of $(x,y)$ be $xb+yc$. We show that these $a$ weights are the Ap\'ery set: the least element of $T$ in each residue modulo $a$.

Indeed, from any nonnegative pair outside $L$, subtract $(\ell,-\eta)$ when $x\ge\ell$, subtract $(-w,h)$ when $y\ge h$, or subtract $(\ell-w,h-\eta)$ when the point is in the removed corner. In each indicated case both resulting coordinates are nonnegative. The weight decreases by $ua$, $va$ or $(u+v)a$, preserving its residue. This process must terminate in $L$. Every residue had a nonnegative starting pair since $b$ is invertible modulo $a$. Therefore the $a$ points in $L$ represent every residue exactly once. For any representation by $a,b,c$, drop its $a$ terms and perform the reduction on its $b,c$ coefficients. The sole weight in $L$ with that residue is at most the original number, proving minimality.

The two possible maximum weights occur at the top-right points of the two rectangles in the decomposition of $L$. Write
\[
W_1=(\ell-1)b+(h-\eta-1)c,\qquad W_2=(\ell-w-1)b+(h-1)c.
\]
We have $W_2>W_1>a$. Since $a\in T$, in any given residue the gaps stop just before the minimal weight, so
\[
F(T)=W_2-a=1078836071896497319661834263595448478109893117.
\]
There are complementary gaps $X=W_1-a$ and $Y=W_2-W_1$: $X$ is one step below its own minimum, and $Y=\eta c-wb$ is $ua$ below the weight of $(\ell-w,0)\in L$, which is the minimum in its residue. Both numbers are positive and $X+Y=F(T)$.

For the count, the least weight $m_r$ in residue $0\le r<a$ leaves $(m_r-r)/a$ gaps there. Summing coordinates in the large rectangle and subtracting the missing corner yields
\begin{align*}
S_x&=\frac{h\ell(\ell-1)-\eta w(2\ell-w-1)}2,\\
S_y&=\frac{\ell h(h-1)-w\eta(2h-\eta-1)}2,\\
g(T)&=\frac{bS_x+cS_y-a(a-1)/2}{a}
=540116018938365429851456760973179421124195455.
\end{align*}
These calculations yield $2g(T)-1-F(T)=1395965980233540041079258350910364138497792$.

\subsection{Residue extensions and the omitted generators}
\begin{SourceRbcTlemma}\label{R12T:glue}
Let $V\subset\N$ be a cofinite additive submonoid with a gap. If $q\ge2$, $d\in V$ and $\gcd(q,d)=1$, put $U=qV+d\N$. Then
\[
F(U)=qF(V)+(q-1)d,\qquad
g(U)=qg(V)+(q-1)(d-1)/2 .
\]
If $x,y$ are complementary positive gaps of $V$, then $qx$ and $qy+(q-1)d$ are complementary positive gaps of $U$.
\end{SourceRbcTlemma}
\begin{proof}
Represent any integer uniquely by $qz+jd$ with $0\le j<q$ and integer $z$. It lies in $U$ exactly when $z\in V$: a nonnegative $d$ coefficient can be reduced modulo $q$ by absorbing copies of $d$ into $V$. The largest gap in class $j$ comes from $z=F(V)$, proving the formula for $F$. The class has $g(V)$ gaps with $z\ge0$ and $\lfloor jd/q\rfloor$ additional nonnegative gaps with $z<0$. The remainders of $jd$ form a permutation of $0,\ldots,q-1$, so these floors sum to $(q-1)(d-1)/2$. Complementary gaps lift in classes $j=0$ and $j=q-1$.
\end{proof}

Let $H=\langle5a,5b,c\rangle=5T+c\N$. Since $c\in T$ and $\gcd(5,c)=1$, Lemma~\ref{R12T:glue} gives $F(H)=5F(T)+4c$. Set $D=17^{64}-1$. The explicit values in the previous subsection give
\[
17F(H)<6^{64}-1,\qquad F(H)<D,\qquad 17F(H)+16D<34^{64}-1.
\]
In particular $D\in H$ and $\gcd(17,D)=1$. Apply Lemma~\ref{R12T:glue} again to $U=17H+D\N$, which has Frobenius number $17F(H)+16D$. Its four generators are all in $S_{64}$. The generator for base $4$ is a multiple of $2^{64}-1$, since $4^{64}-1=(2^{64}-1)(2^{64}+1)$. For the remaining bases, if $n\ge6$ and $17\nmid n$, Fermat's little theorem implies $17\mid n^{64}-1$. This generator divided by $17$ exceeds $F(H)$ by the first bound, so belongs to $H$. The only remaining bases other than 17 are multiples of 17 at least 34, and their generators exceed $F(U)$ by the third bound. Thus $S_{64}=U$.

The two lifts from Lemma~\ref{R12T:glue} give
\[
P=85X,\qquad Q=17(5Y+4c)+16D,
\]
positive gaps of $U$ summing to $F(U)$. Under a single lift the defect $2g-1-F$ is multiplied by $q$, so the defect for $U$ is $85$ times that for $T$. This proves Theorem~\ref{R12T:main}.

\subsection{Discovery, checks, and version scope}
A residue shortest-path search suffices for small $k$ but has $2^k-1$ residues. We tested a mixed-radix construction with an increasing list of generators for which the successive gcd drops and the requisite multiple of each new generator is generated by earlier ones. Exact tests found symmetric semigroups at $k=2,4,8,16,32$. At $64$ the multiple required on adding the base $5$ fails the earlier-generator membership test. This failure alone is not a disproof of symmetry. The certificate above uses the resulting three-generator problem instead. Ap\'ery sets and $L$-shapes for three generators are classical \cite{R12T:AgMar}; their explicit realization here, the four-generator truncation and the application to the conjectures are the specific findings. We give a direct certificate proof rather than import an $L$-shape theorem.

The companion verifier checks the displayed identities, coordinate sums, positivity, omissions and exact final invariants using integer arithmetic. A second implementation independently checks the nonsymmetric three-generator core by Herzog's three-generator symmetry criterion \cite{R12T:Herzog}. The logical proof does not require either script. The exponent-14 counterexample to the older even-exponent identity (v1, Conjecture~4.1) is separate; it is not a counterexample to the revised power-of-two statements. The v1 symmetry statement uses a power-of-two superscript on $\B$ but an even subscript on its center; v2 uses a power of two consistently. The version-specific primary sources are \cite{R12T:BL1,R12T:BL2}. A bounded literature search found no prior exponent-64 disproof; that does not establish worldwide priority or external peer review. No claim is made about asymptotics or the other conjectures in \cite{R12T:BL2}.
\endgroup


\clearpage

\section{A rational cyclic reversal of a cosine rearrangement inequality}
\label{section:R22M}
\noindent{\small\textsf{ROUND 22 \enspace / \enspace Meiji \enspace / \enspace R22M}}\\[3pt]
\noindent\textbf{Target.} Matejicka 2019, Section 3(a).

\noindent\textbf{Editorial status.} Negative resolution: a counterexample, with the additional full arguments preserved. The strict rational counterexample concerns the printed cosine-of-a-power formula in Section 3(a), not the two other inequalities in that source.

\smallskip\hrule\medskip



\begingroup

\setcounter{equation}{0}
\renewcommand{\theequation}{\thesection.\arabic{equation}}
\renewcommand{\theHequation}{R22M.\arabic{equation}}
\paragraph{Overview.}
For a question on cosine values at permuted powers, a cycle on thirteen distinct rational inputs reverses the proposed inequality. The reversal does not arise from the equality case of the identity permutation. Integer comparisons and two alternating Taylor bounds give an exact certificate.


\subsection*{The claim and its reversal}
The open-problem section of Matej\'i\v{c}ka's 2019 paper \cite[\S3(a)]{R22M:M19} asks whether, for $0<x_i\le\pi/2$ and a permutation $(y_i)$ of $(x_i)$, the diagonal sum $\sum_i\cos(x_i^{x_i})$ is smaller than $\sum_i\cos(x_i^{y_i})$. Even replacing the strict inequality by a weak one cannot make this true.
Put
\[
(a_1,\dots,a_{13})=(1,2,4,6,8,10,12,14,16,18,20,22,24),
\qquad a_0=a_{13}.
\]
Take $x_i=a_i/24$ and $y_i=a_{i-1}/24$. All $x_i$ lie in $(0,1]$ and $y$ is a cyclic permutation. We prove
\[
\frac{30467}{10^6}<
D:=\sum_{i=1}^{13}\bigl(\cos(x_i^{x_i})-\cos(x_i^{y_i})\bigr)
<\frac{30651}{10^6}.
\]
This does not contradict the positive result for a swap of two inputs \cite{R22M:M20}. Pairwise comparisons involving a pair and its own swap do not control every cyclic reassignment.

\subsection*{An integer certificate}
Let $M=10^5$. Define $q(a,b)=\lfloor M(a/24)^{b/24}\rfloor$. For positive integers $a,b$, the assertion $q(a,b)=k$ is certified exactly by
\[
 k^{24}24^b\;\le\; M^{24}a^b\;<\;(k+1)^{24}24^b.
\]
Define the partial cosine sum $P_j(t)=\sum_{r=0}^j (-1)^r t^{2r}/(2r)!$. For $k=q(a,b)$, monotonicity of cosine on $[0,1+1/M]$ and its alternating Taylor series give
\[
 L(k):=P_5((k+1)/M)\ \le\ \cos((a/24)^{b/24})\ \le\ U(k):=P_4(k/M).
\]
Successive series terms decrease in absolute value on this interval since $(1+1/M)^2<2$.

\subsection*{Reproducible calculation}
Each row below brackets both power arguments before applying the displayed cosine bounds.
\begin{center}
\begin{tabular}{rrrr}
\toprule
$a_i$ & $a_{i-1}$ & $p_i=q(a_i,a_i)$ & $s_i=q(a_i,a_{i-1})$\\
\midrule
1&24&87597&4166\\
2&1&81295&90164\\
4&2&74183&86129\\
6&4&70710&79370\\
8&6&69336&75983\\
10&8&69435&74690\\
12&10&70710&74915\\
14&12&73021&76376\\
16&14&76314&78936\\
18&16&80592&82548\\
20&18&85904&87219\\
22&20&92333&93005\\
24&22&100000&100000\\
\bottomrule
\end{tabular}
\end{center}
With the tabulated integers, rational additions and multiplications give
\[
30467 < 10^6\sum_i \{L(p_i)-U(s_i)\}\le 10^6D
\le 10^6\sum_i \{U(p_i)-L(s_i)\}<30651.
\]
The companion file \texttt{check.py} checks the integer power bounds, uses exact fractions for the Taylor sums, and reports the outward rounded bounds. No floating-point assumption is used.

\textbf{Scope.} The refutation concerns the displayed cosine-of-a-power inequality in \cite[\S3(a)]{R22M:M19}, arising in the study of the conjectures of \"Ozban \cite{R22M:O}. It does not decide the two other inequalities in the 2019 open-problem section. The source's conjecture-number references are not consistent between its introduction and last section, so we identify our target by its printed formula, not its number. This note does not assert worldwide priority; its certificate was frozen as a nonwinning Round 22 result.
\endgroup


\clearpage

\section{Mode drift in the one-letter plactic centralizer}
\label{section:R24K}
\noindent{\small\textsf{ROUND 24 \enspace / \enspace Keio \enspace / \enspace R24K}}\\[3pt]
\noindent\textbf{Target.} Sagan--Wilson Conjecture 7.3(d).

\noindent\textbf{Editorial status.} Negative resolution: a counterexample, with the additional full arguments preserved. The mode-location conjecture is refuted. The separate asymptotic theorem for every mode is also included; no unrestricted log-concavity assertion is inferred.

\smallskip\hrule\medskip



\begingroup

\setcounter{equation}{0}
\renewcommand{\theequation}{\thesection.\arabic{equation}}
\renewcommand{\theHequation}{R24K.\arabic{equation}}
\paragraph{Overview.}
Sagan and Wilson conjectured a central mode for the binomial coefficients of the counting polynomial of the plactic centralizer of the letter $1$. The first failure occurs at $n=10$. In fact every maximizing index is asymptotic to $n/(2\log 2)$, so that even the predicted linear scale constant is incorrect. The proof embeds almost all shorter packed words into the centralizer by inserting a sublinear suffix of ones. On the exponential scale this transfers the alphabet-size concentration of ordered set partitions. No log-concavity assertion is established.


\subsection{The counterexample}
Write $P(w)$ for the (row) insertion tableau of a word on the positive integers.
Let $c_{n,m}$ count $w\in[m]^n$ for which $P(w1)=P(1w)$.
For $n\ge2$ the coefficients $a_{n,k}$ defined by
\[
 c_{n,m}=\sum_{k=0}^{n-1}a_{n,k}\binom{m}{k}
\]
are nonnegative, with $a_{n,0}=0$, and the coefficients $a_{n,k}$ for $1\le k<n$ are positive \cite[Thm.~4.2]{R24K:SZ}. The assertion of \cite[Conj.~7.3(d)]{R24K:SW} is that $\lceil n/2\rceil$ maximizes these coefficients.
\begin{SourceRceKtheorem}
At $n=10$ the unique maximum is $a_{10,6}=88704$, whereas $a_{10,5}=75095$. All $2\le n<10$ satisfy the proposed mode condition.
\end{SourceRceKtheorem}
\begin{proof}
The first row of $P(w)$ contains only ones precisely when $w$ belongs to this centralizer \cite[Prop.~2.5]{R24K:SZ}. Under RSK, deleting the all-ones top row of the insertion tableau and subtracting one from the other entries gives an arbitrary semistandard tableau of the remaining shape and alphabet $[m-1]$. The recording tableau is unrestricted. Thus, writing $\mu=(\lambda_2,\lambda_3,\dots)$, for $m\ge1$ one has
\[
 c_{n,m}=\sum_{\lambda\vdash n}
 \frac{n!}{\prod_{s\in\lambda}h_{\lambda}(s)}
 \prod_{(i,j)\in\mu}\frac{m-1+j-i}{h_{\mu}(i,j)} .
\]
Here $h_\lambda(i,j)$ is one plus the number of boxes strictly right or strictly below in the same row or column; the formula uses the hook-length and hook-content evaluations. For $n=10$ and $0\le m\le9$ the row is
\[
0,1,252,5670,50820,275010,1086090,3456369,9398925,22677445.
\]
Newton differences give $a_{n,k}=\sum_{j=0}^k(-1)^{k-j}\binom{k}{j}c_{n,j}$. For $1\le k<10$ the values are
\[
1,250,4917,29648,75095,88704,43414,4861,1 .
\]
For $n=2,\ldots,9$, the maximizing sets are, respectively,
$\{1\},\{1,2\},\{2\},\{3\},\{3\},\{4\},\{4\},\{5\}$.
An independent small integer recurrence for readers wishing to avoid hooks starts with the empty first row of multiplicity one. For each new letter $x\in[m]$ replace the leftmost entry strictly greater than $x$ by $x$, or append $x$ if none exists. Aggregate multiplicities of identical rows and retain only all-ones final rows. This computes precisely the same counts.
\end{proof}

\subsection{Why the peak drifts}
A packed word with alphabet size $k$ uses exactly $[k]$. Write $b_{n,k}$ for the number of length-$n$ packed words in this centralizer and $F_{n,k}$ for the analogous unrestricted count. Any centralizing nonempty word contains $1$. Relabeling its other letters in order yields \cite[Thm.~4.3]{R24K:SZ}
\[
 c_{n,m}=\sum_{k=1}^n b_{n,k}\binom{m-1}{k-1}.
\]
Thus $b_{n,k}=a_{n,k-1}+a_{n,k}$ by Pascal's identity, with coefficients off their ranges zero. In particular, $a_{n,k}\le b_{n,k}\le F_{n,k}$, and $\sum b_{n,k}=2\sum a_{n,k}$.

\begin{SourceRceKlemma}\label{R24K:fixed}
Let $F_N=\sum_kF_{N,k}$, and let $B_n=\sum_k b_{n,k}$. Then
$\log B_n=\log F_n+o(n)$.
\end{SourceRceKlemma}
\begin{proof}
For $r=\lceil n^{2/3}\rceil$, $N=n-r$, all but $o(F_N)$ of the packed words of length $N$ have longest weakly increasing subsequence of length at most $r$. To see this, mix uniform words in $[m]^N$ with weights $m^N 2^{-m-1}/F_N$ and then pack by replacing each distinct symbol by its rank. This is a probability law inducing a uniform packed word: each specific packed word with $k$ letters has $\binom{m}{k}$ preimages in $[m]^N$ and $\sum_{m\ge k}\binom mk2^{-m-1}=1$. The probability of $m<r$ is at most $r^N/F_N=o(1)$ because $F_N\ge N!$. For $m\ge r$, a union bound gives probability of even a length-$r$ weakly increasing subsequence at most
\[
 \binom Nr\binom{m+r-1}r m^{-r}\ \le\ (2e^2N/r^2)^r=o(1).
\]
Packing preserves weak inequalities. Increase each letter of a good packed word by one, then append $r$ ones. Before the suffix is inserted, the first row has length at most $r$ by Schensted's theorem and contains no ones. Each inserted one replaces one non-one as long as any remains. Hence the final word centralizes $1$; it is packed, and deleting the suffix and decreasing all letters recovers the input. This proves $B_n\ge(1-o(1))F_N$.

To compare lengths, any length-$j$ packed word is obtained from a length-$(j-1)$ packed word by appending an existing letter, or by first skipping one alphabet rank and appending the missing rank. There are at most $2j$ choices. Consequently $F_n\le(2n)^r F_N$, where $r\log(2n)=o(n)$. Combine the bounds with $B_n\le F_n$.
\end{proof}
\begin{SourceRceKtheorem}
Let $\alpha=1/(2\log2)$. For every $\varepsilon>0$, there is $c_\varepsilon>0$ such that for all sufficiently large $n$,
\[
 \frac{\sum_{|k/n-\alpha|>\varepsilon} a_{n,k}}{\sum_k a_{n,k}}
 \le e^{-c_\varepsilon n}.
\]
The same is true of $b_{n,k}$. Every maximizing index of either row is $\alpha n+o(n)$.
\end{SourceRceKtheorem}
\begin{proof}
A packed word is an ordered partition of labeled positions into nonempty letter classes. Thus $F_n(y)=\sum_k F_{n,k}y^k$ satisfies
\[
 \sum_{n\ge0} F_n(y)\frac{z^n}{n!}=\frac1{1-y(e^z-1)} .
\]
At fixed $y>0$ the unique pole of smallest modulus is the simple pole at $\rho(y)=\log(1+1/y)$; all other poles are $\rho(y)+2\pi i j$, $j\ne0$. Subtracting the principal part and applying Cauchy's estimate on a circle between the first two pole moduli gives
\[
 F_n(y)/n!\sim ((1+y)\rho(y)^{n+1})^{-1}.
\]
Accordingly $n^{-1}\log(F_n(e^t)/F_n(1))\to
L(t)=\log(\rho(1)/\rho(e^t))$, with $L(0)=0$, $L'(0)=\alpha$.
For the upper tail $k/n>\alpha+\varepsilon$, the log of its fraction among all packed words, divided by $n$, has limsup at most $L(t)-t(\alpha+\varepsilon)<0$ on choosing $t>0$ sufficiently small. Use $t<0$ for the lower tail. The resulting exponential bound transfers to $a$ and $b$ by their coefficientwise upper bounds and Lemma~\ref{R24K:fixed}. A maximizing index has fraction at least $1/(n+1)$ and therefore cannot remain in either tail.
\end{proof}

The limiting constant is approximately $0.72134752$. In particular, for all sufficiently large $n$ the proposed central index is not a mode. These arguments decide neither the log-concavity of $a_{n,k}$ in Conjecture~7.3(c) nor the log-concavity of $b_{n,k}$ asked by Sagan and Zhao. No literature-priority claim is made.
\endgroup

\clearpage\appendix
\section{Embedded exact certificate programs}
These listings supply proof-relevant finite computations for the sections identified below. Regression-only programs are explicitly labeled. The mathematical justification for extending any finite certificate to an infinite domain remains in the corresponding proof.
\Needspace{9\baselineskip}
\subsection{Guo--Schlosser arithmetic regression}
\noindent Related to Section~\ref{section:R18T}. This is a finite arithmetic regression only; it does not certify the all-level cyclotomic jet argument. Python standard library.

\begin{Verbatim}[fontsize=\scriptsize,breaklines,breakanywhere]
#!/usr/bin/env python3
"""Finite check for Conjecture 6.1. Uses only Python's standard library."""
from math import isqrt

def isp(p):
    return p >= 2 and all(p % x for x in range(2, isqrt(p)+1))

def split(x, p):
    v = 0
    while x % p == 0:
        x //= p
        v += 1
    return v, x

def sum_mod(p, r, cutoff, m):
    ans = 1
    unit = 1
    v = 0
    for k in range(1, (p**r-1)//cutoff + 1):
        a, num = split(3*k-2, p)
        b, den = split(3*k, p)
        v += a-b
        assert v >= 0
        unit = (unit * num * pow(den,-1,m)) % m
        ans = (ans + (6*k+1) * pow(p, 3*v, m) * pow(unit, 3, m)) % m
    return ans

def gamma23(p, prec):
    m = p**prec
    n = (2*pow(3,-1,m)) % m
    out = (-1)**n
    for k in range(1, n):
        if k % p:
            out = out*k % m
    return out % m

count = 0
for r, max_p in [(1,79), (2,31), (3,13)]:
    for p in range(7, max_p+1, 6):
        if not isp(p):
            continue
        m = p**(3*r)
        g = gamma23(p, 2*r)
        for cutoff in (1,3):
            lo = sum_mod(p,r-1,cutoff,m)
            hi = sum_mod(p,r,cutoff,m)
            assert lo % p**(r-1) == 0
            rhs = (p * g**3 * lo) % m
            assert hi == rhs, (p,r,cutoff,hi,rhs)
            print(f"p={p:2} r={r} cutoff={cutoff} modulus={m:10} residue={hi}")
            count += 1
print(f"Passed {count} finite congruence checks. Not a proof.")
\end{Verbatim}
\Needspace{9\baselineskip}
\subsection{Cubic Franel symbolic and rational certificates}
\noindent Related to Section~\ref{section:R19M}. The orthant positivity, low-degree identities, and rational Gaussian bounds are exact certificates. The final bounded coefficient loop is a separate regression. Requires Python and SymPy.

\begin{Verbatim}[fontsize=\scriptsize,breaklines,breakanywhere]
#!/usr/bin/env python3
"""Exact symbolic certificate and direct integer regression checks."""
from fractions import Fraction as F
from math import comb, factorial
import sympy as s

N,k,d,x,y,n=s.symbols("N k d x y n")
A=N-(k+d)/2
B=N-(k-d)/2
c=(k*(2*N-k)+d*d)/4
t=N*d/2
P=s.Poly(s.expand(N**3*((A*(A-1))**3+(B*(B-1))**3)/2
                   -(N-1)**3*(A*B)**3),d)
H=s.Poly(s.expand(c**3+3*c*t*t-d*(3*c*c*t+t**3)),d)
assert H.nth(0)==s.expand((k*(2*N-k)/4)**3)
R=s.Poly(s.expand(P.as_expr()-P.nth(0)/H.nth(0)*H.as_expr()),d)
assert R.degree()==6 and all(R.nth(j)==0 for j in (0,1,3,5))
a,b,e=[s.cancel(R.nth(j)) for j in (2,4,6)]
U,V,W=[s.Poly(s.cancel(32*k*k*z/N),N,k).as_expr() for z in (a,-b,-e)]
L=N+1
T=(k+1)*(2*N-k+1)
J=s.expand(4*L**2*U-(2*L*T+3*L**2)*V-2*T**2*W)
print("Symbolic positivity certificate: degree, terms, least nonzero coefficient")
for K in range(4,10):
    region={k:K,N:10+y} if K<9 else {k:9+x,N:10+x+y}
    out=[]
    for obj in (V,W,J):
        p=s.Poly(obj.subs(region,simultaneous=True),x,y)
        assert p.nth(0,0)>0 and all(z>0 for z in p.coeffs())
        out.append((p.total_degree(),len(p.terms()), min(p.coeffs())))
    print(K,*out)

def falling_binom(n,r):
    return s.prod(n-i for i in range(r))*s.Rational(1,factorial(r))

small=[6*n,3*n*(n-1)*(n*n-9*n+6)/4,
       n*(n-1)**3*(2*n**3+n*n+20*n-24)/6]
for q,target in enumerate(small,1):
    ans=sum((falling_binom(n+1,j)*falling_binom(n-1,q-j))**3
            -(falling_binom(n,j)*falling_binom(n,q-j))**3
            for j in range(q+1))
    assert s.expand(ans-target)==0
    assert min(s.Poly(target.subs(n,9+x),x).coeffs())>0
assert small[1].subs(n,8)==-84
print("Three small-degree identities and sharp-boundary check: PASS")

t0=F(1269,100)
g=F(157,50)
assert 4*(F(1,5)-F(1,3*5**3))-F(1,239)>g/4
assert g*g*F(85,66)>t0
assert sum(F(3**j,factorial(j)) for j in range(9))>20
assert 400*(1+F(345,1000)+F(345,1000)**2/2)>550
assert (2*t0-15)**2>105 and (25-15)**2<105
B=lambda r: t0**r/F(550*999999)
assert 2*((t0-1)/550+B(1))<F(1,20)
assert F(6,9)*F(20,19)*((t0**3/3-t0**2)/550+B(3)/3)<F(3,4)
assert 2*(F(1,24)+B(3)/15)<F(19,100)
print("Rational inequalities for the Gaussian comparison: PASS")

def binom(n,k):
    return comb(n,k) if 0<=k<=n else 0
checks=0
for n0 in range(9,151):
    v=[[binom(n0+j,i)**3 for i in range(2*n0+1)] for j in (-1,0,1)]
    for k0 in range(2*n0+1):
        u=sum(v[2][j]*v[0][k0-j]-v[1][j]*v[1][k0-j]
              for j in range(k0+1))
        assert u>0 if 0<k0<2*n0 else u==0
        checks+=1
print(f"Direct integer regression, n=9..150: {checks} coefficients PASS.")
\end{Verbatim}
\Needspace{9\baselineskip}
\subsection{Complete overcolored-partition finite certificate}
\noindent Related to Section~\ref{section:R21H}. This is the complete primary certificate through every index required by the stated modular-form bound, at most 1795. The last two rows concern an earlier version and are separate from the maintained nine-case conjecture. Python standard library.

\begin{Verbatim}[fontsize=\scriptsize,breaklines,breakanywhere]
#!/usr/bin/env python3
"""Exact finite certificate using Euler's pentagonal product, modulo 5."""
from math import gcd
from fractions import Fraction
from hashlib import sha256

cases=[(5,2,2,10),(4,4,5,10),(1,2,2,10),(3,1,6,10),(4,1,5,10),(5,3,2,20),(4,5,5,20),(1,3,2,20),(4,1,9,20),(2,2,1,10),(4,4,4,20)]
def raw(r,s,k): return (-2*s*k,(3*s-2*r)*k,(r-s)*k)
def overcolored(r,s,k,n):
    a=[1]+[0]*n
    for j in range(1,n+1):
        for _ in range((r if not j%2 else s)*k):
            for i in range(n,j-1,-1): a[i]+=a[i-j]
            for i in range(j,n+1): a[i]+=a[i-j]
    return a
def euler(a,d):
    L=len(a)-1
    b=a[:]
    h=1
    while d*h*(3*h-1)//2<=L:
        for z in (-1,1):
            v=d*h*(3*h+z)//2
            for j in range(v,L+1): b[j]+=(-1 if h%2 else 1)*a[j-v]
        h+=1
    return [i%5 for i in b]
for r,s,k,t in cases:
    E=tuple(i%25 for i in raw(r,s,k))
    S=sum(d*e for d,e in zip((1,2,4),E))
    P=sorted({(a*a*t+((a*a-1)//24)*S)%25 for a in range(600) if gcd(a,600)==1})
    assert P==sorted({t,25-t})
    assert S%25==0
    assert 20%(600//gcd(24*(-24*t-S),600))==0
    bound=Fraction(36*sum(E),24)-Fraction(S,600)-Fraction(min(P),25)
    B=bound.numerator//bound.denominator
    L=25*B+max(P)
    a=[1]+[0]*L
    for d,e in zip((1,2,4),E):
        for _ in range(e): a=euler(a,d)
    output=[a[25*n+j] for n in range(B+1) for j in P]
    original=overcolored(r,s,k,150)
    assert not any(original[i]%5 for j in P for i in range(j,151,25))
    print((r,s,k),t,E,S,P,B,L,len(output),sum(x!=0 for x in output),original[t],sha256(bytes(a)).hexdigest())
    assert not any(output)
\end{Verbatim}
\Needspace{9\baselineskip}
\subsection{Convexity and sharp-point arithmetic for the Huygens bound}
\noindent Related to Section~\ref{section:R22K}. Exact rational checks reproduce the finite polynomial inequalities and the certified minimizer enclosure. Infinite tails are bounded in the proof. Python standard library.

\begin{Verbatim}[fontsize=\scriptsize,breaklines,breakanywhere]
"""Fraction-only verification of all finite certificates, no third-party dependencies."""
from fractions import Fraction as F
from math import comb, factorial
def coeffs(n):
    e=[1]
    for j in range(1,n+1):
        e.append(-sum(comb(2*j,2*i)*e[i] for i in range(j)))
    return [F(abs(v),factorial(2*j+1)) for j,v in enumerate(e)]
M=120
a=coeffs(M+2)
c=F(3,40)
n=[0 if i%2 else 2*a[i+2] for i in range(M+1)]
d=a[:M+1]
b=[ni-c*di for ni,di in zip(n,d)]
# Bernstein positivity under y=(5/2)t/(1+t)
p=[(k+2)*(k+1)*b[k+2] for k in range(4)]
w=[sum(p[k]*F(5,2)**k*comb(3-k,j-k) for k in range(j+1)) for j in range(4)]
assert min(w)>0
print('Transformed cubic, ascending coefficients:',*[str(v) for v in w],sep='\n  ')
assert F(4,25)*F(13,17)*(1-F(1,3**17))>F(3,25)
assert 1/(1-F(1,3**13))<F(101,100)
assert F(6,25)>c*(1+F(7,5)*F(101,100))
y=F(5,4)
S=sum(b[k]*y**k for k in range(9))
r=y*F(4,9)
T=F(8,3)*F(4,9)**2*r**9/((1-r)*23)
assert S < F(-59,100000) and T<F(26,100000)
print('F_c(5/4) degree-eight sum < -59/100000; upper tail < 26/100000.')
# Exact polynomial evaluations/tail bounds near the point in Conjecture 4.1.
L=F(112840812792,10**11); U=F(112840812794,10**11)
r=F(26,45)
err=3*r**(M+1)*((M+1)/(1-r)+r/(1-r)**2)
def val(co,y,order=0):
    if order: co=[i*co[i] for i in range(1,len(co))]
    acc=F(0)
    for v in reversed(co): acc=acc*y+v
    return acc
def bracket(co,l,u,order=0):return (val(co,l,order),val(co,u,order)+err)
def sign_dg(x):
    y=x*x
    nl,nu=bracket(n,y,y);dl,du=bracket(d,y,y)
    n1l,n1u=bracket(n,y,y,1);d1l,d1u=bracket(d,y,y,1)
    return n1l*dl-nu*d1u,n1u*du-nl*d1l
assert 1<L*L<U*U<F(13,10)
assert sign_dg(L)[1]<0<sign_dg(U)[0]
nl,nu=bracket(n,L*L,U*U);dl,du=bracket(d,L*L,U*U)
assert F(745497347,10**10)<nl/du<=nu/dl<F(745497348,10**10)<c
print('Endpoint derivative signs: negative, positive.')
print('Certified x*: (1.12840812792, 1.12840812794)')
print('Certified lambda: (0.0745497347, 0.0745497348)')
print('All finite rational checks passed.')
\end{Verbatim}
\Needspace{9\baselineskip}
\subsection{Global inverse-Gudermannian coefficient certificate}
\noindent Related to Section~\ref{section:R22S}. The script reconstructs the derivative minorant and all coefficients of the transformed polynomial using exact arithmetic. Python standard library.

\begin{Verbatim}[fontsize=\scriptsize,breaklines,breakanywhere]
#!/usr/bin/env python3
"""Exact certificate for Sun--Chen, Conjecture 7. Python standard library only."""
from fractions import Fraction as F
from math import comb, factorial
from hashlib import sha256

def add(*ps):
    r = {}
    for p in ps:
        for i, a in p.items(): r[i] = r.get(i, 0) + a
    return {i:a for i,a in r.items() if a}

def mul(p, q):
    return add(*({i+j:a*b} for i,a in p.items() for j,b in q.items()))

def scale(p, a):
    return {i:a*v for i,v in p.items()}

def power(p,n):
    q = {0:1}
    for _ in range(n): q = mul(q,p)
    return q

w={0:F(105,47),94:-F(105,47)}
C=add(*(scale(power(w,k),F((-1)**k,factorial(2*k))) for k in range(4)))
M=add({35:24,94:35},scale({0:1,94:-1},70*F(59,94)))
P=add(scale(mul(C,M),59),mul({35:-1},power({0:24,59:35},2)))
H=scale(P,F(78074896,35))
assert all(v.denominator==1 for v in H.values())
h={i:int(v) for i,v in H.items()}
print('Primitive integer coefficients (degree, coefficient):')
print(sorted(h.items()))
D=max(h)
b=[sum(v*comb(D-i,j) for i,v in h.items() if j<=D-i) for j in range(D+1)]
assert b[:4]==[0]*4 and min(b[4:])>0
print('degree:',D)
print('first four transformed coefficients:',b[:4])
print('nonzero transformed coefficients:',len(b[4:]))
print('minimum positive transformed coefficient:',min(b[4:]))
print('minimum positive coefficient attained at j:',4+b[4:].index(min(b[4:])))
print('coefficient at j=4:',b[4])
print('SHA256 of integer array (comma-delimited ASCII, no trailing comma):')
print(sha256(','.join(str(x) for x in b).encode('ascii')).hexdigest())
\end{Verbatim}
\Needspace{9\baselineskip}
\subsection{Singleton-free Bell base differences}
\noindent Related to Section~\ref{section:R23K}. Two exact recurrences reproduce the finite arithmetic inputs. The all-order difference bound and 2-adic isometry are proved analytically in the text. Python standard library.

\begin{Verbatim}[fontsize=\scriptsize,breaklines,breakanywhere]
"""Exact, independent finite checks for the certificate in the note; stdlib only."""
from math import comb
M = 256
N = 400
d = [1]
for n in range(N):
    d.append(sum(comb(n,k) * d[n-k] for k in range(1,n+1)) % M)
# Independent tridiagonal transfer, no rounding.
row = [1]
for n in range(N+1):
    assert row[0] == d[n]
    new = [0] * (len(row)+1)
    for j,w in enumerate(row):
        new[j+1] += w
        new[j] += j*w
        if j: new[j-1] += j*w
    row = [v % M for v in new]
print("Two independent recurrences agree modulo 256 at indices 0..400.")
for a,h in [(1,24),(13,24),(4,12),(7,12),(10,12)]:
    x,y,z = [d[a+i*h] for i in range(3)]
    c = (x,(y-x)%M,(z-2*y+x)%M)
    print(f"a={a:2}, h={h}: D_a, first and second forward differences mod 256 = {c}")
def v2(x):
    assert x!=0
    return (x & -x).bit_length()-1
for a in (1,13):
    c1=(d[a+24]-d[a])%M
    c2=(d[a+48]-2*d[a+24]+d[a])%M
    assert (v2(c1),v2(c2))==(5,7)
assert d[1]==0 and v2(d[13])==7
assert d[4]%8==4 and (d[16]-d[4])%8==0
assert d[7]%4==d[10]%4==2
print("All base congruences in the proof verified.")
\end{Verbatim}
\Needspace{9\baselineskip}
\subsection{Sparse inverse-sine root-count certificate}
\noindent Related to Section~\ref{section:R23M}. The seven-term derivative polynomial and its Taylor coefficients are reconstructed exactly. Requires Python and SymPy.

\begin{Verbatim}[fontsize=\scriptsize,breaklines,breakanywhere]
"""Rebuild the seven-term certificate over the rationals; no sample points."""
from math import comb
from sympy import Poly, Rational as Q, expand, symbols
t,y=symbols("t y")
a,b,c,r=Q(1459,595),Q(864,595),Q(1741,1785),Q(1459,3482)
d=1-c
assert a==1+b and d==Q(44,1785)
D=b+y
deriv=lambda f:f.diff(t)+r*y/t*f.diff(y)
# The quotient rule for x/D, using x^2=(1-t)/c.
N=expand(D+2*(1-t)*deriv(D))
assert expand(N-(b+(1-2*r)*y+2*r*y/t))==0
H=Poly(expand(t*t*(D*N+2*(t-d)*(D*deriv(N)-2*N*deriv(D)))),t,y)
scale=Q(1073078451025,36)
coeff=dict(sorted((9882-3482*i-1459*j,scale*v) for (i,j),v in H.terms()))
expected={0:782041225,1459:-16375689960,2918:62852495616,3482:-45813373270,
          4941:-4231917040,6964:1747495365,8423:1038948064}
assert coeff==expected
print("Terms (exponent, coefficient):",list(coeff.items()))
mom=[sum(v*comb(i,j) for i,v in coeff.items()) for j in range(4)]
assert mom==[0,0,0,46206946595031268620]
print("Four Taylor coefficients at 1:",mom)
sign=[1 if v>0 else -1 for v in coeff.values()]
changes=sum(i!=j for i,j in zip(sign,sign[1:]))
assert changes==4 and coeff[0]>0
print("Sign changes:",changes)
print("Descartes budget: at most 4 positive roots counting multiplicity; 3 at 1 and at least 1 below 1; none above 1.")
print("PASS: exact certificate and derivative normalization.")
\end{Verbatim}
\Needspace{9\baselineskip}
\subsection{Exact cyclic cosine counterexample}
\noindent Related to Section~\ref{section:R22M}. Integer power comparisons and rational alternating-series bounds certify the strict reversal. Python standard library.

\begin{Verbatim}[fontsize=\scriptsize,breaklines,breakanywhere]
"""Exact rational certification; Python 3, standard library."""
from fractions import Fraction as F
from math import factorial
M=10**5
A=[1]+list(range(2,25,2)) # all coordinates a/24
def bracket(a,b):
    # For t=(a/24)**(b/24), find k=floor(M*t) exactly.
    lo,hi=0,M+1
    while hi-lo>1:
        k=(lo+hi)//2
        if k**24*24**b<=M**24*a**b:lo=k
        else:hi=k
    assert lo**24*24**b<=M**24*a**b<hi**24*24**b
    return lo
def C(z,j):
    return sum((-1)**r*z**(2*r)/factorial(2*r) for r in range(j+1))
def cosbox(k):
    # cos decreases on [0,1+1/M]; alternating expansion bounds.
    return C(F(k+1,M),5),C(F(k,M),4)
slo=shi=F(0)
print(' a   b   floor(M*x^x)  floor(M*x^y)')
for i,a in enumerate(A):
    b=A[i-1]; p,q=bracket(a,a),bracket(a,b)
    lo,hi=cosbox(p); ll,hh=cosbox(q)
    slo+=lo-hh;shi+=hi-ll
    print(f'{a:2}  {b:2}       {p:6}       {q:6}')
assert F(30,1000)<slo<=shi<F(31,1000)
print('strict enclosure for diagonal cosine sum minus cyclic cosine sum:', F(30,1000), F(31,1000))
print('bounds rounded outward to 1/1,000,000:',slo*10**6//1,-(-shi*10**6//1))
\end{Verbatim}
\Needspace{9\baselineskip}
\subsection{Plactic-centralizer mode coefficients}
\noindent Related to Section~\ref{section:R24K}. Exact hook and Newton-difference computations reproduce the coefficient row of the counterexample. The asymptotic mode-location theorem has a separate analytic proof. Python standard library.

\begin{Verbatim}[fontsize=\scriptsize,breaklines,breakanywhere]
from math import factorial,prod
def parts(n,b=None):
 if n==0: yield (); return
 for i in range(min(n,b or n),0,-1):
  for tail in parts(n-i,i):yield (i,)+tail
def cells(p):return [(i,j) for i,x in enumerate(p) for j in range(x)]
def hooks(p):return prod(p[i]-j+sum(x>j for x in p[i+1:]) for i,j in cells(p))
def evals(n):
 out=[0]*n
 for p in parts(n):
  mu=p[1:]; f=factorial(n)//hooks(p); h=hooks(mu)
  for m in range(1,n):
   out[m]+=f*(prod(m-1+j-i for i,j in cells(mu))//h)
 tmp=out[:];a=[]
 while tmp:
  a.append(tmp[0]);tmp=[y-x for x,y in zip(tmp,tmp[1:])]
 return out,a
if __name__=='__main__':
 for n in range(2,24):
  ev,a=evals(n)
  print(n,a.index(max(a)),(n+1)//2,a[:10],flush=True)
\end{Verbatim}
\Needspace{9\baselineskip}
\subsection{Independent first-row-state count for the mode counterexample}
\noindent Related to Section~\ref{section:R24K}. An independent integer dynamic program reproduces the counterexample without hook formulas. Python standard library.

\begin{Verbatim}[fontsize=\scriptsize,breaklines,breakanywhere]
from collections import Counter
from bisect import bisect_right
from verify import evals
def count(n,m):
 s=Counter({():1})
 for i in range(n):
  t=Counter()
  for row,weight in s.items():
   for x in range(1,m+1):
    j=bisect_right(row,x)
    t[row[:j]+(x,)+row[j+1:]]+=weight
  s=t
 return sum(v for row,v in s.items() if row and row[-1]==1)
for n in range(2,11):
 vals,coef=evals(n); direct=[count(n,m) for m in range(n)]
 assert vals==direct
 print(n,vals,coef)
\end{Verbatim}
\clearpage
\section{Source register and coverage}
The following register identifies all forty-one original manuscripts and the reconstructed checkpoint. The path is relative to the extracted round archive. Both source contributions for the even Eulerian-maxima topic point to the same principal section.
\begingroup\footnotesize
\begin{longtable}{@{}L{.12\textwidth}L{.08\textwidth}L{.19\textwidth}L{.53\textwidth}@{}}
\toprule ID & Section & Source & Original target\\\midrule\endfirsthead
\toprule ID & Section & Source & Original target\\\midrule\endhead
R13T & \ref{section:R13T} & Round 13, Taisho & Bennett--Nichols Conjecture 2\\
R16M & \ref{section:R16M} & Round 16, Meiji & Tian sub-conjecture (2)\\
R17M & \ref{section:R17M} & Round 17, Meiji & Dilcher--Ulas Conjecture 2 / 5.8\\
R18K & \ref{section:R18K} & Round 18, Keio & Benfield--Lippard Conjecture 10.5\\
R19K & \ref{section:R19K} & Round 19, Keio & Zhi-Wei Sun Conjecture 1.10\\
R20M & \ref{section:R20M} & Round 20, Meiji & Jekl Conjecture 6.1\\
R21T & \ref{section:R21T} & Round 21, Taisho & Galvin--Sharpe Question 1.10(4)\\
R22S & \ref{section:R22S} & Round 22, Showa & Sun--Chen Conjecture 7\\
R23M & \ref{section:R23M} & Round 23, Meiji & Xu--Sun--Chen Conjecture 1\\
R24Hq & \ref{section:R24Hq} & Round 24, Heisei & Karp--Vishnyakova--Zhang, Section 3, Example 8\\
R25M & \ref{section:R25M} & Round 25, Meiji & Sun--Chen Conjecture 10\\
R25K & \ref{section:R25K} & Round 25, Keio & Ambrus--Gargyan log-convexity conjecture\\
R25R & \ref{section:R25K} & Round 25, Reiwa & Ambrus--Gargyan log-convexity conjecture\\
R02K & \ref{section:R02K} & Round 2, Keio & Ulas Conjectures 4.2(2) and 4.4\\
R12K & \ref{section:R12K} & Round 12, Keio & Jenkins equal-shift extension\\
R18H & \ref{section:R18H} & Round 18, Heisei & Au Conjecture 4.3\\
R18T & \ref{section:R18T} & Round 18, Taisho & Guo--Schlosser Conjectures 6.3 and 6.1\\
R18R & \ref{section:R18R} & Round 18, Reiwa & Moll--Ramirez--Villamizar 6.3; Arnoczki--Nyul 3.4\\
R19M & \ref{section:R19M} & Round 19, Meiji & Sun cubic Franel assertion in 1.1(i)\\
R20T & \ref{section:R20T} & Round 20, Taisho & Bach--Remmel {1324,1423} conjecture\\
R20H & \ref{section:R20H} & Round 20, Heisei & Sagan--Swanson Conjecture 7.8\\
R20R & \ref{section:R20R} & Round 20, Reiwa & Alexandersson fixed-peak conjecture\\
R20S & \ref{section:R20S} & Round 20, Showa & Bobby Ocean, MathOverflow 171394\\
R21M & \ref{section:R21M} & Round 21, Meiji & Xu--Sun--Chen Conjecture 2\\
R21H & \ref{section:R21H} & Round 21, Heisei & Thejitha--Fathima Conjecture 7.1\\
R21R & \ref{section:R21R} & Round 21, Reiwa & Sun 2024 Conjecture 5.1(ii),(iii)\\
R21S & \ref{section:R21S} & Round 21, Showa & Chow consecutive-rank question\\
R22K & \ref{section:R22K} & Round 22, Keio & Xu--Chen Conjecture 4.1, remaining lower bound\\
R22H & \ref{section:R22H} & Round 22, Heisei & Qi Remark 5.9\\
R23K & \ref{section:R23K} & Round 23, Keio & Moll--Ramirez--Villamizar Conjecture 7.4\\
R23T & \ref{section:R23T} & Round 23, Taisho & Sandor Conjecture 1\\
R23S & \ref{section:R23S} & Round 23, Showa & Hui-Zuo Xu Conjecture 1\\
R23R & \ref{section:R23R} & Round 23, Reiwa & Wang--Zhang Remark 3.5\\
R24M & \ref{section:R24M} & Round 24, Meiji & Coulter--Do Conjecture 5.7(a), recursive interpretation\\
R24S & \ref{section:R24S} & Round 24, Showa & Ling Zhu Conjecture 4.1\\
R24Hb & \ref{section:R24Hb} & Round 24, Heisei & Karp--Vishnyakova--Zhang Conjecture 1, unimodality clause\\
R24T & \ref{section:R24T} & Round 24, Taisho & Galvin--Sharpe combinatorial explanation request\\
R17R & \ref{section:R17R} & Round 17, Reiwa & Brian Yi Sun Conjecture 4.7\\
R11H & \ref{section:R11H} & Round 11, Heisei & Cazenave Conjecture 2, identifier 20829\\
R12T & \ref{section:R12T} & Round 12, Taisho & Benfield--Lippard Conjectures 10.1 and 10.2\\
R22M & \ref{section:R22M} & Round 22, Meiji & Matejicka 2019, Section 3(a)\\
R24K & \ref{section:R24K} & Round 24, Keio & Sagan--Wilson Conjecture 7.3(d)\\
\bottomrule\end{longtable}
\endgroup
\subsection*{Exact source identities}
\begingroup\scriptsize\setlength{\parindent}{0pt}
\textbf{R13T}: \path{Tournament_Round_13_Final/taisho/paper.tex}\\
\texttt{SHA256 8ed04b59c56339a7ed1bc07be17dc739\allowbreak e3806c89e627bad90b818698d0db8555}\par\medskip
\textbf{R16M}: \path{Tournament_Round_16_Final/meiji/manuscript.tex}\\
\texttt{SHA256 f45be7634253383621ae6aa3b72b6593\allowbreak e5567c90db93c788087a1e7fc2613135}\par\medskip
\textbf{R17M}: \path{Tournament_Round_17_Final/round17/meiji/dilcher_ulas.tex}\\
\texttt{SHA256 f6e8f2639e4b185ee79dc5552abf3ccc\allowbreak 00c1fa032e3bf76145acdc5e663d71b3}\par\medskip
\textbf{R18K}: \path{Tournament_Round_18_Final/round18/keio/shifted_power_gaps.tex}\\
\texttt{SHA256 6ab99a1019906c37a021cc409b4c4d0f\allowbreak 0847a5449df5549b1ef404ee93533a1b}\par\medskip
\textbf{R19K}: \path{Tournament_Round_19_Final/round19/keio/sun_diagonal.tex}\\
\texttt{SHA256 bdc8c2a275cd002176b2b547afa042bd\allowbreak 1b4f81aecb207422e4ff20418728c625}\par\medskip
\textbf{R20M}: \path{Tournament_Round_20_Final/Tournament_Round_20/round20/meiji/jekl.tex}\\
\texttt{SHA256 ab3f0ea5c510f698f1be81d559ba0959\allowbreak 93d3bb7a39f28422b4a749aa3b0c2537}\par\medskip
\textbf{R21T}: \path{Tournament_Round_21_Final/Tournament_Round_21/round21/taisho/hypercomb.tex}\\
\texttt{SHA256 0edab3e851192225a46dfa97af27210d\allowbreak b09c004e1365665485fc7c925df6d82d}\par\medskip
\textbf{R22S}: \path{Tournament_Round_22_Final/Tournament_Round_22/round22/showa/manuscript.tex}\\
\texttt{SHA256 384053fed2e0538f95ecd7c25b057a28\allowbreak 4410f85056ae1c54fbb3cf3cf4edc3d3}\par\medskip
\textbf{R23M}: \path{Tournament_Round_23_Final/Tournament_Round_23/round23/meiji/meiji.tex}\\
\texttt{SHA256 a7945efa42c13d471741ef4292c46020\allowbreak c82e1c31e3d91ebea2b3b4c71c311173}\par\medskip
\textbf{R24Hq}: \path{Tournament_Round_24_Final/Tournament_Round_24/round24/heisei/reciprocal_qgamma.tex}\\
\texttt{SHA256 378c57e0ff40af29788c3d24b0ebc81f\allowbreak 847849bbac245617c4637a599700d446}\par\medskip
\textbf{R25M}: \path{Tournament_Round_25_Final/Tournament_Round_25/round25/meiji/gudermannian.tex}\\
\texttt{SHA256 0ef81d8df4b9150279ecb4787d8bc52b\allowbreak 447b2f3a46ce137f273c458d10f60ff8}\par\medskip
\textbf{R25K}: \path{Tournament_Round_25_Final/Tournament_Round_25/round25/keio/paper.tex}\\
\texttt{SHA256 54a329947d1fd662ea6c9a777b834fd9\allowbreak f0f2d33061d82f14613d1e889a30bac0}\par\medskip
\textbf{R25R}: \path{Tournament_Round_25_Final/Tournament_Round_25/round25/reiwa/eulerian.tex}\\
\texttt{SHA256 c5dd0968dc5e2c7549ab37ac9da18341\allowbreak fe2dd2b6de1e6d90a0bc8e358c26804e}\par\medskip
\textbf{R02K}: \path{Tournament_Round_02/keio/paper.tex}\\
\texttt{SHA256 d25a8f92a556444d0b4c70893948d404\allowbreak ef488efb93871f2e2c3e53210d2b27be}\par\medskip
\textbf{R12K}: \path{Tournament_Round_12/keio/paper.tex}\\
\texttt{SHA256 3a196621ef6fec852503908c22274be6\allowbreak 30151391384966908a9fc79175295e4a}\par\medskip
\textbf{R18H}: \path{Tournament_Round_18_Final/round18/heisei/au43.tex}\\
\texttt{SHA256 0baca950dcc751d5139b1855464cca6c\allowbreak 03f47caa4b9a6d8b290e8b5714992a29}\par\medskip
\textbf{R18T}: \path{Tournament_Round_18_Final/round18/taisho/guo_schlosser_checkpoint.tex}\\
\texttt{SHA256 938309302524b984d8fdbc356d467665\allowbreak 81ce69c013a11575a7ab408d461b79c0}\par\medskip
\textbf{R18R}: \path{Tournament_Round_18_Final/round18/reiwa/checkpoint.tex}\\
\texttt{SHA256 91db7cd13ecee5696ae82967a8c08b85\allowbreak fe78370b0352305451c3d184007a8844}\par\medskip
\textbf{R19M}: \path{Tournament_Round_19_Final/round19/meiji/sun_franel.tex}\\
\texttt{SHA256 509f1d99e955fafab97118aab3858c94\allowbreak afed460d42e3db7142c8cfa6d8b22445}\par\medskip
\textbf{R20T}: \path{Tournament_Round_20_Final/Tournament_Round_20/round20/taisho/paper.tex}\\
\texttt{SHA256 164bfa4def34e675d0c80af7c796dec8\allowbreak 6dd96adac836066b60fa3c4dfd5900f4}\par\medskip
\textbf{R20H}: \path{Tournament_Round_20_Final/Tournament_Round_20/round20/heisei/paper.tex}\\
\texttt{SHA256 974ce7265c59b8545b556d20d20beaed\allowbreak 35f1db79c05e26e1c3da7506ecd10d00}\par\medskip
\textbf{R20R}: \path{Tournament_Round_20_Final/Tournament_Round_20/round20/reiwa/manuscript.tex}\\
\texttt{SHA256 22a42bae4e71f8a9e98c6e14a6021f1a\allowbreak adf53a31043d88ac3a00a0384ca8e933}\par\medskip
\textbf{R20S}: \path{Tournament_Round_20_Final/Tournament_Round_20/round20/showa/ocean_differences.tex}\\
\texttt{SHA256 9616714a97886330b646a056fcacc3fd\allowbreak 5ca2aca8543f36b7f43407e2fb6f99de}\par\medskip
\textbf{R21M}: \path{Tournament_Round_21_Final/Tournament_Round_21/round21/meiji/meiji_r21.tex}\\
\texttt{SHA256 9bd8f53a835ce64eec763caa73630316\allowbreak 251ca8bd66c5b68b267e13c78a5bc27f}\par\medskip
\textbf{R21H}: \path{Tournament_Round_21_Final/Tournament_Round_21/round21/heisei/overcolored.tex}\\
\texttt{SHA256 c0f7d1aed8ab883b093c0d0362b51b47\allowbreak f635f171ddb3b5b7847dc1ce279427ab}\par\medskip
\textbf{R21R}: \path{Tournament_Round_21_Final/Tournament_Round_21/round21/reiwa/manuscript.tex}\\
\texttt{SHA256 4669a23b0eaab7743030391b4a44ed21\allowbreak f439d045cd9f433fa64fff25ab4f92ef}\par\medskip
\textbf{R21S}: \path{Tournament_Round_21_Final/Tournament_Round_21/round21/showa/chow_type_d.tex}\\
\texttt{SHA256 c76c485c54679f7520db032d017845e8\allowbreak bd7845151060eb99388bf966d1fcc5a1}\par\medskip
\textbf{R22K}: \path{Tournament_Round_22_Final/Tournament_Round_22/round22/keio/huygens.tex}\\
\texttt{SHA256 db8dc0900c78aa2179a2eebf1778d272\allowbreak a1938a5fa1a241e0e0cb40edfd30376c}\par\medskip
\textbf{R22H}: \path{Tournament_Round_22_Final/Tournament_Round_22/round22/heisei/qi_threshold.tex}\\
\texttt{SHA256 24f5d3344b06105cd030ddb1f19acb45\allowbreak 132982bfc6d62b2d1d6aaab5c664fb21}\par\medskip
\textbf{R23K}: \path{Tournament_Round_23_Final/Tournament_Round_23/round23/keio/associated_bell.tex}\\
\texttt{SHA256 613b78c4383ef44257f36cc8f92ba605\allowbreak b08dcda5970425a15e4a1403170b5f2e}\par\medskip
\textbf{R23T}: \path{Tournament_Round_23_Final/Tournament_Round_23/round23/taisho/paper.tex}\\
\texttt{SHA256 984cfe545271b3ea115622778bf288c7\allowbreak c23e241f3b2ab95808dad80ad8fd76d0}\par\medskip
\textbf{R23S}: \path{Tournament_Round_23_Final/Tournament_Round_23/round23/showa/xu_sinh.tex}\\
\texttt{SHA256 1713553a7a0b226897a591327a2f7290\allowbreak 45bfab9c47f261b45e30aecbf407e1e7}\par\medskip
\textbf{R23R}: \path{Tournament_Round_23_Final/Tournament_Round_23/round23/reiwa/reiwa.tex}\\
\texttt{SHA256 07a984a0c9dc35a80a2013e6bdcb2b2b\allowbreak 87b7965e453a9ce3edc8307721169f4a}\par\medskip
\textbf{R24M}: \path{Tournament_Round_24_Final/Tournament_Round_24/round24/meiji/manuscript.tex}\\
\texttt{SHA256 3f99dd117321f1fc50f17c51473c1552\allowbreak 8e22df845a7275060af72eb0fcb78da3}\par\medskip
\textbf{R24S}: \path{Tournament_Round_24_Final/Tournament_Round_24/round24/showa/zhu.tex}\\
\texttt{SHA256 ce8b1dd2ca51d662d3d6fd8c6e82699c\allowbreak 9e0bd4b467198d1a9e67b7b7dffbe20c}\par\medskip
\textbf{R24Hb}: \path{Tournament_Round_24_Final/Tournament_Round_24/round24/heisei/bessel_ratio.tex}\\
\texttt{SHA256 d647058e150b0a4ae5a71c0f3ec09dad\allowbreak 354f30f7fd32bc035d3337f284bdb791}\par\medskip
\textbf{R24T}: \path{Tournament_Round_24_Final/Tournament_Round_24/round24/taisho/paper.tex}\\
\texttt{SHA256 7a3d669c1baa8ac1940e9e1e2c7268e2\allowbreak 66e8c1b60b18959c9430fe04a6d1907f}\par\medskip
\textbf{R17R}: \path{Editorial reconstruction: brian_yi_sun_reconstructed.tex}\\
\texttt{SHA256 9c4deba2e29643b7c0431acf71ad1a11\allowbreak 11bb47f0f3c49d0e697643b6a4d54d85}\par\medskip
\textbf{R11H}: \path{Tournament_Round_11/heisei/manuscript.tex}\\
\texttt{SHA256 21522db8b8b2707952366678052c0199\allowbreak b8f5e093a07c4b85b4f35f53977fd17d}\par\medskip
\textbf{R12T}: \path{Tournament_Round_12/taisho/paper.tex}\\
\texttt{SHA256 85d975e07269d55b98d1ea701785be10\allowbreak 4036aabb8e43f96736d4e95843038ba4}\par\medskip
\textbf{R22M}: \path{Tournament_Round_22_Final/Tournament_Round_22/round22/meiji/cycle.tex}\\
\texttt{SHA256 13391b7adc9f5bb337bfce7d9bf9f21b\allowbreak aead1018e65274117dbc639ac6bfadfd}\par\medskip
\textbf{R24K}: \path{Tournament_Round_24_Final/Tournament_Round_24/round24/keio/paper.tex}\\
\texttt{SHA256 a8940808a302f33da6f27f3001b9fb22\allowbreak 5d50c0c55251ba49ec84753b395a2434}\par\medskip
\endgroup
\clearpage
\begin{thebibliography}{999}
\bibitem{R13T:BN}
P. Bennett and J. Nichols, \emph{Some progress on $t$-tone coloring}, arXiv:2510.13382v1 (2025). \url{https://arxiv.org/html/2510.13382v1}
\bibitem{R13T:Keevash}
P. Keevash, \emph{The existence of designs II}, arXiv:1802.05900 (2018), Definition~1.6 and Theorem~1.7. \url{https://arxiv.org/abs/1802.05900}
\bibitem{R13T:FMW}
R. Fabila-Monroy and D. R. Wood, \emph{Decompositions of complete multipartite graphs into complete graphs}, arXiv:1109.3508. \url{https://arxiv.org/abs/1109.3508}
\bibitem{R16M:Tian}
Z.~Zeng, \'A.~Pint\'er, X.~Fu and J.~P.~Tian,
Tian's conjecture on the prime factorization of the binomial
coefficient $\binom{n+1}{2}$,
\emph{Mathematics} 14 (2026), 127; first published
29 December 2025.
\href{https://doi.org/10.3390/math14010127}{doi:10.3390/math14010127}.
\bibitem{R17M:DU}
K.~Dilcher and M.~Ulas,
\emph{Divisibility and Arithmetic Properties of a Class of Sparse Polynomials},
Integers \textbf{22} (2022), A8,
\href{https://math.colgate.edu/~integers/w8/w8.pdf}{publisher PDF}.
Conjecture~2 is numbered~5.8 in \href{https://arxiv.org/abs/2008.13475v2}{arXiv:2008.13475v2}.
\bibitem{R17M:STP}
W.~Su, X.~Tang, and A.~Pott,
\emph{A Note on a Conjecture for Balanced Elementary Symmetric Boolean Functions},
\href{https://arxiv.org/abs/1203.1418v2}{arXiv:1203.1418v2} (2012),
Theorem~1.
\bibitem{R18K:BL}
B.~Benfield and O.~Lippard, \emph{Integers that are not the sum of
positive powers}, arXiv:2404.08193, v2 (31 March 2025),
Conjecture~10.5 and \S10.2.
\url{https://arxiv.org/html/2404.08193v2}.
The same conjecture is numbered 4.5 in v1 (12 April 2024).
\bibitem{R18K:Zenkin}
A.~A. Zenkin, \emph{The generalized Waring problem: A new property
of positive integers}, Mathematical Notes \textbf{58} (1995),
933--937. \url{https://doi.org/10.1007/BF02304770}.
\bibitem{R19K:Sun}
Z.-W. Sun, \emph{On divisibility of binomial coefficients},
J. Aust. Math. Soc. \textbf{93} (2012), 189--201.
\href{https://doi.org/10.1017/S1446788712000171}{doi:10.1017/S1446788712000171}.
\bibitem{R19K:HY}
Y. He and P. Yuan, \emph{Results on some conjectures about binomial
coefficients}, Bull. Math. Soc. Sci. Math. Roumanie \textbf{67}
(2024), 429--448.
\href{https://ssmr.ro/bulletin/pdf/67-4/articol_4.pdf}{Journal PDF}.
\bibitem{R19K:SunOEIS}
Z.-W. Sun, OEIS sequence \href{https://oeis.org/A390107}{A390107},
2025.
\bibitem{R20M:jekl}
J. Jekl, Proper version of Cusa--Huygens' inequality and
Djokvie's type inequality for parabolic trigonometric functions,
\emph{Math. Inequal. Appl.} \textbf{28} (2025), 725--738.
\href{https://doi.org/10.7153/mia-2025-28-42}{doi:10.7153/mia-2025-28-42}.
\bibitem{R21T:vone}
D. Galvin and C. Sharpe, \emph{Independent set sequence
of linear hyperpaths}, arXiv:2409.15555v1 (2024), Sec.~4.
\href{https://arxiv.org/abs/2409.15555v1}{Original preprint}.
\bibitem{R21T:gs}
D. Galvin and C. Sharpe, \emph{Independent set sequence of
some linear hypertrees}, Australas. J. Combin. \textbf{94} (2026),
93--121, Theorem 1.7, Corollary 1.8, Note 1.9, Question 1.10(4).
\href{https://ajc.maths.uq.edu.au/pdf/94/ajc_v94_p093.pdf}{Journal article}.
\bibitem{R21T:eventual}
OpenAI Codex GPT-6-Astra, \emph{Eventual strict
log-concavity of weighted linear hypercombs}, aiXiv:260910.000002
(2026). \href{https://aixiv.science/abs/aixiv.260910.000002}{Preprint}.
\bibitem{R22S:SC}
J.-L. Sun and C.-P. Chen,
Sharp Shafer--Fink, Lazarevi\'c and Cusa type inequalities for the Gudermannian
and inverse Gudermannian functions, \emph{Appl. Anal. Discrete Math.}
\textbf{19} (2025), 130--149.
\href{https://doi.org/10.2298/AADM230809021S}{doi:10.2298/AADM230809021S}.
\bibitem{R23M:XSC}
S.-W. Xu, J.-L. Sun, and C.-P. Chen,
\emph{New Shafer--Fink type inequalities for inverse trigonometric,
inverse hyperbolic and arc lemniscate functions},
Applicable Analysis and Discrete Mathematics \textbf{20} (2026),
Conjecture 2, (27).
\href{https://doi.org/10.2298/AADM250324010X}{doi:10.2298/AADM250324010X}.
\bibitem{R24Hq:KVZ}
D. Karp, A. Vishnyakova and Y. Zhang, Unimodality preservation
by ratios of functional series and integral transforms, \emph{Results Math.}
\textbf{80} (2025), 112.
\href{https://doi.org/10.1007/s00025-025-02431-4}{doi:10.1007/s00025-025-02431-4};
\href{https://arxiv.org/html/2408.01755v2}{arXiv:2408.01755v2}, Section 3.
\bibitem{R24Hq:S}
D. Salazar, Order-moment transport and Hankel determinants
in special-function inequalities,
\href{https://arxiv.org/abs/2606.31647v2}{arXiv:2606.31647v2} (2026).
\bibitem{R25M:dlmf}
NIST Digital Library of Mathematical Functions,
\href{https://dlmf.nist.gov/4.23.E39}{(4.23.39)}
and \href{https://dlmf.nist.gov/4.36.E2}{(4.36.2)}.
\bibitem{R25K:AG}
G. Ambrus and B. G\'argy\'an,
Estimates on the decay of the Laplace--P\'olya integral,
\emph{Bull. London Math. Soc.} \textbf{57} (2025), 3360--3379.
\href{https://doi.org/10.1112/blms.70157}{doi:10.1112/blms.70157}.
\bibitem{R25K:LN}
L. Lesieur and J.-L. Nicolas,
On the Eulerian numbers $M_n=\max_{1\le k\le n}A(n,k)$,
\emph{European J. Combin.} \textbf{13} (1992), 379--399.
\bibitem{R02K:Ulas}
M. Ulas, \emph{Strong arithmetic property of certain Stern
polynomials}, arXiv:1909.10844v1 (24 September 2019), specifically
Conjectures 4.2(2) and 4.4, p.~16; proofs of Theorems 2.3 and 3.1.
\url{https://arxiv.org/abs/1909.10844v1}.
\bibitem{R02K:Publisher}
M. Ulas, same title, Publicationes Mathematicae Debrecen
\textbf{96} (2020), 401--422, DOI 10.5486/PMD.2020.8683.
Publisher metadata: \url{https://publi.math.unideb.hu/paper/2376}.
\bibitem{R02K:Altizio}
D. Altizio, \emph{Zeros of Stern polynomials in the complex
plane}, arXiv:2511.03847v1 (2025).
\url{https://arxiv.org/abs/2511.03847}.
\bibitem{R02K:Corpus}
QEEDDD, \emph{Verifiable Mathematical Reasoning Corpus},
research-draft repository, especially manuscript [071].
\url{https://github.com/QEEDDD/Verifiable-Mathematical-Reasoning-Corpus}.
\bibitem{R12K:singmaster1971}
D.\ Singmaster, How often does an integer occur as a binomial coefficient? \emph{Amer. Math. Monthly} 78 (1971), 385--386. \href{https://doi.org/10.2307/2316907}{doi:10.2307/2316907}.
\bibitem{R12K:jenkins}
H.\ Jenkins, Repeated binomial coefficients and high-degree curves. \emph{Integers} 16 (2016), A69, especially pp.~4, 8--9. \href{https://doi.org/10.5281/zenodo.10623246}{doi:10.5281/zenodo.10623246}.
\bibitem{R12K:mrstt}
K.\ Matom\"aki, M.\ Radziwi\l\l, X.\ Shao, T.\ Tao, and J.\ Ter\"av\"ainen, Singmaster's conjecture in the interior of Pascal's triangle. \emph{Q. J. Math.} 73 (2022), 1137--1177. \href{https://doi.org/10.1093/qmath/haac006}{doi:10.1093/qmath/haac006}.
\bibitem{R12K:lind}
D.\ A.\ Lind, The quadratic field \(\mathbb Q(\sqrt 5)\) and a certain Diophantine equation. \emph{Fibonacci Quart.} 6 (1968), 86--93.
\bibitem{R12K:singmaster1975}
D.\ Singmaster, Repeated binomial coefficients and Fibonacci numbers. \emph{Fibonacci Quart.} 13 (1975), 295--298. \href{https://www.fq.math.ca/Scanned/13-4/singmaster.pdf}{Journal archive}.
\bibitem{R12K:blokhuis}
A.\ Blokhuis, A.\ Brouwer, and B.\ de Weger, Binomial collisions and near collisions. \emph{Integers} 17 (2017), A64. \href{https://arxiv.org/abs/1707.06893}{arXiv:1707.06893}.
\bibitem{R18H:au}
K.~C.~Au, \emph{Discovering hypergeometric series with
harmonic numbers via Wilf--Zeilberger seeds}, arXiv:2602.08721v4,
9 November 2026, \S2.2 and Example VIII (Conjecture 4.3).
\href{https://arxiv.org/html/2602.08721v4}{arXiv full text}.
\bibitem{R18H:schneider}
C.~Schneider, \emph{Parameterized telescoping
proves algebraic independence of sums}, Ann. Comb.\ \textbf{14}
(2010), 533--552. \href{https://arxiv.org/abs/0808.2596}{arXiv:0808.2596}.
\bibitem{R18T:GS}
V. J. W. Guo and M. J. Schlosser, New \(q\)-analogues
of Van Hamme's (E.2) supercongruence and of a supercongruence by Swisher,
\emph{Results Math.} \textbf{78} (2023), 105.
\href{https://doi.org/10.1007/s00025-023-01885-8}{Publisher}.
\bibitem{R18T:LR}
L. Long and R. Ramakrishna, Some supercongruences occurring in
truncated hypergeometric series, \emph{Adv. Math.} \textbf{290}
(2016), 773--808. \href{https://arxiv.org/abs/1403.5232}{arXiv:1403.5232}.
\bibitem{R18R:MRV}
V. H. Moll, J. L. Ram\'{i}rez and D. Villamizar,
\emph{Combinatorial and arithmetical properties of the restricted and associated
Bell and factorial numbers}, J. Combin. \textbf{9} (2018), 693--720,
Conjecture 6.3, p.~711. \url{https://arxiv.org/abs/1707.08138}.
\bibitem{R18R:AN}
T. Arn\'{o}czki and G. Nyul,
\emph{On restricted $r$-Stirling numbers, also known as $r$-Bessel numbers},
Australas. J. Combin. \textbf{89} (2024), 49--60,
Theorems 3.1, 3.2, 3.5 and Conjecture 3.4.
\url{https://ajc.maths.uq.edu.au/pdf/89/ajc_v89_p049.pdf}.
\bibitem{R18R:Menon}
K. V. Menon,
\emph{On the convolution of logarithmically concave sequences},
Proc. Amer. Math. Soc. \textbf{23} (1969), 439--441.
\url{https://doi.org/10.1090/S0002-9939-1969-0246012-5}.
\bibitem{R18R:Godsil}
C. D. Godsil,
\emph{Algebraic matching theory}, Electron. J. Combin.
\textbf{2} (1995), R8. \url{https://doi.org/10.37236/1202}.
\bibitem{R19M:Sun}
Z.-W. Sun, \emph{A family of polynomials and related congruences and
series}, Front. Comb. Number Theory \textbf{2} (2026), 1--22.
\href{https://doi.org/10.3934/fcnt.2026009}{doi:10.3934/fcnt.2026009};
\href{https://arxiv.org/abs/2505.02767}{arXiv:2505.02767}.
\bibitem{R19M:MO}
Z.-W. Sun, author's comment dated 21 September 2024 on
\href{https://mathoverflow.net/questions/479301/on-the-polynomials-sum-k-0n-binomnkkm-qk}
{MathOverflow question 479301}. The cubic assertion is in the comment,
not the main text of that question.
\bibitem{R20T:BR}
Q.~T.~Bach and J.~B.~Remmel, Generating functions for descents over permutations which avoid sets of consecutive patterns, \emph{Australas. J. Combin.} \textbf{64}(1) (2016), 194--231. \href{https://ajc.maths.uq.edu.au/pdf/64/ajc_v64_p194.pdf}{Journal PDF}, especially pp.~207--208.
\bibitem{R20H:ss}
B. E. Sagan and J. P. Swanson, \emph{$q$-Stirling numbers in type B}, European J. Combin. 118 (2024), 103899. \href{https://doi.org/10.1016/j.ejc.2023.103899}{DOI}; \href{https://arxiv.org/abs/2205.14078}{arXiv:2205.14078}.
\bibitem{R20H:k}
A. Kasraoui, \emph{Average values of some $Z$-parameters in a random set partition}, Electron. J. Combin. 18 (2011), P228. \href{https://arxiv.org/abs/1304.4813}{arXiv:1304.4813}.
\bibitem{R20H:j}
S. Janson, \emph{Moment convergence in conditional limit theorems}, J. Appl. Probab. 38 (2001), 421--437. \href{https://www2.math.uu.se/~svantejs/papers/sj132.pdf}{Author's copy}.
\bibitem{R20R:A}
P. Alexandersson, \emph{Real-rooted polynomials: open
problems}, SymCat, Catalan families, conjecture attributed to 2023,
\href{https://www.symmetricfunctions.com/realRootedOpenProblems.htm}{author's problem list}.
\bibitem{R20R:Arm}
D. Armstrong, \emph{Generalized noncrossing partitions
and combinatorics of Coxeter groups}, Mem. Amer. Math. Soc. 202 (2009),
no.~949, \href{https://arxiv.org/abs/math/0611106}{preprint}.
\bibitem{R20R:B}
M. B\'ona et al., \emph{A combinatorial proof of a symmetry
for a refinement of the Narayana numbers}, Electron. J. Combin. 32 (2025),
P2.46, \href{https://arxiv.org/abs/2212.10586}{preprint}.
\bibitem{R20R:WZ}
B. Wang and C.~X.~T. Zhang, \emph{Interlacing property
of a family of generating polynomials over Dyck paths},
Electron. J. Combin. 31 (2024), P1.63,
\href{https://doi.org/10.37236/12375}{doi:10.37236/12375}.
\bibitem{R20S:ocean}
Bobby Ocean, \emph{Positive, Uni-modal, Log-concave Combinatorics}, MathOverflow question 171394, 9 June 2014 (edited 19 June). \url{https://mathoverflow.net/q/171394}.
\bibitem{R21H:TF}
M.~P. Thejitha and S.~N. Fathima, \emph{Overcolored Partition $k$-tuples Restricted by Parity of the Parts}, \href{https://arxiv.org/pdf/2609.03926v4}{arXiv:2609.03926v4}, Conjecture 7.1, p.~29 (2026).
\bibitem{R21H:TFold}
The same authors, \emph{Congruences for Overcolored Partition $k$-tuples Restricted by Parity of the Parts}, \href{https://arxiv.org/abs/2609.32614}{arXiv:2609.32614}, v1, Conjecture 4.1; withdrawn in v2.
\bibitem{R21H:RS}
S.~Radu and J.~A. Sellers, Congruence properties modulo 5 and 7 for the $\mathrm{pod}$ function, \emph{Int. J. Number Theory} \textbf{7} (2011), 2249--2259, \href{https://doi.org/10.1142/S1793042111005064}{doi:10.1142/S1793042111005064}.
\bibitem{R21R:Sun}
Z.-W. Sun, \emph{Problems and results on determinants involving Legendre symbols}, Conjecture 5.1 and equation (5.1), arXiv:2405.03626v8 (2024). \url{https://arxiv.org/html/2405.03626v8}.
\bibitem{R21R:LY}
G. Li and X. Yan, \emph{On Sun's conjectures for truncated Jacobi-symbol determinants via supersingular elliptic curves}, Theorem 1.1 and Remark 1.2, arXiv:2608.08509 (2026). \url{https://doi.org/10.1016/j.ffa.2026.102946}.
\bibitem{R21R:K}
D. Krachun, F. Petrov, Z.-W. Sun and M. Vsemirnov, \emph{On some determinants involving Jacobi symbols}, Finite Fields Appl. 64 (2020), 101672. \url{https://arxiv.org/abs/1812.08080}.
\bibitem{R21R:Trees}
\emph{The number of spanning trees in circulant graphs, its arithmetic properties and asymptotic}, arXiv:1711.00175. \url{https://arxiv.org/abs/1711.00175}.
\bibitem{R21S:chow}
C.-O. Chow, \emph{New proofs of interlacing of zeros of Eulerian polynomials}, J. Math. Anal. Appl. \textbf{510} (2022), 126019. \href{https://doi.org/10.1016/j.jmaa.2022.126019}{doi:10.1016/j.jmaa.2022.126019}.
\bibitem{R21S:sv}
C.~D. Savage and M. Visontai, \emph{The Eulerian polynomials of type D have only real roots}, DMTCS Proc. AS, FPSAC 2013, 527--538. \href{https://dmtcs.episciences.org/articles/2321/download}{Publisher's full text}.
\bibitem{R22K:xc}
S.-W. Xu and C.-P. Chen, Sharp Wilker and Huygens type inequalities for the Gudermannian and inverse Gudermannian functions, \emph{J. Inequal. Appl.} \textbf{2024}, 149. \href{https://doi.org/10.1186/s13660-024-03232-z}{doi:10.1186/s13660-024-03232-z}. See Conjecture 4.1, Theorem 4.2, and the Euler-number bounds in Section 2.
\bibitem{R22H:qi}
F. Qi, \emph{Properties of modified Bessel functions and completely
monotonic degrees of differences between exponential and trigamma
functions}, Math. Inequal. Appl. \textbf{18} (2015), 493--518,
\href{https://doi.org/10.7153/mia-18-37}{doi:10.7153/mia-18-37},
Theorem 1.2, (1.19), and Remark 5.9.
\bibitem{R23K:MRV}
V.~H. Moll, J.~L. Ram\'irez and D.~Villamizar,
Combinatorial and arithmetical properties of the restricted and associated
Bell and factorial numbers, \emph{Journal of Combinatorics}
\textbf{9} (2018), 693--720, Conjecture~7.4, p.~713.
\href{https://arxiv.org/abs/1707.08138}{arXiv:1707.08138}.
\bibitem{R23T:plain1}
J. S\'andor, Extensions of D'Aurizio's
trigonometric inequality, \emph{Notes Number Theory Discrete Math.}
\textbf{23} (2017), no.~2, 81--83; definition of $g$ on p.~82 and
Conjecture~1 on p.~83.
\bibitem{R23S:xu}
H.-Z. Xu, Sharp bounds for the tangent mean by weighted power mean of arithmetic and geometric means, \emph{Open Mathematics}, Conjecture 1 and Lemma 2.
\href{https://doi.org/10.1515/math-2025-0248}{doi:10.1515/math-2025-0248}.
\bibitem{R23S:nw}
M. Nowicka and A. Witkowski, Optimal bounds for the tangent and hyperbolic sine means, \emph{Aequationes Math.} \textbf{94} (2020), 817--827. \href{https://doi.org/10.1007/s00010-020-00705-6}{doi:10.1007/s00010-020-00705-6}.
\bibitem{R23S:nist}
NIST Digital Library of Mathematical Functions, \href{https://dlmf.nist.gov/4.36.E1}{Eq.\ 4.36.1}, the product for $\sinh z$.
\bibitem{R23R:plain1}
B. Wang and C.~X.~T. Zhang, \textit{Some criteria for higher order Tur\'an
inequalities in the spirit of Ma\v{r}\'{i}k's theorem}, Eq.~(3.1) and
\href{https://cfc.nankai.edu.cn/_upload/article/files/7c/59/537dfc2d4ccc940465b4bb926c83/ded5d338-d5c2-4a81-9bed-54e7b8db797d.pdf#page=10}{Remark 3.5 (author-hosted paper)}.
\bibitem{R24M:cd}
X. Coulter and N. Do, \emph{From Weingarten calculus for real Grassmannians to deformations of monotone Hurwitz numbers and Jucys--Murphy elements} (2025), Definitions 4.1 and 5.1, (18)--(19), and Conjectures 5.4 and 5.7. \url{https://arxiv.org/abs/2506.04002}.
\bibitem{R24M:mm}
R. Mickler and A. Moll, \emph{Spectral Theory of the Nazarov--Sklyanin Lax Operator}, SIGMA 19 (2023), 063. \url{https://doi.org/10.3842/SIGMA.2023.063}.
\bibitem{R24M:m}
R. Mickler, \emph{Jack Littlewood--Richardson Coefficients and the Nazarov--Sklyanin Lax Operator}, arXiv:2306.11115v2 (2023), Sections 2.2--3.2, in particular Theorem 3.4. \url{https://arxiv.org/abs/2306.11115}.
\bibitem{R24M:prior}
trureturing, \emph{Deformed Jucys--Murphy operators}, public certificate for 5.4(a) and separate record for 5.8 (2026). \url{https://the-omega-institute.github.io/trureturing-mdbook/Library/Certificates/coulter2025weingarten.html}.
\bibitem{R24S:zhu}
L. Zhu, New Cusa--Huygens type inequalities,
\emph{AIMS Mathematics} \textbf{5} (2020), 5320--5331.
\href{https://doi.org/10.3934/math.2020341}{doi:10.3934/math.2020341},
Conjecture 4.1, p.~5328.
\bibitem{R24Hb:DLMF}
NIST Digital Library of Mathematical Functions,
\href{https://dlmf.nist.gov/10.40}{Section 10.40}.
\bibitem{R24Hb:TT}
Trureturing project,
\href{https://www.trureturing.org/research/karp-vishnyakova-zhang-2025-bessel-ratio-log-concavity/}{Log-concavity of modified Bessel ratios} (2026 record).
\bibitem{R24T:gs}
D. Galvin and C. Sharpe, \emph{Independent set sequence of
some linear hypertrees}, Australas. J. Combin. \textbf{94} (2026),
93--121. \href{https://ajc.maths.uq.edu.au/pdf/94/ajc_v94_p093.pdf}{Journal version}.
\bibitem{R24T:mrv}
V. H. Moll, J. L. Ram\'irez and D. Villamizar,
\emph{Combinatorial and arithmetical properties of the restricted and
associated Bell and factorial numbers}, J. Combin. \textbf{9} (2018),
693--720. \href{https://arxiv.org/abs/1707.08138}{Author preprint}.
\bibitem{R17R:bys:sun}
B. Y. Sun,
\emph{On a ratio monotonicity conjecture of a new kind of numbers},
Journal of Inequalities and Applications \textbf{2018}, article 24
(2018), Conjecture 4.7.
\href{https://doi.org/10.1186/s13660-018-1614-1}
{doi:10.1186/s13660-018-1614-1}.
\bibitem{R17R:bys:archive}
Reiwa, Round 17 research checkpoint, 19 November 2026,
\texttt{CHECKPOINT.md}, \texttt{verify.py}, and \texttt{results.txt},
in \emph{Math Tournament Rounds 01--25},
\texttt{Tournament\_Round\_17\_Final/round17/reiwa}.
Unadjudicated checkpoint; the present article is an editorial
reconstruction, not an original frozen Round 17 manuscript.
\bibitem{R11H:cazenave}
T. Cazenave, \emph{Lemmas and Conjectures in Spectral Graph Theory},
2 November 2026, Section 7, Conjecture 2 (\#20829).
\href{https://arxiv.org/html/2611.03879v1}{arXiv:2611.03879v1}.
\bibitem{R11H:prior}
The Omega Institute, \emph{Cazenave Conjecture 5: the
Laplacian--rank bound}, repository record frozen 13 November 2026.
\href{https://www.trureturing.org/research/op-cazenave2611c5/}{Project report}.
This concerns a different conjecture, not \#20829.
\bibitem{R12T:BL1}
B. Benfield and O. Lippard, \emph{Integers that are not the sum of positive powers}, arXiv:2404.08193v1 (2024), p.~11. \url{https://arxiv.org/pdf/2404.08193v1}
\bibitem{R12T:BL2}
B. Benfield and O. Lippard, \emph{Integers that are not the sum of positive powers}, arXiv:2404.08193v2 (2025), \S\S1, 10.1--10.2. \url{https://arxiv.org/html/2404.08193v2#S10.SS1}
\bibitem{R12T:AgMar}
F. Aguil\'o-Gost and C. Mariju\'an, \emph{Classification of numerical $3$-semigroups by means of $L$-shapes}, Semigroup Forum \textbf{88} (2014), 670--688. \url{https://doi.org/10.1007/s00233-013-9560-1}
\bibitem{R12T:Herzog}
J. Herzog, \emph{Generators and relations of abelian semigroups and semigroup rings}, Manuscripta Math. \textbf{3} (1970), 175--193. \url{https://doi.org/10.1007/BF01273309}
\bibitem{R22M:M19}
L. Matej\'i\v{c}ka, Proof of some new conjectures on algebraic and algebraic-trigonometric inequalities, \emph{Int. J. Open Problems Comput. Math.} \textbf{12}(2) (2019), \href{https://www.ijopcm.org/Vol/2019/2.4.pdf}{open-problem section}, pp. 39--40.
\bibitem{R22M:M20}
L. Matej\'i\v{c}ka, Proof of one open inequality of Laub--Ilani type, \emph{J. Math. Inequal.} \textbf{14} (2020), 83--98, \href{https://doi.org/10.7153/jmi-2020-14-07}{doi:10.7153/jmi-2020-14-07}.
\bibitem{R22M:O}
A. Y. \"Ozban, New algebraic-trigonometric inequalities of Laub--Ilani type, \emph{Bull. Aust. Math. Soc.} \textbf{96} (2017), 87--97, \href{https://doi.org/10.1017/S0004972717000156}{doi:10.1017/S0004972717000156}.
\bibitem{R24K:SW}
B.~E. Sagan and A.~N. Wilson, Centralizers in the plactic monoid, \emph{Semigroup Forum} \textbf{110} (2025), 724--744. \href{https://users.math.msu.edu/users/bsagan/Papers/Old/cpm-pub.pdf}{Published article}.
\bibitem{R24K:SZ}
B.~E. Sagan and C. Zhao, Properties of plactic monoid centralizers, \emph{Semigroup Forum} \textbf{113} (2026), 242--262. \href{https://doi.org/10.1007/s00233-026-10652-4}{doi:10.1007/s00233-026-10652-4}.
\end{thebibliography}
\end{document}
